% !TeX program = pdflatex
% beautybook — decorative book class with coloured theorem boxes and a
% two-tone chapter-heading treatment, adapted from the CTAN/GitHub
% BeautyLaTeX/latex-template package.
%
% Simplified from upstream: the upstream ships four cover styles (cn/en/
% enfig/birkar) and a 36MB set of photographic chapter-corner decorations
% (bookstyle/titleimages/, Pixabay-style stock photos with no separate
% licence grant of their own). This starter uses ONLY the text-only "en"
% cover (no background photo). This handout uses the supplied Klinik AA
% logo (bookstyle/LogoAA.png), with a soft white glow on the cover.
% Local modification (2026-10-29): Indonesian labels and a unified Klinik AA
% palette; see beautybook.cls for template attribution and licensing.
\documentclass[lang=id,12pt]{beautybook}
\RequirePackage[utf8]{inputenc}
% Arial untuk judul materi; times sesudahnya menjaga font sans desain sampul.
\RequirePackage[scaled]{uarial}
\RequirePackage{times}
\RequirePackage[T1]{fontenc}
\RequirePackage{microtype}
\RequirePackage[round]{natbib}

% Font isi: Computer Modern (cmr), sama dengan Tugas 1 Analisis Fungsional.
% Paket times tetap menyediakan font desain lama (Times/Helvetica/Courier).
% OT1 memakai Computer Modern vektor yang tersedia di MiKTeX ini;
% versi T1/EC di instalasi ini berupa bitmap dan tidak mendukung font expansion.
\renewcommand{\encodingdefault}{OT1}
\renewcommand{\rmdefault}{cmr}
\newcommand{\handoutdesignfont}{%
    \renewcommand{\encodingdefault}{T1}% Pertahankan encoding desain lama.
    \renewcommand{\rmdefault}{ptm}% Lokal pada kelompok judul/sampul.
    \fontencoding{T1}\fontfamily{ptm}\selectfont}

% Arial seperti Karakterisasi Dual.tex, khusus judul subbab dan label teorema.
% Pengaturan lokal menjaga paragraf isi tetap memakai Computer Modern.
\newcommand{\handoutheadingfont}{%
    \renewcommand{\encodingdefault}{T1}%
    \renewcommand{\rmdefault}{ua1}% Untuk normalfont pada nama teorema.
    \renewcommand{\sfdefault}{ua1}% Untuk label pita Subbab.
    \fontencoding{T1}\fontfamily{ua1}\selectfont}
\newcommand{\handouttitlefont}[1]{%
    \ifboolexpr{test {\ifstrequal{#1}{part}} or test {\ifstrequal{#1}{chapter}}}
        {\handoutdesignfont}{\handoutheadingfont\bfseries}}

% Pertahankan tata letak judul bawaan, termasuk judul tanpa nomor.
\makeatletter
\patchcmd{\ttl@select}{\begingroup}{\begingroup\handouttitlefont{#1}}{}%
    {\PackageError{handout}{Gagal mengatur font judul}{Periksa versi titlesec.}}
% Label/nama teorema, contoh, latihan, bukti, dan solusi memakai Arial.
\patchcmd{\@begintheorem}{\the\thm@headfont}%
    {\the\thm@headfont\handoutheadingfont}{}%
    {\PackageError{handout}{Gagal mengatur font label teorema}{Periksa versi amsthm.}}
\makeatother
% Judul paragraf dicetak tertunda oleh titlesec, jadi simpan font di formatnya.
% Kelas memakai opsi explicit: judul harus ditulis sebagai #1, tanpa varian *.
\titleformat{\paragraph}[runin]
    {\normalfont\handoutheadingfont\normalsize\bfseries}{\theparagraph}{1em}{#1}
\titleformat{\subparagraph}[runin]
    {\normalfont\handoutheadingfont\normalsize\bfseries}{\thesubparagraph}{1em}{#1}
\AddToHook{cmd/outermarginpagemark/before}{\begingroup\handoutdesignfont}
\AddToHook{cmd/outermarginpagemark/after}{\endgroup}

% Readability adjustments for the A4 handout: the class default is too
% compressed for sustained reading.
\geometry{left=2.2cm,right=2.2cm,top=1.5cm,bottom=2.2cm}
\AtBeginDocument{\fontsize{13}{18}\selectfont}
\linespread{1.06}
\setlength{\parskip}{0.35em}
% Pertahankan jarak vertikal alami saat kotak/rumus berpindah halaman.
\raggedbottom

% Klinik AA: exact HEX values from Logo Guideline.pdf, page 1.
\definecolor{AAIndigo}{HTML}{383277}
\definecolor{AASky}{HTML}{96D2F3}
\definecolor{AARed}{HTML}{E21E34}
\definecolor{AADarkRed}{HTML}{AB1F23}
% Shared palette for the cover, page decorations, theorem boxes, and links.
\colorlet{AAInk}{AAIndigo!65!black}
\colorlet{AABlue}{AAIndigo!55!AASky}
\colorlet{coverbgcolor}{white}
\colorlet{coverfgcolor}{AAInk}
\colorlet{coverbar}{AASky}
\colorlet{bottomcolor}{AAIndigo}
\colorlet{nuanbai}{AASky!6!white}
\colorlet{bg}{AASky!28!white}
\colorlet{fg}{AAIndigo}
\colorlet{outermarginbgcolor}{bg}
\colorlet{outermarginfgcolor}{fg}
% Gradasi latar judul subbab: abu-abu kebiruan menuju biru langit.
\colorlet{sectiontitlebgstart}{AAIndigo!12!gray!26!white}
\colorlet{sectiontitlebgend}{AASky!35!white}
\colorlet{captioncolor}{fg}
\colorlet{framegolden}{fg}
\colorlet{framegray}{AASky!12!white}
\hypersetup{linkcolor=AAIndigo,anchorcolor=AAIndigo,urlcolor=AABlue,citecolor=AAIndigo}
\coverstyle={
    cover-choose=en,
}
% Muat simbol AMS melalui preset yang dipakai pada versi stabil sebelumnya.
% newtxmath tidak dimuat dalam proyek ini; jangan memuat kedua preset sekaligus.
\mathstyle={
    math-font=plain,
}

\mynewtheorem{
    defi={\textbf{Definisi}}[section]{interior style={left color=ReD!8,right color=ReD!5!CyaN!50}, borderline west={1.5mm}{0mm}{ReD}},
    thm={\textbf{Teorema}}[section]{interior style={left color=CyaN!80!black!20,right color=CyaN!80!black!15!CyaN!50}, borderline west={1.5mm}{0mm}{CyaN!80!black}},
    lem={\textbf{Lema}}[section]{interior style={left color=BluE!8,right color=BluE!5!CyaN!50}, borderline west={1.5mm}{0mm}{BluE}},
    prop={\textbf{Proposisi}}[section]{interior style={left color=OrangE!8,right color=OrangE!5!CyaN!50}, borderline west={1.5mm}{0mm}{OrangE}},
    exam={\textbf{Contoh}}[chapter]{interior style={left color=DarkGreen!8,right color=DarkGreen!5!CyaN!50}, borderline west={1.5mm}{0mm}{DarkGreen}},
}
\newtheorem*{remark}{\textbf{Catatan}}

\mynewtcbtheorem{
    definisi={style={handout ribbon={Definisi}{AADarkRed}{defi}}},
    teorema={style={handout ribbon={Teorema}{AAIndigo}{thm}}},
}

% Names used by \autoref for the custom theorem counters.
\renewcommand{\defiautorefname}{Definisi}
\renewcommand{\thmautorefname}{Teorema}
\renewcommand{\lemautorefname}{Lema}
\renewcommand{\propautorefname}{Proposisi}
\renewcommand{\examautorefname}{Contoh}
% Hyperref perlu nomor subbab/bab dalam tujuan tautan penghitung yang direset.
\def\theHdefi{\theHsection.\arabic{defi}}
\def\theHthm{\theHsection.\arabic{thm}}
\def\theHlem{\theHsection.\arabic{lem}}
\def\theHprop{\theHsection.\arabic{prop}}
\def\theHexam{\theHchapter.\arabic{exam}}


% Materi matriks dan dekomposisinya; rujukan buku-buku Fuzhen Zhang.
\DeclareMathOperator{\rank}{rank}
\DeclareMathOperator{\nul}{nul}
\DeclareMathOperator{\tr}{tr}
\DeclareMathOperator{\diag}{diag}
\newcommand{\R}{\mathbb{R}}
\newcommand{\C}{\mathbb{C}}
\newcommand{\F}{\mathbb{F}}
\newcommand{\Z}{\mathbb{Z}}
\newcommand{\Mat}{\mathbb{M}}
\newcommand{\ide}[1]{\par\smallskip\noindent\textbf{Ide utama.} #1\par\smallskip}
\theoremstyle{definition}
\newtheorem{latihan}{Latihan}[chapter]
\theoremstyle{plain}
\providecommand{\latihanautorefname}{Latihan}
\def\theHlatihan{\theHchapter.\arabic{latihan}}
\renewcommand{\bibsection}{\chapter*{\bibname}\addcontentsline{toc}{chapter}{\bibname}}
% Ganti menjadi \pembahasanfalse untuk membagikan soal tanpa solusi.
\newif\ifpembahasan
\pembahasantrue
\setcounter{tocdepth}{1}
\setcounter{secnumdepth}{2}
\hypersetup{pdftitle={Handout ONMIPA-PT: Aljabar Linear -- Matriks},
    pdfauthor={Teosofi Hidayah Agung}}
\begin{document}
\thispagestyle{empty}
\title{Handout ONMIPA-PT Bidang Matematika}
\subtitle{Aljabar Linear}
\edition{Edisi 2026}
\bookseries{}
\author{Teosofi Hidayah Agung}
\pressname{Klinik AA}
\presslogo{bookstyle/LogoAA.png}
\hypersetup{pageanchor=false}
{\handoutdesignfont\makecover}

\makeatletter
\oddoutermargin{\sffamily \leftmark}
\evenoutermargin{\sffamily\@title}
\chapimage{}
\chapstarimage{bookstyle/LogoAA.png}
\partimage{bookstyle/part.png}
\partblankimage{bookstyle/LogoAA.png}
\makeatother


\frontmatter
\pagenumbering{Roman}
\hypersetup{pageanchor=true}

{\thispagestyle{empty}
    \chapter*{Kata Pengantar}
    \addcontentsline{toc}{chapter}{Kata Pengantar}
    Puji syukur ke hadirat Allah SWT atas segala rahmat dan karunia-Nya sehingga buku ``Handout ONMIPA-PT Bidang Matematika: Aljabar Linear'' ini dapat disusun dengan baik oleh penulis dan rekan sejawatnya (GPT Plus).

Buku ini disusun sebagai salah satu bahan pembelajaran dan persiapan dalam menghadapi Olimpiade Nasional Matematika dan Ilmu Pengetahuan Alam Perguruan Tinggi (ONMIPA-PT) Bidang Matematika, yang lomba tersebut pernah penulis ikuti selama 3 kali berturut-turut pada tahun 2024, 2025, dan 2026. Walaupun hanya mentok jadi finalis, setidaknya banyak pengalaman yang dapat penulis ambil dan rasakan hingga mendapatkan niat untuk menyusun buku ini.

Alasan utama penyusunan buku ini adalah untuk merangkum berbagai konsep, definisi, teorema, dan hasil-hasil penting dalam Aljabar Linear yang relevan dengan persiapan ONMIPA-PT. Selain itu, pada dasarnya penulis memang menyukai Aljabar Linear hingga ke tingkat seperti Analisis Fungsional (mata kuliah yang sangat seru sekali hehe). Sangat \textit{recommended} untuk diambil karena campuran antara Aljabar Linear dan Analisis Real yang mungkin saja akan membantu fundamental pembaca untuk menghadapi ONMIPA.

Sembari menunggu \textit{reset token}, penulis menyadari bahwa buku ini masih memiliki berbagai keterbatasan dan jauh dari kesempurnaan, baik dari segi kelengkapan materi, penyajian, maupun penulisan. Oleh karena itu, kritik dan saran yang membangun sangat diharapkan sebagai bahan evaluasi dan perbaikan pada penyusunan berikutnya.

Akhir kata, penulis berharap buku ini dapat memberikan manfaat bagi para pembaca serta menjadi salah satu sarana pendukung dalam proses pembelajaran dan persiapan menghadapi ONMIPA-PT. Semoga buku ini dapat membantu menumbuhkan pemahaman yang lebih mendalam sekaligus meningkatkan minat dan semangat dalam mempelajari matematika.


    \hfill
    \begin{tabular}{lr}
         & Teosofi Hidayah Agung \\
         & 2026
    \end{tabular}
    \clearpage}

\thispagestyle{empty}
{\normalfont\tableofcontents}\let\cleardoublepage\clearpage

\mainmatter
\pagenumbering{arabic}
\part{Matriks}{% !TeX root = ../../ONMIPA Alin.tex
Sebelum memasuki pembahasan pada bagian ini, pembaca
diharapkan telah memahami \textit{operasi baris elementer}
(OBE) dan \textit{operasi kolom elementer} (OKE).
Kedua operasi tersebut merupakan dasar penting dalam
menyederhanakan matriks dan menganalisis sifat-sifatnya.

Sebagai pengingat, terdapat tiga jenis operasi baris
elementer, yaitu
\[
    \begin{aligned}
        R_i & \leftrightarrow R_j
            & & \text{(pertukaran dua baris)}, \\
        R_i & \leftarrow \lambda R_i,\quad \lambda\neq 0
            & & \text{(perkalian dengan skalar)}, \\
        R_i & \leftarrow R_i+\lambda R_j,\quad i\neq j
            & & \text{(penjumlahan kelipatan baris)}.
    \end{aligned}
\]
Operasi kolom elementer didefinisikan secara analog.
Setiap OBE dapat direpresentasikan sebagai perkalian
matriks elementer dari kiri, sedangkan OKE sebagai
perkalian matriks elementer dari kanan, yakni
\[
    A\xrightarrow{\mathrm{OBE}} E_iA
    \qquad\text{dan}\qquad
    A\xrightarrow{\mathrm{OKE}} AF_j,
\]
dengan $E_i$ dan $F_j$ merupakan matriks elementer.
}
% !TeX root = ../../ONMIPA Alin.tex
% Rujukan: Zhang (2011), bagian 1.2, 2.1, 2.3; Zhang (2009),
% soal-soal tentang ruang kolom, kernel, rank, nulitas, dan matriks khusus.
\chapter{Rank dan Nulitas}
\label{ch:rank}

Rank mengukur dimensi ruang hasil suatu matriks, sedangkan nulitas
mengukur dimensi ruang masukan yang dipetakan ke nol. Keduanya
menghubungkan kebebasan kolom, banyaknya parameter bebas dalam solusi,
dan perilaku hasil kali matriks. Bab ini mengembangkan hubungan tersebut
melalui ruang kolom, kernel, faktorisasi, dan eliminasi blok, kemudian
menerapkannya pada persamaan serta pangkat matriks.

Semua matriks berentri dalam $\F=\R$ atau $\C$.
Notasi $\Mat_{m\times n}(\F)$ menyatakan himpunan matriks berukuran
$m\times n$; $\Mat_n(\F)$ digunakan untuk matriks persegi.
Operasi matriks, kebebasan linear, basis, dan dimensi dianggap sudah dikenal.
Determinan belum diperlukan dalam pembuktian-pembuktian bab ini.

\section{Ruang Kolom, Ruang Baris, dan Kernel}
\label{rank:dasar}

\begin{definisi}[][Rank dan nulitas][rank:definisi]
    Untuk $A\in\Mat_{m\times n}(\F)$, definisikan
    \[
        \begin{aligned}
            \mathcal C(A)&=\{Ax:x\in\F^n\}=\operatorname{Im}A,\\
            \mathcal R(A)&=\text{ruang yang direntang oleh baris-baris }A,\\
            \ker A&=\{x\in\F^n:Ax=0\}.
        \end{aligned}
    \]
    \emph{Rank} dan \emph{nulitas} $A$ berturut-turut adalah
    $\rank A=\dim\mathcal C(A)$ dan $\nul A=\dim\ker A$.
    Matriks mempunyai rank kolom penuh jika $\rank A=n$, dan rank baris
    penuh jika $\rank A=m$.
\end{definisi}

Jadi $0\le\rank A\le\min\{m,n\}$, dan $\rank A=0$ tepat ketika $A=0$.
Rank bukan banyaknya kolom tak nol: beberapa kolom tak nol dapat
menghasilkan arah yang sama.

\begin{teorema}[][Kesamaan rank baris dan kolom][rank:baris-kolom]
    Untuk setiap matriks $A$,
    \[
        \dim\mathcal R(A)=\dim\mathcal C(A),\qquad
        \rank A^T=\rank A.
    \]
\end{teorema}
\begin{proof}
    Misalkan $c=\dim\mathcal C(A)>0$. Pilih $c$ kolom bebas dari $A$
    dan tempatkan sebagai kolom-kolom matriks $U$. Setiap kolom $A$
    merupakan kombinasi kolom $U$, sehingga $A=UV$ untuk suatu matriks
    $V$ dengan $c$ baris. Setiap baris $A$ merupakan kombinasi baris $V$;
    akibatnya $\dim\mathcal R(A)\le c$.
    Terapkan argumen yang sama pada $A^T$ untuk memperoleh ketaksamaan
    sebaliknya. Untuk $A=0$, kedua dimensi sama dengan nol.
\end{proof}

\begin{teorema}[][Rank--nulitas][rank:nulitas]
    Jika $A$ mempunyai $n$ kolom, maka
    \[
        \rank A+\nul A=n.
    \]
\end{teorema}
\begin{proof}
    Pilih basis $u_1,\ldots,u_k$ untuk $\ker A$.
    Lengkapi dengan $v_1,\ldots,v_{n-k}$ hingga menjadi basis $\F^n$.
    Vektor $Av_1,\ldots,Av_{n-k}$ merentang $\mathcal C(A)$.
    Jika $\sum_j c_jAv_j=0$, maka $\sum_j c_jv_j\in\ker A$.
    Kebebasan basis yang diperluas memaksa semua $c_j=0$.
    Jadi $\rank A=n-k$ dan $\nul A=k$.
\end{proof}

Dalam rumus ini, $n$ adalah banyaknya \emph{kolom}, yaitu dimensi
domain pemetaan $x\mapsto Ax$. Kernel adalah himpunan vektor, sedangkan
nulitas adalah dimensinya. Secara khusus, $\nul A=0$ ekuivalen dengan
$\ker A=\{0\}$ dan kebebasan semua kolom $A$.
Jika $n>m$, selalu ada solusi tak nol untuk $Ax=0$, sebab
$\nul A=n-\rank A\ge n-m>0$.

\begin{prop}[Operasi yang mempertahankan rank]
    \label{rank:invarian}
    Jika $P\in\Mat_m(\F)$ dan $Q\in\Mat_n(\F)$ invertibel, maka
    \[
        \rank(PAQ)=\rank A.
    \]
    Khususnya, operasi baris dan kolom elementer mempertahankan rank.
\end{prop}
\begin{proof}
    Karena $Q$ memetakan $\F^n$ secara bijektif ke dirinya sendiri,
    $\mathcal C(AQ)=\mathcal C(A)$. Sementara itu,
    $\mathcal C(PA)=P\mathcal C(A)$ mempunyai dimensi yang sama dengan
    $\mathcal C(A)$ karena $P$ invertibel.
\end{proof}

Pada bentuk eselon baris, banyaknya pivot sama dengan rank. Jika yang
dicari adalah basis ruang kolom matriks \emph{asal}, ambil kolom asal
yang indeksnya bersesuaian dengan pivot. Operasi baris mempertahankan
hubungan ketergantungan antarkolom, tetapi dapat mengubah ruang kolomnya.

\begin{exam}[Parameter yang mengubah rank dan nulitas]
    Tentukan rank, nulitas, dan kernel dari
    \[
        A(t)=\begin{pmatrix}1&1&5\\1&3&3\\0&-1&t\end{pmatrix},\qquad t\in\R.
    \]
\end{exam}
\begin{solution}
    Lakukan $R_2\leftarrow R_2-R_1$, kemudian
    $R_3\leftarrow R_3+\tfrac12R_2$:
    \[
        A(t)\sim
        \begin{pmatrix}1&1&5\\0&2&-2\\0&0&t-1\end{pmatrix}.
    \]
    Jika $t\ne1$, terdapat tiga pivot: $\rank A(t)=3$ dan
    $\nul A(t)=0$. Jika $t=1$, ranknya $2$, nulitasnya $1$, dan
    \[
        \ker A(1)=\operatorname{span}\{(-6,1,1)^T\}.
    \]
    Dua kolom pertama matriks asal menjadi basis ruang kolom pada $t=1$.
    Jangan membagi baris dengan $t-1$ sebelum memisahkan kasus $t=1$.
\end{solution}

\subsection{Basis kernel dan variabel bebas}
\label{rank:basis-kernel}

\begin{prop}[Membaca kernel dari bentuk eselon tereduksi]
    Misalkan $A$ mempunyai $n$ kolom dan rank $r$.
    Setelah operasi baris dan pengurutan ulang kolom, tuliskan
    \[
        EA\Pi=\begin{pmatrix}I_r&H\\0&0\end{pmatrix},
    \]
    dengan $E$ invertibel dan $\Pi$ matriks permutasi.
    Maka
    \[
        \ker A=\left\{\Pi\begin{pmatrix}-Hz\\z\end{pmatrix}:
            z\in\F^{n-r}\right\}.
    \]
    Kolom-kolom $\Pi\binom{-H}{I_{n-r}}$ membentuk basis kernel.
\end{prop}
\begin{proof}
    Gunakan perubahan peubah $x=\Pi y$, lalu pisahkan $y$ menjadi
    variabel pivot $y_1$ dan variabel bebas $y_2$.
    Persamaan $Ax=0$ ekuivalen dengan $y_1+Hy_2=0$.
    Jadi $y_2=z$ bebas dan $y_1=-Hz$.
    Pemetaan $z\mapsto\Pi(-Hz,z)^T$ injektif karena $\Pi$ invertibel
    dan komponen terakhir sebelum permutasi adalah $z$ sendiri.
    Bayangan basis standar $\F^{n-r}$ karenanya menjadi basis kernel.
\end{proof}

Jika hanya operasi baris yang dilakukan, $\ker(EA)=\ker A$.
Jika kolom diubah, koordinat solusi harus dikembalikan:
$\ker(AQ)=Q^{-1}\ker A$ untuk $Q$ invertibel.
Jumlah variabel bebas tetap sama, tetapi vektor basis kernelnya
dapat berubah dalam koordinat semula.

\subsection{Nulitas kanan dan nulitas kiri}

Nulitas $A$ biasanya berarti dimensi kernel kanan $\ker A$.
Relasi linear di antara baris-baris $A$ dicatat oleh kernel kiri,
yaitu $\ker A^T$. Untuk $A$ berukuran $m\times n$ dan rank $r$,
\[
    \nul A=n-r,\qquad \nul A^T=m-r,\qquad
    \nul A-\nul A^T=n-m.
\]
Jadi kesamaan rank $A$ dan $A^T$ tidak menjamin kesamaan nulitas
jika matriksnya persegi panjang. Di atas $\C$, kernel kiri juga
sering dinyatakan memakai $A^*=\overline A^{\,T}$ dan pasangan $y^*A=0$;
dimensinya tetap $m-r$.

\begin{exam}[Tiga variabel bebas dan satu relasi baris]
    Tentukan rank, nulitas, dan basis kernel matriks
    \[
        A=\begin{pmatrix}
            1&2&0&-1&3\\0&0&1&2&-1\\2&4&1&0&5
        \end{pmatrix}.
    \]
    Tentukan juga syarat agar $Ax=b$ konsisten.
\end{exam}
\begin{solution}
    Baris ketiga adalah dua kali baris pertama ditambah baris kedua.
    Dua baris pertama bebas, sehingga $\rank A=2$, $\nul A=5-2=3$,
    dan $\nul A^T=3-2=1$.
    Untuk $Ax=0$, ambil $x_2=s$, $x_4=t$, dan $x_5=u$.
    Maka $x_1=-2s+t-3u$ dan $x_3=-2t+u$, sehingga basis kernelnya
    \[
        \begin{pmatrix}-2\\1\\0\\0\\0\end{pmatrix},\quad
        \begin{pmatrix}1\\0\\-2\\1\\0\end{pmatrix},\quad
        \begin{pmatrix}-3\\0\\1\\0\\1\end{pmatrix}.
    \]
    Relasi baris memberi $\ker A^T=\operatorname{span}\{(-2,-1,1)^T\}$.
    Karena dua persamaan pertama dapat diselesaikan untuk setiap
    $b_1,b_2$, syarat konsistensi tepat $b_3=2b_1+b_2$.
    Jika syarat ini dipenuhi, seluruh solusi adalah
    \[
        x=\begin{pmatrix}b_1-2s+t-3u\\s\\b_2-2t+u\\t\\u\end{pmatrix},
        \qquad s,t,u\in\F.
    \]
    Tiga parameter bebas sesuai dengan nulitas $A$.
\end{solution}

Contoh ini memperlihatkan hubungan antara kernel, kebebasan kolom,
dan teorema rank--nulitas.

\section{Bentuk Normal dan Faktorisasi Rank}
\label{rank:faktorisasi}

\begin{teorema}[][Bentuk normal terhadap operasi baris dan kolom][rank:normal]
    Jika $A\in\Mat_{m\times n}(\F)$ mempunyai rank $r$, terdapat matriks
    invertibel $P,Q$ sehingga
    \[
        PAQ=\begin{pmatrix}I_r&0\\0&0\end{pmatrix}=:D_r.
    \]
    Dua matriks dengan ukuran sama mempunyai rank sama jika dan hanya
    jika dapat dihubungkan oleh perkalian matriks invertibel di kiri dan kanan.
\end{teorema}
\begin{proof}
    Reduksi baris dan permutasi kolom mengubah $A$ menjadi
    $\begin{pmatrix}I_r&H\\0&0\end{pmatrix}$.
    Kurangkan kombinasi kolom-kolom pivot dari kolom-kolom sisanya untuk
    meniadakan $H$. Semua operasi tersebut invertibel dan menghasilkan
    $D_r$. Pernyataan terakhir mengikuti bentuk normal ini dan
    \autoref{rank:invarian}.
\end{proof}

Di sini $P$ dan $Q$ boleh dipilih secara independen. Hubungan ini disebut
\emph{ekuivalensi matriks}; syaratnya berbeda dari keserupaan
$B=S^{-1}AS$, yang mewajibkan kedua faktor saling invers.

\begin{teorema}[][Faktorisasi rank penuh][rank:full-factor]
    Jika $\rank A=r>0$, maka $A$ dapat ditulis sebagai
    \[
        A=UV,\qquad U\in\Mat_{m\times r}(\F),\quad
        V\in\Mat_{r\times n}(\F),
    \]
    dengan $\rank U=\rank V=r$.
\end{teorema}
\begin{proof}
    Ambil $U$ dari $r$ kolom bebas milik $A$, lalu isi kolom ke-$j$ dari
    $V$ dengan koefisien kolom ke-$j$ dari $A$ terhadap basis tersebut.
    Pada indeks kolom yang dipilih, $V$ memuat kolom-kolom $I_r$.
    Maka $V$ mempunyai rank baris penuh dan $A=UV$.
\end{proof}

\begin{prop}[Invers satu sisi dan pencoretan faktor]
    \label{rank:inverse-sisi}
    Matriks $U$ dengan rank kolom penuh mempunyai invers kiri $L$, yaitu
    $LU=I$. Matriks $V$ dengan rank baris penuh mempunyai invers kanan
    $R$, yaitu $VR=I$. Akibatnya, untuk ukuran yang sesuai,
    \[
        \rank(UAV)=\rank A.
    \]
\end{prop}
\begin{proof}
    Dari $PUQ=\binom{I}{0}$, diperoleh
    $L=Q\begin{pmatrix}I&0\end{pmatrix}P$ dengan $LU=I$.
    Transpose memberikan pernyataan untuk invers kanan.
    Pemetaan $x\mapsto Ux$ injektif, sehingga
    $\dim U\mathcal C(A)=\dim\mathcal C(A)$.
    Karena $V$ surjektif, $\mathcal C(AV)=\mathcal C(A)$.
    Gabungkan keduanya.
\end{proof}

\begin{prop}[Submatriks invertibel terbesar]
    Rank $A$ adalah ukuran terbesar submatriks persegi invertibel
    yang dapat dipilih dari $A$, dengan nilai $0$ jika $A=0$.
\end{prop}
\begin{proof}
    Menghapus kolom tidak menambah dimensi ruang kolom; menghapus baris
    juga tidak menaikkan rank, dengan argumen yang sama pada transpose.
    Jadi submatriks persegi invertibel berukuran $s$ memaksa $s\le\rank A$.
    Sebaliknya, pilih $r=\rank A>0$ kolom bebas dan bentuk matriks $U$.
    Rank baris $U$ juga $r$, sehingga dapat dipilih $r$ baris bebasnya.
    Irisan baris dan kolom yang terpilih adalah matriks persegi dengan
    rank penuh, sehingga invertibel oleh rank--nulitas.
\end{proof}

\subsection{Membaca rank dari pola entri}

Jika $U=[u_1\ \cdots\ u_r]$ dan baris-baris $V$ adalah
$v_1^T,\ldots,v_r^T$, maka
\[
    A=\sum_{j=1}^r u_jv_j^T.
\]
Setiap $u_jv_j^T$ mempunyai rank satu jika kedua vektornya tak nol.
Sebaliknya, jika $A$ merupakan jumlah $s$ matriks rank satu, ruang kolomnya
direntang oleh paling banyak $s$ vektor. Jadi rank $A$ adalah banyaknya
suku rank satu paling sedikit yang diperlukan untuk menyatakan $A$.

\begin{exam}[Matriks besar dengan rank dua]
    Untuk $n\ge2$, tentukan rank matriks $A=(i+j-1)_{i,j=1}^n$.
\end{exam}
\begin{solution}
    Definisikan $u=(1,2,\ldots,n)^T$ dan $e=(1,\ldots,1)^T$.
    Kolom ke-$j$ dari $A$ adalah $u+(j-1)e$, sehingga
    $\mathcal C(A)\subseteq\operatorname{span}\{u,e\}$.
    Dua kolom pertama bebas karena $u$ dan $e$ bebas.
    Maka $\rank A=2$, berapa pun besarnya $n$.
\end{solution}

\begin{exam}[Pola kuadratik tanpa menghitung determinan]
    Untuk $n\ge3$, matriks $B=(i^2+ij+j^2)_{i,j=1}^n$ mempunyai rank $3$.
\end{exam}
\begin{solution}
    Faktorkan
    \[
        B=\begin{pmatrix}
            1&1&1\\1&2&2^2\\\vdots&\vdots&\vdots\\1&n&n^2
        \end{pmatrix}
        \begin{pmatrix}
            1^2&2^2&\cdots&n^2\\1&2&\cdots&n\\1&1&\cdots&1
        \end{pmatrix}=UV.
    \]
    Kolom $U$ bebas: polinom berderajat paling tinggi dua yang nol di
    $1,2,3$ harus identik nol. Alasan sama membuktikan baris $V$ bebas.
    Dengan \autoref{rank:inverse-sisi}, $\rank(UV)=3$.
    Faktorisasi memberi batas atas; kebebasan kedua faktor memastikan
    bahwa batas tersebut benar-benar tercapai.
\end{solution}

\section{Rank dan Nulitas Hasil Kali}
\label{rank:produk}

\begin{teorema}[][Kehilangan dimensi pada hasil kali][rank:produk-tepat]
    Untuk $A\in\Mat_{m\times n}(\F)$ dan $B\in\Mat_{n\times p}(\F)$,
    \[
        \rank(AB)=\rank B-\dim\bigl(\mathcal C(B)\cap\ker A\bigr).
    \]
\end{teorema}
\begin{proof}
    Batasi pemetaan $x\mapsto Ax$ pada domain $\mathcal C(B)$.
    Citranya adalah $\mathcal C(AB)$, sedangkan kernelnya
    $\mathcal C(B)\cap\ker A$. Terapkan rank--nulitas pada pemetaan ini.
\end{proof}

\begin{prop}[Nulitas hasil kali]
    \label{rank:nulitas-produk}
    Untuk $A$ berukuran $m\times n$ dan $B$ berukuran $n\times p$,
    \[
        \nul(AB)=\nul B+\dim\bigl(\mathcal C(B)\cap\ker A\bigr).
    \]
    Akibatnya,
    \[
        \nul B\le\nul(AB)\le\min\{p,\nul A+\nul B\}.
    \]
    Jika $B$ mempunyai rank baris penuh, maka
    $\nul(AB)=\nul A+\nul B$.
\end{prop}
\begin{proof}
    Matriks $B$ dan $AB$ sama-sama mempunyai $p$ kolom.
    Kurangkan rumus pada \autoref{rank:produk-tepat} dari $p$.
    Untuk batasnya, gunakan
    $0\le\dim(\mathcal C(B)\cap\ker A)\le\nul A$.
    Jika $B$ mempunyai rank baris penuh, $\mathcal C(B)=\F^n$,
    sehingga irisannya dengan $\ker A$ adalah seluruh $\ker A$.
\end{proof}

\begin{exam}[Kernel tambahan pada hasil kali]
    Misalkan
    \[
        A=\begin{pmatrix}1&1&0\\0&0&0\end{pmatrix},\qquad
        B=\begin{pmatrix}1&0&1&0\\0&1&1&0\\0&0&0&0\end{pmatrix}.
    \]
    Tentukan nulitas $A$, $B$, dan $AB$ melalui ruang kolom dan kernel.
\end{exam}
\begin{solution}
    Rank $A$ adalah $1$ dan rank $B$ adalah $2$, sehingga
    $\nul A=3-1=2$ dan $\nul B=4-2=2$.
    Selanjutnya,
    \[
        \mathcal C(B)=\operatorname{span}\{e_1,e_2\},\qquad
        \ker A=\operatorname{span}\{e_1-e_2,e_3\}.
    \]
    Irisannya berdimensi satu. Jadi $\nul(AB)=2+1=3$ dan $\rank(AB)=1$.
    Selain arah dalam $\ker B$, ada satu arah bebas tambahan yang
    keluaran $B$-nya dimatikan oleh $A$.
    Arah $e_3\in\ker A$ tidak menyumbang kehilangan tambahan karena
    tidak termasuk ruang kolom $B$.
\end{solution}

\ide{Hasil kali $AB$ kehilangan rank ketika sebagian keluaran $B$
    masuk ke kernel $A$. Rank kedua faktor saja belum menentukan
    besar kehilangan tersebut.}

\begin{teorema}[][Ketaksamaan Sylvester][rank:sylvester]
    Dengan ukuran seperti di atas,
    \[
        \max\{0,\rank A+\rank B-n\}
        \le\rank(AB)\le\min\{\rank A,\rank B\}.
    \]
    Pada batas bawah, $n$ adalah dimensi \emph{ruang perantara}:
    banyak kolom $A$, sekaligus banyak baris $B$.
\end{teorema}
\begin{proof}
    Gunakan $\mathcal C(AB)\subseteq\mathcal C(A)$ dan
    \autoref{rank:produk-tepat}. Untuk batas bawah,
    \[
        \dim(\mathcal C(B)\cap\ker A)\le\dim\ker A=n-\rank A.
    \]
\end{proof}

\begin{prop}[Kapan batas hasil kali tercapai?]
    \label{rank:kesamaan-produk}
    Untuk ukuran yang sesuai,
    \begin{align*}
        \rank(AB)=\rank B
        &\iff \mathcal C(B)\cap\ker A=\{0\},\\
        \rank(AB)=\rank A
        &\iff \mathcal C(B)+\ker A=\F^n,\\
        \rank(AB)=\rank A+\rank B-n
        &\iff \ker A\subseteq\mathcal C(B).
    \end{align*}
\end{prop}
\begin{proof}
    Pernyataan pertama dan ketiga langsung mengikuti rumus rank hasil
    kali. Untuk yang kedua, $A\mathcal C(B)=\mathcal C(A)$ tepat ketika
    setiap $x\in\F^n$ dapat ditulis $x=y+z$ dengan
    $y\in\mathcal C(B)$ dan $z\in\ker A$.
\end{proof}

Khususnya, $AB=0$ ekuivalen dengan $\mathcal C(B)\subseteq\ker A$,
sehingga $\rank A+\rank B\le n$.
Ketaksamaan rank tersebut saja tidak cukup untuk menyimpulkan $AB=0$:
ambil $A=B=\diag(1,0)$.

\begin{exam}[Seluruh kemungkinan rank hasil kali]
    Misalkan $A,B\in\Mat_8(\R)$, $\rank A=5$, dan $\rank B=6$.
    Tentukan semua nilai yang mungkin bagi $\rank(AB)$.
\end{exam}
\begin{solution}
    Sylvester memberi $3\le\rank(AB)\le5$.
    Untuk membuktikan bahwa $3,4,5$ semuanya mungkin, tetapkan
    $A=\diag(I_5,0_3)$. Untuk setiap $k\in\{3,4,5\}$, pilih ruang
    kolom $B$ yang direntang oleh $k$ vektor koordinat dari
    $e_1,\ldots,e_5$ dan $6-k$ vektor dari $e_6,e_7,e_8$.
    Jadikan keenam vektor tersebut kolom pertama $B$, lalu tambahkan
    dua kolom nol. Maka $\rank B=6$ dan tepat $k$ arah bertahan pada $AB$.
    Secara umum, cara ini merealisasikan setiap bilangan bulat antara
    $\max\{0,a+b-n\}$ dan $\min\{a,b\}$ untuk dua matriks persegi
    berukuran $n$ dengan rank $a,b$.
\end{solution}

\begin{teorema}[][Ketaksamaan Frobenius][rank:frobenius]
    Jika $AB$ dan $BC$ terdefinisi, maka
    \[
        \rank(ABC)\ge\rank(AB)+\rank(BC)-\rank B.
    \]
\end{teorema}
\begin{proof}
    Karena $\mathcal C(BC)\subseteq\mathcal C(B)$,
    \begin{align*}
        \rank(ABC)
        &=\rank(BC)-\dim(\mathcal C(BC)\cap\ker A)\\
        &\ge\rank(BC)-\dim(\mathcal C(B)\cap\ker A)\\
        &=\rank(BC)+\rank(AB)-\rank B.
    \end{align*}
\end{proof}

Batas Frobenius mempertahankan informasi tentang faktor tengah $B$.
Jika hanya ukuran ruang perantara yang dipakai, informasi itu dapat
hilang ketika Sylvester diterapkan secara berulang.

\section{Rank Jumlah dan Perubahan Rank Satu}
\label{rank:jumlah}

\begin{prop}[Matriks gabungan]
    \label{rank:gabungan}
    Untuk $A,B$ dengan jumlah baris sama,
    \[
        \rank[A\ B]=\rank A+\rank B
        -\dim(\mathcal C(A)\cap\mathcal C(B)).
    \]
    Untuk $A,B$ dengan jumlah kolom sama, berlaku
    \[
        \ker\begin{pmatrix}A\\B\end{pmatrix}=\ker A\cap\ker B.
    \]
\end{prop}
\begin{proof}
    Ruang kolom $[A\ B]$ adalah $\mathcal C(A)+\mathcal C(B)$.
    Gunakan identitas
    $\dim(U+V)=\dim U+\dim V-\dim(U\cap V)$, yang diperoleh dengan
    memperluas basis $U\cap V$ menjadi basis masing-masing ruang.
    Pernyataan kernel mengikuti kesetaraan
    $\binom{A}{B}x=0\iff Ax=Bx=0$.
\end{proof}

\begin{teorema}[][Batas rank jumlah][rank:batas-jumlah]
    Jika $A,B\in\Mat_{m\times n}(\F)$, maka
    \[
        |\rank A-\rank B|\le\rank(A+B)
        \le\min\{m,n,\rank A+\rank B\}.
    \]
    Ketaksamaan yang sama berlaku untuk $A-B$.
\end{teorema}
\begin{proof}
    Inklusi $\mathcal C(A+B)\subseteq\mathcal C(A)+\mathcal C(B)$
    memberi batas atas. Terapkan batas ini pada
    $A=(A+B)-B$ dan $B=(A+B)-A$ untuk memperoleh batas bawah.
\end{proof}

\begin{teorema}[][Kapan rank bersifat aditif?][rank:aditif]
    Untuk dua matriks berukuran sama,
    \[
        \rank(A+B)=\rank A+\rank B
    \]
    jika dan hanya jika kedua syarat berikut dipenuhi:
    \[
        \mathcal C(A)\cap\mathcal C(B)=\{0\},\qquad
        \mathcal R(A)\cap\mathcal R(B)=\{0\}.
    \]
\end{teorema}
\begin{proof}
    Dari $A+B=[A\ B]\binom{I}{I}$ dan
    $A+B=[I\ I]\binom{A}{B}$, rank jumlah tidak melebihi rank kedua
    matriks gabungan. Jika ranknya $\rank A+\rank B$, rumus dimensi
    ruang kolom dan versi transposenya memaksa kedua irisan menjadi nol.

    Sebaliknya, irisan ruang baris yang nol memberi
    $\rank\binom{A}{B}=\rank A+\rank B$.
    Pemetaan $(u,v)\mapsto u+v$ injektif pada
    $\mathcal C(A)\times\mathcal C(B)$ karena irisan ruang kolom nol.
    Maka penerapan pemetaan tersebut pada citra $\binom{A}{B}$ tidak
    menurunkan dimensi, sehingga rank $A+B$ sama dengan jumlah ranknya.
\end{proof}

Satu syarat saja tidak cukup. Misalnya,
\[
    A=\begin{pmatrix}1&0\\0&0\end{pmatrix},\qquad
    B=\begin{pmatrix}0&0\\1&0\end{pmatrix}
\]
mempunyai ruang kolom yang beririsan nol, tetapi
$\rank(A+B)=1<2=\rank A+\rank B$ karena ruang barisnya berimpit.

\subsection{Mengganggu matriks dengan satu hasil kali luar}

Karena $\rank(uv^T)\le1$, selalu berlaku
\[
    \bigl|\rank(A+uv^T)-\rank A\bigr|\le1.
\]
Untuk matriks invertibel, kasus penurunan rank dapat ditentukan tepat.

\begin{prop}[Perubahan rank satu pada matriks invertibel]
    \label{rank:update}
    Jika $A\in\Mat_n(\F)$ invertibel dan $u,v\in\F^n$, maka
    \[
        \rank(A+uv^T)=
        \begin{cases}
            n,&1+v^TA^{-1}u\ne0,\\
            n-1,&1+v^TA^{-1}u=0.
        \end{cases}
    \]
    Pada kasus kedua, $\ker(A+uv^T)=\operatorname{span}\{A^{-1}u\}$.
\end{prop}
\begin{proof}
    Tuliskan $w=A^{-1}u$. Persamaan $(A+uv^T)x=0$ ekuivalen dengan
    $x=-w(v^Tx)$. Jadi semua vektor kernel berada pada
    $\operatorname{span}\{w\}$. Mengalikan dengan $v^T$ memberi
    $(1+v^Tw)v^Tx=0$.
    Jika koefisiennya tak nol, $x=0$.
    Jika koefisiennya nol, $w\ne0$ dan $(I+wv^T)w=0$;
    kernel berdimensi satu. Gunakan rank--nulitas.
\end{proof}

\begin{exam}[Diagonal konstan dan satu matriks rank satu]
    Misalkan $J_n$ adalah matriks yang semua entrinya $1$, $n\ge2$.
    Tentukan rank $tI_n+J_n$ untuk setiap $t\in\F$.
\end{exam}
\begin{solution}
    Tuliskan $J_n=ee^T$ dengan $e=(1,\ldots,1)^T$.
    Jika $t\ne0$, terapkan \autoref{rank:update} pada $A=tI_n$;
    skalar penentunya adalah $1+n/t$.
    Kasus $t=0$ dihitung langsung. Hasilnya
    \[
        \rank(tI_n+J_n)=
        \begin{cases}
            1,&t=0,\\
            n-1,&t=-n,\\
            n,&t\notin\{0,-n\}.
        \end{cases}
    \]
    Pada $t=-n$ kernelnya $\operatorname{span}\{e\}$; pada $t=0$
    kernelnya $\{x:e^Tx=0\}$.
    Ini juga menyelesaikan matriks dengan entri diagonal $a$ dan entri
    luar diagonal $b$: matriksnya adalah $(a-b)I_n+bJ_n$.
\end{solution}

\section{Eliminasi Blok dan Persamaan Matriks}
\label{rank:blok}

Operasi baris blok dilakukan dengan mengalikan blok dari kiri;
operasi kolom blok mengalikan dari kanan. Urutan perkalian harus dijaga
karena dua matriks umumnya tidak komutatif.

\begin{teorema}[][Rumus rank komplemen Schur][rank:schur]
    Misalkan
    $M=\begin{pmatrix}P&Q\\R&S\end{pmatrix}$, dengan $P$ persegi
    berukuran $k$ dan invertibel. Maka
    \[
        \rank M=k+\rank(S-RP^{-1}Q).
    \]
\end{teorema}
\begin{proof}
    Identitas
    \[
        \begin{pmatrix}I&0\\-RP^{-1}&I\end{pmatrix}
        \begin{pmatrix}P&Q\\R&S\end{pmatrix}
        \begin{pmatrix}I&-P^{-1}Q\\0&I\end{pmatrix}
        =\begin{pmatrix}P&0\\0&S-RP^{-1}Q\end{pmatrix}
    \]
    menggunakan dua faktor invertibel di kiri dan kanan. Rank matriks
    diagonal blok adalah jumlah rank blok-bloknya karena ruang kolomnya
    menempati kelompok koordinat yang terpisah.
\end{proof}

\begin{prop}[Segitiga blok]
    \label{rank:segitiga-blok}
    Untuk blok-blok berukuran sesuai,
    \[
        \rank\begin{pmatrix}A&X\\0&B\end{pmatrix}
        \ge\rank A+\rank B.
    \]
    Kesamaan selalu berlaku jika $X=0$, tetapi belum tentu jika $X\ne0$.
\end{prop}
\begin{proof}
    Pilih $\rank A$ kolom bebas dari blok kolom pertama dan
    $\rank B$ kolom dari blok kolom kedua yang bagian bawahnya bebas.
    Gabungannya bebas: komponen bawah terlebih dahulu memaksa koefisien
    kelompok kedua nol, lalu komponen atas memaksa sisanya nol.
    Untuk melihat kemungkinan ketaksamaan ketat, ambil
    $A=B=(0)$ dan $X=(1)$.
\end{proof}

\begin{exam}[Dua hasil kali dengan ukuran berbeda]
    Jika $A$ berukuran $m\times n$ dan $B$ berukuran $n\times m$, buktikan
    \[
        m-\rank(I_m-AB)=n-\rank(I_n-BA).
    \]
\end{exam}
\begin{solution}
    Hitung rank matriks yang sama dengan dua pilihan pivot blok:
    \[
        \rank\begin{pmatrix}I_m&A\\B&I_n\end{pmatrix}
        =m+\rank(I_n-BA)=n+\rank(I_m-AB).
    \]
    Jadi kedua matriks mempunyai nulitas sama, walaupun ukurannya
    berbeda. Secara langsung, $x\mapsto Bx$ memberi isomorfisme dari
    $\ker(I_m-AB)$ ke $\ker(I_n-BA)$ dengan invers $y\mapsto Ay$.
    Pernyataan ini tidak berarti $\rank(AB)=\rank(BA)$.
\end{solution}

\subsection{Konsistensi dan banyaknya parameter bebas}

\begin{prop}[Persamaan dengan satu faktor tetap]
    \label{rank:ax-b}
    Untuk $A\in\Mat_{m\times n}(\F)$ dan $C\in\Mat_{m\times q}(\F)$,
    persamaan $AX=C$ mempunyai solusi jika dan hanya jika
    \[
        \mathcal C(C)\subseteq\mathcal C(A)
        \iff \rank[A\ C]=\rank A.
    \]
    Jika konsisten, himpunan solusinya berbentuk $X_0+\mathcal N$, dengan
    $\mathcal N=\{Z:AZ=0\}$ berdimensi $q\,\nul A=q(n-\rank A)$.
\end{prop}
\begin{proof}
    Persamaan ini setara dengan $q$ persamaan kolom $Ax_j=c_j$.
    Setiap kolom $C$ harus berada pada $\mathcal C(A)$.
    Jika $X_0$ adalah satu solusi, seluruh solusi adalah $X_0+Z$,
    dengan setiap kolom $Z$ dipilih secara bebas dari $\ker A$.
    Tiap kolom menyumbang $n-\rank A$ parameter bebas.
\end{proof}

Untuk $q=1$, diperoleh kriteria konsistensi sistem linear $Ax=c$.
Jika konsisten, solusi tunggal tepat ketika $\nul A=0$, atau ekuivalen
$\rank A=n$.
Untuk $C\ne0$, himpunan solusi umumnya merupakan ruang afin
(translasi suatu subruang), bukan subruang yang melalui nol.

\begin{teorema}[][Persamaan dengan dua faktor tetap][rank:axb]
    Ambil $A\in\Mat_{m\times n}(\F)$,
    $B\in\Mat_{p\times q}(\F)$, dan $C\in\Mat_{m\times q}(\F)$.
    Persamaan $AXB=C$, dengan peubah $X\in\Mat_{n\times p}(\F)$,
    konsisten jika dan hanya jika
    \[
        \mathcal C(C)\subseteq\mathcal C(A),\qquad
        \mathcal R(C)\subseteq\mathcal R(B).
    \]
    Jika konsisten, banyak parameter bebasnya adalah
    \[
        np-(\rank A)(\rank B).
    \]
\end{teorema}
\begin{proof}
    Syarat inklusi perlu karena perkalian di kanan menghasilkan
    kombinasi kolom, sedangkan perkalian di kiri menghasilkan kombinasi baris.
    Untuk membuktikan kecukupan dan menghitung dimensinya, pilih
    \[
        PAQ=D_a,\qquad RBS=D_b,\qquad a=\rank A,\quad b=\rank B,
    \]
    dalam bentuk normal \autoref{rank:normal}.
    Dengan $Y=Q^{-1}XR^{-1}$ dan $C'=PCS$, persamaannya menjadi
    \[
        D_aYD_b=C'.
    \]
    Kedua syarat inklusi tepat menyatakan bahwa entri $C'$ di luar
    $a$ baris pertama dan $b$ kolom pertama semuanya nol.
    Perkalian $D_aYD_b$ mengambil blok $a\times b$ di pojok kiri atas
    $Y$ dan mengisi bagian lainnya dengan nol.
    Jadi $ab$ entri $Y$ ditentukan oleh $C'$, sedangkan $np-ab$ entri
    sisanya bebas. Perubahan peubah antara $X$ dan $Y$ bijektif.
\end{proof}

Hasil ini merupakan penerapan bentuk normal rank dan kriteria
faktorisasi ruang kolom dan baris.
Kondisi di atas juga dapat diperiksa melalui
$\rank[A\ C]=\rank A$ dan
$\rank\binom{B}{C}=\rank B$.

Pandang $\mathcal L(X)=AXB$ sebagai pemetaan linear dari
$\Mat_{n\times p}(\F)$ ke $\Mat_{m\times q}(\F)$.
Domainnya berdimensi $np$, dan pembuktian di atas menunjukkan
\[
    \rank\mathcal L=(\rank A)(\rank B),\qquad
    \nul\mathcal L=np-(\rank A)(\rank B).
\]
Di sini rank dan nulitas pemetaan berarti dimensi citra dan kernelnya.
Nulitas tersebut adalah dimensi ruang matriks $Z$ yang memenuhi $AZB=0$,
sekaligus banyaknya parameter bebas pada setiap persamaan $AXB=C$
yang konsisten.

\begin{exam}[Sembilan peubah, empat kendala bebas]
    Misalkan
    \[
        A=\begin{pmatrix}1&2&0\\0&0&1\end{pmatrix},\qquad
        B=\begin{pmatrix}1&0\\0&1\\1&1\end{pmatrix},\qquad
        C\in\Mat_2(\R).
    \]
    Tentukan apakah $AXB=C$ selalu dapat diselesaikan dan berapa banyak
    parameter bebas dalam solusinya.
\end{exam}
\begin{solution}
    Matriks $A$ mempunyai rank baris penuh $2$, dan $B$ mempunyai
    rank kolom penuh $2$. Maka kedua syarat inklusi selalu dipenuhi,
    dan banyak parameter bebasnya $3\cdot3-2\cdot2=5$.
    Secara konstruktif, gunakan
    \[
        R=\begin{pmatrix}1&0\\0&0\\0&1\end{pmatrix},\qquad
        L=\begin{pmatrix}1&0&0\\0&1&0\end{pmatrix}.
    \]
    Karena $AR=I_2$ dan $LB=I_2$, matriks $X_0=RCL$ adalah satu solusi.
    Semua solusi diperoleh dengan menambahkan $Z$ yang memenuhi $AZB=0$.
\end{solution}

\section{Transpose Konjugat dan Pangkat Matriks}
\label{rank:pangkat}

\subsection{Menggunakan kepositifan untuk menemukan kernel}

Untuk matriks kompleks, tuliskan $A^*=\overline A^{\,T}$.
Jika $A$ real, $A^*=A^T$.

\begin{teorema}[][Rank matriks Gram][rank:gram]
    Untuk setiap matriks kompleks $A$,
    \[
        \ker(A^*A)=\ker A,\qquad
        \rank(A^*A)=\rank A=\rank(AA^*).
    \]
\end{teorema}
\begin{proof}
    Jelas $Ax=0$ mengakibatkan $A^*Ax=0$.
    Sebaliknya, jika $A^*Ax=0$, maka
    \[
        0=x^*A^*Ax=(Ax)^*(Ax)=\sum_i |(Ax)_i|^2,
    \]
    sehingga $Ax=0$.
    Rank--nulitas memberi kesamaan rank pertama.
    Konjugasi mempertahankan kebebasan linear, sehingga
    $\rank A^*=\rank A$. Terapkan hasil pertama pada $A^*$ untuk
    mendapatkan $\rank(AA^*)=\rank A$.
\end{proof}

\begin{remark}
    Di atas $\C$, jangan mengganti $A^*$ dengan $A^T$ secara sembarang.
    Untuk $A=(1,i)^T$, diperoleh $A^TA=0$, padahal $\rank A=1$.
    Sebaliknya, $A^*A=(2)$ mempunyai rank satu sebagaimana mestinya.
\end{remark}

Jika $A$ berukuran $m\times n$ dan mempunyai rank $r$, teorema Gram memberi
\[
    \nul(A^*A)=n-r=\nul A,\qquad \nul(AA^*)=m-r.
\]
Kedua matriks Gram mempunyai rank sama, sedangkan nulitasnya mengikuti
ukuran masing-masing matriks.

\subsection{Barisan rank, nulitas, dan stabilisasi kernel}

\begin{teorema}[][Penurunan rank pangkat][rank:deret-pangkat]
    Misalkan $A\in\Mat_n(\F)$ dan $r_k=\rank(A^k)$, dengan $r_0=n$.
    Maka
    \[
        n=r_0\ge r_1\ge r_2\ge\cdots\ge0,
        \qquad r_{k-1}-r_k\ge r_k-r_{k+1}\quad(k\ge1).
    \]
    Jika $r_s=r_{s+1}$ untuk suatu $s\ge0$, semua rank berikutnya
    sama dengan $r_s$.
\end{teorema}
\begin{proof}
    Inklusi
    \[
        \ker A^k\subseteq\ker A^{k+1},\qquad
        \mathcal C(A^{k+1})\subseteq\mathcal C(A^k)
    \]
    memberi penurunan rank. Terapkan Frobenius pada
    $A\cdot A^{k-1}\cdot A$:
    \[
        r_{k+1}\ge2r_k-r_{k-1}.
    \]
    Ini tepat ketaksamaan untuk selisih rank.
    Selisih-selisih tersebut nonnegatif dan tidak meningkat, sehingga
    sesudah satu selisih nol, semua selisih berikutnya juga nol.
\end{proof}

Jadi bukan setiap barisan bilangan bulat menurun dapat menjadi barisan
rank pangkat. Misalnya, untuk matriks berukuran $8$ dengan rank $6$,
harus berlaku $\rank A^2\ge4$; nilai $3$ mustahil.
Stabilisasi pasti terjadi paling lambat pada $A^n$, sebab rank tidak
dapat turun ketat lebih dari $n$ kali.

Dalam bahasa nulitas, tuliskan $d_k=\nul(A^k)=n-r_k$, dengan $d_0=0$.
Maka
\[
    0=d_0\le d_1\le d_2\le\cdots\le n,\qquad
    d_{k+1}-d_k\le d_k-d_{k-1}.
\]
Kernel bertambah besar, tetapi pertambahan dimensinya tidak meningkat.
Karena $\ker A^k\subseteq\ker A^{k+1}$, kesamaan dua nulitas berurutan
berarti kedua kernel sudah sama; setelah itu seluruh kernel berikutnya
juga sama. Selain itu, nulitas hasil kali memberi
\[
    \nul A\le\nul(A^k)\le\min\{n,k\,\nul A\}\qquad(k\ge1).
\]
Batas kanan diperoleh dengan menerapkan
\autoref{rank:nulitas-produk} secara berulang.

\begin{prop}[Rank tidak turun setelah dikuadratkan]
    \label{rank:stabil-awal}
    Untuk $A\in\Mat_n(\F)$, pernyataan berikut ekuivalen:
    \begin{enumerate}[label=(\roman*),nosep]
        \item $\rank A^2=\rank A$;
        \item $\mathcal C(A)\cap\ker A=\{0\}$;
        \item $\F^n=\mathcal C(A)\oplus\ker A$;
        \item dalam suatu basis, matriks $A$ berbentuk
            $\begin{pmatrix}D&0\\0&0\end{pmatrix}$ dengan $D$ invertibel
            berukuran $\rank A$.
    \end{enumerate}
    Untuk $A=0$, blok $D$ kosong dan pernyataan terakhir diartikan
    sebagai matriks nol.
\end{prop}
\begin{proof}
    Ekuivalensi (i)--(ii) berasal dari rumus rank hasil kali.
    Karena dimensi ruang kolom dan kernel berjumlah $n$, (ii) ekuivalen
    dengan (iii). Pada (iii), kedua subruang invarian terhadap $A$.
    Pembatasan $A$ pada $\mathcal C(A)$ injektif, sehingga invertibel
    pada ruang berdimensi hingga itu; pada $\ker A$, pemetaannya nol.
    Gabungan basis kedua subruang memberikan (iv).
    Sebaliknya, bentuk (iv) langsung memberikan (i).
\end{proof}

\begin{prop}[Batas rank dan nulitas matriks nilpoten]
    \label{rank:nilpoten}
    Jika $A^k=0$ untuk suatu bilangan bulat $k\ge1$, maka
    \[
        \rank A\le n-\left\lceil\frac nk\right\rceil,
        \qquad \nul A\ge\left\lceil\frac nk\right\rceil.
    \]
    Khususnya, $A^2=0$ mengakibatkan
    $\rank A\le\lfloor n/2\rfloor$ dan $\nul A\ge\lceil n/2\rceil$.
\end{prop}
\begin{proof}
    Sylvester berulang memberi
    $\rank A^k\ge k\rank A-(k-1)n$.
    Karena ruas kiri nol, $\rank A\le n-n/k$.
    Rank bilangan bulat, sehingga batasnya
    $\lfloor n-n/k\rfloor=n-\lceil n/k\rceil$.
    Untuk $k=2$, bukti langsungnya adalah
    $\mathcal C(A)\subseteq\ker A$.
\end{proof}

\begin{exam}[Batas yang dapat dicapai]
    Berapakah rank terbesar dan nulitas terkecil matriks $10\times10$
    yang memenuhi $A^3=0$?
\end{exam}
\begin{solution}
    Batas di atas memberi $\rank A\le10-\lceil10/3\rceil=6$.
    Definisikan matriks pergeseran
    \[
        N_3=\begin{pmatrix}0&1&0\\0&0&1\\0&0&0\end{pmatrix}.
    \]
    Matriks $A=\diag(N_3,N_3,N_3,0)$ berukuran $10$, memenuhi $A^3=0$,
    dan mempunyai rank $2+2+2=6$. Jadi rank terbesar adalah $6$ dan
    nulitas terkecil adalah $10-6=4$.
    Untuk ukuran umum, rantai pergeseran dengan panjang paling tinggi
    $k$ juga memperlihatkan ketajaman batas pada \autoref{rank:nilpoten}.
\end{solution}

\section{Identitas Matriks, Rank, dan Nulitas}
\label{rank:identitas}

\begin{teorema}[][Idempoten dan pasangan komplemennya][rank:idempoten]
    Jika $E^2=E$, maka
    \[
        \mathcal C(E)=\ker(I-E),\qquad
        \ker E=\mathcal C(I-E),\qquad
        \rank E+\rank(I-E)=n.
    \]
    Sebaliknya, jika $A+B=I_n$ dan $\rank A+\rank B=n$, maka
    $A^2=A$, $B^2=B$, dan $AB=BA=0$.
\end{teorema}
\begin{proof}
    Jika $y=Ex$, maka $Ey=E^2x=y$. Sebaliknya, $Ey=y$ berarti
    $y\in\mathcal C(E)$. Ini membuktikan kesamaan pertama.
    Kesamaan kedua mengikuti $E(I-E)=0$ dan
    $x\in\ker E\Rightarrow x=(I-E)x$.
    Rank--nulitas memberi jumlah rank yang dinyatakan.

    Untuk arah sebaliknya, $A+B=I$ memberi
    $\mathcal C(A)+\mathcal C(B)=\F^n$. Syarat jumlah dimensi
    memaksa irisannya nol. Karena $B=I-A$, berlaku $AB=BA$;
    setiap kolom $AB$ berada pada kedua ruang kolom, sehingga $AB=BA=0$.
    Maka $A=A(A+B)=A^2$ dan $B=(A+B)B=B^2$.
\end{proof}

\begin{prop}[Identitas kuadratik dengan dua akar berbeda]
    \label{rank:kuadratik}
    Jika $(A-aI)(A-bI)=0$, dengan $a\ne b$, maka
    \[
        \rank(A-aI)+\rank(A-bI)=n.
    \]
    Khususnya, jika $A^2=I$, maka
    $\rank(A-I)+\rank(A+I)=n$.
\end{prop}
\begin{proof}
    Definisikan
    \[
        E=\frac{A-bI}{a-b},\qquad F=\frac{aI-A}{a-b}.
    \]
    Kedua matriks berjumlah $I$ dan memenuhi $EF=FE=0$.
    Akibatnya $E^2=E$ dan $F=I-E$. Terapkan
    \autoref{rank:idempoten}, lalu gunakan bahwa perkalian dengan
    skalar tak nol mempertahankan rank.
\end{proof}

Sebagai contoh, jika $A\in\Mat_3(\F)$ memenuhi $A^2=I$ dan
$A\ne\pm I$, kedua rank $A-I$ dan $A+I$ positif serta berjumlah $3$.
Jadi ranknya $1$ dan $2$, dalam suatu urutan.

\begin{prop}[Identitas rank untuk $A-ABA$]
    \label{rank:aba}
    Jika $A\in\Mat_{m\times n}(\F)$ dan
    $B\in\Mat_{n\times m}(\F)$, maka
    \[
        \rank(A-ABA)=\rank A+\rank(I_n-BA)-n.
    \]
\end{prop}
\begin{proof}
    Untuk setiap $x\in\ker A$, berlaku $(I_n-BA)x=x$.
    Maka $\ker A\subseteq\mathcal C(I_n-BA)$, sehingga irisan kedua
    ruang itu tepat $\ker A$. Terapkan rumus rank hasil kali pada
    $A(I_n-BA)$:
    \[
        \rank(A-ABA)=\rank(I_n-BA)-\dim\ker A.
    \]
\end{proof}

Untuk idempoten $E$, hubungan ruang kolom dan kernel di atas juga memberi
$\nul E=\rank(I-E)$ dan $\nul E+\nul(I-E)=n$.
Untuk $(A-aI)(A-bI)=0$ dengan $a\ne b$, rank--nulitas memberikan
\[
    \nul(A-aI)+\nul(A-bI)=n.
\]
Jadi solusi kedua sistem homogen itu bersama-sama menyediakan $n$
arah bebas: kernelnya beririsan nol dan berjumlah langsung menjadi $\F^n$.
Irisannya nol karena $Ax=ax$ dan $Ax=bx$ memaksa $(a-b)x=0$.

\subsection{Memilih pendekatan pada soal rank dan nulitas}

\begin{enumerate}[label=\arabic*.,leftmargin=*,itemsep=3pt]
    \item Untuk pola entri, cari faktorisasi $UV$ atau tulis kolom
        sebagai kombinasi beberapa vektor tetap. Buktikan juga batas
        bawahnya; faktorisasi saja baru memberi batas atas.
    \item Untuk nulitas, tentukan banyak kolom lebih dahulu. Jika kernel
        mudah dideskripsikan, hitung dimensinya langsung; jika rank lebih
        mudah ditemukan, gunakan $\nul A=n-\rank A$.
    \item Untuk hasil kali, selidiki $\mathcal C(B)\cap\ker A$;
        gunakan Sylvester jika hanya rank kedua faktor yang diketahui.
    \item Untuk jumlah, periksa ruang kolom dan ruang baris sekaligus.
        Rank umumnya tidak aditif.
    \item Untuk blok, cari blok invertibel yang bisa dijadikan pivot.
        Untuk parameter, pisahkan dahulu nilai yang membuat pivot gagal.
    \item Untuk identitas seperti $A^2=0$ atau $A^2=A$, ubah identitas
        menjadi inklusi atau dekomposisi ruang kolom dan kernel.
    \item Untuk rank pangkat, periksa apakah penurunan rank mengecil
        dan apakah kernel sudah stabil.
\end{enumerate}

Hubungan dengan determinan akan dibahas pada \autoref{ch:determinan}.
Rank penuh matriks persegi berarti matriks tersebut invertibel;
nanti kondisi yang sama dapat dikenali melalui determinan tak nol.

\section{Latihan Pilihan}
\label{rank:latihan}

Soal berikut menggabungkan perhitungan struktural dan pembuktian.

\begin{latihan}\label{rank:lat-siklus}
    Misalkan $e_1,\ldots,e_n$ basis standar $\F^n$, $n\ge2$.
    Matriks $C_t$ mempunyai kolom
    \[
        e_1-e_2,\ e_2-e_3,\ldots,e_{n-1}-e_n,\ e_n-te_1.
    \]
    Tentukan rank dan kernel $C_t$ untuk semua $t\in\F$.
\end{latihan}

\begin{latihan}\label{rank:lat-cancel}
    Misalkan $\rank(AB)=\rank B$. Buktikan bahwa, untuk setiap
    $X,Y$ berukuran sesuai,
    \[
        ABX=ABY\iff BX=BY.
    \]
    Buktikan pula versi kanan: jika $\rank(AB)=\rank A$, maka
    $XAB=YAB\iff XA=YA$.
\end{latihan}

\begin{latihan}\label{rank:lat-produk-identitas}
    Untuk $A_1,\ldots,A_k\in\Mat_n(\F)$, buktikan
    \[
        \rank(A_1A_2\cdots A_k-I)
        \le\sum_{j=1}^k\rank(A_j-I).
    \]
    Untuk kasus $A_1=\cdots=A_k=A$, temukan batas atas yang lebih
    kuat dengan memfaktorkan $A^k-I$.
\end{latihan}

\begin{latihan}\label{rank:lat-konstruksi}
    Diberikan $A\in\Mat_{m\times n}(\F)$ dengan rank $r$.
    Untuk setiap bilangan bulat $k$ dengan $r\le k\le n$, buktikan
    bahwa terdapat $B\in\Mat_n(\F)$ sedemikian sehingga
    \[
        AB=0,\qquad \rank A+\rank B=k.
    \]
\end{latihan}

\begin{latihan}\label{rank:lat-kernel-sama}
    Misalkan $A,B\in\Mat_n(\F)$. Buktikan bahwa $\ker A=\ker B$
    jika dan hanya jika terdapat matriks invertibel $P$ sehingga $A=PB$.
    Apakah kesamaan $\rank A=\rank B$ saja sudah cukup untuk
    menjamin keberadaan $P$ tersebut?
\end{latihan}

\begin{latihan}\label{rank:lat-gabungan}
    Benar atau salah: untuk setiap $A,B\in\Mat_n(\F)$,
    \[
        \rank\begin{pmatrix}A\\B\end{pmatrix}=\rank[A\ B]?
    \]
    Jika salah, berikan contoh eksplisit terkecil yang memungkinkan.
\end{latihan}

\begin{latihan}\label{rank:lat-tiga-ruang}
    Untuk matriks $A,B,C$ dengan jumlah baris sama, buktikan
    \[
        \rank[A\ B\ C]\le\rank[A\ B]+\rank[B\ C]-\rank B.
    \]
    Tentukan syarat tepat terjadinya kesamaan dalam bahasa ruang kolom.
\end{latihan}

\begin{latihan}\label{rank:lat-proyeksi}
    Untuk $0\ne x\in\C^n$, definisikan
    \[
        P=I_n-\frac{xx^*}{x^*x}.
    \]
    Tentukan rank, ruang kolom, dan kernel $P$.
\end{latihan}

\begin{latihan}\label{rank:lat-stabil}
    Buktikan bahwa $\rank A^2=\rank A$ jika dan hanya jika terdapat
    matriks invertibel $B$ dengan $A^2=BA$.
\end{latihan}

\begin{latihan}\label{rank:lat-dua-idempoten}
    Misalkan $E^2=E$ dan $F^2=F$. Buktikan bahwa $E+F$ idempoten
    jika dan hanya jika $EF=FE=0$. Jika kondisi ini dipenuhi,
    buktikan pula $\rank(E+F)=\rank E+\rank F$.
\end{latihan}

\begin{latihan}\label{rank:lat-blok-identitas}
    Untuk $A,B\in\Mat_n(\F)$, buktikan
    \[
        \rank\begin{pmatrix}A&I_n\\0&B\end{pmatrix}=n+\rank(BA).
    \]
    Gunakan hasil tersebut untuk membuktikan ketaksamaan Sylvester
    dengan bantuan \autoref{rank:segitiga-blok}.
\end{latihan}

\begin{latihan}\label{rank:lat-nilpoten}
    Diberikan $A\in\Mat_{11}(\F)$ dengan $A^3=0$ dan $\rank A=7$.
    Tentukan $\rank A^2$, $\nul A$, dan $\nul A^2$, lalu konstruksikan
    satu matriks yang memenuhi seluruh syarat tersebut.
\end{latihan}

\begin{latihan}\label{rank:lat-irisan-kernel}
    Matriks $A$ dan $B$ sama-sama mempunyai $n$ kolom, dengan
    $\nul A=a$ dan $\nul B=b$.
    Buktikan bahwa
    \[
        \max\{0,a+b-n\}\le\dim(\ker A\cap\ker B)\le\min\{a,b\}.
    \]
    Nyatakan dimensi irisan tersebut sebagai nulitas sebuah matriks blok,
    lalu tentukan syarat agar $Ax=0$ dan $Bx=0$ hanya mempunyai solusi nol.
\end{latihan}

\begin{latihan}\label{rank:lat-nulitas-produk}
    Matriks $A$ berukuran $3\times5$ dengan $\nul A=3$, sedangkan
    $B$ berukuran $5\times4$ dengan $\nul B=1$.
    Tentukan semua nilai yang mungkin bagi $\nul(AB)$ dan berikan
    konstruksi untuk setiap nilai tersebut.
\end{latihan}

\ifpembahasan
\section{Pembahasan Latihan}
\label{rank:pembahasan}

\subsection*{Latihan \ref{rank:lat-siklus}: relasi kolom sepanjang siklus}
Persamaan $C_tc=0$ memberikan
\[
    c_1-tc_n=0,\qquad -c_{j-1}+c_j=0\quad(2\le j\le n).
\]
Jadi semua $c_j$ sama, dan $(1-t)c_1=0$.
Untuk $t\ne1$, kernelnya nol dan ranknya $n$.
Untuk $t=1$, kernelnya $\operatorname{span}\{(1,\ldots,1)^T\}$
dan ranknya $n-1$. Hubungan kolom yang menyebabkan penurunan rank
adalah jumlah seluruh kolom sama dengan nol.

\subsection*{Latihan \ref{rank:lat-cancel}: pencoretan pada ruang kolom}
Dari \autoref{rank:kesamaan-produk},
$\mathcal C(B)\cap\ker A=\{0\}$.
Jika $AB(X-Y)=0$, setiap kolom $B(X-Y)$ terletak pada irisan tersebut,
sehingga $B(X-Y)=0$. Arah sebaliknya langsung.
Untuk versi kanan, transpose persamaan dan gunakan
$\rank(B^TA^T)=\rank A^T$. Pencoretan ini sah pada ruang yang relevan
walaupun $A$ atau $B$ tidak invertibel pada seluruh ruang.

\subsection*{Latihan \ref{rank:lat-produk-identitas}: ekspansi teleskopik}
Gunakan identitas
\[
    A_1\cdots A_k-I
    =\sum_{j=1}^k A_1\cdots A_{j-1}(A_j-I),
\]
dengan hasil kali kosong untuk $j=1$ diartikan $I$.
Subaditivitas rank dan batas rank hasil kali memberi
\[
    \rank(A_1\cdots A_k-I)
    \le\sum_{j=1}^k\rank\bigl(A_1\cdots A_{j-1}(A_j-I)\bigr)
    \le\sum_{j=1}^k\rank(A_j-I).
\]
Jika semua $A_j=A$, batas umum tersebut dapat diperkuat:
\[
    A^k-I=(A-I)(I+A+\cdots+A^{k-1})
    \quad\Longrightarrow\quad
    \rank(A^k-I)\le\rank(A-I).
\]
Kesamaan dapat terjadi, misalnya pada $A=0$.
Identitas khusus untuk faktor-faktor yang sama memberi informasi
lebih kuat daripada hanya menjumlahkan batas rank setiap suku.

\subsection*{Latihan \ref{rank:lat-konstruksi}: memilih citra di dalam kernel}
Karena $\dim\ker A=n-r\ge k-r$, pilih $k-r$ vektor bebas
$u_1,\ldots,u_{k-r}$ di dalam $\ker A$.
Ambil $B$ yang kolom pertamanya adalah vektor-vektor itu dan semua
kolom sisanya nol. Maka $AB=0$ dan $\rank B=k-r$.
Untuk $k=r$, cukup pilih $B=0$.
Konstruksi ini juga menjelaskan mengapa syarat $k\le n$ diperlukan.

\subsection*{Latihan \ref{rank:lat-kernel-sama}: memperluas isomorfisme}
Jika $A=PB$ dengan $P$ invertibel, jelas $Ax=0\iff Bx=0$.
Sebaliknya, asumsikan kernelnya sama dan definisikan
\[
    T:\mathcal C(B)\longrightarrow\mathcal C(A),\qquad T(Bx)=Ax.
\]
Pemetaan ini terdefinisi baik: $Bx=By$ mengakibatkan
$x-y\in\ker B=\ker A$, sehingga $Ax=Ay$.
Ia surjektif menurut definisi ruang kolom dan injektif oleh kesamaan kernel.
Pilih basis ruang kolom $B$, ambil bayangannya sebagai basis ruang
kolom $A$, lalu perluas keduanya menjadi basis $\F^n$.
Pemetaan yang memasangkan kedua basis lengkap memberikan matriks
invertibel $P$ dengan $PBx=Ax$ untuk semua $x$.

Kesamaan rank saja tidak cukup. Matriks
$A=\diag(1,0)$ dan $B=\diag(0,1)$ sama-sama mempunyai rank satu tetapi
mempunyai kernel berbeda, sehingga tidak dapat memenuhi $A=PB$
dengan $P$ invertibel.

\subsection*{Latihan \ref{rank:lat-gabungan}: ruang baris dan kolom berbeda}
Pernyataannya salah. Ambil
\[
    A=\begin{pmatrix}1&0\\0&0\end{pmatrix},\qquad
    B=\begin{pmatrix}0&1\\0&0\end{pmatrix}.
\]
Gabungan horizontal mempunyai rank satu, sedangkan gabungan vertikal
mempunyai rank dua. Ukuran $2$ adalah yang terkecil: pada ukuran $1$,
kedua gabungan sama-sama mempunyai rank nol jika kedua skalar nol,
dan rank satu jika setidaknya satu skalar tak nol.

\subsection*{Latihan \ref{rank:lat-tiga-ruang}: satu irisan yang pasti ada}
Tuliskan
$U=\mathcal C(A)+\mathcal C(B)$ dan
$V=\mathcal C(B)+\mathcal C(C)$.
Karena $\mathcal C(B)\subseteq U\cap V$, diperoleh
\begin{align*}
    \rank[A\ B\ C]
    &=\dim(U+V)=\dim U+\dim V-\dim(U\cap V)\\
    &\le\rank[A\ B]+\rank[B\ C]-\rank B.
\end{align*}
Kesamaan berlaku tepat ketika
\[
    \bigl(\mathcal C(A)+\mathcal C(B)\bigr)
    \cap\bigl(\mathcal C(B)+\mathcal C(C)\bigr)=\mathcal C(B).
\]

\subsection*{Latihan \ref{rank:lat-proyeksi}: satu arah dihilangkan}
Karena $x^*x>0$, rumus $P$ terdefinisi. Untuk setiap $y$,
$Py=y-x(x^*y)/(x^*x)$.
Jadi $Py=0$ tepat ketika $y$ kelipatan $x$.
Selanjutnya, $x^*Py=0$; sebaliknya, jika $x^*z=0$, maka $Pz=z$.
Dengan demikian
\[
    \ker P=\operatorname{span}\{x\},\qquad
    \mathcal C(P)=\{z:x^*z=0\},\qquad \rank P=n-1.
\]

\subsection*{Latihan \ref{rank:lat-stabil}: faktor invertibel dari dekomposisi}
Jika $A^2=BA$ dengan $B$ invertibel, perkalian di kiri tidak
mengubah rank, sehingga $\rank A^2=\rank A$.
Sebaliknya, \autoref{rank:stabil-awal} memberi
\[
    A=S\begin{pmatrix}D&0\\0&0\end{pmatrix}S^{-1}
\]
dengan $D$ invertibel. Pilih
$B=S\diag(D,I_{n-\rank A})S^{-1}$.
Matriks ini invertibel dan $BA=A^2$.
Untuk $A=0$, dapat langsung dipilih $B=I_n$.

\subsection*{Latihan \ref{rank:lat-dua-idempoten}: memaksa dua hasil kali nol}
Dari $(E+F)^2=E+F$ diperoleh $EF+FE=0$.
Kalikan dari kiri dan kanan dengan $E$ secara terpisah:
\[
    EF+EFE=0,\qquad EFE+FE=0.
\]
Maka $EF=FE$, sehingga $2EF=0$ dan $EF=FE=0$ di atas $\R$ atau $\C$.
Arah sebaliknya diperoleh dengan mengembangkan $(E+F)^2$.

Jika $EF=FE=0$, setiap $z\in\mathcal C(E)\cap\mathcal C(F)$ memenuhi
$Ez=z$ karena $E$ idempoten, tetapi juga $Ez=0$ karena $EF=0$.
Jadi irisannya nol. Selain itu,
$(E+F)E=E$ dan $(E+F)F=F$, sehingga
$\mathcal C(E+F)=\mathcal C(E)+\mathcal C(F)$.
Rumus dimensi memberi $\rank(E+F)=\rank E+\rank F$.

\subsection*{Latihan \ref{rank:lat-blok-identitas}: menukar blok kolom}
Pertukaran kedua blok kolom mempertahankan rank dan menghasilkan
\[
    \begin{pmatrix}I&A\\B&0\end{pmatrix}.
\]
Komplemen Schur terhadap $I$ adalah $-BA$, sehingga ranknya
$n+\rank(BA)$.
Di sisi lain, \autoref{rank:segitiga-blok} memberi batas bawah
$\rank A+\rank B$ untuk matriks semula. Maka
\[
    \rank(BA)\ge\rank A+\rank B-n.
\]
Ini adalah Sylvester untuk $BA$; tukar nama $A,B$ untuk bentuk $AB$.

\subsection*{Latihan \ref{rank:lat-nilpoten}: mengekang selisih rank}
Tuliskan $s=\rank A^2$. Barisan ranknya dimulai dengan $11,7,s,0$.
Selisih-selisihnya harus tidak meningkat:
\[
    4\ge7-s\ge s\ge0.
\]
Jadi $s\ge3$ dan $s\le7/2$; karena $s$ bilangan bulat, $s=3$.
Dengan rank--nulitas, $\nul A=11-7=4$ dan $\nul A^2=11-3=8$.
Dengan $N_3$ seperti pada contoh sebelumnya dan
$N_2=\begin{pmatrix}0&1\\0&0\end{pmatrix}$, ambil
\[
    A=\diag(N_3,N_3,N_3,N_2).
\]
Matriks ini berukuran $11$, ranknya $2+2+2+1=7$,
rank kuadratnya $1+1+1=3$, dan pangkat tiganya nol.

\subsection*{Latihan \ref{rank:lat-irisan-kernel}: dua sistem homogen sekaligus}
Ambil $U=\ker A$ dan $V=\ker B$ sebagai subruang $\F^n$.
Rumus dimensi memberi
\[
    \dim(U\cap V)=a+b-\dim(U+V)\ge a+b-n.
\]
Dimensi tidak negatif, dan $U\cap V$ termuat dalam masing-masing
ruang, sehingga juga tidak melebihi $\min\{a,b\}$.
Selanjutnya,
\[
    \ker\begin{pmatrix}A\\B\end{pmatrix}=U\cap V.
\]
Jadi sistem gabungan hanya mempunyai solusi nol tepat ketika
$\nul\binom{A}{B}=0$, atau ekuivalen $\rank\binom{A}{B}=n$.
Syarat $a+b\le n$ perlu untuk itu, tetapi belum cukup;
misalnya, dua kernel yang sama-sama berupa garis yang sama di $\F^2$
tetap beririsan tak nol.

\subsection*{Latihan \ref{rank:lat-nulitas-produk}: memperhatikan ruang perantara}
Rank kedua faktor adalah $\rank A=5-3=2$ dan $\rank B=4-1=3$.
Di ruang perantara $\F^5$, ruang $\mathcal C(B)$ dan $\ker A$
sama-sama berdimensi $3$. Dimensi irisannya sedikitnya $3+3-5=1$
dan paling banyak $3$. Karena itu,
\[
    \nul(AB)=\nul B+\dim(\mathcal C(B)\cap\ker A)\in\{2,3,4\}.
\]
Untuk merealisasikan ketiganya, tetapkan
\[
    A=\begin{pmatrix}1&0&0&0&0\\0&1&0&0&0\\0&0&0&0&0\end{pmatrix}.
\]
Dengan $e_j$ basis standar $\F^5$, pilih secara berturut-turut
\[
    B_2=[e_1\ e_2\ e_3\ 0],\quad
    B_3=[e_1\ e_3\ e_4\ 0],\quad
    B_4=[e_3\ e_4\ e_5\ 0].
\]
Masing-masing mempunyai rank $3$ dan nulitas $1$.
Hasil kali $AB_2,AB_3,AB_4$ mempunyai rank $2,1,0$ dan
nulitas $2,3,4$, sesuai konstruksi yang diminta.
\fi

% !TeX root = ../../ONMIPA Alin.tex
% Sumber utama: Fuzhen Zhang, Linear Algebra: Challenging Problems for
% Students, edisi kedua, Bab 2; mulai hlm. cetak 21 / halaman PDF 40.
\chapter{Determinan}
\label{ch:determinan}

Bab ini membahas sifat dasar determinan dan mengembangkan teknik
perhitungannya melalui identitas, pola entri, rekurensi, dan matriks blok.
Operasi matriks, operasi baris elementer, dan pengertian invers dianggap
sudah dikenal. Semua matriks berentri dalam $\F=\R$ atau $\C$, kecuali
dinyatakan lain. Notasi $|A|$ berarti $\det A$; untuk suatu skalar
kompleks $z$, notasi $|z|$ tetap berarti modulus.

\section{Definisi dan Ekspansi}
\label{det:ekspansi}

\begin{defi}[Minor, kofaktor, dan determinan]
      Untuk $A=(a_{ij})\in\Mat_n(\F)$, tuliskan $A_{ij}$ untuk submatriks
      yang diperoleh dengan menghapus baris ke-$i$ dan kolom ke-$j$.
      Determinan didefinisikan secara rekursif oleh
      \[
            \det(a_{11})=a_{11},\qquad
            \det A=\sum_{j=1}^n(-1)^{1+j}a_{1j}\det A_{1j}\quad(n\ge2).
      \]
      Bilangan $M_{ij}=\det A_{ij}$ disebut \emph{minor} dari entri $a_{ij}$,
      sedangkan $C_{ij}=(-1)^{i+j}M_{ij}$ disebut \emph{kofaktornya}.
      Untuk konsistensi pada ukuran satu, determinan matriks kosong ditetapkan $1$.
\end{defi}

Definisi di atas menggunakan baris pertama. Untuk membuktikan bahwa
baris atau kolom lain memberikan nilai yang sama, diperlukan hubungan
dengan permutasi.

\begin{defi}[Tanda permutasi]
      Tuliskan $S_n$ untuk himpunan seluruh permutasi $\{1,\ldots,n\}$.
      Untuk $\sigma\in S_n$, definisikan banyaknya inversi
      \[
            \nu(\sigma)=\#\{(i,j):i<j,\ \sigma(i)>\sigma(j)\},
            \qquad \operatorname{sgn}(\sigma)=(-1)^{\nu(\sigma)}.
      \]
      Permutasi disebut genap jika tandanya $1$, dan ganjil jika tandanya $-1$.
      Satu pertukaran dua posisi membalik tanda permutasi.
\end{defi}

\begin{thm}[Rumus Leibniz]
      \label{det:leibniz}
      Untuk $A\in\Mat_n(\F)$,
      \[
            \det A=\sum_{\sigma\in S_n}\operatorname{sgn}(\sigma)
            \prod_{i=1}^n a_{i,\sigma(i)}.
      \]
      Setiap suku mengambil tepat satu entri dari setiap baris dan setiap kolom.
\end{thm}
\begin{proof}
      Gunakan induksi pada $n$. Kelompokkan permutasi menurut nilai
      $\sigma(1)=j$. Setelah posisi pertama dan nilai $j$ dihapus,
      permutasi sisanya menjadi permutasi pada $n-1$ unsur.
      Pasangan inversi yang melibatkan posisi pertama berjumlah $j-1$;
      karena $(-1)^{j-1}=(-1)^{1+j}$, kontribusi kelompok ini adalah
      $(-1)^{1+j}a_{1j}\det A_{1j}$.
      Jumlah seluruh kelompok tepat menghasilkan definisi rekursif.
\end{proof}

\begin{thm}[Ekspansi Laplace]
      \label{det:laplace}
      Untuk setiap baris $i$ dan setiap kolom $j$,
      \[
            \det A=\sum_{k=1}^n a_{ik}C_{ik}
            =\sum_{k=1}^n a_{kj}C_{kj}.
      \]
\end{thm}
\begin{proof}
      Kelompokkan suku Leibniz menurut entri yang diambil dari baris $i$.
      Pada kelompok yang mengambil entri $(i,k)$, menghapus baris $i$ dan
      kolom $k$ menyisakan determinan $A_{ik}$ dengan faktor tanda
      $(-1)^{i+k}$. Ini memberikan ekspansi baris.
      Pengelompokan menurut kolom memberikan ekspansi kolom dengan alasan sama.
\end{proof}

Rumus Leibniz terutama berguna untuk membaca struktur: jika tidak ada
permutasi yang memilih seluruh entri tak nol, determinannya nol.
Ekspansi Laplace berguna ketika terdapat baris atau kolom dengan sedikit
entri tak nol. Mengembangkan seluruh $n!$ suku biasanya bukan cara
perhitungan yang efektif.

\section{Sifat Fundamental}
\label{det:sifat}

\begin{thm}[Multilinear dan berselang-seling]
      \label{det:multilinear}
      Determinan linear terhadap \emph{masing-masing} kolom ketika kolom
      lain dipertahankan. Menukar dua kolom mengubah tanda determinan;
      jika dua kolom sama, determinannya nol. Pernyataan yang sama berlaku
      untuk baris.
\end{thm}
\begin{proof}
      Dalam setiap suku Leibniz terdapat tepat satu faktor dari kolom yang
      sedang diubah; distributivitas memberi linearitas terhadap kolom itu.
      Menukar dua kolom memasangkan setiap permutasi dengan permutasi
      bertanda berlawanan. Jika kolom yang ditukar sama, matriks tidak
      berubah, sehingga $\det A=-\det A$ dan $\det A=0$.
      Argumen baris identik.
\end{proof}

Akibatnya, tiga operasi elementer mempunyai pengaruh berikut:
\[
      \begin{array}{c|c}
            \text{Operasi pada }A\text{ menghasilkan }B & \det B  \\ \hline
            R_i\leftrightarrow R_j                      & -\det A \\
            R_i\leftarrow cR_i                          & c\det A \\
            R_i\leftarrow R_i+cR_j,\quad i\ne j         & \det A
      \end{array}
\]
Pada operasi kedua, $c\ne0$ diperlukan jika operasi hendak dibalik;
identitas determinannya sendiri juga benar untuk $c=0$.
Aturan yang sama berlaku untuk kolom. Aturan penambahan baris mengikuti
multilinearitas: suku tambahan mempunyai dua baris yang sama.

\begin{lem}[Karakterisasi determinan]
      \label{det:unik}
      Jika fungsi $f(v_1,\ldots,v_n)$ multilinear dan berselang-seling
      terhadap $n$ vektor kolom dalam $\F^n$, maka
      \[
            f(v_1,\ldots,v_n)=f(e_1,\ldots,e_n)\det[v_1\ \cdots\ v_n].
      \]
\end{lem}
\begin{proof}
      Ekspansikan setiap $v_j$ sebagai kombinasi vektor koordinat $e_i$.
      Multilinearitas menghasilkan jumlah nilai $f(e_{i_1},\ldots,e_{i_n})$.
      Suku dengan indeks berulang nol, sedangkan suku tanpa indeks berulang
      merupakan permutasi dari $f(e_1,\ldots,e_n)$ dengan tanda permutasinya.
      Jumlah koefisien yang tersisa adalah rumus Leibniz.
\end{proof}

\begin{thm}[Determinan hasil kali]
      \label{det:produk}
      Untuk $A,B\in\Mat_n(\F)$, berlaku $\det(AB)=\det A\det B$.
\end{thm}
\begin{proof}
      Dengan $A$ tetap, definisikan
      $f(b_1,\ldots,b_n)=\det[Ab_1\ \cdots\ Ab_n]$.
      Fungsi ini multilinear dan berselang-seling, serta
      $f(e_1,\ldots,e_n)=\det A$.
      Terapkan \autoref{det:unik} pada kolom-kolom $B$.
\end{proof}

\begin{prop}[Sifat-sifat dasar determinan]
      \label{det:dasar}
      Untuk $A,B\in\Mat_n(\F)$, skalar $c$, dan matriks invertibel $S$:
      \begin{enumerate}[label=(d\arabic*),nosep]
            \item $A$ singular jika dan hanya jika $\det A=0$.
            \item Baris-baris atau kolom-kolom $A$ bergantung linear jika dan
                  hanya jika $\det A=0$.
            \item Jika $A$ segitiga, $\det A=\prod_{i=1}^n a_{ii}$.
            \item $\det(A^T)=\det A$.
            \item $\det(cA)=c^n\det A$.
            \item $\det(AB)=\det A\det B$.
            \item $\det(S^{-1}AS)=\det A$.
      \end{enumerate}
\end{prop}
\begin{proof}
      Pada matriks segitiga, hanya permutasi identitas yang dapat memberi
      suku tak nol, sehingga (d3) berlaku. Mengganti permutasi dengan
      inversnya dalam rumus Leibniz membuktikan (d4); multilinearitas
      memberi (d5). Sifat (d6) sudah dibuktikan.

      Untuk (d1)--(d2), reduksi baris ke bentuk eselon mempertahankan
      kenolannya determinan. Jika terdapat baris nol, determinannya nol
      dan kolom-kolomnya tidak bebas. Jika semua pivot ada, matriks eselon
      segitiga mempunyai diagonal tak nol, sehingga determinannya tak nol
      dan matriks dapat direduksi menjadi $I_n$; semua OBE dapat dibalik,
      maka $A$ invertibel. Transpose memberikan pernyataan untuk baris.

      Akhirnya, $\det(S^{-1})\det S=\det I=1$. Dengan (d6),
      \[
            \det(S^{-1}AS)=\det(S^{-1})\det A\det S=\det A.
      \]
\end{proof}

\begin{exam}[Linear pada kolom, bukan pada matriks]
      Misalkan $A=[a\ r_2\ r_3\ r_4]$, $B=[b\ r_2\ r_3\ r_4]$,
      $\det A=4$, dan $\det B=1$. Maka
      \[
            \det(A+B)=2^3\det[a+b\ r_2\ r_3\ r_4]=8(4+1)=40.
      \]
      Jika $C=[r_4\ r_3\ r_2\ a+b]$, urutan empat kolom dibalik melalui
      $\binom42=6$ pertukaran, sehingga $\det C=(-1)^6(4+1)=5$.
      Pada umumnya $\det(A+B)\ne\det A+\det B$.
\end{exam}

\subsection{Tanda permutasi dan kenolan struktural}
Jika urutan $n$ kolom dibalik, faktor tandanya adalah
$(-1)^{n(n-1)/2}$. Karena itu, untuk matriks dengan entri nol
di sebelah kiri diagonal sekundernya,
\[
      a_{ij}=0\text{ jika }i+j<n+1
      \quad\Longrightarrow\quad
      \det A=(-1)^{n(n-1)/2}\prod_{i=1}^n a_{i,n+1-i}.
\]
Pembalikan kolom mengubah matriks tersebut menjadi segitiga bawah.
Argumen ini juga menerangkan tanda determinan matriks pembalik urutan.

\begin{exam}[Determinan nol tanpa ekspansi]
      Determinan matriks berikut nol untuk semua nilai entrinya:
      \[
            \begin{vmatrix}
                  a_1 & a_2 & a_3 & a_4 & a_5 \\b_1&b_2&b_3&b_4&b_5\\
                  c_1 & c_2 & 0   & 0   & 0   \\d_1&d_2&0&0&0\\e_1&e_2&0&0&0
            \end{vmatrix}=0.
      \]
      Dalam suku Leibniz yang tidak nol, tiga baris terakhir harus memilih
      kolom dari himpunan $\{1,2\}$. Tiga kolom yang berbeda tidak dapat
      dipilih dari dua kemungkinan. Jadi setiap suku Leibniz nol.
\end{exam}

\section{Adjugat dan Perubahan Entri}
\label{det:adjugat}

\begin{defi}[Adjugat]
      Matriks adjugat $A$, ditulis $\operatorname{adj}(A)$, adalah transpose
      matriks kofaktor:
      \[
            (\operatorname{adj}(A))_{ij}=C_{ji}.
      \]
      Adjugat berbeda dari transpose konjugat $A^*=\overline A^{\,T}$.
\end{defi}

\begin{thm}[Identitas adjugat]
      \label{det:identitas-adj}
      Untuk setiap matriks persegi, termasuk yang singular,
      \[
            A\operatorname{adj}(A)=\operatorname{adj}(A)A=(\det A)I_n.
      \]
      Jika $\det A\ne0$, maka $A^{-1}=(\det A)^{-1}\operatorname{adj}(A)$.
\end{thm}
\begin{proof}
      Entri $(i,j)$ dari $A\operatorname{adj}(A)$ adalah
      $\sum_k a_{ik}C_{jk}$. Jika $i=j$, jumlah ini adalah ekspansi Laplace
      dan bernilai $\det A$. Jika $i\ne j$, jumlah tersebut adalah
      determinan matriks yang baris ke-$j$-nya diganti dengan baris ke-$i$.
      Kofaktor pada baris $j$ tidak berubah karena baris itu dihapus ketika
      minor dihitung. Matriks pengganti mempunyai dua baris sama, sehingga
      determinannya nol. Identitas dari sisi lain dibuktikan dengan kolom.
\end{proof}

\begin{thm}[Perubahan rank satu]
      \label{det:rankone}
      Untuk $u,v\in\F^n$ dan setiap $A\in\Mat_n(\F)$,
      \[
            \boxed{\det(A+uv^T)=\det A+v^T\operatorname{adj}(A)u.}
      \]
      Jika $A$ invertibel, rumus ini menjadi
      \[
            \det(A+uv^T)=\det A\,(1+v^TA^{-1}u).
      \]
\end{thm}
\begin{proof}
      Kolom ke-$j$ dari $A+uv^T$ adalah $a_j+v_ju$.
      Ekspansikan determinan secara multilinear. Semua suku yang memilih
      $u$ dari dua kolom atau lebih nol, karena dua kolom itu sebanding.
      Suku tanpa $u$ adalah $\det A$; jumlah suku dengan tepat satu $u$ ialah
      \[
            \sum_{j=1}^n v_j\sum_{i=1}^n u_iC_{ij}
            =v^T\operatorname{adj}(A)u.
      \]
      Untuk $A$ invertibel, gunakan
      $\operatorname{adj}(A)=(\det A)A^{-1}$.
\end{proof}

Khususnya, $\det(I_n+uv^T)=1+v^Tu$. Dengan $v^T$ diganti oleh
$-y^*$, diperoleh $\det(I_n-xy^*)=1-y^*x$. Rumus adjugat di atas tetap sah
ketika $A$ singular; rumus yang mengandung $A^{-1}$ tidak boleh
dipakai pada kasus tersebut.

\begin{prop}[Determinan sebagai polinom]
      \label{det:principal}
      Jika $A[S,S]$ menyatakan submatriks utama pada himpunan indeks
      $S\subseteq\{1,\ldots,n\}$, maka
      \[
            \det(I_n+tA)=\sum_{S\subseteq\{1,\ldots,n\}}
            t^{|S|}\det A[S,S].
      \]
      Suku untuk $S=\varnothing$ bernilai $1$. Khususnya, koefisien $t$
      adalah $\tr(A)$ dan koefisien $t^n$ adalah $\det A$.
\end{prop}
\begin{proof}
      Pada kolom ke-$j$, pilih $e_j$ atau $ta_j$ lalu gunakan multilinearitas.
      Jika $ta_j$ dipilih tepat untuk $j\in S$, kolom di luar $S$ memaksa
      permutasi Leibniz menetapkan setiap indeks di luar $S$.
      Jumlah suku yang tersisa adalah $t^{|S|}\det A[S,S]$.
\end{proof}

Untuk $n\ge2$, koefisien $t^2$ adalah
\[
      \sum_{i<j}(a_{ii}a_{jj}-a_{ij}a_{ji})
      =\frac12\bigl((\tr A)^2-\tr(A^2)\bigr).
\]
Dengan demikian, beberapa informasi determinan dapat diperoleh tanpa
menghitung keseluruhan polinom.

\begin{prop}[Rumus diferensiasi Jacobi]
      Jika setiap entri $A(t)$ dapat diturunkan terhadap parameter real $t$,
      maka
      \[
            \frac{d}{dt}\det A(t)
            =\sum_{i,j}C_{ij}(t)a'_{ij}(t)
            =\tr\bigl(\operatorname{adj}(A(t))A'(t)\bigr).
      \]
      Jika $A(t)$ invertibel, bentuknya menjadi
      $\det A(t)\,\tr(A(t)^{-1}A'(t))$.
\end{prop}
\begin{proof}
      Turunkan rumus Leibniz dengan aturan turunan hasil kali.
      Kelompokkan suku menurut entri yang diturunkan; koefisien
      $a'_{ij}(t)$ adalah kofaktor $C_{ij}(t)$.
      Identitas trace mengikuti definisi adjugat.
\end{proof}
Sebagai akibat, $\left.\frac{d}{dt}\det(I_n+tA)\right|_{t=0}=\tr(A)$.
Rumus tanpa invers tetap berlaku pada nilai $t$ yang membuat $A(t)$ singular.

\section{Pola Entri dan Parameter}
\label{det:parameter}

\begin{thm}[Diagonal ditambah matriks rank satu]
      \label{det:diagrankone}
      Untuk $D=\diag(d_1,\ldots,d_n)$ dan $u,v\in\F^n$,
      \[
            \det(D+uv^T)=\prod_{i=1}^n d_i
            +\sum_{i=1}^n u_iv_i\prod_{j\ne i}d_j.
      \]
      Tidak disyaratkan semua $d_i$ tak nol.
\end{thm}
\begin{proof}
      Adjugat matriks diagonal adalah matriks diagonal dengan entri ke-$i$
      $\prod_{j\ne i}d_j$. Terapkan \autoref{det:rankone}.
\end{proof}

\begin{exam}[Entri di luar diagonal sama]
      Tentukan determinan matriks dengan entri diagonal $p_1,\ldots,p_n$
      dan seluruh entri di luar diagonal sama dengan $a$.
\end{exam}
\begin{solution}
      Dengan $\boldsymbol{1}=(1,\ldots,1)^T$, matriksnya adalah
      \[
            M=\diag(p_1-a,\ldots,p_n-a)+a\boldsymbol{1}\boldsymbol{1}^T.
      \]
      Maka
      \[
            \det M=\prod_{i=1}^n(p_i-a)
            +a\sum_{i=1}^n\prod_{j\ne i}(p_j-a).
      \]
      Jika $f(x)=\prod_i(p_i-x)$ dan $f_i(a)=\prod_{j\ne i}(p_j-a)$,
      hasil ini sama dengan $f(a)+a\sum_i f_i(a)$.
      Karena $f(a)=(p_n-a)f_n(a)$, bentuk ekuivalennya ialah
      $a\sum_{i=1}^{n-1}f_i(a)+p_nf_n(a)$.

      Khususnya, jika semua entri diagonal adalah $\alpha$ dan semua entri
      di luar diagonal adalah $\beta$, maka
      \[
            \boxed{\det M=(\alpha-\beta)^{n-1}
                  \bigl(\alpha+(n-1)\beta\bigr).}
      \]
      Untuk $n\ge2$, kedua kemungkinan
      singularnya adalah $\alpha=\beta$ atau $\alpha=-(n-1)\beta$.
\end{solution}

\begin{exam}[Dua konstanta di atas dan di bawah diagonal]
      \label{det:zhang26}
      Misalkan
      \[
            M=\begin{pmatrix}
                  p_1    & a      & a      & \cdots & a      \\
                  b      & p_2    & a      & \cdots & a      \\
                  b      & b      & p_3    & \cdots & a      \\
                  \vdots & \vdots & \vdots & \ddots & \vdots \\
                  b      & b      & b      & \cdots & p_n
            \end{pmatrix},\qquad f(x)=\prod_{i=1}^n(p_i-x).
      \]
      Untuk $a\ne b$, buktikan
      \[
            \det M=\frac{bf(a)-af(b)}{b-a}.
      \]
\end{exam}
\begin{solution}
      \ide{Kurangkan satu matriks rank satu agar sisanya segitiga.}
      Tuliskan $M=L+a\boldsymbol{1}\boldsymbol{1}^T$, dengan $L$ segitiga
      bawah, diagonal $p_i-a$, dan semua entri di bawah diagonal $b-a$.
      Mula-mula anggap $p_i\ne a$ untuk setiap $i$, sehingga $L$ invertibel.
      Jika $y=L^{-1}\boldsymbol{1}$ dan $s_i=y_1+\cdots+y_i$, $s_0=0$,
      persamaan baris ke-$i$ memberikan
      \[
            (p_i-a)y_i+(b-a)s_{i-1}=1.
      \]
      Definisikan $r_i=1-(b-a)s_i$. Karena
      \[
            r_i=r_{i-1}-(b-a)y_i
            =\frac{p_i-b}{p_i-a}r_{i-1},\qquad r_0=1,
      \]
      diperoleh $r_n=f(b)/f(a)$ dan $s_n=(1-r_n)/(b-a)$.
      Sekarang terapkan rumus perubahan rank satu:
      \begin{align*}
            \det M
             & =\det L\,(1+a\boldsymbol{1}^TL^{-1}\boldsymbol{1}) \\
             & =f(a)(1+as_n)
            =\frac{bf(a)-af(b)}{b-a}.
      \end{align*}
      Setelah dikalikan $b-a$, kedua ruas adalah polinom dalam
      $p_1,\ldots,p_n$. Identitas yang benar untuk semua $p_i\ne a$
      tetap benar ketika sebagian $p_i=a$, melalui identitas polinom
      atau kontinuitas. Jadi syarat sementara tersebut dapat dihapus.
      Kasus $a=b$ ditangani oleh contoh sebelumnya, bukan dengan membagi oleh nol.
\end{solution}

\section{Determinan Vandermonde}
\label{det:vandermonde}

\begin{thm}[Vandermonde]
      \label{det:rumus-vandermonde}
      Untuk $x_1,\ldots,x_n\in\F$,
      \[
            V_n(x_1,\ldots,x_n)=
            \begin{vmatrix}
                  1         & 1         & \cdots & 1         \\x_1&x_2&\cdots&x_n\\
                  x_1^2     & x_2^2     & \cdots & x_n^2     \\
                  \vdots    & \vdots    &        & \vdots    \\
                  x_1^{n-1} & x_2^{n-1} & \cdots & x_n^{n-1}
            \end{vmatrix}
            =\prod_{1\le i<j\le n}(x_j-x_i).
      \]
\end{thm}
\begin{proof}
      Pandang $V_n$ sebagai polinom dalam $x_n$, dengan koefisien berupa polinom
      dalam $x_1,\ldots,x_{n-1}$. Derajatnya paling tinggi $n-1$.
      Jika $x_n=x_i$ untuk $i<n$, dua kolom sama, sehingga polinom tersebut
      memiliki faktor $x_n-x_i$. Koefisien $x_n^{n-1}$ adalah
      $V_{n-1}(x_1,\ldots,x_{n-1})$, hasil ekspansi pada entri $(n,n)$.
      Karena itu
      \[
            V_n=V_{n-1}\prod_{i=1}^{n-1}(x_n-x_i).
      \]
      Induksi dengan $V_1=1$ menyelesaikan bukti sebagai identitas polinom,
      termasuk saat beberapa $x_i$ berimpit.
\end{proof}

\begin{prop}[Basis polinom lain]
      Jika $p_i(x)$ berderajat tepat $i$ dengan koefisien utama $c_i$,
      untuk $i=0,\ldots,n-1$, maka
      \[
            \det\bigl(p_{i-1}(x_j)\bigr)_{i,j=1}^n
            =\left(\prod_{i=0}^{n-1}c_i\right)
            \prod_{1\le i<j\le n}(x_j-x_i).
      \]
\end{prop}
\begin{proof}
      Tuliskan setiap $p_i$ sebagai kombinasi $1,x,\ldots,x^i$.
      Matriks evaluasinya merupakan hasil kali matriks koefisien segitiga
      bawah dengan matriks Vandermonde. Diagonal matriks koefisien adalah
      $c_0,\ldots,c_{n-1}$; gunakan determinan hasil kali.
\end{proof}

\begin{exam}[Pangkat dan koefisien binomial]
      Untuk $V=(j^{i-1})_{i,j=1}^n$, rumus Vandermonde memberi
      \[
            \det V=\prod_{1\le i<j\le n}(j-i)
            =\prod_{j=2}^n(j-1)!=\prod_{k=1}^{n-1}k!.
      \]
      Sebagai pengembangan, jika $B_{ij}=\binom{j-1}{i-1}$, maka
      $\det B=1$. Memang, polinom
      \[
            p_k(x)=\frac{x(x-1)\cdots(x-k+1)}{k!},\qquad p_0(x)=1,
      \]
      mempunyai koefisien utama $1/k!$. Evaluasi pada $0,1,\ldots,n-1$
      menghasilkan faktor Vandermonde $\prod_{k=1}^{n-1}k!$, yang tepat
      membatalkan hasil kali koefisien utama.
\end{exam}

\section{Determinan melalui Rekurensi}
\label{det:rekurensi}

\begin{prop}[Matriks tridiagonal]
      Definisikan $D_n=\det T_n$, dengan
      \[
            T_n=\begin{pmatrix}
                  p      & q      & 0      & \cdots & 0      \\
                  r      & p      & q      & \ddots & \vdots \\
                  0      & r      & p      & \ddots & 0      \\
                  \vdots & \ddots & \ddots & \ddots & q      \\
                  0      & \cdots & 0      & r      & p
            \end{pmatrix}.
      \]
      Dengan $D_0=1$ dan $D_1=p$, berlaku
      \[
            D_n=pD_{n-1}-qrD_{n-2}\qquad(n\ge2).
      \]
\end{prop}
\begin{proof}
      Ekspansi pada baris pertama memberi $pD_{n-1}$ dari entri $(1,1)$.
      Minor entri $(1,2)$ mempunyai kolom pertama $(r,0,\ldots,0)^T$,
      sehingga determinannya $rD_{n-2}$. Tanda kofaktor posisi $(1,2)$
      negatif, sehingga suku kedua adalah $-qrD_{n-2}$.
\end{proof}

\begin{exam}[Dua akar rekurensi yang dapat berimpit]
      Ambil $p=a+b$, $q=ab$, dan $r=1$ pada $T_n$.
      Tentukan $D_n$, termasuk ketika $a=b$.
\end{exam}
\begin{solution}
      Rekurensinya
      $D_n=(a+b)D_{n-1}-abD_{n-2}$, dengan $D_0=1$ dan $D_1=a+b$.
      Jika $a\ne b$, barisan
      \[
            D_n=\frac{a^{n+1}-b^{n+1}}{a-b}
            =\sum_{k=0}^n a^{n-k}b^k
      \]
      memenuhi rekurensi dan kedua nilai awal, sehingga merupakan jawabannya.
      Bentuk jumlah menunjukkan langsung bahwa untuk $a=b$,
      \[
            D_n=(n+1)a^n\qquad(n\ge1),
      \]
      dan $D_0=1$. Kasus akar berimpit ini tidak memerlukan pembagian oleh $a-b$.
\end{solution}

\begin{exam}[Kasus Fibonacci]
      Untuk $p=q=1$ dan $r=-1$, rekurensinya $D_n=D_{n-1}+D_{n-2}$,
      dengan $D_0=D_1=1$. Dengan konvensi $F_0=0$, $F_1=1$, diperoleh
      $D_n=F_{n+1}$. Khususnya determinan berukuran $5$ bernilai $F_6=8$.
\end{exam}

\begin{exam}[Determinan matriks pendamping]
      Untuk $a_0,\ldots,a_{n-1}\in\F$, misalkan
      \[
            C=\begin{pmatrix}
                  0      & 1      & 0      & \cdots   & 0        \\
                  0      & 0      & 1      & \ddots   & \vdots   \\
                  \vdots & \vdots & \ddots & \ddots   & 0        \\
                  0      & 0      & \cdots & 0        & 1        \\
                  -a_0   & -a_1   & \cdots & -a_{n-2} & -a_{n-1}
            \end{pmatrix}.
      \]
      Buktikan $\det(tI_n-C)=t^n+a_{n-1}t^{n-1}+\cdots+a_1t+a_0$.
\end{exam}
\begin{solution}
      Tuliskan determinan yang dicari sebagai $P_n(t)$.
      Ekspansi kolom pertama pada $tI-C$ memberi dua kontribusi:
      $t$ kali determinan berukuran $n-1$ dengan koefisien
      $a_1,\ldots,a_{n-1}$, dan kontribusi $a_0$ dari baris terakhir.
      Pada kontribusi kedua, minor segitiga mempunyai diagonal $-1$,
      sehingga tandanya $(-1)^{n+1}(-1)^{n-1}=1$.
      Dengan demikian,
      \[
            P_n(t)=tP_{n-1}(t;a_1,\ldots,a_{n-1})+a_0.
      \]
      Induksi, dimulai dari $P_1(t)=t+a_0$, memberi polinom yang diminta.
      Perhitungan ini murni merupakan ekspansi determinan.
\end{solution}

\section{Determinan Matriks Blok}
\label{det:blok}

Ukuran blok harus sesuai, tetapi blok tidak dapat diperlakukan sebagai
skalar komutatif. Misalnya, bentuk $\det(AD-BC)$ untuk matriks
$\left(\begin{smallmatrix}A&B\\C&D\end{smallmatrix}\right)$ memerlukan
hipotesis tambahan dan urutan faktor yang tepat.

\begin{prop}[Segitiga blok]
      Jika $A$ dan $D$ persegi, maka
      \[
            \det\begin{pmatrix}A&B\\0&D\end{pmatrix}
            =\det\begin{pmatrix}A&0\\C&D\end{pmatrix}
            =\det A\det D.
      \]
\end{prop}
\begin{proof}
      Pada bentuk pertama, setiap suku Leibniz tak nol harus memasangkan
      baris-baris blok bawah dengan kolom-kolom blok kanan.
      Akibatnya, baris blok atas dipasangkan dengan kolom blok kiri.
      Permutasi terurai menjadi dua permutasi dalam masing-masing blok;
      penjumlahannya menghasilkan $\det A\det D$. Bentuk kedua mengikuti
      transpose. Tidak diperlukan invertibilitas blok diagonal.
\end{proof}

\begin{thm}[Komplemen Schur]
      \label{det:schur}
      Untuk $M=\begin{pmatrix}A&B\\C&D\end{pmatrix}$ dengan $A,D$ persegi,
      \[
            \det M=\det A\det(D-CA^{-1}B)\quad\text{jika }A\text{ invertibel}.
      \]
      Jika $D$ invertibel, rumus alternatifnya adalah
      $\det M=\det D\det(A-BD^{-1}C)$.
\end{thm}
\begin{proof}
      Kalikan dari kiri dengan matriks segitiga blok berdeterminan $1$:
      \[
            \begin{pmatrix}I&0\\-CA^{-1}&I\end{pmatrix}
            \begin{pmatrix}A&B\\C&D\end{pmatrix}
            =\begin{pmatrix}A&B\\0&D-CA^{-1}B\end{pmatrix}.
      \]
      Gunakan determinan hasil kali dan segitiga blok.
      Untuk rumus kedua, gunakan pengali
      $\left(\begin{smallmatrix}I&-BD^{-1}\\0&I\end{smallmatrix}\right)$.
\end{proof}

\begin{thm}[Kondisi komutativitas yang cukup]
      \label{det:blokkomutatif}
      Jika $A,B,C,D\in\Mat_n(\F)$ dan $AC=CA$, maka
      \[
            \det\begin{pmatrix}A&B\\C&D\end{pmatrix}=\det(AD-CB).
      \]
\end{thm}
\begin{proof}
      Jika $A$ invertibel, rumus Schur memberi
      \begin{align*}
            \det M
             & =\det A\det(D-CA^{-1}B)          \\
             & =\det(AD-ACA^{-1}B)=\det(AD-CB).
      \end{align*}
      Kesamaan terakhir memakai $AC=CA$.
      Untuk $A$ singular, ganti $A$ dengan $A+tI$.
      Matriks ini tetap komutatif dengan $C$ dan invertibel untuk semua
      $t$ kecuali sejumlah hingga nilai, sebab $\det(A+tI)$ adalah
      polinom monik berderajat $n$.
      Kedua determinan menjadi polinom dalam $t$ yang sama pada tak hingga
      banyak nilai, sehingga identik. Substitusi $t=0$ menyelesaikan kasus singular.
\end{proof}

Urutan $CB$ pada kesimpulan harus dipertahankan.
Latihan~\ref{det:lat-blok-syarat} menunjukkan bahwa menukarnya menjadi
$BC$, atau menghapus syarat $AC=CA$, dapat mengubah hasil.

\begin{prop}[Identitas Sylvester untuk determinan]
      Untuk $A\in\Mat_{m\times n}(\F)$ dan $B\in\Mat_{n\times m}(\F)$,
      \[
            \det(I_m+AB)=\det(I_n+BA).
      \]
\end{prop}
\begin{proof}
      Hitung determinan
      $\left(\begin{smallmatrix}I_m&A\\-B&I_n\end{smallmatrix}\right)$
      dengan komplemen Schur terhadap $I_m$, lalu terhadap $I_n$.
      Kedua hasilnya adalah ruas kanan dan ruas kiri identitas tersebut.
\end{proof}

Identitas ini mengurangi ukuran masalah jika salah satu ukuran jauh
lebih kecil. Contohnya, jika $U,V$ berukuran $n\times2$, maka
\[
      \det(I_n+UV^T)=\det(I_2+V^TU).
\]
Dengan demikian, perubahan rank paling tinggi dua dapat dihitung
melalui satu determinan $2\times2$, berapa pun besarnya $n$.

\begin{exam}[Blok dengan pola simetris]
      Untuk sembarang $A,B\in\Mat_n(\F)$, buktikan
      \[
            \det\begin{pmatrix}A&B\\B&A\end{pmatrix}
            =\det(A+B)\det(A-B).
      \]
\end{exam}
\begin{solution}
      Ambil $T=\left(\begin{smallmatrix}I&I\\I&-I\end{smallmatrix}\right)$,
      sehingga $T^{-1}=T/2$. Perkalian blok memberi
      \[
            T^{-1}\begin{pmatrix}A&B\\B&A\end{pmatrix}T
            =\begin{pmatrix}A+B&0\\0&A-B\end{pmatrix}.
      \]
      Determinan tidak berubah oleh keserupaan, lalu rumus blok diagonal
      memberikan hasilnya. Syarat $AB=BA$ tidak diperlukan.
\end{solution}

\begin{exam}[Determinan blok real yang tidak negatif]
      Jika $A,B\in\Mat_n(\R)$, buktikan
      \[
            \det\begin{pmatrix}A&B\\-B&A\end{pmatrix}\ge0.
      \]
\end{exam}
\begin{solution}
      Lakukan perhitungan di $\C$ dengan
      $T=\left(\begin{smallmatrix}I&I\\iI&-iI\end{smallmatrix}\right)$.
      Matriks ini invertibel dan
      \[
            T^{-1}\begin{pmatrix}A&B\\-B&A\end{pmatrix}T
            =\begin{pmatrix}A+iB&0\\0&A-iB\end{pmatrix}.
      \]
      Karena $A,B$ real, kedua determinan blok diagonal saling konjugat.
      Oleh karena itu,
      \[
            \det\begin{pmatrix}A&B\\-B&A\end{pmatrix}
            =\det(A+iB)\,\overline{\det(A+iB)}
            =|\det(A+iB)|^2\ge0.
      \]
      Jika $A,B$ boleh kompleks, faktorisasi menjadi $\det(A+iB)\det(A-iB)$
      masih benar, tetapi kesimpulan tidak negatif tidak berlaku umum.
      Sebagai contoh, pada ukuran satu, $A=(i)$ dan $B=(0)$ memberi determinan $-1$.
\end{solution}

\clearpage
\section{Latihan Pilihan}
\label{det:latihan}

Soal-soal berikut melatih penggunaan sifat dan teknik determinan.
Sertakan identitas atau struktur yang digunakan dalam setiap jawaban.

\begin{latihan}\label{det:lat-tigadeterminan}
      Tentukan ketiga determinan berikut, dengan $x\in\F$:
      \[
            \begin{vmatrix}1&2&3\\8&9&4\\7&6&5\end{vmatrix},\qquad
            \begin{vmatrix}1+x&2+x&3+x\\8+x&9+x&4+x\\7+x&6+x&5+x\end{vmatrix},
            \qquad
            \begin{vmatrix}x&x^2&x^3\\x^8&x^9&x^4\\x^7&x^6&x^5\end{vmatrix}.
      \]
\end{latihan}

\begin{latihan}\label{det:lat-divisibilitas}
      Jika seluruh entri $A\in\Mat_n(\Z)$ bernilai $1$ atau $-1$,
      buktikan bahwa $2^{n-1}$ membagi $\det A$.
\end{latihan}

\begin{latihan}\label{det:lat-paritas}
      Misalkan $A$ berukuran $n\times n$.
      \begin{enumerate}[label=(\alph*)]
            \item Jika $A^T=-A$ dan $n$ ganjil, buktikan $\det A=0$.
            \item Jika $A\in\Mat_n(\R)$ memenuhi $A^2+I_n=0$, buktikan $n$ genap.
            \item Apakah kesimpulan (b) tetap berlaku untuk matriks kompleks?
      \end{enumerate}
\end{latihan}

\begin{latihan}\label{det:lat-ortogonal}
      Misalkan $A\in\Mat_n(\C)$ memenuhi $AA^T=I_n$ dan $\det A<0$.
      Tentukan $\det(A+I_n)$ tanpa menggunakan nilai eigen.
\end{latihan}

\begin{latihan}\label{det:lat-blok-syarat}
      Untuk $M=\left(\begin{smallmatrix}A&B\\C&D\end{smallmatrix}\right)$
      dengan blok real $2\times2$, berikan contoh tandingan untuk:
      \begin{enumerate}[label=(\alph*)]
            \item mengganti $\det(AD-CB)$ oleh $\det(AD-BC)$ meskipun $AC=CA$;
            \item menggunakan $\det M=\det(AD-CB)$ tanpa syarat $AC=CA$.
      \end{enumerate}
\end{latihan}

\begin{latihan}\label{det:lat-dependensi}
      Misalkan $A\ne0$ dan $B_1,B_2,B_3,B_4\in\Mat_2(\R)$, dengan
      $A\in\Mat_2(\R)$. Jika
      \[
            \det(A+B_i)=\det A+\det B_i,\qquad i=1,2,3,4,
      \]
      buktikan bahwa $B_1,B_2,B_3,B_4$ bergantung linear atas $\R$.
\end{latihan}

\begin{latihan}\label{det:lat-permutasiblok}
      Tentukan kedua determinan berikut dalam bentuk terfaktor:
      \[
            \begin{vmatrix}
                  0 & 0 & a_1 & b_1 \\0&0&a_2&b_2\\a_3&b_3&0&0\\a_4&b_4&0&0
            \end{vmatrix},\qquad
            \begin{vmatrix}
                  a_1 & 0 & 0 & b_1 \\0&a_2&b_2&0\\0&b_3&a_3&0\\b_4&0&0&a_4
            \end{vmatrix}.
      \]
\end{latihan}

\ifpembahasan
      \clearpage
      \section{Pembahasan Latihan}
      \label{det:pembahasan}

      \paragraph{Latihan~\ref{det:lat-tigadeterminan}.}
      Untuk determinan kedua, kurangkan baris pertama dari baris kedua
      dan ketiga:
      \[
            \det\begin{pmatrix}
                  1+x & 2+x & 3+x \\7&7&1\\6&4&2
            \end{pmatrix}
            =10(1+x)-8(2+x)-14(3+x)=-12(x+4).
      \]
      Kasus $x=0$ langsung memberi determinan pertama, yaitu $-48$.
      Untuk determinan ketiga, ekspansi enam suku memberi
      \begin{align*}
            D(x) & =x^{15}+x^{13}+x^{17}-x^{19}-x^{15}-x^{11} \\
                 & =-x^{11}(1-x^2)(1-x^6).
      \end{align*}
      Seluruh langkah merupakan identitas polinom, sehingga hasilnya juga
      berlaku untuk $x=0$ tanpa pemeriksaan pembagian tersendiri.

      \paragraph{Latihan~\ref{det:lat-divisibilitas}.}
      Untuk setiap $i\ge2$, lakukan $R_i\leftarrow R_i-R_1$.
      Setiap entri pada $n-1$ baris tersebut menjadi $0$, $2$, atau $-2$.
      Keluarkan faktor $2$ dari masing-masing baris. Maka
      $\det A=2^{n-1}\det B$ untuk suatu matriks bilangan bulat $B$.
      Karena $\det B\in\Z$, keterbagian yang diminta terbukti.

      \paragraph{Latihan~\ref{det:lat-paritas}.}
      (a) Transpose dan perkalian skalar memberi
      \[
            \det A=\det(A^T)=\det(-A)=(-1)^n\det A.
      \]
      Jika $n$ ganjil, $2\det A=0$, sehingga determinannya nol.

      (b) Persamaan $A^2=-I_n$ memberi $(\det A)^2=(-1)^n$.
      Karena $A$ real, ruas kiri tidak negatif, sehingga $n$ tidak mungkin
      ganjil. (c) Untuk matriks kompleks kesimpulan itu gagal: matriks
      $1\times1$ dengan entri $i$ sudah memenuhi $A^2+I_1=0$.

      \paragraph{Latihan~\ref{det:lat-ortogonal}.}
      Dari $AA^T=I$ diperoleh $(\det A)^2=1$, sehingga $\det A=\pm1$,
      meskipun entri $A$ boleh kompleks. Syarat negatif memilih $\det A=-1$.
      Selain itu, $A^{-1}=A^T$, maka
      \[
            \det(A+I)=\det A\det(I+A^{-1})
            =-\det(I+A^T)=-\det(I+A).
      \]
      Jadi $\det(A+I)=0$.

      \paragraph{Latihan~\ref{det:lat-blok-syarat}.}
      Untuk (a), ambil
      \[
            A=\begin{pmatrix}1&0\\0&2\end{pmatrix},\quad
            C=\begin{pmatrix}3&0\\0&4\end{pmatrix},\quad
            B=\begin{pmatrix}0&1\\1&0\end{pmatrix},\quad
            D=\begin{pmatrix}1&1\\0&1\end{pmatrix}.
      \]
      Jelas $AC=CA$. Tetapi
      \[
            AD-CB=\begin{pmatrix}1&-2\\-4&2\end{pmatrix},\qquad
            AD-BC=\begin{pmatrix}1&-3\\-3&2\end{pmatrix}.
      \]
      Determinan pertama $-6$, sama dengan $\det M$ menurut
      \autoref{det:blokkomutatif}, sedangkan determinan kedua $-7$.

      Untuk (b), pilih
      \[
            A=\begin{pmatrix}1&1\\0&1\end{pmatrix},\quad
            B=\begin{pmatrix}0&1\\1&0\end{pmatrix},\quad
            C=\begin{pmatrix}1&0\\0&2\end{pmatrix},\quad D=I_2.
      \]
      Maka $AC\ne CA$. Komplemen Schur memberi
      \[
            D-CA^{-1}B=\begin{pmatrix}2&-1\\-2&1\end{pmatrix},
            \qquad \det M=1\cdot0=0.
      \]
      Namun $AD-CB=\left(\begin{smallmatrix}1&0\\-2&1\end{smallmatrix}\right)$
      berdeterminan $1$. Jadi hipotesis komutativitas tidak dapat dihapus.

      \paragraph{Latihan~\ref{det:lat-dependensi}.}
      Tuliskan $A=\left(\begin{smallmatrix}a&b\\c&d\end{smallmatrix}\right)$
      dan $X=\left(\begin{smallmatrix}x&y\\z&w\end{smallmatrix}\right)$.
      Ekspansi langsung memberikan
      \[
            \det(A+X)-\det A-\det X=dx-cy-bz+aw.
      \]
      Jadi keempat vektor koordinat $\operatorname{vec}(B_i)$, dengan urutan
      koordinat $(x,y,z,w)^T$, memenuhi persamaan linear homogen yang sama:
      \[
            (d,-c,-b,a)\operatorname{vec}(B_i)=0.
      \]
      Baris $(d,-c,-b,a)$ tidak nol karena $A\ne0$.
      Jika $K$ adalah matriks $4\times4$ dengan kolom-kolom
      $\operatorname{vec}(B_i)$, maka $(d,-c,-b,a)K=0$ memaksa $K$ singular.
      Kolom-kolom $K$ bergantung linear; hubungan linear yang sama berlaku
      untuk $B_1,B_2,B_3,B_4$.

      \paragraph{Latihan~\ref{det:lat-permutasiblok}.}
      Pada determinan pertama, tukarkan dua kelompok kolom berukuran dua.
      Diperlukan $2\cdot2=4$ pertukaran, sehingga tandanya positif.
      Matriks menjadi blok diagonal dan determinannya
      \[
            (a_1b_2-a_2b_1)(a_3b_4-a_4b_3).
      \]
      Pada determinan kedua, urutkan baris \emph{dan} kolom menjadi
      $(1,4,2,3)$. Kedua faktor tanda saling meniadakan, dan terbentuk dua
      blok diagonal
      $\left(\begin{smallmatrix}a_1&b_1\\b_4&a_4\end{smallmatrix}\right)$
      serta $\left(\begin{smallmatrix}a_2&b_2\\b_3&a_3\end{smallmatrix}\right)$.
      Hasilnya
      \[
            (a_1a_4-b_1b_4)(a_2a_3-b_2b_3).
      \]
\fi

% !TeX root = ../../ONMIPA Alin.tex
% Sumber utama: Zhang (2009), Bab 3; Zhang (2011), bagian 1.3, 2.4, 3.3.
\chapter{Nilai Eigen dan Vektor Eigen}
\label{ch:eigen}

Bab ini menghubungkan akar polinom karakteristik dengan ruang eigen,
rank, dan nulitas. Pembahasan berlanjut ke keserupaan, polinom minimal,
Cayley--Hamilton, serta persamaan matriks. Diagonalisasi dan bentuk-bentuk
dekomposisi dikembangkan tersendiri pada \autoref{ch:dekomposisi}.

Matriks berentri dalam $\F=\R$ atau $\C$; medan yang digunakan disebutkan
secara eksplisit ketika seluruh akar suatu polinom diperlukan.

\section{Nilai Eigen dan Ruang Eigen}
\label{eigen:dasar}

\begin{definisi}[][Nilai eigen dan vektor eigen][eigen:definisi]
    Untuk $A\in\Mat_n(\F)$, skalar $\lambda\in\F$ disebut
    \emph{nilai eigen} jika terdapat $0\ne v\in\F^n$ dengan
    $Av=\lambda v$. Vektor $v$ disebut \emph{vektor eigen} yang
    bersesuaian dengan $\lambda$. Ruang eigennya adalah
    \[
        E_\lambda(A)=\ker(A-\lambda I).
    \]
    Polinom karakteristik $A$ didefinisikan oleh
    $\chi_A(t)=\det(tI-A)$.
\end{definisi}

Ruang eigen memuat vektor nol, tetapi vektor eigen selalu tak nol.
Dengan hasil pada bab rank dan determinan,
\[
    \lambda\text{ nilai eigen}
    \iff\chi_A(\lambda)=0
    \iff\rank(A-\lambda I)<n.
\]
Menghitung nilai eigen berarti menentukan akar polinom karakteristik;
menghitung vektor eigen berarti menyelesaikan sistem homogen untuk
setiap akar tersebut.

\begin{definisi}[][Multiplisitas aljabar dan geometri][eigen:multiplicitas]
    Multiplisitas aljabar $a_\lambda$ adalah banyaknya faktor
    $t-\lambda$ dalam $\chi_A(t)$. Multiplisitas geometri adalah
    \[
        g_\lambda=\dim E_\lambda(A)=\nul(A-\lambda I).
    \]
\end{definisi}

\begin{teorema}[][Batas dimensi ruang eigen][eigen:batas-geometri]
    Untuk setiap nilai eigen $\lambda$, berlaku
    $1\le g_\lambda\le a_\lambda$.
    Ruang-ruang eigen untuk nilai eigen yang berbeda berjumlah langsung.
\end{teorema}
\begin{proof}
    Pilih basis $v_1,\ldots,v_g$ untuk $E_\lambda(A)$ dan perluas
    menjadi basis $\F^n$. Dalam basis ini, matriks $A$ berbentuk
    \[
        \begin{pmatrix}\lambda I_g&C\\0&D\end{pmatrix}.
    \]
    Determinan segitiga blok memberi
    $\chi_A(t)=(t-\lambda)^g\chi_D(t)$, sehingga $g\le a_\lambda$.

    Untuk sifat jumlah langsung, ambil nilai eigen berbeda
    $\lambda_1,\ldots,\lambda_s$ dan relasi
    $x_1+\cdots+x_s=0$ dengan $x_i\in E_{\lambda_i}(A)$.
    Untuk suatu $i$, terapkan $\prod_{j\ne i}(A-\lambda_jI)$.
    Semua suku kecuali $x_i$ lenyap, sehingga
    \[
        \left(\prod_{j\ne i}(\lambda_i-\lambda_j)\right)x_i=0.
    \]
    Koefisiennya tak nol; maka setiap $x_i=0$.
\end{proof}

Khususnya, vektor-vektor eigen untuk nilai eigen berbeda bebas linear.
Namun, beberapa vektor eigen untuk \emph{nilai eigen yang sama}
masih perlu diperiksa kebebasannya.

\begin{prop}[Trace, determinan, dan rank]
    \label{eigen:trace-det}
    Jika nilai eigen $A\in\Mat_n(\C)$ adalah
    $\lambda_1,\ldots,\lambda_n$, dihitung dengan multiplisitas aljabar,
    maka
    \[
        \chi_A(t)=\prod_{i=1}^n(t-\lambda_i),\qquad
        \tr A=\sum_{i=1}^n\lambda_i,\qquad
        \det A=\prod_{i=1}^n\lambda_i.
    \]
    Banyaknya nilai eigen tak nol, dengan multiplisitas, tidak melebihi
    $\rank A$.
\end{prop}
\begin{proof}
    Faktorkan polinom karakteristik di atas $\C$.
    Koefisien $t^{n-1}$ dalam $\det(tI-A)$ adalah
    $-\sum_i a_{ii}=-\tr A$, sedangkan suku konstannya $(-1)^n\det A$.
    Membandingkan koefisien memberi dua identitas.
    Jika $0$ adalah nilai eigen dengan multiplisitas aljabar $a_0$,
    maka $\rank A=n-g_0\ge n-a_0$.
    Jika $0$ bukan nilai eigen, $A$ invertibel dan ranknya $n$.
\end{proof}

Rank dapat lebih besar daripada banyaknya nilai eigen tak nol.
Misalnya, $\begin{pmatrix}0&1\\0&0\end{pmatrix}$ mempunyai rank satu,
tetapi kedua nilai eigennya nol. Ini membedakan informasi kernel dari
informasi akar polinom karakteristik.


\begin{exam}[Nilai eigen matriks rank satu]
    Misalkan $A=uv^T$, dengan $u,v\in\F^n$ tak nol dan $n\ge2$.
    Tuliskan $c=v^Tu$. Tentukan polinom karakteristik, ruang eigen,
    dan pangkat $A$.
\end{exam}
\begin{solution}
    Karena $Ax=u(v^Tx)$, diperoleh $\rank A=1$ dan
    $\ker A=\ker v^T$ berdimensi $n-1$.
    Trace $A$ adalah $c$, sehingga
    \[
        \chi_A(t)=t^{n-1}(t-c).
    \]
    Jika $c\ne0$, ruang eigen untuk $c$ adalah $\operatorname{span}\{u\}$,
    sedangkan ruang eigen untuk $0$ adalah $\ker v^T$.
    Jika $c=0$, satu-satunya nilai eigen adalah $0$, dan ruang eigennya
    tetap berdimensi $n-1$. Vektor $u$ kini juga berada pada kernel.
    Identitas $A^2=u(v^Tu)v^T=cA$ memberi
    $A^k=c^{k-1}A$ untuk $k\ge1$, dengan faktor berpangkat nol diartikan $1$.
\end{solution}

\section{Keserupaan dan Perubahan Basis}
\label{eigen:similaritas}

\begin{definisi}[][Keserupaan matriks][eigen:similar]
    Matriks $A,B\in\Mat_n(\F)$ disebut serupa jika
    $B=P^{-1}AP$ untuk suatu matriks invertibel $P$.
\end{definisi}

Kolom-kolom $P$ menyatakan basis baru dalam koordinat lama.
Persamaan $AP=PB$ berarti kolom ke-$j$ dari $B$ menyimpan koordinat
bayangan vektor basis ke-$j$ terhadap basis baru tersebut.
Keserupaan mempertahankan polinom karakteristik, trace, determinan,
dan rank, sebab
\[
    tI-B=P^{-1}(tI-A)P.
\]
Kesamaan nilai eigen beserta multiplisitasnya belum cukup untuk
menjamin keserupaan.

\section{Polinom Matriks dan Cayley--Hamilton}
\label{eigen:polinom}

Untuk $f(t)=a_0+a_1t+\cdots+a_dt^d$, definisikan
$f(A)=a_0I+a_1A+\cdots+a_dA^d$.
Semua polinom dalam satu matriks $A$ saling komutatif.
Jika $Av=\lambda v$, maka $f(A)v=f(\lambda)v$.

\begin{prop}[Nilai eigen polinom matriks]
    \label{eigen:spectral-mapping}
    Untuk $A\in\Mat_n(\C)$ dan polinom $f$, nilai eigen $f(A)$
    tepat merupakan bilangan $f(\lambda)$ dengan $\lambda$ nilai eigen $A$.
    Khususnya, $f(A)$ invertibel jika dan hanya jika
    $f(\lambda)\ne0$ untuk setiap nilai eigen $A$.
\end{prop}
\begin{proof}
    Satu arah mengikuti $f(A)v=f(\lambda)v$.
    Jika $f$ tidak konstan, faktorkan
    $f(t)-\mu=c\prod_j(t-\alpha_j)$ di atas $\C$.
    Matriks $f(A)-\mu I=c\prod_j(A-\alpha_jI)$ singular tepat ketika
    salah satu faktornya singular. Jadi $\mu=f(\alpha_j)$ untuk
    suatu nilai eigen $\alpha_j$ dari $A$.
    Untuk $f$ konstan, pernyataannya langsung.
\end{proof}

\begin{teorema}[][Cayley--Hamilton][eigen:ch]
    Setiap matriks persegi memenuhi polinom karakteristiknya:
    jika $\chi_A(t)=t^n+c_{n-1}t^{n-1}+\cdots+c_0$, maka
    \[
        A^n+c_{n-1}A^{n-1}+\cdots+c_1A+c_0I=0.
    \]
\end{teorema}
\begin{proof}
    Ekspansikan adjugat sebagai polinom matriks
    $\operatorname{adj}(tI-A)=\sum_{j=0}^{n-1}B_jt^j$.
    Identitas adjugat memberi
    \[
        (tI-A)\sum_{j=0}^{n-1}B_jt^j=\chi_A(t)I.
    \]
    Perbandingan koefisien menghasilkan
    \[
        B_{n-1}=I,\qquad B_{j-1}-AB_j=c_jI\ (1\le j<n),
        \qquad -AB_0=c_0I.
    \]
    Kalikan persamaan untuk $c_j$ dari kiri dengan $A^j$, kemudian
    jumlahkan bersama persamaan konstanta dan $A^n=A^nB_{n-1}$.
    Suku-suku yang memuat $B_j$ saling menghapus, menyisakan
    $\chi_A(A)=0$.
\end{proof}

Identitas ini diperoleh melalui koefisien polinom. Substitusi formal
$t=A$ ke dalam simbol determinan $\det(tI-A)$ bukan pembuktian:
semula entri matriks tersebut adalah polinom skalar, sedangkan
substitusi itu mengganti jenis entrinya dengan matriks.

\begin{exam}[Pangkat melalui rekurensi]
    Untuk $F=\begin{pmatrix}1&1\\1&0\end{pmatrix}$, tentukan $F^k$.
\end{exam}
\begin{solution}
    Karena $\chi_F(t)=t^2-t-1$, Cayley--Hamilton memberi $F^2=F+I$.
    Dengan bilangan Fibonacci $f_0=0$, $f_1=1$, dan
    $f_{k+1}=f_k+f_{k-1}$, induksi memberikan
    \[
        F^k=f_kF+f_{k-1}I
        =\begin{pmatrix}f_{k+1}&f_k\\f_k&f_{k-1}\end{pmatrix}
        \qquad(k\ge1).
    \]
    Jadi setiap pangkat tereduksi menjadi kombinasi $I$ dan $F$.
    Selain itu $F(F-I)=I$, sehingga $F^{-1}=F-I$.
\end{solution}

\subsection{Polinom minimal dan akar berulang}

\begin{definisi}[][Polinom minimal][eigen:minimal]
    Polinom minimal $m_A(t)$ adalah polinom monik berderajat terkecil
    yang memenuhi $m_A(A)=0$.
\end{definisi}

\begin{prop}[Sifat polinom minimal]
    \label{eigen:sifat-minimal}
    Polinom minimal ada dan unik. Jika $f(A)=0$, maka $m_A\mid f$;
    khususnya $m_A\mid\chi_A$.
    Akar-akar $m_A$ tepat nilai eigen $A$, dan matriks serupa mempunyai
    polinom minimal yang sama.
\end{prop}
\begin{proof}
    Keberadaan mengikuti Cayley--Hamilton.
    Bagi $f$ oleh $m_A$: $f=qm_A+r$ dengan $\deg r<\deg m_A$.
    Maka $r(A)=0$, sehingga minimalitas memaksa $r=0$.
    Ini sekaligus membuktikan keunikan polinom monik minimal.

    Jika $Av=\lambda v$ dengan $v\ne0$, maka
    $0=m_A(A)v=m_A(\lambda)v$, sehingga $m_A(\lambda)=0$.
    Sebaliknya, jika $m_A(t)=(t-\lambda)q(t)$ tetapi $A-\lambda I$
    invertibel, $m_A(A)=0$ memaksa $q(A)=0$, bertentangan dengan minimalitas.
    Terakhir, $f(P^{-1}AP)=P^{-1}f(A)P$ menunjukkan bahwa matriks
    serupa mempunyai polinom penganulir yang sama.
\end{proof}


Kriteria polinom minimal untuk diagonalisasi dan penggunaan proyektor
akan dibahas pada \autoref{ch:dekomposisi}.

\section{Hasil Kali dan Persamaan Matriks}
\label{eigen:hasilkali}

\begin{teorema}[][Nilai eigen $AB$ dan $BA$][eigen:abba]
    Untuk $A\in\Mat_{m\times n}(\C)$ dan $B\in\Mat_{n\times m}(\C)$,
    \[
        t^n\chi_{AB}(t)=t^m\chi_{BA}(t).
    \]
    Jadi $AB$ dan $BA$ mempunyai nilai eigen tak nol yang sama,
    termasuk multiplisitas aljabarnya.
\end{teorema}
\begin{proof}
    Untuk $t\ne0$, hitung determinan
    $\begin{pmatrix}tI_m&A\\B&I_n\end{pmatrix}$ dengan dua komplemen Schur.
    Diperoleh
    \[
        \det(tI_m-AB)=t^m\det(I_n-t^{-1}BA).
    \]
    Kalikan dengan $t^n$ untuk mendapatkan identitas yang dinyatakan.
    Kedua ruas adalah polinom yang sama pada semua $t\ne0$,
    sehingga identitas berlaku juga pada $t=0$.
\end{proof}

Jika $ABx=\lambda x$ dengan $\lambda\ne0$, maka $Bx\ne0$ dan
$BA(Bx)=\lambda Bx$. Bahkan $B$ memberi isomorfisme
$E_\lambda(AB)\to E_\lambda(BA)$ dengan invers $A/\lambda$.
Kesamaan multiplisitas aljabar pada teorema memerlukan identitas polinom,
sedangkan argumen ruang eigen ini menjelaskan multiplisitas geometrinya.

Untuk $m=n$, polinom karakteristiknya sama, tetapi matriksnya belum
tentu serupa. Ambil
\[
    A=\begin{pmatrix}1&0\\0&0\end{pmatrix},\qquad
    B=\begin{pmatrix}0&1\\0&0\end{pmatrix}.
\]
Maka $AB=B\ne0$, sedangkan $BA=0$; rank keduanya berbeda.
Jika salah satu faktor invertibel, keserupaan berlaku, misalnya
$BA=A^{-1}(AB)A$ ketika $A$ invertibel.

\begin{prop}[Persamaan Sylvester dengan spektrum terpisah]
    \label{eigen:sylvester}
    Jika $A\in\Mat_m(\C)$ dan $B\in\Mat_n(\C)$ tidak mempunyai
    nilai eigen bersama, maka untuk setiap $C\in\Mat_{m\times n}(\C)$,
    persamaan $AX-XB=C$ mempunyai tepat satu solusi.
\end{prop}
\begin{proof}
    Jika $AX=XB$, maka $A^kX=XB^k$, sehingga
    $\chi_A(A)X=X\chi_A(B)$.
    Ruas kiri nol oleh Cayley--Hamilton; matriks $\chi_A(B)$
    invertibel oleh \autoref{eigen:spectral-mapping} dan pemisahan spektrum.
    Jadi $X=0$.
    Pemetaan linear $X\mapsto AX-XB$ pada ruang berdimensi $mn$
    dengan demikian injektif, sehingga bijektif.
\end{proof}

\section{Latihan Pilihan}
\label{eigen:latihan}

Latihan berikut menghubungkan ruang eigen, polinom, dan persamaan matriks.

\begin{latihan}\label{eigen:lat-polinom}
    Misalkan $A\in\Mat_m(\C)$ dan $B\in\Mat_n(\C)$.
    Buktikan bahwa $\chi_B(A)$ invertibel jika dan hanya jika
    $A$ dan $B$ tidak mempunyai nilai eigen bersama.
\end{latihan}

\begin{latihan}\label{eigen:lat-dua-rankone}
    Misalkan $a=(a_1,\ldots,a_n)^T\in\R^n$, $n\ge2$,
    $\sum_i a_i=0$, dan $e=(1,\ldots,1)^T$.
    Tentukan semua nilai eigen dan ruang eigennya untuk $M=aa^T+ee^T$.
    Tangani juga kasus $a=0$ dan $\|a\|^2=n$.
\end{latihan}

\begin{latihan}\label{eigen:lat-persamaan}
    Misalkan
    \[
        A=\begin{pmatrix}1&1\\0&2\end{pmatrix},\qquad
        B=\begin{pmatrix}2&0\\0&3\end{pmatrix},\qquad
        C=(c_{ij})\in\Mat_2(\C).
    \]
    Tentukan syarat konsistensi dan seluruh solusi $AX-XB=C$.
    Hubungkan parameter bebasnya dengan nilai eigen bersama $A$ dan $B$.
\end{latihan}

\begin{latihan}\label{eigen:lat-jumlah-vektor}
    Misalkan $u,v$ vektor eigen $A$ untuk nilai eigen yang berbeda.
    Buktikan bahwa $u+v$ bukan vektor eigen $A$.
\end{latihan}

\begin{latihan}\label{eigen:lat-trace-komutator}
    Untuk $A,B\in\Mat_n(\C)$, buktikan $\tr(AB)=\tr(BA)$.
    Simpulkan bahwa $AB-BA=I$ tidak mungkin terjadi.
    Selanjutnya, jika $AB-BA=A$, buktikan $A$ singular.
\end{latihan}

\begin{latihan}\label{eigen:lat-adjugat}
    Misalkan $A\in\Mat_n(\C)$ singular, $n\ge2$, dengan
    $\chi_A(t)=t^n+c_{n-1}t^{n-1}+\cdots+c_1t$.
    Buktikan
    \[
        \chi_{\operatorname{adj}A}(t)
        =t^{n-1}\bigl(t-(-1)^{n-1}c_1\bigr).
    \]
\end{latihan}

\ifpembahasan
\section{Pembahasan Latihan}
\label{eigen:pembahasan}

\subsection*{Latihan \ref{eigen:lat-polinom}: akar karakteristik sebagai penghalang invers}
Dengan \autoref{eigen:spectral-mapping}, $\chi_B(A)$ invertibel
tepat ketika $\chi_B(\lambda)\ne0$ untuk semua nilai eigen $\lambda$ dari $A$.
Nol-nol $\chi_B$ tepat merupakan nilai eigen $B$.
Jadi syarat tersebut ekuivalen dengan tidak adanya nilai eigen bersama.
Ukuran $A$ dan $B$ tidak harus sama karena $\chi_B$ dievaluasi sebagai
polinom skalar pada matriks $A$.

\subsection*{Latihan \ref{eigen:lat-dua-rankone}: dua arah yang ortogonal}
Karena $a^Te=0$, diperoleh $Me=ne$ dan $Ma=\|a\|^2a$.
Untuk $x$ ortogonal terhadap $a$ dan $e$, berlaku $Mx=0$.
Jika $a\ne0$ dan $\|a\|^2\ne n$, ruang eigennya adalah
\[
    E_n=\operatorname{span}\{e\},\quad
    E_{\|a\|^2}=\operatorname{span}\{a\},\quad
    E_0=\operatorname{span}\{a,e\}^{\perp}.
\]
Dimensinya $1,1,n-2$; bila $n=2$, ruang terakhir nol dan $0$ bukan nilai eigen.
Jika $a\ne0$ dan $\|a\|^2=n$, kedua ruang pertama bergabung menjadi
$E_n=\operatorname{span}\{a,e\}$, berdimensi dua.
Jika $a=0$, $M=ee^T$ mempunyai nilai eigen $n$ pada
$\operatorname{span}\{e\}$ dan $0$ pada $e^\perp$.

\subsection*{Latihan \ref{eigen:lat-persamaan}: nilai eigen bersama menyisakan parameter}
Tuliskan $X=\begin{pmatrix}x&y\\z&w\end{pmatrix}$.
Perhitungan memberikan
\[
    AX-XB=\begin{pmatrix}-x+z&-2y+w\\0&-w\end{pmatrix}.
\]
Jadi konsistensi tepat ketika $c_{21}=0$. Jika konsisten, semua solusi adalah
\[
    X=\begin{pmatrix}
        t-c_{11}&-(c_{12}+c_{22})/2\\t&-c_{22}
    \end{pmatrix},\qquad t\in\C.
\]
Kedua matriks mempunyai satu nilai eigen bersama, yaitu $2$, dengan
ruang eigen masing-masing berdimensi satu. Satu pasangan arah ini
memberikan kernel berdimensi satu pada pemetaan $X\mapsto AX-XB$.

\subsection*{Latihan \ref{eigen:lat-jumlah-vektor}: kebebasan dua vektor eigen}
Tuliskan $Au=\lambda u$ dan $Av=\mu v$, dengan $\lambda\ne\mu$.
Kedua vektor bebas linear. Jika $A(u+v)=\nu(u+v)$, maka
$(\lambda-\nu)u+(\mu-\nu)v=0$, sehingga
$\lambda=\nu=\mu$, suatu kontradiksi.

\subsection*{Latihan \ref{eigen:lat-trace-komutator}: trace menghalangi invertibilitas}
Ekspansi entri memberi
$\tr(AB)=\sum_{i,j}a_{ij}b_{ji}=\tr(BA)$.
Jadi trace suatu komutator $AB-BA$ selalu nol,
sedangkan $\tr I=n\ne0$.
Sekarang anggap $AB-BA=A$ dan $A$ invertibel.
Mengalikan dari kiri dengan $A^{-1}$ memberi
$B-A^{-1}BA=I$. Trace ruas kiri nol karena
$\tr(A^{-1}BA)=\tr B$, tetapi trace ruas kanan $n$.
Kontradiksi ini membuktikan $A$ singular.

\subsection*{Latihan \ref{eigen:lat-adjugat}: satu nilai eigen yang mungkin tak nol}
Jika $\rank A\le n-2$, semua minor berukuran $n-1$ nol,
sehingga $\operatorname{adj}A=0$.
Jika $\rank A=n-1$, identitas $A\operatorname{adj}A=0$
menempatkan semua kolom adjugat pada kernel $A$ yang berdimensi satu.
Jadi dalam kedua kasus $\rank(\operatorname{adj}A)\le1$.
Matriks berank paling besar satu mempunyai polinom karakteristik
$t^{n-1}(t-\tr M)$: multiplisitas aljabar nol minimal $n-1$,
dan trace menentukan akar yang tersisa.
Koefisien $t$ dalam $\det(tI-A)$, melalui ekspansi minor utama,
adalah $(-1)^{n-1}$ kali jumlah minor utama berukuran $n-1$.
Jumlah itu tepat $\tr(\operatorname{adj}A)$.
Dengan demikian $\tr(\operatorname{adj}A)=(-1)^{n-1}c_1$,
yang menghasilkan rumus yang diminta.
\fi

% !TeX root = ../../ONMIPA Alin.tex
% Sumber utama: Zhang (2009), Bab 3; Zhang (2011), bagian 3.1--3.4.
\chapter{Dekomposisi Matriks}
\label{ch:dekomposisi}

Bab sebelumnya menyediakan nilai eigen, ruang eigen, dan polinom minimal.
Sekarang kita memilih basis yang menyederhanakan matriks: mula-mula
basis vektor eigen, lalu basis rantai Jordan, dan akhirnya basis
ortonormal pada dekomposisi Schur serta spektral.

\section{Diagonalisasi}
\label{eigen:diagonalisasi}

\begin{definisi}[][Matriks yang dapat didiagonalkan][dekomp:def-diagonal]
    Matriks $A\in\Mat_n(\F)$ dapat didiagonalkan di atas $\F$ jika
    $A=PDP^{-1}$ dengan $D$ diagonal dan $P$ invertibel, keduanya
    berentri dalam $\F$.
\end{definisi}

\begin{teorema}[][Kriteria diagonalisasi][eigen:kriteria-diagonal]
    Untuk $A\in\Mat_n(\F)$, pernyataan berikut ekuivalen:
    \begin{enumerate}[label=(\roman*),nosep]
        \item $A$ dapat didiagonalkan di atas $\F$;
        \item terdapat basis $\F^n$ yang seluruh anggotanya vektor eigen $A$;
        \item $\chi_A$ terurai menjadi faktor linear di atas $\F$ dan
            $g_\lambda=a_\lambda$ untuk setiap nilai eigen $\lambda$.
    \end{enumerate}
\end{teorema}
\begin{proof}
    Jika $P=[v_1\ \cdots\ v_n]$ dan
    $D=\diag(\lambda_1,\ldots,\lambda_n)$, persamaan $AP=PD$
    ekuivalen dengan $Av_i=\lambda_iv_i$ untuk setiap $i$.
    Ini membuktikan (i)--(ii).
    Ruang-ruang eigen berbeda berjumlah langsung dan dimensinya
    tidak melebihi multiplisitas aljabarnya.
    Karena $\sum_\lambda a_\lambda=n$ ketika polinomnya terurai,
    basis vektor eigen ada tepat ketika setiap batas tersebut tercapai.
\end{proof}

Adanya $n$ nilai eigen berbeda cukup untuk diagonalisasi, tetapi
tidak perlu: matriks identitas mempunyai satu nilai eigen dan tetap
diagonal. Sebaliknya, akar berulang mengharuskan pemeriksaan dimensi
ruang eigen; akar berulang sendiri tidak memutuskan jawabannya.

\begin{exam}[Parameter pada ruang eigen berulang]
    Tentukan kapan
    \[
        A_t=\begin{pmatrix}1&t&0\\0&1&0\\0&0&2\end{pmatrix}
    \]
    dapat didiagonalkan.
\end{exam}
\begin{solution}
    Untuk setiap $t$, $\chi_{A_t}(x)=(x-1)^2(x-2)$.
    Persamaan $(A_t-I)v=0$ memberi $tv_2=0$ dan $v_3=0$.
    Jadi $g_1=2$ jika $t=0$, tetapi $g_1=1$ jika $t\ne0$.
    Karena $g_2=1$, matriks dapat didiagonalkan tepat ketika $t=0$.
    Contoh ini juga menunjukkan bahwa polinom karakteristik yang sama
    dapat dimiliki oleh matriks yang tidak serupa.
\end{solution}

\begin{exam}[Pola diagonal dan luar diagonal konstan]
    Misalkan $A$ berukuran $n\ge2$, dengan entri diagonal $a$ dan
    semua entri luar diagonal $b$. Tentukan ruang eigen dan pangkat $A$.
\end{exam}
\begin{solution}
    Dengan $e=(1,\ldots,1)^T$ dan $J_n=ee^T$,
    \[
        A=(a-b)I+bJ_n,\qquad J_ne=ne.
    \]
    Pada ruang $H=\{x:e^Tx=0\}$, berlaku $J_nx=0$.
    Karena $\F^n=\operatorname{span}\{e\}\oplus H$, matriks $A$
    bekerja sebagai perkalian dengan $a+(n-1)b$ pada ruang pertama
    dan $a-b$ pada ruang kedua. Jika $b\ne0$, kedua nilai ini berbeda
    dan multiplisitasnya $1$ serta $n-1$. Jika $b=0$, $A=aI$ dan
    ruang eigen untuk $a$ adalah seluruh $\F^n$.

    Definisikan $P=J_n/n$ dan $Q=I-P$.
    Karena $P^2=P$, $Q^2=Q$, dan $PQ=QP=0$,
    \[
        A^k=\bigl(a+(n-1)b\bigr)^kP+(a-b)^kQ\qquad(k\ge1).
    \]
    Untuk $k=0$, $A^0=P+Q=I$.
    Jika kedua skalar tak nol, invers diperoleh dengan mengganti
    pangkat $k$ oleh $-1$.
\end{solution}


\begin{exam}[Menyusun matriks perubahan basis]
    Diagonalkan
    \[
        A=\begin{pmatrix}1&2&2\\2&1&2\\2&2&1\end{pmatrix}.
    \]
\end{exam}
\begin{solution}
    Karena $A=-I+2J_3$, nilai eigennya $5,-1,-1$.
    Pilih $v_1=(1,1,1)^T$, $v_2=(1,-1,0)^T$, dan $v_3=(1,1,-2)^T$.
    Ketiganya bebas dan memenuhi
    \[
        P=\begin{pmatrix}1&1&1\\1&-1&1\\1&0&-2\end{pmatrix},\qquad
        AP=P\diag(5,-1,-1).
    \]
    Jadi $P^{-1}AP=\diag(5,-1,-1)$.
    Jika tiap kolom dinormalisasi, diperoleh matriks ortogonal dan
    diagonalisasi yang sekaligus mempertahankan hasil kali dalam.
\end{solution}

Medan skalar tetap penting. Matriks real
$R=\begin{pmatrix}0&-1\\1&0\end{pmatrix}$ mempunyai polinom
$t^2+1$, sehingga tidak mempunyai vektor eigen real.
Di atas $\C$, nilai eigennya $i,-i$ berbeda dan $R$ dapat didiagonalkan.

\section{Polinom Minimal dan Proyektor}
\label{dekomp:polinom}

Polinom minimal dari \autoref{eigen:minimal} memberikan kriteria
untuk menentukan apakah suatu basis vektor eigen lengkap tersedia.

\begin{teorema}[][Kriteria polinom untuk diagonalisasi][eigen:squarefree]
    Matriks $A$ dapat didiagonalkan di atas $\F$ jika dan hanya jika
    $m_A$ merupakan hasil kali faktor-faktor linear yang berbeda di atas $\F$.
    Lebih umum, cukup ada polinom penganulir $q(A)=0$ yang merupakan
    hasil kali faktor linear berbeda.
\end{teorema}
\begin{proof}
    Jika $A$ diagonal dalam suatu basis, hasil kali
    $\prod_\lambda(t-\lambda)$ atas nilai eigen berbeda menganulir $A$.
    Polinom minimal membagi hasil kali itu dan memuat semua akarnya.

    Sebaliknya, misalkan $q(t)=\prod_{i=1}^s(t-\lambda_i)$ dengan
    $\lambda_i$ berbeda dan $q(A)=0$.
    Bentuk polinom Lagrange
    \[
        \ell_i(t)=\prod_{j\ne i}\frac{t-\lambda_j}{\lambda_i-\lambda_j}.
    \]
    Karena $\sum_i\ell_i(t)=1$ dan $(t-\lambda_i)\ell_i(t)$
    adalah kelipatan $q(t)$, setiap $x\in\F^n$ dapat ditulis
    \[
        x=\sum_i\ell_i(A)x,\qquad
        (A-\lambda_iI)\ell_i(A)x=0.
    \]
    Jadi ruang-ruang eigen merentang seluruh ruang dan berjumlah langsung.
    Gabungan basisnya merupakan basis vektor eigen.
\end{proof}

Akibatnya, $A^2=A$ dan $A^3=A$ menjamin diagonalisasi di atas $\R$
atau $\C$, sebab polinom $t(t-1)$ dan $t(t-1)(t+1)$ mempunyai akar
berbeda. Jika $A^k=I$, $A$ dapat didiagonalkan di atas $\C$, karena
$t^k-1$ tidak mempunyai akar berulang.
Sebaliknya, $A^2=0$ belum menjamin diagonalisasi; matriks nilpoten
yang dapat didiagonalkan harus merupakan matriks nol.

\subsection{Proyektor dan reduksi pangkat}

Jika nilai eigen berbeda dari matriks yang dapat didiagonalkan adalah
$\lambda_1,\ldots,\lambda_s$, definisikan $P_i=\ell_i(A)$.
Dengan memeriksa aksinya pada setiap ruang eigen, diperoleh
\[
    P_i^2=P_i,\qquad P_iP_j=0\ (i\ne j),\qquad
    \sum_iP_i=I,\qquad A=\sum_i\lambda_iP_i.
\]
Maka $f(A)=\sum_i f(\lambda_i)P_i$ untuk setiap polinom $f$.
Proyektor ini belum tentu ortogonal; ortogonalitas akan muncul pada
dekomposisi spektral matriks normal.

\begin{exam}[Identitas kuadratik menentukan semua pangkat]
    Misalkan $A^2-3A+2I=0$. Tentukan $A^k$ dan $A^{-1}$.
\end{exam}
\begin{solution}
    Polinom $(t-1)(t-2)$ mempunyai dua akar berbeda.
    Proyektor yang bersesuaian adalah $P_1=2I-A$ dan $P_2=A-I$;
    salah satunya boleh nol jika hanya satu nilai eigen yang muncul.
    Untuk setiap $k\ge0$,
    \[
        A^k=P_1+2^kP_2=(2-2^k)I+(2^k-1)A.
    \]
    Dari $A(3I-A)=2I$, diperoleh $A^{-1}=(3I-A)/2$.
    Tidak diperlukan perhitungan entri atau vektor eigen satu per satu.
\end{solution}

\section{Diagonalisasi Simultan}
\label{eigen:komutatif}

\begin{prop}[Ruang eigen yang invarian]
    Jika $AB=BA$, maka setiap ruang eigen $E_\lambda(A)$ invarian
    terhadap $B$: $B(E_\lambda(A))\subseteq E_\lambda(A)$.
\end{prop}
\begin{proof}
    Untuk $Ax=\lambda x$, berlaku $A(Bx)=B(Ax)=\lambda Bx$.
    Vektor $Bx$ boleh nol, sehingga pernyataan yang tepat memakai
    ruang eigen, bukan selalu vektor eigen tak nol.
\end{proof}

\begin{teorema}[][Diagonalisasi simultan][eigen:simultan]
    Dua matriks yang masing-masing dapat didiagonalkan di atas $\F$
    dapat didiagonalkan oleh satu perubahan basis yang sama jika dan
    hanya jika keduanya komutatif.
\end{teorema}
\begin{proof}
    Jika keduanya diagonal dalam satu basis, jelas keduanya komutatif.
    Sebaliknya, misalkan $AB=BA$.
    Dekomposisikan $\F^n$ menjadi jumlah langsung ruang eigen $A$.
    Pada setiap ruang ini, pembatasan $B$ tetap dianulir oleh polinom
    minimal $B$, yang mempunyai faktor linear berbeda.
    Maka pembatasan itu dapat didiagonalkan.
    Pilih basis vektor eigen $B$ di setiap ruang eigen $A$, lalu gabungkan;
    setiap anggotanya merupakan vektor eigen untuk kedua matriks.
\end{proof}

\begin{teorema}[][Komutativitas dan spektrum sederhana][eigen:centralizer]
    Jika $A$ mempunyai $n$ nilai eigen berbeda dan $AB=BA$, maka
    terdapat tepat satu polinom $p$ berderajat kurang dari $n$ sehingga
    $B=p(A)$. Khususnya, $B$ juga dapat didiagonalkan.
\end{teorema}
\begin{proof}
    Setiap ruang eigen $A$ berdimensi satu, sehingga $Bv_i=\mu_iv_i$
    pada basis vektor eigen $A$.
    Interpolasi Lagrange memberikan satu polinom berderajat kurang
    dari $n$ dengan $p(\lambda_i)=\mu_i$.
    Maka $p(A)$ dan $B$ mempunyai aksi sama pada suatu basis.
    Keunikannya mengikuti fakta bahwa polinom berderajat kurang dari
    $n$ tidak dapat mempunyai $n$ akar berbeda kecuali polinom nol.
\end{proof}

\section{Bentuk Jordan dan Rantai Vektor Eigen}
\label{eigen:jordan}

Ketika jumlah dimensi ruang eigen kurang dari $n$, vektor eigen biasa
belum membentuk basis. Penggantinya adalah rantai vektor yang
berturut-turut dipetakan ke anggota rantai sebelumnya.

\begin{definisi}[][Blok dan rantai Jordan][eigen:def-jordan]
    Blok Jordan berukuran $s$ untuk $\lambda$ adalah
    \[
        J_s(\lambda)=\lambda I_s+N_s
        =\begin{pmatrix}
            \lambda&1&&0\\&\lambda&\ddots&\\&&\ddots&1\\0&&&\lambda
        \end{pmatrix},
    \]
    dengan $N_s$ mempunyai entri $1$ pada superdiagonal dan nol
    di tempat lain. Untuk $s=1$, bloknya hanya $(\lambda)$.
    Rantai Jordan $v_1,\ldots,v_s$ memenuhi
    \[
        (A-\lambda I)v_1=0,\qquad
        (A-\lambda I)v_j=v_{j-1}\quad(2\le j\le s),\qquad v_1\ne0.
    \]
\end{definisi}

Semua anggota rantai bebas linear. Jika $\sum_jc_jv_j=0$, terapkan
$(A-\lambda I)^{s-1}$ untuk memperoleh $c_sv_1=0$;
lanjutkan dengan pangkat yang lebih rendah untuk mendapatkan semua
koefisien nol. Vektor $v_j$ disebut vektor eigen tergeneralisasi karena
$(A-\lambda I)^jv_j=0$.

\begin{teorema}[][Dekomposisi Jordan][eigen:teorema-jordan]
    Setiap $A\in\Mat_n(\C)$ dapat ditulis
    \[
        A=PJP^{-1},\qquad
        J=\diag\bigl(J_{s_1}(\lambda_1),\ldots,J_{s_r}(\lambda_r)\bigr),
    \]
    dengan $P$ invertibel. Nilai $\lambda_i$ boleh berulang.
    Ukuran blok untuk setiap nilai eigen ditentukan secara unik,
    hingga urutan blok.
\end{teorema}

Secara struktural, dekomposisi ini mula-mula memisahkan ruang menjadi
ruang-ruang eigen tergeneralisasi $\ker(A-\lambda I)^n$.
Pemisahan tersebut mengikuti identitas Bezout untuk faktor-faktor
polinom karakteristik yang saling prima.
Pada ruang untuk $\lambda$, pemetaan $A-\lambda I$ nilpoten dan
dapat disusun dalam basis rantai. Setiap rantai menghasilkan satu
blok Jordan. Pembuktian lengkap keberadaan basis rantai
tidak dirinci di sini.
Untuk perhitungan, ukuran rantainya dapat dibaca dari nulitas,
seperti dibuktikan berikut.

\begin{teorema}[][Membaca blok Jordan dari nulitas][eigen:jordan-nulitas]
    Tetapkan satu nilai eigen $\lambda$ dan tuliskan
    $d_k=\nul(A-\lambda I)^k$, dengan $d_0=0$.
    Jika ukuran blok Jordan untuk $\lambda$ adalah $s_1,\ldots,s_r$, maka
    \[
        d_k=\sum_{j=1}^r\min\{k,s_j\}.
    \]
    Akibatnya,
    \begin{align*}
        \#\{\text{blok berukuran setidaknya }k\}&=d_k-d_{k-1},\\
        \#\{\text{blok berukuran tepat }k\}&=2d_k-d_{k-1}-d_{k+1}.
    \end{align*}
\end{teorema}
\begin{proof}
    Pada satu blok untuk $\lambda$, matriks $A-\lambda I$
    berubah menjadi $N_s$ dan $\nul N_s^k=\min\{k,s\}$.
    Pada blok untuk nilai eigen $\mu\ne\lambda$, matriks
    $J_s(\mu)-\lambda I$ invertibel, sehingga tidak menyumbang kernel.
    Nulitas matriks diagonal blok adalah jumlah nulitas blok-bloknya.
    Satu blok menyumbang $1$ pada $d_k-d_{k-1}$ tepat ketika $s\ge k$;
    selisih dua hitungan berurutan memberi jumlah blok berukuran tepat $k$.
\end{proof}

Dengan demikian, $g_\lambda=d_1$ adalah banyaknya blok untuk $\lambda$,
$a_\lambda=\lim_k d_k$ adalah jumlah ukurannya, dan ukuran blok
terbesar adalah indeks pertama $k$ dengan $d_k=d_{k+1}$.
Ini memberi makna struktural pada stabilisasi kernel yang dibahas
dalam \autoref{rank:deret-pangkat}.

\begin{prop}[Polinom minimal dan blok terbesar]
    \label{eigen:minimal-jordan}
    Jika ukuran blok terbesar untuk nilai eigen $\lambda$ adalah
    $h_\lambda$, maka
    \[
        m_A(t)=\prod_\lambda(t-\lambda)^{h_\lambda}.
    \]
\end{prop}
\begin{proof}
    Satu blok $J_s(\lambda)$ dianulir oleh $(t-\lambda)^s$
    dan tidak oleh pangkat yang lebih kecil karena $N_s^{s-1}\ne0$.
    Faktor polinom yang tidak nol di $\lambda$ memberikan matriks
    segitiga dengan diagonal tak nol, sehingga invertibel pada blok itu.
    Suatu polinom menganulir seluruh matriks tepat ketika menganulir
    setiap blok, sehingga pangkat yang diperlukan untuk tiap faktor
    adalah ukuran blok terbesarnya.
\end{proof}

\begin{exam}[Membangun perubahan basis dari rantai]
    Tentukan bentuk Jordan dan matriks perubahan basis untuk
    \[
        A=\begin{pmatrix}3&-1&0\\1&1&0\\-1&1&3\end{pmatrix}.
    \]
\end{exam}
\begin{solution}
    Polinom karakteristiknya $(t-2)^2(t-3)$.
    Ambil $v_1=(1,1,0)^T$, $v_2=(1,0,1)^T$, dan $w=(0,0,1)^T$.
    Perkalian langsung memberikan
    \[
        (A-2I)v_1=0,\qquad (A-2I)v_2=v_1,\qquad Aw=3w.
    \]
    Kolom-kolom tersebut bebas, sehingga
    \[
        P=\begin{pmatrix}1&1&0\\1&0&0\\0&1&1\end{pmatrix},\qquad
        P^{-1}AP=\begin{pmatrix}2&1&0\\0&2&0\\0&0&3\end{pmatrix}.
    \]
    Urutan $v_1,v_2$ membuat entri $1$ berada di superdiagonal.
    Matriks ini tidak dapat didiagonalkan karena ruang eigen untuk $2$
    hanya berdimensi satu.
\end{solution}

\begin{exam}[Ukuran blok dari beberapa nulitas]
    Diketahui $A\in\Mat_8(\C)$ memenuhi
    \[
        \chi_A(t)=(t-2)^6(t+1)^2,
    \]
    serta
    \[
        \nul(A-2I)=3,\quad\nul(A-2I)^2=5,\quad
        \nul(A-2I)^3=6,\quad\nul(A+I)=1.
    \]
    Tentukan bentuk Jordan dan polinom minimalnya.
\end{exam}
\begin{solution}
    Untuk $\lambda=2$, barisan nulitas dimulai dengan $0,3,5,6,6$.
    Selisihnya $3,2,1,0$, sehingga terdapat tiga blok, dua berukuran
    setidaknya dua, dan satu berukuran setidaknya tiga.
    Ukurannya tepat $3,2,1$.
    Untuk $\lambda=-1$, multiplisitas aljabarnya dua dan jumlah bloknya
    satu, sehingga bloknya berukuran dua. Maka
    \[
        J=\diag\bigl(J_3(2),J_2(2),(2),J_2(-1)\bigr),\qquad
        m_A(t)=(t-2)^3(t+1)^2.
    \]
\end{solution}

\subsection{Pangkat dan polinom pada satu blok}

Karena $\lambda I$ dan $N_s$ komutatif serta $N_s^s=0$,
\[
    J_s(\lambda)^k
    =\sum_{j=0}^{\min\{k,s-1\}}\binom{k}{j}\lambda^{k-j}N_s^j
    \qquad(k\ge0),
\]
dengan faktor berpangkat nol diartikan $1$.
Untuk polinom $f$, ekspansi Taylor yang berhingga memberi
\[
    f(J_s(\lambda))
    =\sum_{j=0}^{s-1}\frac{f^{(j)}(\lambda)}{j!}N_s^j.
\]
Jadi akar berulang pada polinom minimal membuat turunan $f$ ikut
menentukan $f(A)$, selain nilai $f$ pada nilai eigen.

\begin{prop}[Nulitas polinom matriks]
    \label{eigen:nulitas-polinom}
    Misalkan $f$ polinom tak nol dan $q_\lambda$ adalah orde akarnya
    di $\lambda$, dengan $q_\lambda=0$ jika $f(\lambda)\ne0$.
    Setiap blok $J_s(\lambda)$ menyumbang
    $\min\{q_\lambda,s\}$ pada nulitas $f(A)$.
\end{prop}
\begin{proof}
    Tuliskan $f(t)=(t-\lambda)^{q_\lambda}g(t)$ dengan
    $g(\lambda)\ne0$. Pada blok tersebut,
    $f(J_s(\lambda))=N_s^{q_\lambda}g(J_s(\lambda))$.
    Faktor kedua invertibel karena segitiga dengan diagonal
    $g(\lambda)$, sehingga nulitasnya sama dengan nulitas
    $N_s^{q_\lambda}$.
\end{proof}

Pada contoh ukuran delapan di atas, ambil
$f(t)=(t-2)^2(t+1)$. Blok-blok untuk $2$ menyumbang nulitas
$2+2+1=5$, dan blok untuk $-1$ menyumbang $1$.
Jadi $\nul f(A)=6$ dan $\rank f(A)=2$ tanpa perlu mengetahui entri $A$.

\section{Dekomposisi Schur dan Spektral}
\label{eigen:spektral}

Gunakan hasil kali dalam standar $\langle x,y\rangle=x^*y$ pada
$\C^n$, dengan $A^*=\overline A^{\,T}$ dan $\|x\|^2=x^*x$.
Matriks $U$ disebut \emph{uniter} jika $U^*U=UU^*=I$;
kolom-kolomnya merupakan basis ortonormal.
Versi realnya disebut matriks ortogonal, dengan $U^TU=I$.

\begin{teorema}[][Dekomposisi Schur][eigen:schur]
    Untuk setiap $A\in\Mat_n(\C)$, terdapat matriks uniter $U$
    sehingga $U^*AU=T$ segitiga atas. Entri diagonal $T$ adalah
    nilai eigen $A$, dihitung dengan multiplisitas aljabar.
\end{teorema}
\begin{proof}
    Pilih satu vektor eigen satuan $u_1$, lalu perluas menjadi basis
    ortonormal, misalnya dengan Gram--Schmidt. Jika $U_1$ mempunyai
    basis tersebut sebagai kolom-kolomnya, maka
    \[
        U_1^*AU_1=\begin{pmatrix}\lambda_1&w\\0&B\end{pmatrix}.
    \]
    Dengan induksi pada ukuran, terdapat uniter $V$ sehingga
    $V^*BV$ segitiga atas. Ambil $U=U_1\diag(1,V)$.
    Determinan matriks segitiga menunjukkan bahwa entri diagonalnya
    adalah akar polinom karakteristik, termasuk multiplisitasnya.
\end{proof}

Schur berlaku untuk semua matriks kompleks, termasuk yang tidak dapat
didiagonalkan. Dari segitiga atas, untuk setiap polinom $f$ diperoleh
\[
    \tr f(A)=\sum_{i=1}^n f(\lambda_i),\qquad
    \det f(A)=\prod_{i=1}^n f(\lambda_i).
\]
Secara khusus, $\tr A^k=\sum_i\lambda_i^k$ untuk $k\ge1$,
tanpa perlu mengasumsikan diagonalisasi.

\begin{definisi}[][Matriks normal dan Hermitian][eigen:normal]
    Matriks $A$ disebut \emph{normal} jika $A^*A=AA^*$,
    dan \emph{Hermitian} jika $A^*=A$.
    Matriks Hermitian disebut semidefinit positif jika $x^*Ax\ge0$
    untuk semua $x$, dan definit positif jika $x^*Ax>0$ untuk semua $x\ne0$.
\end{definisi}

\begin{teorema}[][Dekomposisi spektral][eigen:teorema-spektral]
    Matriks $A\in\Mat_n(\C)$ normal jika dan hanya jika
    \[
        A=U\diag(\lambda_1,\ldots,\lambda_n)U^*
    \]
    untuk suatu matriks uniter $U$.
    Khususnya, matriks Hermitian mempunyai nilai eigen real dan
    basis ortonormal yang terdiri atas vektor eigen.
\end{teorema}
\begin{proof}
    Jika $A=UDU^*$ dengan $D$ diagonal, maka
    $A^*A=UD^*DU^*=UDD^*U^*=AA^*$.

    Sebaliknya, gunakan Schur: $T=U^*AU$ segitiga atas dan normal.
    Kesamaan entri $(1,1)$ pada $TT^*=T^*T$ memberi
    \[
        \sum_{j=1}^n|t_{1j}|^2=|t_{11}|^2.
    \]
    Maka semua entri baris pertama di luar diagonal nol.
    Blok sisanya tetap normal dan segitiga atas; induksi menunjukkan
    bahwa $T$ diagonal.
    Jika $A=A^*$, untuk vektor eigen $v\ne0$ berlaku
    $\lambda=v^*Av/(v^*v)\in\R$.
\end{proof}

Untuk matriks simetris real, vektor eigen dapat dipilih real, lalu
diortonormalkan di setiap ruang eigen. Jadi terdapat matriks
ortogonal real $Q$ dengan $A=QDQ^T$.
Kondisi simetris kompleks $A^T=A$ berbeda dari Hermitian;
contoh pada latihan menunjukkan bahwa matriks simetris kompleks
bahkan dapat gagal didiagonalkan.

\subsection{Merekonstruksi matriks dari ruang eigennya}

Jika kolom-kolom $U$ adalah $u_1,\ldots,u_n$, bentuk spektral memberi
\[
    A=\sum_{i=1}^n\lambda_i u_iu_i^*.
\]
Dengan mengelompokkan nilai eigen yang sama, dapat ditulis
$A=\sum_\lambda\lambda P_\lambda$, dengan $P_\lambda$ proyektor
ortogonal ke $E_\lambda(A)$. Proyektor ini memenuhi
$P_\lambda^*=P_\lambda$, $P_\lambda^2=P_\lambda$, dan
$P_\lambda P_\mu=0$ jika $\lambda\ne\mu$.

\begin{exam}[Matriks simetris dari dua vektor eigen]
    Konstruksikan matriks simetris real $3\times3$ yang mempunyai
    nilai eigen $1,1,-1$, dengan $(1,1,1)^T$ dan $(2,2,1)^T$
    sebagai vektor eigen untuk $1$.
\end{exam}
\begin{solution}
    Kedua vektor bebas dan merentang bidang $x_1=x_2$.
    Komplemen ortogonalnya direntang oleh
    $u=(1,-1,0)^T/\sqrt2$, yang harus menjadi vektor eigen untuk $-1$.
    Matriks yang dicari adalah
    \[
        A=1(I-uu^T)-uu^T=I-2uu^T
        =\begin{pmatrix}0&1&0\\1&0&0\\0&0&1\end{pmatrix}.
    \]
    Ruang eigennya menentukan aksi pada seluruh ruang, sehingga
    matriks yang memenuhi syarat ini unik.
\end{solution}

\begin{prop}[Tanda bentuk kuadratik dan nilai eigen]
    \label{eigen:kuadratik}
    Jika $A$ Hermitian, maka $A$ semidefinit positif tepat ketika
    semua nilai eigennya nonnegatif, dan definit positif tepat ketika
    semua nilai eigennya positif. Selain itu,
    \[
        \lambda_{\min}\|x\|^2\le x^*Ax\le\lambda_{\max}\|x\|^2.
    \]
\end{prop}
\begin{proof}
    Tuliskan $A=UDU^*$ dan $y=U^*x$.
    Maka $x^*Ax=\sum_i\lambda_i|y_i|^2$ dan
    $\|y\|^2=\|x\|^2$. Ketaksamaan mengikuti dari nilai eigen terkecil
    dan terbesar. Syarat perlu untuk kepositifan diperoleh dengan
    mengambil $x$ berupa masing-masing vektor eigen satuan.
\end{proof}

Pada matriks normal, rank tepat sama dengan banyaknya nilai eigen
tak nol, termasuk multiplisitasnya, karena bentuk diagonal diperoleh
melalui perkalian invertibel. Pada matriks Hermitian, informasi
spektral juga menentukan tanda bentuk kuadratik dan memungkinkan
pemilihan basis ortonormal.

\section{Memilih Dekomposisi}
\label{eigen:peta-dekomposisi}

Keempat bentuk berikut menjawab kebutuhan yang berbeda.
Pada baris Jordan dan Schur, medan yang dipakai adalah $\C$.

\begin{center}
    \small
    \renewcommand{\arraystretch}{1.35}
    \begin{tabular}{>{\raggedright\arraybackslash}p{.18\linewidth}
        >{\raggedright\arraybackslash}p{.29\linewidth}
        >{\raggedright\arraybackslash}p{.42\linewidth}}
        \textbf{Bentuk} & \textbf{Syarat dan rumus} & \textbf{Informasi utama}\\\hline
        Diagonalisasi & Basis vektor eigen; $A=PDP^{-1}$ &
            Pangkat dan polinom matriks dihitung pada entri diagonal.\\
        Jordan & Semua matriks persegi kompleks; $A=PJP^{-1}$ &
            Ukuran rantai, polinom minimal, dan nulitas pangkat.\\
        Schur & Semua matriks persegi kompleks; $A=UTU^*$ &
            Bentuk segitiga dengan perubahan basis ortonormal.\\
        Spektral & Matriks normal; $A=UDU^*$ &
            Basis eigen ortonormal dan proyektor ortogonal.\\\hline
    \end{tabular}
\end{center}

\begin{enumerate}[label=\arabic*.,leftmargin=*,itemsep=3pt]
    \item Jika diberikan persamaan polinom untuk $A$, periksa apakah
        polinom itu mempunyai akar berulang sebelum mencari vektor eigen.
    \item Jika akar karakteristik berulang, hitung
        $\nul(A-\lambda I)$ dan bandingkan dengan multiplisitas aljabarnya.
    \item Jika diperlukan struktur yang lebih rinci, hitung nulitas
        pangkat $A-\lambda I$ untuk menentukan blok Jordan.
    \item Jika matriks Hermitian atau simetris real, manfaatkan
        ortogonalitas ruang eigen dan bentuk $\sum_i\lambda_i u_iu_i^*$.
    \item Untuk identitas trace yang melibatkan pangkat, Schur memberi
        akses ke nilai eigen tanpa mengasumsikan diagonalisasi.
    \item Jika dua matriks komutatif, periksa ruang eigen yang invarian;
        spektrum sederhana sering mengubah persoalan menjadi interpolasi.
\end{enumerate}

\section{Latihan Pilihan}
\label{dekomp:latihan}

Latihan berikut menggunakan diagonalisasi, bentuk Jordan, dan dekomposisi spektral.

\begin{latihan}\label{eigen:lat-komutatif}
    Misalkan $A$ invertibel dan mempunyai $n$ nilai eigen berbeda.
    Jika $AB=BA^{-1}$, buktikan bahwa $B^2$ dapat didiagonalkan.
    Apakah $B$ sendiri harus dapat didiagonalkan?
\end{latihan}

\begin{latihan}\label{eigen:lat-trace}
    Jika $A\in\Mat_n(\C)$ memenuhi $\tr A^k=0$ untuk
    $k=1,\ldots,n$, buktikan bahwa $A^n=0$.
    Buktikan pula arah sebaliknya.
\end{latihan}

\begin{latihan}\label{eigen:lat-real}
    Misalkan $A\in\Mat_2(\R)$ memenuhi $A^2=-I$.
    Buktikan bahwa $A$ serupa di atas $\R$ dengan
    $\begin{pmatrix}0&-1\\1&0\end{pmatrix}$.
    Jelaskan perbedaan diagonalisasinya di atas $\R$ dan $\C$.
\end{latihan}

\begin{latihan}\label{eigen:lat-rekonstruksi}
    Matriks simetris real $A$ mempunyai nilai eigen $1,2,3$.
    Diketahui $u=(-1,-1,1)^T$ dan $v=(1,-2,-1)^T$ merupakan
    vektor eigen untuk $1$ dan $2$. Tentukan vektor eigen untuk $3$
    dan hitung matriks $A$.
\end{latihan}

\begin{latihan}\label{eigen:lat-rank-diagonal}
    Buktikan bahwa $A\in\Mat_n(\C)$ dapat didiagonalkan jika dan
    hanya jika untuk setiap nilai eigen $\lambda$ berlaku
    \[
        \rank(A-\lambda I)=\rank(A-\lambda I)^2.
    \]
\end{latihan}

\begin{latihan}\label{eigen:lat-minimal-sama}
    Konstruksikan dua matriks $6\times6$ yang mempunyai polinom
    karakteristik dan polinom minimal sama, tetapi tidak serupa.
    Jelaskan mengapa setiap polinom yang menganulir salah satunya
    juga menganulir yang lain.
\end{latihan}

\begin{latihan}\label{eigen:lat-simetris-kompleks}
    Matriks $A=\begin{pmatrix}1&i\\i&-1\end{pmatrix}$ memenuhi $A^T=A$.
    Tentukan polinom karakteristik, polinom minimal, dan dimensi
    ruang eigennya. Apakah $A$ dapat didiagonalkan di atas $\C$?
\end{latihan}

\begin{latihan}\label{eigen:lat-frobenius}
    Misalkan $\lambda_1,\ldots,\lambda_n$ adalah nilai eigen
    $A=(a_{ij})\in\Mat_n(\C)$, dihitung dengan multiplisitas.
    Buktikan bahwa
    \[
        \sum_{i=1}^n|\lambda_i|^2\le\sum_{i,j=1}^n|a_{ij}|^2,
    \]
    dengan kesamaan tepat ketika $A$ normal.
\end{latihan}

\begin{latihan}\label{eigen:lat-komutator}
    Misalkan $A\in\Mat_n(\C)$ dapat didiagonalkan dan
    $\mathcal T:\Mat_n(\C)\to\Mat_n(\C)$ didefinisikan oleh
    $\mathcal T(X)=AX-XA$.
    Tentukan nilai eigen $\mathcal T$ dan buktikan bahwa pemetaan ini
    dapat didiagonalkan. Jika multiplisitas nilai eigen berbeda dari
    $A$ adalah $m_1,\ldots,m_s$, tentukan rank dan nulitas $\mathcal T$.
\end{latihan}

\begin{latihan}\label{eigen:lat-siklik}
    Vektor $v$ disebut siklik untuk $A\in\Mat_n(\C)$ jika
    $v,Av,\ldots,A^{n-1}v$ merupakan basis $\C^n$.
    \begin{enumerate}[label=(\alph*),nosep]
        \item Jika ada vektor siklik, buktikan $m_A=\chi_A$ dan bahwa
            setiap ruang eigen $A$ berdimensi satu.
        \item Jika $A$ mempunyai $n$ nilai eigen berbeda, konstruksikan
            satu vektor siklik.
    \end{enumerate}
\end{latihan}

\begin{latihan}\label{eigen:lat-blok}
    Untuk $A,B\in\Mat_n(\C)$, definisikan
    $M=\begin{pmatrix}A&B\\B&A\end{pmatrix}$.
    Nyatakan $\chi_M$ melalui $\chi_{A+B}$ dan $\chi_{A-B}$.
    Buktikan bahwa $M$ dapat didiagonalkan jika dan hanya jika
    $A+B$ dan $A-B$ masing-masing dapat didiagonalkan.
    Apakah diperlukan asumsi $AB=BA$?
\end{latihan}

\ifpembahasan
\section{Pembahasan Latihan}
\label{dekomp:pembahasan}

\subsection*{Latihan \ref{eigen:lat-komutatif}: kuadrat yang komutatif}
Dari $AB=BA^{-1}$ diperoleh $A^{-1}B=BA$.
Maka
\[
    AB^2=BA^{-1}B=B^2A.
\]
Karena $A$ mempunyai spektrum sederhana, \autoref{eigen:centralizer}
menjamin $B^2$ diagonal dalam basis vektor eigen $A$.
Matriks $B$ sendiri belum tentu dapat didiagonalkan. Contohnya,
\[
    A=\begin{pmatrix}2&0\\0&1/2\end{pmatrix},\qquad
    B=\begin{pmatrix}0&1\\0&0\end{pmatrix}.
\]
Relasi $AB=BA^{-1}$ berlaku, tetapi $B$ nilpoten tak nol.

\subsection*{Latihan \ref{eigen:lat-trace}: memisahkan nilai eigen tak nol}
Andaikan terdapat nilai eigen tak nol yang berbeda
$\mu_1,\ldots,\mu_s$, dengan multiplisitas aljabar $r_1,\ldots,r_s$.
Dari Schur dan hipotesis, untuk $k=1,\ldots,s$ berlaku
\[
    \sum_{j=1}^s r_j\mu_j^k=0.
\]
Matriks koefisiennya adalah matriks Vandermonde dengan setiap kolom
ke-$j$ dikalikan $\mu_j$. Determinannya tak nol karena semua $\mu_j$
berbeda dan tak nol. Akibatnya semua $r_j=0$, bertentangan dengan
multiplisitas yang positif.
Jadi seluruh nilai eigen $A$ nol, $\chi_A(t)=t^n$, dan $A^n=0$
oleh Cayley--Hamilton.
Sebaliknya, $A^n=0$ memaksa setiap nilai eigen nol; Schur memberi
$\tr A^k=0$ untuk semua $k\ge1$.

\subsection*{Latihan \ref{eigen:lat-real}: basis dari satu vektor dan bayangannya}
Pilih $v\ne0$ real. Vektor $v$ dan $Av$ bebas: jika $Av=cv$
dengan $c\in\R$, maka $-v=A^2v=c^2v$, yang mustahil.
Dalam basis $(v,Av)$, aksi $A$ adalah
\[
    v\longmapsto Av,\qquad Av\longmapsto-v.
\]
Jadi dengan $P=[v\ Av]$ diperoleh
$P^{-1}AP=\begin{pmatrix}0&-1\\1&0\end{pmatrix}$.
Di atas $\R$ tidak ada nilai eigen, sehingga diagonalisasi mustahil.
Di atas $\C$, kedua nilai eigen $i,-i$ berbeda, sehingga diagonalisasi ada.

\subsection*{Latihan \ref{eigen:lat-rekonstruksi}: proyektor ortogonal berbobot}
Vektor $u$ dan $v$ ortogonal. Vektor $w=(1,0,1)^T$ ortogonal terhadap
keduanya, sehingga menjadi vektor eigen untuk $3$.
Karena $\|u\|^2=3$, $\|v\|^2=6$, dan $\|w\|^2=2$,
\[
    A=\frac{uu^T}{3}+2\frac{vv^T}{6}+3\frac{ww^T}{2}
    =\frac16\begin{pmatrix}13&-2&5\\-2&10&2\\5&2&13\end{pmatrix}.
\]
Masing-masing suku mengatur aksi pada satu arah dalam basis ortonormal.

\subsection*{Latihan \ref{eigen:lat-rank-diagonal}: mendeteksi blok berukuran dua atau lebih}
Rank yang sama ekuivalen dengan
$\nul(A-\lambda I)=\nul(A-\lambda I)^2$.
Menurut \autoref{eigen:jordan-nulitas}, selisih kedua nulitas adalah
banyaknya blok Jordan untuk $\lambda$ yang berukuran setidaknya dua.
Jadi syarat tersebut untuk setiap nilai eigen berarti semua blok
berukuran satu, tepat syarat diagonalisasi.

\subsection*{Latihan \ref{eigen:lat-minimal-sama}: ukuran blok masih dapat berbeda}
Ambil
\[
    A=\diag\bigl(J_3(0),J_3(0)\bigr),\qquad
    B=\diag\bigl(J_3(0),J_2(0),J_1(0)\bigr).
\]
Keduanya mempunyai $\chi(t)=t^6$ dan $m(t)=t^3$.
Namun, $\nul A=2$ dan $\nul B=3$, sehingga keduanya tidak serupa.
Untuk kedua matriks, $f$ menganulirnya tepat ketika $t^3$ membagi $f$.
Jadi kesamaan seluruh polinom penganulir belum menentukan ukuran
semua blok Jordan.

\subsection*{Latihan \ref{eigen:lat-simetris-kompleks}: simetris tidak selalu Hermitian}
Perkalian langsung memberi $A^2=0$, sedangkan $A\ne0$.
Karena trace dan determinannya nol, $\chi_A(t)=t^2$.
Polinom minimalnya juga $t^2$, sebab $A$ tidak dianulir oleh $t$.
Rank $A$ adalah satu dan $\ker A=\operatorname{span}\{(-i,1)^T\}$.
Jadi satu-satunya nilai eigen adalah $0$, dengan multiplisitas
aljabar dua dan geometri satu; matriks tidak dapat didiagonalkan.
Tidak ada pertentangan dengan teorema spektral karena $A^*\ne A$.

\subsection*{Latihan \ref{eigen:lat-frobenius}: entri di atas diagonal Schur}
Untuk matriks persegi berukuran sama,
$\tr(XY)=\sum_{i,j}x_{ij}y_{ji}=\tr(YX)$.
Dengan dekomposisi Schur $A=UTU^*$,
\[
    \sum_{i,j}|a_{ij}|^2=\tr(A^*A)=\tr(T^*T)
    =\sum_i|\lambda_i|^2+\sum_{i<j}|t_{ij}|^2.
\]
Ini memberi ketaksamaan. Kesamaan berlaku tepat ketika semua entri
di atas diagonal $T$ nol, yaitu ketika $A$ dapat didiagonalkan
secara uniter. Oleh teorema spektral, ini ekuivalen dengan normalitas $A$.

\subsection*{Latihan \ref{eigen:lat-komutator}: basis matriks satuan}
Tuliskan $A=PDP^{-1}$, dengan $D=\diag(\lambda_1,\ldots,\lambda_n)$.
Untuk matriks satuan $E_{ij}$, yang hanya mempunyai entri $1$ pada
posisi $(i,j),$ berlaku
\[
    DE_{ij}-E_{ij}D=(\lambda_i-\lambda_j)E_{ij}.
\]
Karena $PE_{ij}P^{-1}$ untuk semua pasangan $(i,j)$ membentuk basis
$\Mat_n(\C)$, matriks-matriks ini merupakan basis vektor eigen
$\mathcal T$, dengan nilai eigen $\lambda_i-\lambda_j$.
Nilai eigen nol muncul tepat ketika kedua indeks berasal dari
nilai eigen $A$ yang sama. Dengan demikian
\[
    \nul\mathcal T=\sum_{j=1}^s m_j^2,\qquad
    \rank\mathcal T=n^2-\sum_{j=1}^s m_j^2.
\]

\subsection*{Latihan \ref{eigen:lat-siklik}: vektor siklik dan Vandermonde}
Jika $v,Av,\ldots,A^{n-1}v$ bebas, tidak ada polinom tak nol
berderajat kurang dari $n$ yang menganulir $A$: menerapkannya pada $v$
akan menghasilkan relasi linear tak trivial.
Maka $\deg m_A\ge n$; karena $m_A\mid\chi_A$ dan $\deg\chi_A=n$,
diperoleh $m_A=\chi_A$.
Dalam bentuk Jordan, pangkat setiap faktor dalam $m_A$ adalah ukuran
blok terbesar, sedangkan pangkatnya dalam $\chi_A$ adalah jumlah
ukuran semua blok untuk nilai eigen itu. Kesamaan keduanya memaksa
tepat satu blok per nilai eigen, sehingga setiap ruang eigen berdimensi satu.

Sebaliknya, jika nilai eigen $\lambda_1,\ldots,\lambda_n$ berbeda,
pilih basis vektor eigen $v_1,\ldots,v_n$ dan ambil $v=\sum_i v_i$.
Koordinat $v,Av,\ldots,A^{n-1}v$ terhadap basis ini membentuk matriks
Vandermonde dengan baris ke-$i$ sama dengan
$(1,\lambda_i,\ldots,\lambda_i^{n-1})$.
Determinan tak nol, sehingga $v$ siklik.

\subsection*{Latihan \ref{eigen:lat-blok}: memisahkan jumlah dan selisih}
Matriks
\[
    U=\frac1{\sqrt2}\begin{pmatrix}I&I\\I&-I\end{pmatrix}
\]
uniter dan memenuhi
\[
    U^*MU=\begin{pmatrix}A+B&0\\0&A-B\end{pmatrix}.
\]
Jadi $\chi_M(t)=\chi_{A+B}(t)\chi_{A-B}(t)$.
Jika kedua blok dapat didiagonalkan, gabungan perubahan basisnya
mendiagonalkan matriks diagonal blok, dan karenanya juga $M$.
Sebaliknya, jika $M$ dapat didiagonalkan, suatu polinom dengan faktor
linear berbeda menganulir $M$ dan kedua bloknya. Maka kedua blok
dapat didiagonalkan oleh \autoref{eigen:squarefree}.
Tidak diperlukan $AB=BA$; identitas blok di atas berlaku tanpa asumsi itu.

\fi

% !TeX root = ../../ONMIPA Alin.tex
% Sumber utama: Zhang (2009), Bab 4, Definitions and Facts serta soal-soal terpilih.
\chapter{Matriks Khusus}
\label{ch:matriks-khusus}

Syarat seperti $A^*=A$, $A^2=A$, atau $A^*A=I$ membatasi bentuk
spektrum, kernel, dan hasil kali suatu matriks. Dalam soal ONMIPA,
mengenali struktur tersebut sering menghindarkan perhitungan
determinan atau eliminasi yang panjang.
Pembahasan memanfaatkan dekomposisi pada
\autoref{ch:dekomposisi} sebagai alat pembuktian.

Kecuali disebutkan lain, matriks persegi berukuran $n$ dan berentri
kompleks. Notasi $A^*=\overline A^{\,T}$ berarti transpose konjugat,
sedangkan $\|x\|^2=x^*x$. Untuk matriks real, $A^*=A^T$.

\section{Hermitian, Skew-Hermitian, dan Normal}
\label{khusus:hermitian}

\begin{definisi}[][Tiga struktur terhadap transpose konjugat][khusus:def-hermitian]
    Matriks $A$ disebut \emph{Hermitian} jika $A^*=A$,
    \emph{skew-Hermitian} jika $A^*=-A$, dan \emph{normal} jika
    $A^*A=AA^*$. Matriks Hermitian real adalah matriks simetris.
\end{definisi}

Matriks Hermitian dan skew-Hermitian selalu normal.
Teorema spektral (\autoref{eigen:teorema-spektral}) menjamin bahwa
matriks normal dapat didiagonalkan dengan perubahan basis uniter.
Jika $Av=\lambda v$ dan $v\ne0$, maka
$v^*Av=\lambda\|v\|^2$. Karena $v^*Av$ real untuk $A$ Hermitian,
nilai eigennya real. Untuk $A$ skew-Hermitian, $iA$ Hermitian,
sehingga nilai eigen $A$ murni imajiner, termasuk kemungkinan nol.

\begin{prop}[Bagian Hermitian dari suatu matriks]
    \label{khusus:bagian-hermitian}
    Setiap matriks kompleks mempunyai representasi unik
    \[
        A=H+iK,\qquad H=\frac{A+A^*}{2},\qquad
        K=\frac{A-A^*}{2i},
    \]
    dengan $H,K$ Hermitian. Matriks $A$ normal jika dan hanya jika
    $HK=KH$.
\end{prop}
\begin{proof}
    Rumus langsung memberi $H^*=H$, $K^*=K$, dan $A=H+iK$.
    Jika representasi demikian ada, $A^*=H-iK$ menentukan $H,K$
    secara unik. Selanjutnya,
    \[
        A^*A-AA^*=2i(HK-KH),
    \]
    sehingga normalitas ekuivalen dengan komutativitas kedua bagian.
\end{proof}

\begin{prop}[Kernel dan bayangan matriks normal]
    \label{khusus:kernel-normal}
    Jika $A$ normal, maka
    \[
        \ker A=\ker A^*,\qquad
        \operatorname{Im}A=(\ker A)^\perp,\qquad
        \C^n=\ker A\oplus\operatorname{Im}A.
    \]
    Jumlah langsung terakhir bersifat ortogonal.
\end{prop}
\begin{proof}
    Normalitas memberi
    $\|Ax\|^2=x^*A^*Ax=x^*AA^*x=\|A^*x\|^2$.
    Jadi kedua kernel sama. Untuk matriks apa pun,
    $y\perp\operatorname{Im}A$ ekuivalen dengan $A^*y=0$;
    akibatnya $\operatorname{Im}A=(\ker A^*)^\perp$.
\end{proof}

Perlu dibedakan: $A^T=A$ pada matriks kompleks tidak menjamin
$A^*=A$. Contohnya,
$A=\begin{pmatrix}1&i\\i&-1\end{pmatrix}$ simetris dan memenuhi
$A^2=0$, tetapi $A\ne0$. Matriks normal dengan semua nilai eigen nol
justru harus nol, sebab diagonal pada dekomposisi uniternya nol.

\section{Matriks Semidefinit dan Definit Positif}
\label{khusus:positif}

\begin{definisi}[][Positivitas dan urutan Loewner][khusus:def-positif]
    Matriks Hermitian $A$ disebut \emph{semidefinit positif}, ditulis
    $A\succeq0$, jika $x^*Ax\ge0$ untuk setiap $x$.
    Matriks ini disebut \emph{definit positif}, ditulis $A\succ0$,
    jika $x^*Ax>0$ untuk setiap $x\ne0$.
    Untuk $A,B$ Hermitian, $A\succeq B$ berarti $A-B\succeq0$.
\end{definisi}

Urutan ini ditentukan oleh bentuk kuadrat, bukan oleh perbandingan
tiap entri. Matriks $\begin{pmatrix}1&-1\\-1&1\end{pmatrix}$
semidefinit positif walaupun mempunyai entri negatif. Sebaliknya,
$\begin{pmatrix}1&2\\2&1\end{pmatrix}$ berentri positif tetapi
mempunyai nilai eigen $-1$.

\begin{teorema}[][Kriteria spektral dan representasi Gram][khusus:gram]
    Untuk $A$ Hermitian, pernyataan berikut ekuivalen:
    \begin{enumerate}[label=(\roman*),nosep]
        \item $A\succeq0$;
        \item semua nilai eigen $A$ tak negatif;
        \item $A=B^*B$ untuk suatu matriks $B$.
    \end{enumerate}
    Selain itu, $A\succ0$ jika dan hanya jika semua nilai eigennya
    positif. Jika $A=B^*B$, maka $\ker A=\ker B$ dan
    $\rank A=\rank B$.
\end{teorema}
\begin{proof}
    Tuliskan $A=UDU^*$ dengan $U$ uniter dan $D$ diagonal real.
    Dengan $y=U^*x$, bentuk kuadratnya menjadi
    $x^*Ax=\sum_i d_i|y_i|^2$. Ini memberi kriteria spektral.
    Jika $d_i\ge0$, pilih
    $B=\diag(\sqrt{d_1},\ldots,\sqrt{d_n})U^*$.
    Sebaliknya, $A=B^*B$ memberi $x^*Ax=\|Bx\|^2\ge0$.
    Kesamaan kernel mengikuti identitas yang sama; kesamaan rank
    mengikuti teorema rank--nulitas.
\end{proof}

Untuk setiap $X$, kongruensi mempertahankan semidefinit positif:
$A\succeq0$ mengakibatkan $X^*AX\succeq0$.
Jika $X$ persegi dan invertibel, implikasinya dua arah, juga untuk
definit positif. Ini berbeda dari keserupaan $X^{-1}AX$:
kongruensi mengubah koordinat suatu bentuk kuadrat.

\begin{teorema}[][Kriteria minor utama][khusus:sylvester]
    Misalkan $A$ Hermitian.
    \begin{enumerate}[label=(\roman*),nosep]
        \item $A\succ0$ jika dan hanya jika semua minor utama
            terdepan $\det A[1\!:\!k,1\!:\!k]$ positif,
            untuk $k=1,\ldots,n$.
        \item $A\succeq0$ jika dan hanya jika semua minor utama
            $\det A[S,S]$ tak negatif, untuk setiap himpunan indeks
            tak kosong $S\subseteq\{1,\ldots,n\}$.
    \end{enumerate}
\end{teorema}
\begin{proof}
    Setiap submatriks utama dari matriks definit positif juga
    definit positif, sehingga determinannya positif. Untuk arah
    sebaliknya pada (i), gunakan induksi pada ukuran. Tulis
    \[
        A=\begin{pmatrix}B&b\\b^*&c\end{pmatrix}.
    \]
    Hipotesis induksi memberi $B\succ0$. Karena
    $\det A=\det B\,(c-b^*B^{-1}b)>0$, komplemen Schurnya positif.
    Melengkapi kuadrat memberi
    \[
        \begin{pmatrix}x\\t\end{pmatrix}^{\!*}
        A\begin{pmatrix}x\\t\end{pmatrix}
        =(x+B^{-1}bt)^*B(x+B^{-1}bt)
          +(c-b^*B^{-1}b)|t|^2,
    \]
    yang positif untuk $(x,t)\ne0$.

    Untuk (ii), arah perlu kembali mengikuti positivitas submatriks
    utama dan kriteria spektral. Sebaliknya, ambil $\varepsilon>0$.
    Setiap minor utama terdepan dari $A+\varepsilon I$ adalah
    \[
        \det\bigl(A[1\!:\!k,1\!:\!k]+\varepsilon I_k\bigr)
        =\sum_{S\subseteq\{1,\ldots,k\}}
           \varepsilon^{k-|S|}\det A[S,S],
    \]
    dengan determinan kosong bernilai $1$; lihat
    \autoref{det:principal}. Semua suku tak negatif dan suku
    $\varepsilon^k$ positif. Bagian (i) memberi $A+\varepsilon I\succ0$.
    Untuk setiap $x$, lewatkan $\varepsilon\to0$ pada
    $x^*(A+\varepsilon I)x\ge0$ untuk memperoleh $x^*Ax\ge0$.
\end{proof}

Pada kasus semidefinit, memeriksa minor utama terdepan saja
\emph{tidak cukup}. Matriks $\diag(0,-1)$ mempunyai kedua minor
terdepan nol, tetapi tidak semidefinit positif.
Minor utama berukuran dua juga memberi batas yang sering berguna:
\[
    A\succeq0\quad\Longrightarrow\quad
    a_{ii}\ge0,\qquad |a_{ij}|^2\le a_{ii}a_{jj}.
\]
Jadi, jika $a_{ii}=0$, seluruh baris dan kolom ke-$i$ bernilai nol.

\begin{prop}[Positivitas matriks blok]
    \label{khusus:schur}
    Misalkan $A\succ0$, $C=C^*$, dan
    $M=\begin{pmatrix}A&B\\B^*&C\end{pmatrix}$. Maka
    \[
        M\succeq0\iff C-B^*A^{-1}B\succeq0,
        \qquad
        M\succ0\iff C-B^*A^{-1}B\succ0.
    \]
\end{prop}
\begin{proof}
    Untuk $T=\begin{pmatrix}I&-A^{-1}B\\0&I\end{pmatrix}$,
    perkalian langsung memberi
    $T^*MT=\diag(A,C-B^*A^{-1}B)$.
    Gunakan invariansi positivitas di bawah kongruensi invertibel.
\end{proof}

\begin{exam}[Parameter pada matriks simetris]
    Tentukan semua $\lambda\in\R$ agar
    \[
        A_\lambda=\begin{pmatrix}
            1&\lambda&-1\\\lambda&4&2\\-1&2&4
        \end{pmatrix}
    \]
    definit positif atau semidefinit positif. Tentukan pula ranknya
    pada batas interval semidefinit positif.
\end{exam}
\begin{solution}
    Minor utama terdepan berturut-turut adalah
    $1$, $4-\lambda^2$, dan $4(1-\lambda)(\lambda+2)$.
    Ketiganya positif tepat ketika $-2<\lambda<1$.
    Untuk semidefinit positif, minor utama berukuran dua yang lain
    adalah $3$ dan $12$; minor berukuran satu semuanya positif.
    Jadi syarat lengkapnya
    $4-\lambda^2\ge0$ dan $(1-\lambda)(\lambda+2)\ge0$,
    yaitu $-2\le\lambda\le1$.
    Pada kedua ujung interval, determinan nol tetapi minor utama
    pada indeks $\{1,3\}$ bernilai $3$; ranknya tepat dua.
    Di bagian dalam interval, ranknya tiga.
\end{solution}

\section{Akar Kuadrat, Hasil Kali, dan Urutan Matriks}
\label{khusus:akar}

\begin{teorema}[][Akar kuadrat semidefinit positif][khusus:akar-unik]
    Setiap $A\succeq0$ mempunyai tepat satu akar kuadrat
    semidefinit positif, ditulis $A^{1/2}$. Jika
    $A=U\diag(\lambda_1,\ldots,\lambda_n)U^*$, maka
    \[
        A^{1/2}=U\diag(\sqrt{\lambda_1},\ldots,\sqrt{\lambda_n})U^*.
    \]
    Jika $A\succ0$, akar tersebut invertibel; inversnya ditulis $A^{-1/2}$.
\end{teorema}
\begin{proof}
    Rumus spektral membuktikan keberadaan. Jika $X\succeq0$ dan
    $X^2=A$, maka $XA=AX$. Jadi setiap ruang eigen $E_\lambda(A)$
    invarian terhadap $X$. Pembatasan $X$ pada ruang ini Hermitian
    semidefinit positif dan kuadratnya $\lambda I$.
    Semua nilai eigennya harus $\sqrt\lambda$, sehingga pembatasan
    tersebut tepat $\sqrt\lambda I$. Ini menentukan $X$ secara unik.
\end{proof}

Syarat semidefinit positif diperlukan untuk keunikan: $I$ dan $-I$
keduanya mempunyai kuadrat $I$, tetapi hanya $I$ yang semidefinit positif.

\begin{exam}[Akar kuadrat tanpa menghitung vektor eigen satu per satu]
    Misalkan $J_n$ matriks yang seluruh entrinya $1$, dan
    $A=\alpha I+\beta J_n$, dengan $\alpha,\beta\in\R$.
    Tentukan syarat $A\succeq0$ dan rumus $A^{1/2}$.
\end{exam}
\begin{solution}
    Proyektor ortogonal $P=J_n/n$ dan $Q=I-P$ memenuhi
    $PQ=0$. Untuk $n\ge2$,
    $A=\alpha Q+(\alpha+n\beta)P$, sehingga syaratnya
    $\alpha\ge0$ dan $\alpha+n\beta\ge0$. Pada kondisi itu,
    \[
        A^{1/2}=\sqrt\alpha\left(I-\frac{J_n}{n}\right)
           +\sqrt{\alpha+n\beta}\,\frac{J_n}{n}.
    \]
    Untuk $n=1$, cukup $\alpha+\beta\ge0$ dan akar skalarnya.
\end{solution}

\begin{prop}[Kapan hasil kali tetap positif?]
    \label{khusus:produk-positif}
    Jika $A,B\succeq0$, maka
    \[
        AB\succeq0\iff AB=BA.
    \]
    Jika $A\succ0$ dan $B$ Hermitian, $AB$ dapat didiagonalkan
    dan seluruh nilai eigennya real.
\end{prop}
\begin{proof}
    Jika $AB\succeq0$, matriks itu Hermitian, maka
    $AB=(AB)^*=BA$. Sebaliknya, dua matriks Hermitian yang komutatif
    mempunyai basis eigen ortonormal bersama: diagonalkan $B$
    pada tiap ruang eigen $A$ yang invarian terhadap $B$.
    Diagonal $AB$ dalam basis itu berupa hasil kali bilangan tak negatif.

    Jika $A\succ0$, gunakan keserupaan
    \[
        A^{-1/2}(AB)A^{1/2}=A^{1/2}BA^{1/2},
    \]
    yang ruas kanannya Hermitian. Terapkan teorema spektral.
\end{proof}

Jika hanya $A\succeq0$ dan $B$ Hermitian, nilai eigen $AB$ masih real:
$AB$ dan $A^{1/2}BA^{1/2}$ mempunyai polinom karakteristik sama,
melalui identitas hasil kali pada \autoref{eigen:abba}.
Namun, diagonalisasi dapat gagal. Ambil
$A=\diag(1,0)$ dan $B=\begin{pmatrix}0&1\\1&0\end{pmatrix}$;
hasil kalinya $\begin{pmatrix}0&1\\0&0\end{pmatrix}$ tidak dapat didiagonalkan.

\begin{teorema}[][Invers membalik urutan][khusus:invers-order]
    Jika $A\succeq B\succ0$, maka $B^{-1}\succeq A^{-1}$.
\end{teorema}
\begin{proof}
    Matriks $C=B^{-1/2}AB^{-1/2}$ memenuhi $C\succeq I$.
    Nilai eigennya minimal $1$, sehingga $C^{-1}\preceq I$.
    Karena $C^{-1}=B^{1/2}A^{-1}B^{1/2}$, kongruensi dengan
    $B^{-1/2}$ menghasilkan $A^{-1}\preceq B^{-1}$.
\end{proof}

\section{Ketaksamaan Trace}
\label{khusus:trace}

Identitas $\tr(XY)=\tr(YX)$ membolehkan pergeseran siklik pada
trace suatu hasil kali, tetapi tidak membolehkan pertukaran faktor
secara sembarang. Sebagai contoh, $\tr(ABAB)$ tidak selalu sama
dengan $\tr(A^2B^2)$.

\begin{teorema}[][Cauchy--Schwarz untuk matriks][khusus:cs]
    Untuk $A,B\in\Mat_{m\times n}(\C)$, berlaku
    \[
        |\tr(A^*B)|^2\le \tr(A^*A)\tr(B^*B).
    \]
    Kesamaan berlaku tepat ketika $A,B$ bergantung linear.
    Khusus untuk $H$ Hermitian dengan rank $r>0$,
    \[
        (\tr H)^2\le r\tr(H^2).
    \]
    Pada ketaksamaan terakhir, kesamaan berlaku tepat ketika semua
    nilai eigen tak nol $H$ sama.
\end{teorema}
\begin{proof}
    Karena $\tr(A^*B)=\sum_{i,j}\overline{a_{ij}}b_{ij}$,
    ketaksamaan pertama beserta kasus kesamaannya adalah
    Cauchy--Schwarz untuk vektor seluruh entri kedua matriks.
    Untuk yang kedua, terapkan Cauchy--Schwarz pada $r$ nilai eigen
    real tak nol $\lambda_1,\ldots,\lambda_r$:
    $(\sum_i\lambda_i)^2\le r\sum_i\lambda_i^2$.
\end{proof}

\begin{teorema}[][Trace mendeteksi komutativitas][khusus:trace-komutator]
    Untuk $A,B$ Hermitian,
    \[
        \tr(A^2B^2)-\tr\bigl((AB)^2\bigr)
        =\tfrac12\tr\bigl((AB-BA)^*(AB-BA)\bigr)\ge0.
    \]
    Kesamaan berlaku jika dan hanya jika $AB=BA$.
\end{teorema}
\begin{proof}
    Tuliskan $C=AB-BA$ sehingga $C^*=BA-AB$.
    Ekspansi $C^*C$ diikuti pergeseran siklik dalam trace memberi
    \[
        \tr(C^*C)=2\tr(A^2B^2)-2\tr(ABAB).
    \]
    Karena $\tr(C^*C)=\sum_{i,j}|c_{ij}|^2$, ruas kanan real
    tak negatif dan nol tepat ketika $C=0$.
\end{proof}

\begin{prop}[Trace hasil kali matriks semidefinit positif]
    \label{khusus:trace-positif}
    Jika $A,B\succeq0$, maka $\tr(AB)\ge0$, dengan kesamaan
    jika dan hanya jika $AB=0$.
\end{prop}
\begin{proof}
    Ambil $C=B^{1/2}A^{1/2}$. Pergeseran siklik memberi
    $\tr(AB)=\tr(C^*C)\ge0$. Jika trace nol, $C=0$ dan
    $BA=B^{1/2}CA^{1/2}=0$, sehingga $AB=(BA)^*=0$.
    Arah sebaliknya langsung.
\end{proof}

\begin{exam}[Hasil kali Hermitian tidak harus berspektrum real]
    Ambil $A=\diag(1,-1)$ dan
    $B=\begin{pmatrix}0&1\\1&0\end{pmatrix}$.
    Keduanya Hermitian, tetapi
    $AB=\begin{pmatrix}0&1\\-1&0\end{pmatrix}$ mempunyai nilai eigen
    $i,-i$. Di sini
    \[
        \tr\bigl((AB)^2\bigr)=-2,\qquad \tr(A^2B^2)=2.
    \]
    Jadi hipotesis semidefinit positif pada salah satu faktor
    dalam hasil tentang spektrum real tidak boleh dihilangkan.
\end{exam}

\section{Determinan dan Hasil Kali Hadamard}
\label{khusus:hadamard}

\begin{teorema}[][Ketaksamaan determinan Hadamard][khusus:det-hadamard]
    Jika $A=(a_{ij})\succeq0$, maka
    \[
        0\le\det A\le\prod_{i=1}^n a_{ii}.
    \]
    Kesamaan pada batas atas berlaku tepat ketika $A$ diagonal
    atau sedikitnya satu entri diagonalnya nol.
\end{teorema}
\begin{proof}
    Jika $a_{ii}=0$, baris dan kolom ke-$i$ nol, sehingga
    determinan dan hasil kali diagonal sama-sama nol.
    Sekarang anggap semua entri diagonal positif dan gunakan induksi.
    Tulis $A=\begin{pmatrix}a&b^*\\b&C\end{pmatrix}$.
    Komplemen Schur $S=C-bb^*/a$ semidefinit positif; maka
    \[
        \det A=a\det S
        \le a\prod_j\left(c_{jj}-\frac{|b_j|^2}{a}\right)
        \le a\prod_j c_{jj}.
    \]
    Semua faktor di tengah tak negatif. Jika batas atas tercapai,
    nilainya positif, sehingga $b_j=0$ untuk setiap $j$.
    Selanjutnya hipotesis induksi memaksa $C$ diagonal.
    Sebaliknya, matriks diagonal jelas mencapai kesamaan.
\end{proof}

Untuk matriks persegi $B$ dengan kolom $b_1,\ldots,b_n$, terapkan
teorema ini pada $B^*B$ untuk memperoleh
\[
    |\det B|\le\prod_{j=1}^n\|b_j\|.
\]
Jika semua kolom tak nol, kesamaan berlaku tepat ketika
kolom-kolomnya saling ortogonal. Jika ada kolom nol, kedua ruas nol.
Sementara itu, rata-rata aritmetika--geometri pada nilai eigen
$A\succeq0$ menghasilkan
\[
    \det A\le\left(\frac{\tr A}{n}\right)^n,
\]
dengan kesamaan tepat ketika $A$ kelipatan skalar dari $I$.

\begin{definisi}[][Hasil kali Hadamard][khusus:def-hadamard]
    Untuk dua matriks berukuran sama, hasil kali Hadamard adalah
    $A\circ B=(a_{ij}b_{ij})$, yaitu perkalian entri yang seletak.
    Operasi ini berbeda dari perkalian matriks biasa $AB$.
\end{definisi}

\begin{teorema}[][Teorema hasil kali Schur][khusus:schur-product]
    Jika $A,B\succeq0$, maka $A\circ B\succeq0$.
    Jika keduanya definit positif, $A\circ B$ juga definit positif.
\end{teorema}
\begin{proof}
    Dari dekomposisi spektral, tulis $A=\sum_k u_ku_k^*$,
    dengan akar nilai eigen dimasukkan ke dalam $u_k$.
    Jika $D_k=\diag(u_k)$, perkalian entri memberi
    \[
        A\circ B=\sum_k D_kBD_k^*\succeq0.
    \]
    Untuk kasus definit positif, pilih $\alpha>0$ sehingga
    $A\succeq\alpha I$. Bagian pertama memberi
    \[
        A\circ B\succeq\alpha(I\circ B)
           =\alpha\diag(b_{11},\ldots,b_{nn})\succ0.
    \]
\end{proof}

\begin{exam}[Dua faktor singular, hasil kali Hadamard invertibel]
    Matriks
    \[
        A=\begin{pmatrix}1&1&0\\1&1&0\\0&0&1\end{pmatrix},\qquad
        B=\begin{pmatrix}1&0&1\\0&1&0\\1&0&1\end{pmatrix}
    \]
    keduanya semidefinit positif dan mempunyai rank dua. Hal ini
    terlihat dari
    $x^*Ax=|x_1+x_2|^2+|x_3|^2$ dan
    $x^*Bx=|x_1+x_3|^2+|x_2|^2$.
    Namun, $A\circ B=I_3$. Jadi singularitas kedua faktor tidak
    memaksa singularitas hasil kali Hadamard.
\end{exam}

\section{Matriks Uniter dan Ortogonal}
\label{khusus:uniter}

\begin{definisi}[][Uniter dan ortogonal][khusus:def-uniter]
    Matriks kompleks $U$ disebut \emph{uniter} jika $U^*U=I$.
    Matriks real $Q$ disebut \emph{ortogonal} jika $Q^TQ=I$.
    Kolom-kolom matriks ini membentuk basis ortonormal.
\end{definisi}

Untuk matriks persegi, $U^*U=I$ ekuivalen dengan $U^{-1}=U^*$
dan $UU^*=I$. Jadi $U$ normal, mempertahankan hasil kali dalam,
dan $\|Ux\|=\|x\|$. Jika $Uv=\lambda v$, persamaan norma memberi
$|\lambda|=1$.
Sebaliknya, matriks \emph{normal} yang seluruh nilai eigennya
bermodulus satu pasti uniter, melalui dekomposisi spektral.
Tanpa normalitas, kesimpulan ini gagal: $\begin{pmatrix}1&1\\0&1\end{pmatrix}$
mempunyai nilai eigen $1,1$ tetapi tidak uniter.

\begin{prop}[Spektrum matriks ortogonal real]
    \label{khusus:ortogonal-real}
    Nilai eigen matriks ortogonal real terdiri atas $1$, $-1$,
    dan pasangan kompleks konjugat bermodulus satu.
    Determinannya $1$ atau $-1$. Untuk ukuran ganjil, $Q$
    mempunyai nilai eigen yang sama dengan $\det Q$.
\end{prop}
\begin{proof}
    Polinom karakteristik real memasangkan setiap nilai eigen
    tak real dengan konjugatnya. Hasil kali setiap pasangan adalah
    $|\lambda|^2=1$; nilai eigen real yang bermodulus satu hanya
    $1$ dan $-1$. Pada ukuran ganjil, banyaknya nilai eigen real
    ganjil. Jika determinan $-1$, ada sejumlah ganjil nilai eigen
    $-1$. Jika determinan $1$, banyaknya $-1$ genap, sehingga
    sedikitnya ada satu nilai eigen $1$.
\end{proof}

\begin{exam}[Refleksi terhadap komplemen ortogonal suatu vektor]
    Untuk $0\ne u\in\C^n$, definisikan
    \[
        P=\frac{uu^*}{u^*u},\qquad H=I-2P.
    \]
    Karena $P^*=P$ dan $P^2=P$, diperoleh $H^*=H$ serta $H^2=I$.
    Jadi $H$ Hermitian sekaligus uniter. Pada $\operatorname{span}\{u\}$
    matriks $H$ bekerja sebagai $-I$, sedangkan pada $u^\perp$
    sebagai $I$. Akibatnya $\det H=-1$ dan $\tr H=n-2$.
    Untuk $u$ real, ini merupakan matriks refleksi ortogonal.
\end{exam}

\section{Idempoten, Involusi, dan Nilpoten}
\label{khusus:polinomial}

\begin{teorema}[][Struktur idempoten dan involusi][khusus:idempoten-involusi]
    Untuk matriks di atas $\R$ atau $\C$:
    \begin{enumerate}[label=(\roman*),nosep]
        \item Jika $E^2=E$, maka $E$ serupa dengan $\diag(I_r,0)$,
            dengan $r=\rank E=\tr E$.
        \item Jika $Q^2=I$, maka $Q$ serupa dengan $\diag(I_p,-I_q)$,
            dengan $p+q=n$. Akibatnya
            \[
                \tr Q=p-q,\quad\det Q=(-1)^q,\quad
                \rank(Q-I)=q,\quad\rank(Q+I)=p.
            \]
        \item Hubungan $Q=2E-I$ merupakan korespondensi antara
            idempoten dan involusi.
    \end{enumerate}
\end{teorema}
\begin{proof}
    Polinom $t(t-1)$ dan $(t-1)(t+1)$ mempunyai akar-akar berbeda.
    Kriteria polinom minimal pada \autoref{eigen:squarefree}
    menjamin diagonalisasi, dengan nilai eigen masing-masing
    di dalam $\{0,1\}$ dan $\{-1,1\}$. Rumus-rumus mengikuti
    bentuk diagonalnya. Terakhir,
    $(2E-I)^2=I$ ekuivalen dengan $E^2=E$.
\end{proof}

Secara geometris, idempoten memproyeksikan ruang pada
$\operatorname{Im}E$ sepanjang $\ker E$, seperti dalam
\autoref{rank:idempoten}. Proyeksi ini ortogonal tepat ketika
$E$ Hermitian. Jika Hermitian, eigenspace untuk $0$ dan $1$
ortogonal. Sebaliknya, proyeksi ortogonal mempunyai matriks
$\diag(I_r,0)$ dalam suatu basis ortonormal, sehingga Hermitian.

\begin{prop}[Faktorisasi rank penuh suatu idempoten]
    \label{khusus:faktor-idempoten}
    Misalkan $E=UV$, dengan $U\in\Mat_{n\times r}(\F)$
    berank kolom penuh dan $V\in\Mat_{r\times n}(\F)$ berank
    baris penuh, $r>0$. Maka $E^2=E$ jika dan hanya jika $VU=I_r$.
\end{prop}
\begin{proof}
    Persamaan $E^2=E$ berarti $U(VU-I_r)V=0$.
    Kalikan dengan invers kiri $U$ dan invers kanan $V$ untuk
    memperoleh $VU=I_r$. Arah sebaliknya langsung dari perkalian.
\end{proof}

\begin{definisi}[][Matriks nilpoten][khusus:def-nilpoten]
    Matriks $N$ disebut \emph{nilpoten} jika $N^k=0$ untuk suatu
    $k\ge1$. Bilangan terkecil demikian disebut indeks nilpotensi.
\end{definisi}

Matriks kompleks nilpoten tepat ketika seluruh nilai eigennya nol.
Arah perlu mengikuti pemetaan spektral; arah cukup mengikuti
Cayley--Hamilton, sebab $\chi_N(t)=t^n$ mengakibatkan $N^n=0$.
Jadi indeks nilpotensi paling besar $n$.
Jika $N^k=0$, identitas geometri hingga memberi
\[
    (I-N)^{-1}=I+N+\cdots+N^{k-1}.
\]
Matriks normal yang nilpoten harus nol. Matriks nilpoten tak nol
karena itu selalu gagal menjadi normal dan gagal didiagonalkan.

\begin{prop}[Jumlah nilpoten yang komutatif]
    \label{khusus:nilpoten-komutatif}
    Jika $A^p=0$, $B^q=0$, dan $AB=BA$, maka
    $(A+B)^{p+q-1}=0$ dan $AB$ nilpoten.
\end{prop}
\begin{proof}
    Pada ekspansi binomial, setiap suku $A^jB^{p+q-1-j}$ nol:
    jika $j<p$, pangkat $B$ sedikitnya $q$.
    Selain itu, $(AB)^p=A^pB^p=0$.
\end{proof}

Tanpa komutativitas, keduanya dapat gagal. Untuk matriks satuan
$A=E_{12}$ dan $B=E_{21}$ berukuran dua, $A^2=B^2=0$,
tetapi $(A+B)^2=I$ dan $AB=\diag(1,0)$.

\section{Matriks Permutasi}
\label{khusus:permutasi}

\begin{definisi}[][Matriks suatu permutasi][khusus:def-permutasi]
    Untuk permutasi $\pi$ pada $\{1,\ldots,n\}$, matriks permutasi
    $P_\pi$ ditentukan oleh $P_\pi e_j=e_{\pi(j)}$.
    Dengan konvensi ini, kolom ke-$j$ adalah $e_{\pi(j)}$.
\end{definisi}

Matriks ini mempunyai tepat satu entri $1$ di tiap baris dan kolom;
entri lainnya nol. Selain itu,
$P_\pi P_\sigma=P_{\pi\circ\sigma}$ dan
$P_\pi^{-1}=P_{\pi^{-1}}=P_\pi^T$, sehingga matriks permutasi ortogonal.

\begin{teorema}[][Siklus menentukan spektrum dan rank][khusus:siklus]
    Misalkan siklus-siklus saling lepas dari $\pi$ mempunyai panjang
    $\ell_1,\ldots,\ell_c$, termasuk siklus panjang satu. Maka
    \begin{align*}
        \chi_{P_\pi}(t)&=\prod_{j=1}^c(t^{\ell_j}-1),&
        \det P_\pi&=(-1)^{n-c},\\
        \nul(P_\pi-I)&=c,&\rank(P_\pi-I)&=n-c.
    \end{align*}
    Untuk setiap bilangan bulat $k\ge1$,
    \[
        \nul(P_\pi^k-I)=\sum_{j=1}^c\gcd(\ell_j,k),
        \qquad
        \tr(P_\pi^k)=\sum_{\ell_j\mid k}\ell_j.
    \]
    Pangkat positif terkecil yang menghasilkan $I$ adalah
    $\operatorname{KPK}(\ell_1,\ldots,\ell_c)$.
\end{teorema}
\begin{proof}
    Susun basis standar mengikuti urutan siklus. Matriksnya menjadi
    diagonal blok, dengan satu blok per siklus. Pada siklus panjang
    $\ell$, vektor $e,P e,\ldots,P^{\ell-1}e$ membentuk basis dan
    $P^\ell e=e$. Jadi polinom minimal blok itu berderajat $\ell$
    dan sama dengan $t^\ell-1$, sekaligus polinom karakteristiknya.
    Determinan blok adalah $(-1)^{\ell-1}$.

    Persamaan $Px=x$ memaksa koordinat konstan pada setiap siklus,
    sehingga kernel $P-I$ berdimensi banyaknya siklus.
    Pangkat ke-$k$ dari satu siklus panjang $\ell$ terpecah menjadi
    $\gcd(\ell,k)$ siklus; ini memberi rumus nulitas untuk $P^k-I$.
    Trace menghitung titik tetap: satu siklus menyumbang seluruh
    $\ell$ titiknya tepat ketika $\ell\mid k$, dan tidak menyumbang
    titik tetap pada kasus lainnya. Rumus orde mengikuti syarat
    bahwa setiap panjang siklus harus membagi pangkat tersebut.
\end{proof}

Semua matriks permutasi dapat didiagonalkan di atas $\C$ karena
normal. Di atas $\R$, diagonalisasi mungkin tepat ketika setiap
siklus mempunyai panjang satu atau dua. Siklus dengan panjang
lebih dari dua menyumbang akar satuan yang tak real.

\begin{exam}[Spektrum dari siklus permutasi]
    Misalkan $P=P_\pi$ untuk permutasi
    \[
        \pi=(1\ 2\ 3\ 4)(5\ 6\ 7)(8\ 9).
    \]
    Tentukan orde dan determinan $P$, serta rank $P-I$ dan $P^6-I$.
\end{exam}
\begin{solution}
    Panjang siklusnya adalah $4,3,2$. Jadi orde $P$ adalah
    $\operatorname{KPK}(4,3,2)=12$. Selain itu,
    \[
        \det P=(-1)^{9-3}=1,\qquad \rank(P-I)=9-3=6.
    \]
    Selanjutnya
    \[
        \nul(P^6-I)=\gcd(4,6)+\gcd(3,6)+\gcd(2,6)=7,
    \]
    sehingga $\rank(P^6-I)=2$.
    Polinom karakteristiknya $(t^4-1)(t^3-1)(t^2-1)$:
    multiplisitas $1$ adalah tiga, sedangkan multiplisitas $-1$
    adalah dua. Semua titik pada siklus panjang tiga dan dua
    tetap di bawah $\pi^6$, sehingga $\tr(P^6)=5$.
\end{solution}

\section{Latihan Pilihan}
\label{khusus:latihan}

Pada soal-soal kompleks, tanda $*$ selalu berarti transpose
konjugat. Gunakan struktur matriks sebelum mencoba menghitung
polinom karakteristiknya secara langsung.

\begin{latihan}\label{khusus:lat-trace-hermitian}
    Buktikan bahwa untuk $A\in\Mat_n(\C)$,
    \[
        A=A^*\iff\tr(A^2)=\tr(A^*A).
    \]
    Apakah cukup mengganti syarat di kanan dengan
    $\tr(A^2)\ge0$?
\end{latihan}

\begin{latihan}\label{khusus:lat-nonhermitian}
    Berikan matriks real $2\times2$ yang semua nilai eigennya
    positif dan semua minor utamanya positif, tetapi bentuk
    kuadrat $x^TAx$ bernilai negatif untuk suatu $x\ne0$.
    Jelaskan hipotesis mana pada kriteria minor utama yang gagal.
\end{latihan}

\begin{latihan}\label{khusus:lat-komutan-kuadrat}
    Misalkan $A\succeq0$ dan $B$ matriks berukuran sama.
    Buktikan $A^2B=BA^2$ jika dan hanya jika $AB=BA$.
    Apakah kesimpulan masih berlaku jika $A$ hanya Hermitian?
\end{latihan}

\begin{latihan}\label{khusus:lat-order-kuadrat}
    Tentukan apakah $A\succeq B\succeq0$ selalu mengakibatkan
    $A^2\succeq B^2$. Jika tidak, berikan contoh $2\times2$ dan
    buktikan bahwa implikasinya benar apabila $AB=BA$.
\end{latihan}

\begin{latihan}\label{khusus:lat-variasi}
    Jika $A\succ0$ dan $x\in\C^n$, buktikan
    \[
        x^*A^{-1}x=
        \max_{y\in\C^n}\bigl(2\operatorname{Re}(x^*y)-y^*Ay\bigr).
    \]
    Tentukan vektor yang mencapai maksimum dan gunakan hasilnya
    untuk membuktikan kembali bahwa invers membalik urutan.
\end{latihan}

\begin{latihan}\label{khusus:lat-blok}
    Untuk $A,B$ Hermitian, buktikan
    \[
        \begin{pmatrix}A&B\\B&A\end{pmatrix}\succeq0
        \iff A+B\succeq0\ \text{dan}\ A-B\succeq0.
    \]
    Apakah $AB=BA$ diperlukan? Tentukan pula rank matriks blok
    tersebut dalam bentuk rank $A+B$ dan $A-B$.
\end{latihan}

\begin{latihan}\label{khusus:lat-rank-hadamard}
    Untuk $A,B\in\Mat_{m\times n}(\C)$, buktikan
    \[
        \rank(A\circ B)\le\rank A\,\rank B.
    \]
    Perlihatkan bahwa batas ini dapat tercapai dengan
    $\rank A=\rank B=2$ pada matriks berukuran $4\times4$.
\end{latihan}

\begin{latihan}\label{khusus:lat-lyapunov}
    Misalkan $A\succ0$ dan $B\succeq0$.
    Buktikan bahwa $AX+XA=B$ mempunyai tepat satu solusi $X$,
    dan solusi itu semidefinit positif.
    Jika $B\succ0$, buktikan $X\succ0$.
\end{latihan}

\begin{latihan}\label{khusus:lat-involusi}
    Misalkan $Q^2=I$. Buktikan bahwa pernyataan berikut ekuivalen:
    $Q$ normal, $Q$ Hermitian, dan $Q$ uniter.
    Berikan contoh involusi real $2\times2$ yang tidak normal.
\end{latihan}

\begin{latihan}\label{khusus:lat-idempoten}
    Misalkan $E^2=E$, $F^2=F$, dan $EF=FE$. Buktikan bahwa
    $EF$ dan $G=E+F-EF$ idempoten, dengan
    \[
        \operatorname{Im}(EF)=\operatorname{Im}E\cap\operatorname{Im}F,
        \qquad
        \operatorname{Im}G=\operatorname{Im}E+\operatorname{Im}F.
    \]
    Turunkan kedua identitas
    \begin{align*}
        \rank E+\rank F&=\rank(EF)+\rank G,\\
        \rank(E-F)&=\rank E+\rank F-2\rank(EF).
    \end{align*}
\end{latihan}

\begin{latihan}\label{khusus:lat-siklik}
    Misalkan $P$ matriks permutasi untuk siklus $(1\ 2\ \cdots\ n)$,
    dengan $n\ge2$, dan $J_n$ matriks yang seluruh entrinya $1$.
    Buktikan
    \[
        I+P+\cdots+P^{n-1}=J_n.
    \]
    Tentukan rank jumlah tersebut, rank $I-P$, dan buktikan
    bahwa $I,P,\ldots,P^{n-1}$ bebas linear dalam ruang matriks.
\end{latihan}

\begin{latihan}\label{khusus:lat-determinan-trace}
    Misalkan $A\succeq0$ dan $\tr A=s$, dengan $s\ge0$ tetap.
    Tentukan nilai maksimum $\det A$ beserta semua kasus kesamaannya.
    Selanjutnya, jika seluruh entri diagonal $A$ adalah $1$,
    buktikan $\det A\le1$ dan tentukan kapan kesamaan terjadi.
\end{latihan}

\ifpembahasan
\section{Pembahasan Latihan}
\label{khusus:pembahasan}

\subsection*{Latihan \ref{khusus:lat-trace-hermitian}: pisahkan bagian Hermitian}
Tuliskan $A=H+iK$ dengan $H,K$ Hermitian.
Karena $\tr(HK)=\tr(KH)$ real, diperoleh
\[
    \tr(A^2)=\tr(H^2)-\tr(K^2)+2i\tr(HK),\qquad
    \tr(A^*A)=\tr(H^2)+\tr(K^2).
\]
Kesamaan bagian real memaksa $\tr(K^2)=0$.
Karena $K=K^*$, ini sama dengan $\sum_{i,j}|k_{ij}|^2=0$,
sehingga $K=0$ dan $A$ Hermitian. Arah sebaliknya langsung.
Syarat $\tr(A^2)\ge0$ saja tidak cukup: matriks
$A=\begin{pmatrix}1&1\\0&1\end{pmatrix}$ mempunyai
$\tr(A^2)=2$ tetapi tidak Hermitian.

\subsection*{Latihan \ref{khusus:lat-nonhermitian}: spektrum positif belum cukup}
Ambil $A=\begin{pmatrix}1&3\\0&1\end{pmatrix}$.
Kedua nilai eigennya $1$, kedua minor utama berukuran satu
bernilai $1$, dan determinannya $1$.
Namun, untuk $x=(1,-1)^T$, diperoleh $x^TAx=-1$.
Hipotesis simetris, atau Hermitian pada kasus kompleks, tidak
terpenuhi. Kriteria minor utama tidak dapat dipakai tanpa hipotesis itu.

\subsection*{Latihan \ref{khusus:lat-komutan-kuadrat}: kuadrat injektif pada spektrum positif}
Diagonalkan $A=UDU^*$ dengan $D=\diag(\lambda_1,\ldots,\lambda_n)$,
$\lambda_i\ge0$, dan tuliskan $C=U^*BU$.
Persamaan $D^2C=CD^2$ berarti
$(\lambda_i^2-\lambda_j^2)c_{ij}=0$ untuk semua $i,j$.
Karena kedua nilai eigen tak negatif, $\lambda_i^2=\lambda_j^2$
ekuivalen dengan $\lambda_i=\lambda_j$.
Jadi $DC=CD$, yang memberikan $AB=BA$.
Arah sebaliknya berlaku untuk sembarang $A$.
Jika hanya Hermitian, ambil $A=\diag(1,-1)$ dan
$B=\begin{pmatrix}0&1\\1&0\end{pmatrix}$.
Matriks $A^2=I$ komutatif dengan $B$, tetapi $AB\ne BA$.

\subsection*{Latihan \ref{khusus:lat-order-kuadrat}: urutan tidak otomatis bertahan}
Ambil
\[
    B=\begin{pmatrix}1&0\\0&0\end{pmatrix},\qquad
    A=\begin{pmatrix}2&1\\1&1\end{pmatrix}.
\]
Matriks $B$ semidefinit positif dan
$A-B=\begin{pmatrix}1&1\\1&1\end{pmatrix}\succeq0$.
Namun,
\[
    A^2-B^2=\begin{pmatrix}4&3\\3&2\end{pmatrix}
\]
mempunyai determinan $-1$, sehingga tidak semidefinit positif.
Jika $AB=BA$, matriks $A-B$ dan $A+B$ komutatif dan keduanya
semidefinit positif. Oleh \autoref{khusus:produk-positif},
$A^2-B^2=(A-B)(A+B)\succeq0$.

\subsection*{Latihan \ref{khusus:lat-variasi}: melengkapi kuadrat}
Untuk $z=y-A^{-1}x$, ekspansi langsung memberi
\[
    2\operatorname{Re}(x^*y)-y^*Ay=x^*A^{-1}x-z^*Az.
\]
Karena $A\succ0$, maksimum dicapai tepat saat $z=0$, yaitu
$y=A^{-1}x$. Jika $A\succeq B\succ0$, untuk setiap $y$ berlaku
\[
    2\operatorname{Re}(x^*y)-y^*Ay
    \le 2\operatorname{Re}(x^*y)-y^*By.
\]
Mengambil maksimum memberi $x^*A^{-1}x\le x^*B^{-1}x$ untuk
setiap $x$, yakni $A^{-1}\preceq B^{-1}$.

\subsection*{Latihan \ref{khusus:lat-blok}: jumlah dan selisih dalam basis ortonormal}
Matriks
\[
    U=\frac1{\sqrt2}\begin{pmatrix}I&I\\I&-I\end{pmatrix}
\]
uniter dan memenuhi
\[
    U^*\begin{pmatrix}A&B\\B&A\end{pmatrix}U
       =\begin{pmatrix}A+B&0\\0&A-B\end{pmatrix}.
\]
Kongruensi ini membuktikan ekuivalensi positivitas tanpa
memerlukan $AB=BA$. Karena $U$ invertibel, rank matriks blok
adalah $\rank(A+B)+\rank(A-B)$.

\subsection*{Latihan \ref{khusus:lat-rank-hadamard}: uraikan menjadi matriks rank satu}
Jika salah satu faktor nol, hasilnya langsung. Untuk
$r=\rank A>0$ dan $s=\rank B>0$, gunakan faktorisasi rank penuh
untuk menulis
$A=\sum_{i=1}^r u_iv_i^T$ dan $B=\sum_{j=1}^s x_jy_j^T$.
Dengan perkalian entri,
\[
    A\circ B=\sum_{i=1}^r\sum_{j=1}^s
       (u_i\circ x_j)(v_i\circ y_j)^T.
\]
Setiap suku berank paling besar satu, sehingga rank jumlahnya
paling besar $rs$. Untuk ketajaman, ambil
\[
    A=\begin{pmatrix}1&1&0&0\\1&1&0&0\\0&0&1&1\\0&0&1&1\end{pmatrix},
    \qquad
    B=\begin{pmatrix}1&0&1&0\\0&1&0&1\\1&0&1&0\\0&1&0&1\end{pmatrix}.
\]
Kedua matriks berank dua, tetapi $A\circ B=I_4$ berank empat.

\subsection*{Latihan \ref{khusus:lat-lyapunov}: uji pada vektor eigen solusi}
Nilai eigen $A$ positif, sedangkan nilai eigen $-A$ negatif.
Teorema persamaan Sylvester (\autoref{eigen:sylvester}) memberi
tepat satu solusi untuk $AX-X(-A)=B$.
Mengambil transpose konjugat menunjukkan $X^*$ juga solusi;
keunikan memaksa $X=X^*$.
Untuk vektor eigen $Xv=\mu v$, $v\ne0$, diperoleh
\[
    v^*Bv=v^*(AX+XA)v=2\mu\,v^*Av.
\]
Karena $v^*Av>0$ dan $v^*Bv\ge0$, setiap $\mu\ge0$.
Jadi $X\succeq0$. Jika $B\succ0$, ruas kiri positif untuk
setiap vektor eigen, sehingga setiap $\mu>0$ dan $X\succ0$.

\subsection*{Latihan \ref{khusus:lat-involusi}: spektrum hanya pada dua titik}
Nilai eigen involusi ada di $\{-1,1\}$.
Jika $Q$ normal, teorema spektral memberi $Q=UDU^*$ dengan
$D$ diagonal real; maka $Q$ Hermitian.
Jika $Q$ Hermitian, $Q^*Q=Q^2=I$, sehingga uniter.
Matriks uniter selalu normal, menutup rangkaian implikasi.
Contoh involusi yang tidak normal adalah
$Q=\begin{pmatrix}1&1\\0&-1\end{pmatrix}$:
$Q^2=I$, tetapi $Q^T\ne Q$, sehingga ekuivalensi di atas
menjamin $Q$ tidak normal.

\subsection*{Latihan \ref{khusus:lat-idempoten}: empat ruang eigen bersama}
Idempoten dapat didiagonalkan, dan $E,F$ komutatif.
Diagonalisasi simultan (\autoref{eigen:simultan}) memecah ruang
menjadi empat ruang $V_{ab}$, $a,b\in\{0,1\}$, tempat $E$
bekerja sebagai $aI$ dan $F$ sebagai $bI$.
Pada $V_{ab}$, matriks $EF$ bekerja sebagai $abI$,
sedangkan $G$ sebagai $(a+b-ab)I$.
Kedua skalar terakhir bernilai $0$ atau $1$, sehingga kedua
matriks idempoten. Bayangan $EF$ adalah $V_{11}$, sedangkan
bayangan $G$ adalah $V_{10}\oplus V_{01}\oplus V_{11}$.
Ini sekaligus membuktikan pernyataan tentang irisan dan jumlah.
Dengan $d_{ab}=\dim V_{ab}$,
\begin{align*}
    \rank E&=d_{10}+d_{11},&\rank F&=d_{01}+d_{11},\\
    \rank(EF)&=d_{11},&\rank G&=d_{10}+d_{01}+d_{11}.
\end{align*}
Selain itu, $E-F$ bernilai $1$ pada $V_{10}$, $-1$ pada
$V_{01}$, dan nol pada dua ruang lainnya.
Jadi $\rank(E-F)=d_{10}+d_{01}$, yang memberi kedua identitas rank.

\subsection*{Latihan \ref{khusus:lat-siklik}: orbit satu vektor standar}
Untuk setiap $e_j$, vektor $e_j,Pe_j,\ldots,P^{n-1}e_j$
menelusuri semua vektor basis standar tepat satu kali.
Jadi $(I+P+\cdots+P^{n-1})e_j=(1,\ldots,1)^T$ dan jumlah
matriksnya $J_n$, dengan rank satu.
Persamaan $Px=x$ mengharuskan semua koordinat $x$ sama,
sehingga $\rank(I-P)=n-1$.
Jika $\sum_{k=0}^{n-1}a_kP^k=0$, terapkan pada $e_1$.
Karena $e_1,Pe_1,\ldots,P^{n-1}e_1$ merupakan basis,
semua $a_k=0$. Maka pangkat-pangkat tersebut bebas linear.

\subsection*{Latihan \ref{khusus:lat-determinan-trace}: dua cara membatasi determinan}
Tuliskan nilai eigen $A$ sebagai $\lambda_i\ge0$ dengan
$\sum_i\lambda_i=s$. Rata-rata aritmetika--geometri memberi
\[
    \det A=\prod_i\lambda_i\le(s/n)^n.
\]
Jika $s>0$, kesamaan terjadi tepat ketika semua
$\lambda_i=s/n$, yakni $A=(s/n)I$.
Jika $s=0$, semua nilai eigen nol dan $A=0$, sehingga rumus
serta karakterisasi yang sama tetap berlaku.
Jika entri diagonal semuanya $1$, ketaksamaan Hadamard memberi
$\det A\le1$. Karena tidak ada entri diagonal nol, kesamaan
memaksa $A$ diagonal, yaitu $A=I$.
\fi


\part{Ruang Vektor}{``Vektor adalah objek yang memiliki besaran dan arah'', kalimat yang sering didengar ketika belajar fisika. Namun, obsesi untuk mengabstraksi suatu objek dalam matematika menimbulkan pergeseran makna untuk sebuah objek.\\

``Ruang'' diartikan sebagai himpunan yang memiliki struktur tertentu biasanya dapat berupa operasi, relasi, atau sifat-sifat aturan tertentu. Jadi secara singkatnya, jika kita membicarakan vektor pada konteks ini adalah anggota dari ruang vektor.}
% !TeX root = ../../ONMIPA Alin.tex
% Acuan: Zhang (2009), Bab 1, khususnya Soal 1.1--1.7, 1.13--1.14,
% 1.22, 1.38, 1.41, 1.43, 1.44, dan 1.51; Zhang (2011), bagian 1.1.
% Pembahasan dibatasi sampai subruang; tidak membahas basis dan dimensi.
\chapter{Ruang Linear}
\label{ch:ruang-linear}

\textbf{Ruang vektor} atau bisa disebut \textbf{ruang linear}  merupakan himpunan yang dilengkapi dengan dua operasi, yaitu penjumlahan dan perkalian skalar. Menurut penulis, materi ini cukup penting dan bisa dipelajari untuk awal-awal karena tidak terlalu me-\textit{reference} definisi-definisi lain. Materi ini juga menjadi dasar untuk memahami materi-materi selanjutnya, seperti kebebasan linear, basis, dimensi, dan transformasi linear.

\section{Definisi Ruang Vektor}
\label{linear:aksioma}
Anggota $V$ disebut \emph{vektor}, sedangkan anggota medan $\F$
disebut \emph{skalar}. Contoh medan adalah $\R$, $\C$, dan
bilangan rasional $\mathbb Q$. Setiap skalar tak nol dalam suatu
medan mempunyai invers perkalian.

\begin{definisi}[][Ruang linear atas suatu medan][linear:def-ruang]
    Himpunan tak kosong $V$ disebut \emph{ruang linear atas $\F$}
    jika dilengkapi operasi
    \[
        +:V\times V\longrightarrow V,
        \qquad
        \cdot:\F\times V\longrightarrow V,
    \]
    yang memenuhi aksioma berikut untuk setiap $u,v,w\in V$
    dan $\alpha,\beta\in\F$:
    \begin{enumerate}[label=(\roman*),nosep]
        \item $u+v=v+u$;
        \item $(u+v)+w=u+(v+w)$;
        \item terdapat $0_V\in V$ sehingga $v+0_V=v$;
        \item untuk setiap $v$ terdapat $-v\in V$ sehingga $v+(-v)=0_V$;
        \item $\alpha(u+v)=\alpha u+\alpha v$;
        \item $(\alpha+\beta)v=\alpha v+\beta v$;
        \item $(\alpha\beta)v=\alpha(\beta v)$;
        \item $1v=v$, dengan $1$ identitas perkalian dalam $\F$.
    \end{enumerate}
\end{definisi}

Penulisan kedua operasi sebagai pemetaan ke $V$ sudah mencakup
\emph{ketertutupan}: $u+v$ dan $\alpha v$ harus tetap berada di $V$.
Jika ketertutupan ditulis sebagai dua syarat tersendiri, daftar di
atas menjadi sepuluh syarat. Kedua cara penulisan itu ekuivalen.

Vektor nol $0_V$ dan skalar nol $0\in\F$ mempunyai peran berbeda.
Pada $\R^n$, vektor nol adalah $(0,\ldots,0)^T$; pada ruang matriks,
ia adalah matriks nol; pada ruang fungsi, ia adalah fungsi yang
bernilai nol di setiap titik. Setelah jenis objeknya jelas, keduanya
sering ditulis dengan simbol $0$.

Untuk $v_1,\ldots,v_r\in V$ dan $\alpha_1,\ldots,\alpha_r\in\F$,
ekspresi $\alpha_1v_1+\cdots+\alpha_rv_r$ disebut
\emph{kombinasi linear} dari vektor-vektor tersebut.
Penjumlahan ini mempunyai hingga banyak suku dan selalu
menghasilkan anggota $V$ karena ketertutupan kedua operasi.

\begin{exam}[Medan skalar harus disebutkan]
    Tentukan apakah $\C$ merupakan ruang linear atas $\C$ dan atas
    $\R$. Apakah $\R$ merupakan ruang linear atas $\C$ dengan
    penjumlahan dan perkalian biasa?
\end{exam}
\begin{solution}
    Himpunan $\C$ merupakan ruang linear atas kedua medan tersebut:
    penjumlahan kompleks dan perkalian oleh skalar yang diperbolehkan
    tetap menghasilkan bilangan kompleks, dan aksiomanya mengikuti
    sifat aritmetika kompleks.
    Namun, $\R$ tidak tertutup terhadap perkalian oleh skalar kompleks,
    sebab $i\cdot1=i\notin\R$. Dengan operasi yang disebutkan,
    $\R$ bukan ruang linear atas $\C$.
\end{solution}

Secara umum, ruang linear atas $\C$ juga dapat dipandang sebagai
ruang linear atas $\R$ dengan membatasi skalar yang digunakan.
Arah sebaliknya memerlukan definisi perkalian oleh skalar kompleks
dan tidak berlaku secara otomatis.

\section{Contoh \emph{Well-Known}}
\label{linear:contoh}

\subsection{Koordinat, matriks, polinom, fungsi, dan barisan}

\begin{enumerate}[label=(\alph*)]
    \item \textbf{Vektor koordinat.} Himpunan $\F^n$ merupakan
          ruang linear dengan penjumlahan dan perkalian skalar per
          koordinat. Untuk perkalian $Ax$, vektor $x$ ditulis sebagai kolom.
    \item \textbf{Matriks berukuran tetap.} Himpunan
          $\Mat_{m\times n}(\F)$ merupakan ruang linear atas $\F$.
          Ukuran harus tetap agar penjumlahan terdefinisi.
          Perkalian dua matriks bukan bagian dari aksioma ruang linear.
    \item \textbf{Polinom.} Himpunan $\F[t]$ berisi semua polinom
          dengan koefisien dalam $\F$. Untuk $n\ge0$, tuliskan
          \[
              \mathcal P_{\le n}(\F)
              =\{a_0+a_1t+\cdots+a_nt^n:a_j\in\F\}.
          \]
          Polinom nol ikut termuat. Penjumlahan dan perkalian skalar
          tidak menaikkan derajat melampaui $n$, sehingga himpunan
          ini juga merupakan ruang linear. Notasi $\le n$ menegaskan
          bahwa batas derajatnya termasuk $n$.
    \item \textbf{Fungsi.} Untuk himpunan tak kosong $S$,
          himpunan semua fungsi $S\to\F$, ditulis $\F^S$,
          mempunyai operasi titik demi titik:
          \[
              (f+g)(s)=f(s)+g(s),\qquad
              (\alpha f)(s)=\alpha f(s).
          \]
          Aksioma berlaku karena dapat diperiksa pada setiap $s\in S$.
          Himpunan $C([a,b],\R)$ dari fungsi real kontinu pada
          $[a,b]$, $a<b$, juga merupakan ruang linear, sebab jumlah
          dan kelipatan skalar fungsi kontinu tetap kontinu.
    \item \textbf{Barisan.} Himpunan $\F^{\mathbb N}$ dari semua
          barisan $(a_1,a_2,\ldots)$ merupakan ruang linear dengan
          operasi per suku. Tidak ada syarat kekonvergenan dalam
          definisi ruang ini.
\end{enumerate}

Pada contoh-contoh tersebut, objeknya berbeda tetapi mekanisme
pembuktiannya sama: operasi diwarisi dari aritmetika dalam $\F$,
lalu diterapkan pada setiap koordinat, entri, koefisien, atau nilai fungsi.

\subsection{Bilangan positif sebagai vektor}

\begin{exam}[Penjumlahan berupa perkalian]
    Pada $V=\R_{>0}$, definisikan
    \[
        u\oplus v=uv,\qquad \alpha\odot u=u^\alpha
        \quad(u,v>0,\ \alpha\in\R).
    \]
    Buktikan bahwa $V$ merupakan ruang linear atas $\R$.
    Tentukan vektor nol dan invers penjumlahan dari $u$.
\end{exam}
\begin{solution}
    Kedua operasi tertutup karena hasilnya tetap positif.
    Sifat komutatif dan asosiatif $\oplus$ mengikuti perkalian real.
    Identitas penjumlahannya adalah $1$, dan invers penjumlahan $u$
    adalah $u^{-1}$, sebab $u\oplus1=u$ dan $u\oplus u^{-1}=1$.
    Hukum pangkat memberi
    \begin{align*}
        \alpha\odot(u\oplus v)
         & =(uv)^\alpha=u^\alpha v^\alpha
        =(\alpha\odot u)\oplus(\alpha\odot v), \\
        (\alpha+\beta)\odot u
         & =u^{\alpha+\beta}=u^\alpha u^\beta
        =(\alpha\odot u)\oplus(\beta\odot u),  \\
        (\alpha\beta)\odot u
         & =(u^\beta)^\alpha
        =\alpha\odot(\beta\odot u),
    \end{align*}
    serta $1\odot u=u$. Jadi seluruh aksioma terpenuhi.
    Di ruang ini, $0\odot u=1$ memang benar: bilangan $1$
    adalah vektor nol untuk operasi $\oplus$.
\end{solution}

Fungsi $u\mapsto\log u$ membantu membaca struktur tadi:
$\log(u\oplus v)=\log u+\log v$ dan
$\log(\alpha\odot u)=\alpha\log u$.
Dengan operasi biasa, $\R_{>0}$ tidak menjadi ruang linear atas
$\R$, sebab perkalian oleh skalar nol atau negatif keluar dari himpunan.

\subsection{Vektor nol dapat berada di titik lain}

\begin{exam}[Menggeser operasi pada ruang koordinat]
    Tetapkan $a\in\F^n$. Pada himpunan $\F^n$, definisikan
    \[
        u\oplus v=u+v-a,\qquad
        \alpha\odot u=a+\alpha(u-a).
    \]
    Tunjukkan bahwa operasi ini membentuk ruang linear atas $\F$.
\end{exam}
\begin{solution}
    Pemetaan $T(u)=u-a$ bijektif dari $\F^n$ ke dirinya dan memenuhi
    \[
        T(u\oplus v)=T(u)+T(v),\qquad
        T(\alpha\odot u)=\alpha T(u).
    \]
    Maka setiap identitas aksioma, setelah dikenai $T$, menjadi
    identitas pada ruang koordinat biasa. Karena $T$ injektif,
    identitas sebelum dikenai $T$ juga berlaku.
    Vektor nolnya adalah $a$, sebab $u\oplus a=u$.
    Invers penjumlahan dari $u$ adalah $2a-u$.
\end{solution}

Kedua contoh tersebut menunjukkan bahwa himpunan saja belum
menentukan suatu ruang linear. Medan dan kedua operasi harus diketahui.

\section{Definisi dan Kriteria Subruang}
\label{linear:subruang}

\begin{definisi}[][Subruang linear][linear:def-subruang]
    Misalkan $V$ ruang linear atas $\F$. Himpunan $W\subseteq V$
    disebut \emph{subruang linear} dari $V$ jika $W$ sendiri
    merupakan ruang linear atas $\F$ dengan penjumlahan dan
    perkalian skalar yang diwarisi dari $V$.
\end{definisi}

Frasa ``diwarisi dari $V$'' penting: kita tidak boleh mengganti
operasi atau medan skalar ketika menguji apakah $W$ subruang.
Misalnya, ruang bilangan positif dengan operasi $\oplus,\odot$
di atas bukan subruang dari $\R$ yang menggunakan operasi biasa.

\begin{teorema}[][Uji subruang][linear:uji-subruang]
    Misalkan $W\subseteq V$. Pernyataan berikut ekuivalen:
    \begin{enumerate}[label=(\roman*),nosep]
        \item $W$ merupakan subruang dari $V$;
        \item $0_V\in W$, $u+v\in W$ untuk semua $u,v\in W$,
              dan $\alpha u\in W$ untuk semua $\alpha\in\F$, $u\in W$;
        \item $W\ne\varnothing$ dan
              $\alpha u+\beta v\in W$ untuk semua $u,v\in W$
              serta $\alpha,\beta\in\F$.
    \end{enumerate}
\end{teorema}
\begin{proof}
    Jika $W$ subruang, vektor nolnya sama dengan vektor nol $V$:
    identitas $w+0_W=w$ dalam $V$ dan pencoretan memberi $0_W=0_V$.
    Ketertutupan kedua operasi mengikuti definisi, sehingga (i)
    mengakibatkan (ii). Dari (ii), perkalian skalar dan penjumlahan
    berurutan memberi (iii).

    Sekarang anggap (iii). Ambil satu $w\in W$.
    Dengan $\alpha=\beta=0$ diperoleh $0_V\in W$.
    Pilihan $\alpha=\beta=1$ memberi ketertutupan penjumlahan,
    sedangkan $\beta=0$ memberi ketertutupan perkalian skalar.
    Khususnya, $-u=(-1)u\in W$ untuk setiap $u\in W$.
    Aksioma asosiatif, komutatif, distributif, dan identitas skalar
    sudah berlaku di $V$, sehingga juga berlaku pada anggota $W$.
    Dengan demikian $W$ merupakan ruang linear.
\end{proof}

Syarat tak kosong tidak boleh dihilangkan. Pernyataan
``untuk setiap $u,v\in W$'' otomatis benar apabila $W$ kosong,
tetapi himpunan kosong tidak mempunyai vektor nol.
Subruang $\{0_V\}$ dan $V$ selalu ada. Subruang $W$ disebut
\emph{sejati} jika $W\ne V$; dengan pengertian ini,
$\{0_V\}$ juga sejati apabila $V\ne\{0_V\}$.

\begin{prop}[Subruang di dalam subruang]
    Jika $U$ subruang dari $W$ dan $W$ subruang dari $V$,
    semuanya atas medan yang sama, maka $U$ subruang dari $V$.
\end{prop}
\begin{proof}
    Inklusi $U\subseteq W\subseteq V$ berlaku dan kedua operasi
    yang digunakan sama. Vektor nol berada di $U$, dan $U$
    tertutup terhadap kedua operasi; terapkan uji subruang.
\end{proof}

\begin{exam}[Garis homogen dan garis bergeser]
    Untuk $c\in\R$, tentukan kapan
    $W_c=\{(x,y)\in\R^2:2x-3y=c\}$ merupakan subruang.
\end{exam}
\begin{solution}
    Jika $c\ne0$, vektor $(0,0)$ tidak berada di $W_c$, sehingga
    $W_c$ bukan subruang.
    Jika $c=0$ dan $u=(x_1,y_1),v=(x_2,y_2)\in W_0$, maka
    \[
        2(\alpha x_1+\beta x_2)-3(\alpha y_1+\beta y_2)
        =\alpha(2x_1-3y_1)+\beta(2x_2-3y_2)=0.
    \]
    Karena $(0,0)\in W_0$, uji subruang berlaku.
    Jadi syaratnya tepat $c=0$.
\end{solution}

\begin{exam}[Dua jenis ketertutupan yang berbeda]
    Himpunan $\Z^2\subset\R^2$ tertutup terhadap penjumlahan,
    tetapi tidak terhadap semua skalar real: $(1,0)\in\Z^2$,
    sedangkan $\tfrac12(1,0)\notin\Z^2$.
    Sebaliknya,
    \[
        S=\{(x,y)\in\R^2:xy=0\}
    \]
    tertutup terhadap perkalian skalar, tetapi tidak terhadap penjumlahan:
    \[
        (1,0),(0,1)\in S,\qquad (1,0)+(0,1)\notin S.
    \]
    Kedua himpunan memuat nol, tetapi keduanya bukan subruang real.
\end{exam}

\section{Mengenali Subruang dari Syarat Pembentuknya}
\label{linear:kendala}

\subsection{Persamaan homogen pada vektor dan matriks}

\begin{prop}[Himpunan solusi sistem homogen]
    \label{linear:kernel}
    Untuk $A\in\Mat_{m\times n}(\F)$, himpunan
    \[
        W=\{x\in\F^n:Ax=0\}
    \]
    merupakan subruang dari $\F^n$.
    Lebih umum, untuk $A\in\Mat_{m\times n}(\F)$ dan
    $B\in\Mat_{p\times q}(\F)$ tetap,
    \[
        \{X\in\Mat_{n\times p}(\F):AXB=0\}
    \]
    merupakan subruang dari $\Mat_{n\times p}(\F)$.
\end{prop}
\begin{proof}
    Vektor atau matriks nol memenuhi persamaan masing-masing.
    Jika $Au=Av=0$, maka
    $A(\alpha u+\beta v)=\alpha Au+\beta Av=0$.
    Untuk versi matriks, jika $AXB=AYB=0$, maka
    \[
        A(\alpha X+\beta Y)B=\alpha AXB+\beta AYB=0.
    \]
    Kedua kesimpulan mengikuti uji subruang.
\end{proof}

Himpunan pada bagian pertama adalah $\ker A$ yang telah muncul
pada bab rank dan nulitas. Untuk sistem tak homogen $Ax=b$
dengan $b\ne0$, himpunan solusinya bukan subruang dari $\F^n$,
karena vektor nol tidak memenuhi persamaan. Jika tidak ada solusi,
himpunannya kosong dan juga bukan subruang.

\begin{exam}[Parameter dalam suatu subruang]
    Buktikan bahwa
    \[
        W=\{(x_1,x_2,x_3,x_4)\in\F^4:
        x_3=x_1+x_2,\ x_4=x_1-x_2\}
    \]
    merupakan subruang. Tentukan semua anggotanya dalam bentuk parameter.
\end{exam}
\begin{solution}
    Kedua syarat adalah persamaan linear homogen, sehingga
    \autoref{linear:kernel} dapat diterapkan.
    Dengan $s=x_1$ dan $t=x_2$, seluruh anggota $W$ berbentuk
    \[
        (s,t,s+t,s-t),\qquad s,t\in\F.
    \]
    Ini juga memberi pemeriksaan langsung: kombinasi dengan skalar
    $\alpha,\beta$ dari anggota berparameter $(s_1,t_1)$ dan
    $(s_2,t_2)$ mempunyai parameter
    $(\alpha s_1+\beta s_2,\alpha t_1+\beta t_2)$, sehingga tetap di $W$.
\end{solution}

\subsection{Kendala pada polinom dan fungsi}

Evaluasi di titik tetap menghormati kedua operasi:
$(\alpha p+\beta q)(a)=\alpha p(a)+\beta q(a)$.
Karena itu, syarat seperti $p(a)=0$, $p(a)=p(b)$, atau
$2p(0)-p(1)=0$ menghasilkan subruang polinom.
Demikian pula, jika turunan terdefinisi, sifat
$(\alpha p+\beta q)'=\alpha p'+\beta q'$ dapat digunakan.

\begin{exam}[Mengubah syarat polinom menjadi syarat koefisien]
    Dalam $\mathcal P_{\le3}(\R)$, tentukan seluruh polinom
    yang memenuhi $p(0)=p(1)$ dan $p'(0)=0$.
    Buktikan bahwa himpunannya merupakan subruang.
\end{exam}
\begin{solution}
    Tuliskan $p(t)=a+bt+ct^2+dt^3$.
    Syarat $p'(0)=0$ memberi $b=0$, sedangkan $p(0)=p(1)$
    memberi $b+c+d=0$. Jadi $c=-d$ dan
    \[
        p(t)=a+d(t^3-t^2),\qquad a,d\in\R.
    \]
    Polinom nol diperoleh saat $a=d=0$.
    Penjumlahan dan perkalian skalar pada polinom hanya menjumlahkan
    dan mengalikan kedua parameter, sehingga himpunan tersebut subruang.
\end{solution}

Berbeda dari derajat \emph{paling besar} $n$, syarat derajat
\emph{tepat} $n$ tidak membentuk subruang, sebab polinom nol
tidak memenuhi syarat itu. Syarat $p(0)=1$ juga langsung gagal
pada polinom nol. Sementara itu, polinom yang tak negatif pada
seluruh $\R$ memuat polinom nol, tetapi tidak tertutup terhadap
skalar negatif: polinom $1$ memenuhi syarat, sedangkan $-1$ tidak.

\begin{prop}[Fungsi genap dan fungsi ganjil]
    \label{linear:paritas}
    Di dalam ruang semua fungsi $\R\to\F$, masing-masing himpunan
    \begin{align*}
        E & =\{f:f(-t)=f(t)\text{ untuk setiap }t\in\R\}, \\
        O & =\{f:f(-t)=-f(t)\text{ untuk setiap }t\in\R\}
    \end{align*}
    merupakan subruang atas $\F$.
\end{prop}
\begin{proof}
    Fungsi nol termasuk kedua himpunan. Jika $f,g\in E$,
    maka untuk setiap $t$,
    \[
        (\alpha f+\beta g)(-t)
        =\alpha f(t)+\beta g(t)=(\alpha f+\beta g)(t).
    \]
    Jadi $\alpha f+\beta g\in E$. Jika $f,g\in O$,
    perhitungan yang sama memberi
    $(\alpha f+\beta g)(-t)=-(\alpha f+\beta g)(t)$.
\end{proof}

Syarat itu harus berlaku untuk \emph{setiap} $t$, bukan hanya
pada sejumlah titik yang dipilih. Jika fungsi juga dipersyaratkan
kontinu, pembuktian tetap berlaku karena kombinasi linear hingga
dari fungsi kontinu tetap kontinu.

\subsection{Ketergantungan pada medan skalar}

Untuk matriks simetris, transpose memenuhi
$(\alpha A+\beta B)^T=\alpha A^T+\beta B^T$.
Jadi matriks simetris dalam $\Mat_n(\F)$ membentuk subruang atas
$\F$, baik ketika $\F=\R$ maupun ketika $\F=\C$.

Untuk transpose konjugat, skalar ikut dikonjugatkan:
\[
    (\alpha A+\beta B)^*=
    \overline\alpha A^*+\overline\beta B^*.
\]
Akibatnya, matriks Hermitian membentuk subruang \emph{real}
dari $\Mat_n(\C)$ yang dipandang sebagai ruang atas $\R$.
Namun, himpunan itu bukan subruang kompleks: $I$ Hermitian,
sedangkan $(iI)^*=-iI\ne iI$.

\begin{exam}[Subruang kompleks yang seluruh anggotanya Hermitian]
    Misalkan $W$ subruang kompleks dari $\Mat_n(\C)$ dan setiap
    $A\in W$ Hermitian. Buktikan $W=\{0\}$.
\end{exam}
\begin{solution}
    Untuk setiap $A\in W$, ketertutupan terhadap skalar kompleks
    memberi $iA\in W$. Karena $A$ dan $iA$ keduanya Hermitian,
    \[
        iA=(iA)^*=-iA^*=-iA.
    \]
    Jadi $2iA=0$ dan $A=0$. Karena $0\in W$, diperoleh $W=\{0\}$.
\end{solution}

\subsection{Syarat nonlinear perlu diperiksa langsung}

Tidak semua syarat yang tampak homogen menghasilkan subruang.
Contohnya, matriks $2\times2$ berank paling besar satu tertutup
terhadap perkalian skalar, tetapi
\[
    \rank\begin{pmatrix}1&0\\0&0\end{pmatrix}
    =\rank\begin{pmatrix}0&0\\0&1\end{pmatrix}=1,
    \qquad \rank I_2=2.
\]
Himpunan tersebut tidak tertutup terhadap penjumlahan.
Sebaliknya, suatu persamaan nonlinear dapat saja mendefinisikan
subruang: dalam $\R^2$, syarat $x^2+y^2=0$ hanya dipenuhi oleh
$(0,0)$, sehingga himpunan solusinya merupakan subruang.
Yang harus diperiksa adalah himpunan solusinya, bukan sekadar
bentuk persamaan yang dipakai untuk menuliskannya.

\section{Irisan dan Gabungan Subruang}
\label{linear:irisan-gabungan}

\begin{teorema}[][Irisan sebarang keluarga subruang][linear:irisan]
    Jika setiap $W_j$, $j\in J$, merupakan subruang dari $V$,
    maka $\bigcap_{j\in J}W_j$ merupakan subruang dari $V$.
\end{teorema}
\begin{proof}
    Vektor nol berada di setiap $W_j$, sehingga berada di irisan.
    Jika $u,v$ berada di irisan, maka keduanya berada di setiap
    $W_j$. Ketertutupan setiap $W_j$ memberi
    $\alpha u+\beta v\in W_j$ untuk setiap $j$, sehingga kombinasi
    itu berada di irisan. Untuk keluarga kosong, irisan relatif
    terhadap semesta $V$ diartikan sebagai $V$ sendiri.
\end{proof}

Sebagai contoh, polinom yang memenuhi $p(0)=p(1)=p(2)=0$
merupakan irisan tiga subruang yang masing-masing ditentukan
oleh satu syarat evaluasi. Hasil ini juga berlaku untuk tak hingga
banyak syarat: irisan himpunan $\{p\in\R[t]:p(a)=0\}$ atas
semua $a\in\R$ adalah subruang yang hanya berisi polinom nol.
Untuk fungsi pada \autoref{linear:paritas}, berlaku $E\cap O=\{0\}$,
sebab fungsi yang sekaligus genap dan ganjil memenuhi $f(t)=-f(t)$.

\begin{teorema}[][Gabungan dua subruang][linear:gabungan-dua]
    Jika $U,W$ subruang dari $V$, maka
    \[
        U\cup W\text{ subruang}
        \iff U\subseteq W\ \text{atau}\ W\subseteq U.
    \]
\end{teorema}
\begin{proof}
    Jika salah satu termuat di dalam yang lain, gabungannya adalah
    subruang yang lebih besar.
    Sebaliknya, andaikan tidak ada inklusi di antara keduanya.
    Pilih $u\in U\setminus W$ dan $w\in W\setminus U$.
    Jika gabungannya subruang, $u+w$ harus berada di $U\cup W$.
    Jika $u+w\in U$, maka $w=(u+w)-u\in U$, kontradiksi.
    Jika $u+w\in W$, maka $u=(u+w)-w\in W$, juga kontradiksi.
\end{proof}

Karena itu, gabungan dua garis berbeda yang melalui titik asal
di $\R^2$ bukan subruang. Gabungannya memuat dua arah, tetapi
tidak memuat seluruh hasil penjumlahan vektor dari arah-arah tersebut.

\begin{prop}[Gabungan subruang yang meningkat]
    \label{linear:rantai}
    Jika $W_1\subseteq W_2\subseteq\cdots$ adalah barisan
    subruang dari $V$, maka $\bigcup_{j\ge1}W_j$ merupakan subruang.
\end{prop}
\begin{proof}
    Vektor nol berada di gabungan. Jika $u\in W_r$ dan $v\in W_s$,
    ambil $k=\max(r,s)$. Kedua vektor berada di $W_k$, sehingga
    $\alpha u+\beta v\in W_k$ dan karenanya berada di gabungan.
\end{proof}

Contohnya,
$\R[t]=\bigcup_{n\ge0}\mathcal P_{\le n}(\R)$.
Setiap polinom mempunyai derajat berhingga, tetapi tidak ada satu
batas derajat yang berlaku untuk seluruh anggota $\R[t]$.

\subsection{Gabungan hingga dan medan tak hingga}

\begin{teorema}[][Gabungan hingga subruang sejati][linear:penutup-hingga]
    Misalkan $V$ ruang linear atas medan tak hingga, khususnya
    $\R$ atau $\C$. Jika $W_1,\ldots,W_m$ subruang sejati dari $V$,
    maka
    \[
        V\ne\bigcup_{j=1}^m W_j.
    \]
\end{teorema}
\begin{proof}
    Gunakan induksi pada $m$. Untuk $m=1$, pernyataan mengikuti
    definisi subruang sejati.
    Untuk langkah induksi, pilih
    $x\notin W_1\cup\cdots\cup W_{m-1}$.
    Jika $x\notin W_m$, pembuktian selesai.
    Jika $x\in W_m$, pilih $y\notin W_m$ dan pertimbangkan
    vektor $y+tx$, $t\in\F$.

    Tidak satu pun vektor itu berada di $W_m$: jika $y+tx\in W_m$,
    maka $y=(y+tx)-tx\in W_m$, kontradiksi.
    Untuk setiap $j<m$, paling banyak ada satu nilai $t$ yang
    membuat $y+tx\in W_j$. Sebab, jika dua nilai berbeda $s,t$
    memenuhi syarat itu, maka $(s-t)x\in W_j$, sehingga
    $x\in W_j$, kontradiksi.
    Hanya ada hingga banyak nilai parameter yang terlarang.
    Karena $\F$ tak hingga, kita dapat memilih $t$ lain dan
    memperoleh vektor di luar seluruh $W_j$.
\end{proof}

\begin{prop}[Satu subruang harus memuat semuanya]
    \label{linear:termuat-gabungan}
    Di atas $\R$ atau $\C$, jika subruang $U$ memenuhi
    $U\subseteq W_1\cup\cdots\cup W_m$, dengan semua $W_j$
    subruang, maka $U\subseteq W_j$ untuk sedikitnya satu $j$.
    Akibatnya, gabungan hingga subruang merupakan subruang tepat
    ketika salah satu anggotanya memuat semua anggota lainnya.
\end{prop}
\begin{proof}
    Tulis $U=\bigcup_{j=1}^m(U\cap W_j)$.
    Jika setiap $U\cap W_j$ sejati di dalam $U$, teorema sebelumnya
    memberi kontradiksi. Maka $U=U\cap W_j$ untuk suatu $j$.
    Untuk pernyataan terakhir, terapkan hasil ini pada subruang
    $U=\bigcup_j W_j$; arah sebaliknya langsung.
\end{proof}

Sifat hingga pada keluarga subruang penting. Semua garis yang
melalui titik asal di $\R^2$, jika digabungkan, menutupi seluruh
$\R^2$, walaupun masing-masing garis merupakan subruang sejati.
Pada penggunaan teorema ini, medan skalar juga harus diperiksa;
pembuktian pemilihan parameter tadi memakai ketersediaan tak
hingga banyak skalar.

\section{Latihan Pilihan}
\label{linear:latihan}

Nyatakan medan dan operasi yang digunakan sebelum menguji suatu
himpunan. Untuk jawaban negatif, berikan kegagalan aksioma atau
contoh konkret yang melanggar ketertutupan.

\begin{latihan}\label{linear:lat-operasi}
    Tetapkan $p,q\in\R^n$, $n\ge1$. Pada himpunan $\R^n$, definisikan
    \[
        u\oplus v=u+v+p,\qquad
        \alpha\odot u=\alpha u+(\alpha-1)q.
    \]
    Tentukan syarat perlu dan cukup pada $p,q$ agar kedua operasi
    itu membentuk ruang linear atas $\R$. Tentukan pula vektor nol
    dan invers penjumlahan setiap vektor.
\end{latihan}

\begin{latihan}\label{linear:lat-kompleks}
    Misalkan $V$ ruang linear kompleks, dan $W\subseteq V$
    merupakan subruang ketika skalar dibatasi pada $\R$.
    \begin{enumerate}[label=(\alph*),nosep]
        \item Buktikan bahwa $W$ subruang kompleks jika dan hanya
              jika $iw\in W$ untuk setiap $w\in W$.
        \item Tentukan apakah
              \[
                  \mathcal H=\left\{
                  \begin{pmatrix}z&w\\-\overline w&\overline z\end{pmatrix}
                  :z,w\in\C\right\}
              \]
              merupakan subruang real dan apakah merupakan subruang
              kompleks dari $\Mat_2(\C)$.
    \end{enumerate}
\end{latihan}

\begin{latihan}\label{linear:lat-kuadratik}
    Untuk $\lambda\in\R$, definisikan
    \[
        W_\lambda=\{(x,y)\in\R^2:x^2+\lambda y^2=0\}.
    \]
    Tentukan semua $\lambda$ yang membuat $W_\lambda$ subruang real.
    Kemudian jawab pertanyaan yang sama untuk
    $\widetilde W_\lambda=\{(x,y)\in\C^2:x^2+\lambda y^2=0\}$
    sebagai subruang kompleks, dengan parameter $\lambda$ tetap real.
\end{latihan}

\begin{latihan}\label{linear:lat-polinom}
    Di dalam $\mathcal P_{\le4}(\R)$, pertimbangkan
    \begin{align*}
        U & =\{p:p(0)=p(1)=p(2)\}, \\
        W & =\{p:p(0)p(1)=0\}.
    \end{align*}
    Tentukan himpunan mana yang merupakan subruang.
    Tuliskan semua polinom dalam $U$ dalam bentuk parameter,
    tanpa menghitung basis atau dimensi.
\end{latihan}

\begin{latihan}\label{linear:lat-komutatif}
    Misalkan $N=\begin{pmatrix}0&1\\0&0\end{pmatrix}$ dan
    \[
        U=\{X\in\Mat_2(\R):XN=NX\},\qquad
        W=\{X\in\Mat_2(\R):XN=-NX\}.
    \]
    Buktikan keduanya merupakan subruang dan tentukan semua
    anggotanya secara eksplisit. Apakah $W$ tertutup terhadap
    perkalian dua matriks? Jelaskan kaitannya dengan definisi subruang.
\end{latihan}

\begin{latihan}\label{linear:lat-trace-kuadrat}
    Tentukan apakah masing-masing himpunan berikut merupakan
    subruang dari $\Mat_2(\R)$:
    \[
        S=\{X:\tr(X^2)=0\},\qquad
        T=\{X:X^T=X,\ \tr(X^2)=0\}.
    \]
    Semua operasi menggunakan penjumlahan dan perkalian skalar biasa.
\end{latihan}

\begin{latihan}\label{linear:lat-rekurensi}
    Untuk $c\in\R$, misalkan $W_c$ terdiri atas semua barisan real
    $(a_n)_{n\ge1}$ yang memenuhi
    \[
        a_{n+2}-a_{n+1}-a_n=c\qquad(n\ge1).
    \]
    Tentukan kapan $W_c$ merupakan subruang dari $\R^{\mathbb N}$.
    Untuk $W_0$, buktikan bahwa setiap pilihan $a_1,a_2$ menentukan
    tepat satu barisan, dan bahwa syarat $a_1=a_2=0$ hanya
    menghasilkan barisan nol.
\end{latihan}

\begin{latihan}\label{linear:lat-afine}
    Misalkan $W\subseteq V$, dengan $V$ ruang linear real,
    memenuhi $0_V\in W$ serta
    \[
        \alpha u+(1-\alpha)v\in W
        \quad\text{untuk setiap }u,v\in W,\ \alpha\in\R.
    \]
    Buktikan $W$ subruang.
    Apakah masih cukup jika ketertutupan tersebut hanya diketahui
    untuk $\alpha=\tfrac12$? Apakah syarat $0_V\in W$ boleh dihilangkan?
\end{latihan}

\begin{latihan}\label{linear:lat-barisan}
    Di dalam $\R^{\mathbb N}$, definisikan
    \begin{align*}
        C_{00} & =\{(a_n):a_n=0\text{ untuk semua }n\text{ yang cukup besar}\},       \\
        E      & =\{(a_n):a_n\text{ konstan untuk semua }n\text{ yang cukup besar}\}.
    \end{align*}
    Buktikan bahwa keduanya subruang dan bahwa
    $C_{00}\subsetneq E\subsetneq\R^{\mathbb N}$.
    Nyatakan $C_{00}$ sebagai gabungan meningkat subruang-subruang
    yang masing-masing mengizinkan hanya hingga banyak suku awal tak nol.
\end{latihan}

\begin{latihan}\label{linear:lat-garis}
    Buktikan bahwa setiap subruang sejati tak nol dari $\R^2$
    berbentuk $\{tu:t\in\R\}$ untuk suatu $u\ne0$.
    Gunakan hasil ini untuk menjelaskan mengapa gabungan tak hingga
    subruang sejati dapat sama dengan seluruh ruang.
\end{latihan}

\begin{latihan}\label{linear:lat-hindari}
    Diberikan vektor-vektor tak nol $a_1,\ldots,a_m\in\R^n$,
    dengan $m,n\ge1$. Buktikan bahwa ada $x\in\R^n$ sehingga
    \[
        a_j^Tx\ne0\qquad\text{untuk semua }j=1,\ldots,m.
    \]
    Berikan dua pembuktian: satu menggunakan teorema gabungan hingga,
    dan satu lagi dengan memilih $x=(1,t,\ldots,t^{n-1})^T$ secara tepat.
\end{latihan}

\begin{latihan}\label{linear:lat-nol-bersama}
    Buktikan kedua pernyataan berikut.
    \begin{enumerate}[label=(\alph*)]
        \item Misalkan $W$ subruang dari $\R^n$ dan setiap $w\in W$
              mempunyai sedikitnya satu koordinat nol. Maka terdapat
              satu indeks $j$ yang sama sehingga $w_j=0$ untuk semua $w\in W$.
        \item Misalkan $U$ subruang dari $\R[t]$, dan
              $a_1,\ldots,a_m$ bilangan real tetap yang berbeda, $m\ge1$.
              Jika setiap $p\in U$ mempunyai akar di antara
              $a_1,\ldots,a_m$, maka ada satu $a_j$ yang menjadi
              akar semua polinom dalam $U$.
    \end{enumerate}
\end{latihan}

\ifpembahasan
    \section{Pembahasan Latihan}
    \label{linear:pembahasan}

    \subsection*{Latihan \ref{linear:lat-operasi}: dua cara menentukan vektor nol}
    Identitas untuk operasi $\oplus$ harus memenuhi $u\oplus e=u$,
    yakni $u+e+p=u$. Jadi $e=-p$.
    Jika kedua operasi membentuk ruang linear, identitas $0\odot u=e$
    harus berlaku. Namun, $0\odot u=-q$, sehingga $p=q$ merupakan
    syarat perlu.
    Jika $p=q$, ambil $a=-p$ pada contoh pergeseran operasi:
    \[
        u\oplus v=u+v-a,\qquad
        \alpha\odot u=a+\alpha(u-a).
    \]
    Contoh tersebut membuktikan kecukupan. Dengan demikian syaratnya
    tepat $p=q$; vektor nolnya $-p$ dan invers penjumlahan dari $u$
    adalah $-u-2p$.

    \subsection*{Latihan \ref{linear:lat-kompleks}: cukup menguji perkalian oleh $i$}
    Jika $W$ subruang kompleks, tentu $iw\in W$.
    Sebaliknya, tulis $\alpha=a+bi$ dengan $a,b\in\R$.
    Karena $W$ subruang real dan $iw\in W$, diperoleh
    \[
        \alpha w=aw+b(iw)\in W.
    \]
    Ketertutupan penjumlahan sudah diketahui, sehingga $W$ subruang kompleks.

    Untuk $\mathcal H$, matriks nol diperoleh dari $z=w=0$.
    Penjumlahan menjumlahkan parameter $z,w$, dan perkalian oleh
    skalar real mempertahankan hubungan konjugasi pada kedua entri
    baris bawah. Jadi $\mathcal H$ subruang real.
    Matriks $I_2$ berada di $\mathcal H$, tetapi $iI_2$ tidak:
    entri kiri atasnya memaksa $z=i$, sehingga entri kanan bawah
    seharusnya $\overline z=-i$, bukan $i$.
    Maka $\mathcal H$ bukan subruang kompleks.

    \subsection*{Latihan \ref{linear:lat-kuadratik}: medan mengubah himpunan solusi}
    Untuk skalar real, jika $\lambda>0$, persamaan hanya mempunyai
    solusi $(0,0)$. Jika $\lambda=0$, solusinya seluruh pasangan
    $(0,y)$; kedua himpunan merupakan subruang.
    Jika $\lambda<0$, letakkan $c=\sqrt{-\lambda}>0$.
    Vektor $(c,1)$ dan $(-c,1)$ memenuhi persamaan, tetapi jumlahnya
    $(0,2)$ tidak, sebab $0^2+\lambda\,2^2=4\lambda\ne0$.
    Jadi $W_\lambda$ subruang real tepat ketika $\lambda\ge0$.

    Di atas $\C$, untuk setiap $\lambda\ne0$ terdapat $c\ne0$
    dengan $c^2=-\lambda$. Kedua vektor yang sama kembali memberi
    kegagalan ketertutupan penjumlahan.
    Karena kasus $\lambda=0$ tetap menghasilkan $\{(0,y):y\in\C\}$,
    $\widetilde W_\lambda$ subruang kompleks tepat ketika $\lambda=0$.

    \subsection*{Latihan \ref{linear:lat-polinom}: tiga nilai yang sama dan satu perkalian}
    Syarat pada $U$ ekuivalen dengan dua persamaan linear homogen
    $p(0)-p(1)=0$ dan $p(1)-p(2)=0$, sehingga $U$ subruang.
    Jika $a=p(0)$, polinom $p(t)-a$ mempunyai akar $0,1,2$.
    Teorema faktor memberi
    \[
        p(t)=a+t(t-1)(t-2)(bt+c),\qquad a,b,c\in\R,
    \]
    karena derajat $p$ paling besar empat. Sebaliknya, setiap polinom
    berbentuk ini mempunyai nilai $a$ pada ketiga titik.

    Himpunan $W$ bukan subruang. Polinom $p(t)=t$ dan $q(t)=t-1$
    masing-masing berada di $W$, tetapi $p+q=2t-1$ mempunyai nilai
    $-1$ di $0$ dan $1$ di $1$, sehingga tidak berada di $W$.

    \subsection*{Latihan \ref{linear:lat-komutatif}: persamaan linear pada entri}
    Dengan $X=\begin{pmatrix}a&b\\c&d\end{pmatrix}$,
    \[
        XN=\begin{pmatrix}0&a\\0&c\end{pmatrix},\qquad
        NX=\begin{pmatrix}c&d\\0&0\end{pmatrix}.
    \]
    Maka
    \[
        U=\left\{\begin{pmatrix}a&b\\0&a\end{pmatrix}:a,b\in\R\right\},
        \qquad
        W=\left\{\begin{pmatrix}a&b\\0&-a\end{pmatrix}:a,b\in\R\right\}.
    \]
    Matriks nol berada di kedua himpunan, dan kombinasi linear
    mempertahankan bentuk masing-masing. Jadi keduanya subruang.
    Matriks $D=\diag(1,-1)$ berada di $W$, tetapi $D^2=I_2$ tidak
    berada di $W$. Ini tidak bertentangan dengan definisi subruang,
    karena definisi tersebut tidak mensyaratkan ketertutupan perkalian matriks.

    \subsection*{Latihan \ref{linear:lat-trace-kuadrat}: simetris memberi jumlah kuadrat}
    Ambil $A=\begin{pmatrix}0&1\\0&0\end{pmatrix}$ dan
    $B=\begin{pmatrix}0&0\\1&0\end{pmatrix}$.
    Karena $A^2=B^2=0$, keduanya berada di $S$.
    Namun, $(A+B)^2=I_2$, sehingga $\tr((A+B)^2)=2\ne0$.
    Jadi $S$ bukan subruang.
    Untuk matriks simetris real $X=\begin{pmatrix}a&b\\b&c\end{pmatrix}$,
    \[
        \tr(X^2)=a^2+2b^2+c^2.
    \]
    Nilainya nol tepat ketika $a=b=c=0$. Jadi $T=\{0\}$ merupakan subruang.

    \subsection*{Latihan \ref{linear:lat-rekurensi}: ketertutupan diperiksa pada setiap indeks}
    Jika $c\ne0$, barisan nol tidak memenuhi persamaan, sehingga
    $W_c$ bukan subruang. Untuk $c=0$, barisan nol termasuk $W_0$.
    Jika $(a_n),(b_n)\in W_0$ dan $d_n=\alpha a_n+\beta b_n$, maka
    \[
        d_{n+2}-d_{n+1}-d_n
        =\alpha(a_{n+2}-a_{n+1}-a_n)
        +\beta(b_{n+2}-b_{n+1}-b_n)=0
    \]
    untuk setiap $n$. Jadi $W_0$ subruang.
    Setelah $a_1,a_2$ dipilih, rumus $a_{n+2}=a_{n+1}+a_n$
    mendefinisikan setiap suku berikutnya secara tunggal.
    Induksi menunjukkan bahwa dua barisan dengan suku awal sama
    harus identik. Jika kedua suku awal nol, induksi juga memberi
    $a_n=0$ untuk setiap $n$.

    \subsection*{Latihan \ref{linear:lat-afine}: peran nol dan seluruh skalar}
    Dengan memilih $v=0_V$, syarat soal memberi $\alpha u\in W$
    untuk setiap skalar real $\alpha$.
    Untuk $u,v\in W$, pilihan $\alpha=\tfrac12$ memberi
    $(u+v)/2\in W$. Ketertutupan skalar yang baru dibuktikan
    kemudian memberi $u+v\in W$. Maka $W$ subruang.

    Jika hanya titik tengah yang diketahui berada di $W$, kesimpulan
    dapat gagal. Himpunan $\mathbb Q\subset\R$ memuat nol dan tertutup
    terhadap rata-rata dua anggotanya, tetapi
    $\sqrt2\cdot1\notin\mathbb Q$.
    Syarat memuat nol juga penting: garis
    $W=\{(x,1):x\in\R\}$ tertutup terhadap semua kombinasi
    $\alpha u+(1-\alpha)v$, tetapi bukan subruang karena tidak memuat $(0,0)$.

    \subsection*{Latihan \ref{linear:lat-barisan}: ambil batas indeks yang lebih besar}
    Jika dua barisan nol setelah indeks masing-masing $r$ dan $s$,
    kombinasi linearnya nol setelah $\max(r,s)$.
    Jika dua barisan konstan setelah indeks-indeks itu, dengan nilai
    akhir $a$ dan $b$, kombinasi linearnya konstan dengan nilai
    $\alpha a+\beta b$ setelah $\max(r,s)$.
    Kedua himpunan memuat barisan nol, sehingga keduanya subruang.
    Barisan yang akhirnya nol tentu akhirnya konstan, tetapi
    $(1,1,1,\ldots)\in E\setminus C_{00}$.
    Barisan $(1,2,3,\ldots)$ tidak berada di $E$, sehingga kedua
    inklusi pada soal bersifat sejati.
    Dengan
    \[
        V_k=\{(a_n):a_n=0\text{ untuk setiap }n>k\},\qquad k\ge1,
    \]
    setiap $V_k$ subruang, $V_k\subseteq V_{k+1}$, dan
    $C_{00}=\bigcup_{k\ge1}V_k$.

    \subsection*{Latihan \ref{linear:lat-garis}: dua arah berbeda memenuhi seluruh bidang}
    Misalkan $W$ subruang sejati tak nol. Pilih $u\in W\setminus\{0\}$.
    Ketertutupan skalar memberi $\{tu:t\in\R\}\subseteq W$.
    Andaikan ada $v\in W$ yang bukan kelipatan $u$.
    Matriks $[u\ v]$ berukuran dua mempunyai determinan tak nol:
    jika determinannya nol, hubungan antar kedua koordinat
    akan memaksa $v$ kelipatan $u$.
    Karena itu, untuk setiap $x\in\R^2$, sistem
    $[u\ v](\alpha,\beta)^T=x$ mempunyai solusi.
    Jadi setiap $x$ adalah $\alpha u+\beta v\in W$, bertentangan
    dengan $W$ sejati. Maka $W=\{tu:t\in\R\}$.
    Setiap vektor tak nol terletak pada garis yang dibentuk oleh
    kelipatan dirinya. Gabungan seluruh garis demikian, bersama
    vektor nol yang sudah termuat pada tiap garis, adalah $\R^2$.

    \subsection*{Latihan \ref{linear:lat-hindari}: hindari hingga banyak kendala}
    Untuk setiap $j$, definisikan $H_j=\{x:a_j^Tx=0\}$.
    Ini subruang sejati: karena $a_j\ne0$, terdapat koordinat tak nol,
    dan vektor satuan pada koordinat tersebut tidak berada di $H_j$.
    Teorema \ref{linear:penutup-hingga} menjamin ada
    $x\notin\bigcup_j H_j$, sehingga semua $a_j^Tx\ne0$.

    Untuk cara kedua, tulis $a_j=(a_{j1},\ldots,a_{jn})^T$ dan bentuk
    \[
        f_j(t)=a_{j1}+a_{j2}t+\cdots+a_{jn}t^{n-1}.
    \]
    Setiap $f_j$ polinom tak nol, sehingga hanya mempunyai hingga
    banyak akar real. Pilih $t$ di luar gabungan seluruh akar
    $f_1,\ldots,f_m$. Vektor $x=(1,t,\ldots,t^{n-1})^T$
    memenuhi $a_j^Tx=f_j(t)\ne0$ untuk setiap $j$.

    \subsection*{Latihan \ref{linear:lat-nol-bersama}: satu koordinat atau akar bersama}
    Untuk (a), definisikan subruang koordinat
    $H_j=\{x\in\R^n:x_j=0\}$.
    Hipotesis soal menyatakan $W\subseteq\bigcup_{j=1}^n H_j$.
    Proposisi \ref{linear:termuat-gabungan} memberi $W\subseteq H_j$
    untuk sedikitnya satu indeks $j$. Itulah koordinat yang selalu nol.

    Untuk (b), ambil $K_j=\{p\in\R[t]:p(a_j)=0\}$, yang merupakan
    subruang karena evaluasi menghormati penjumlahan dan perkalian skalar.
    Hipotesis menyatakan $U\subseteq\bigcup_{j=1}^m K_j$.
    Proposisi yang sama memberi $U\subseteq K_j$ untuk suatu $j$,
    sehingga $p(a_j)=0$ untuk semua $p\in U$.
    Perhatikan perubahan kuantornya: pilihan koordinat atau akar
    pada hipotesis boleh bergantung pada vektor; kesimpulannya memberi
    satu pilihan yang berlaku serentak untuk seluruh subruang.
\fi

% !TeX root = ../../ONMIPA Alin.tex
% Acuan: Zhang (2009), Bab 1, khususnya Soal 1.15--1.36 dan 1.43--1.50;
% Zhang (2011), bagian 1.1, khususnya latihan tentang dimensi dan komplemen.
\chapter{Kebebaslinearan}
\label{ch:kebebaslinearan}

Kebebaslinearan membedakan daftar vektor yang memuat informasi
berulang dari daftar yang setiap anggotanya diperlukan. Dari gagasan
ini muncul basis dan dimensi: basis memberikan representasi yang
unik, sedangkan dimensi mengukur banyaknya parameter bebas.
Kecuali dinyatakan lain, medan skalar adalah $\F=\R$ atau $\C$.
Dimensi selalu dihitung terhadap medan skalar yang sedang digunakan.

\section{Kombinasi Linear dan Ruang Pembangun}

\begin{definisi}[][Ruang yang dibangun oleh suatu himpunan][bebas:span]
    Untuk $S\subseteq V$, ruang pembangun dari $S$ adalah
    \[
        \operatorname{span}_{\F}S
        =\left\{\sum_{j=1}^{r}a_jv_j:
        r\ge0,\ a_j\in\F,\ v_j\in S\right\}.
    \]
    Jumlah kosong bernilai $0$, sehingga
    $\operatorname{span}\varnothing=\{0\}$.
    Jika $\operatorname{span}S=V$, maka $S$ disebut
    \emph{himpunan pembangun} bagi $V$.
\end{definisi}

Walaupun $S$ tak hingga, setiap kombinasi linear tetap memakai
hingga banyak anggota. Deret tak hingga tidak termasuk definisi ini.

\begin{prop}[Subruang terkecil yang memuat himpunan]
    Ruang $\operatorname{span}S$ merupakan subruang terkecil dari
    $V$ yang memuat $S$. Akibatnya, untuk $S,T\subseteq V$,
    \[
        S\subseteq\operatorname{span}T
        \quad\Longrightarrow\quad
        \operatorname{span}S\subseteq\operatorname{span}T.
    \]
\end{prop}
\begin{proof}
    Vektor nol termasuk $\operatorname{span}S$. Kombinasi linear
    dari dua jumlah hingga kembali berupa jumlah hingga anggota
    $S$, sehingga uji subruang berlaku. Setiap subruang yang memuat
    $S$ harus memuat semua kombinasi linear anggotanya. Ini sekaligus
    membuktikan sifat terkecil dan implikasi yang dinyatakan.
\end{proof}

Dua daftar membangun ruang yang sama jika setiap vektor pada
masing-masing daftar dapat dituliskan dari daftar lainnya.
Menambahkan $v$ ke daftar pembangun tidak memperbesar ruang
tepat ketika $v$ sudah berada dalam ruang yang dibangun sebelumnya.

\begin{exam}[Daftar panjang dengan dua arah]
    Untuk $n\ge2$, tentukan ruang yang dibangun oleh
    \[
        v_k=(k,k+1,\ldots,k+n-1)^T,\qquad k=1,2,\ldots.
    \]
\end{exam}
\begin{solution}
    Tuliskan $\boldsymbol{1}=(1,\ldots,1)^T$ dan
    $d=(0,1,\ldots,n-1)^T$. Karena $v_k=k\boldsymbol{1}+d$,
    semua $v_k$ berada dalam $\operatorname{span}\{\boldsymbol{1},d\}$.
    Sebaliknya, $\boldsymbol{1}=v_2-v_1$ dan $d=2v_1-v_2$.
    Jadi ruang yang dicari adalah
    \[
        \operatorname{span}\{\boldsymbol{1},d\}
        =\operatorname{span}\{v_1,v_2\}.
    \]
    Banyaknya vektor yang diberikan tidak menentukan banyaknya
    arah yang sungguh-sungguh berbeda.
\end{solution}

\section{Kebebasan dan Ketergantungan Linear}

\begin{definisi}[][Bebas linear dan bergantung linear][bebas:def]
    Daftar $v_1,\ldots,v_r$ disebut \emph{bebas linear} jika
    \[
        a_1v_1+\cdots+a_rv_r=0
        \quad\Longrightarrow\quad a_1=\cdots=a_r=0.
    \]
    Jika ada hubungan tersebut dengan koefisien tidak semuanya
    nol, daftar itu disebut \emph{bergantung linear}.
    Himpunan tak hingga disebut bebas linear jika setiap
    subhimpunan hingganya bebas linear.
\end{definisi}

Daftar kosong bebas linear. Daftar yang memuat vektor nol atau
memuat vektor yang sama dua kali bergantung linear. Setiap
subdaftar dari daftar bebas linear juga bebas linear; setiap
daftar yang memuat subdaftar bergantung linear tetap bergantung linear.

\begin{teorema}[][Kriteria ketergantungan linear][bebas:kriteria]
    Daftar tak kosong $v_1,\ldots,v_r$ bergantung linear jika dan
    hanya jika sedikitnya satu $v_j$ merupakan kombinasi linear
    dari vektor-vektor lainnya.
    Jika $v_1\ne0$, hal ini ekuivalen dengan adanya $k\ge2$ sehingga
    \[
        v_k\in\operatorname{span}\{v_1,\ldots,v_{k-1}\}.
    \]
\end{teorema}
\begin{proof}
    Dari hubungan tak trivial, pilih koefisien $a_j\ne0$ dan
    selesaikan hubungan itu terhadap $v_j$. Sebaliknya, pindahkan
    ekspresi suatu $v_j$ ke satu ruas; koefisien $v_j$ adalah $1$,
    sehingga hubungannya tak trivial.
    Untuk pernyataan terakhir, pilih indeks terbesar $k$ dengan
    $a_k\ne0$. Karena $v_1\ne0$, tidak mungkin $k=1$.
    Maka $v_k$ merupakan kombinasi vektor-vektor sebelumnya.
\end{proof}

Frasa \emph{sedikitnya satu} tidak boleh diganti dengan
\emph{setiap}. Daftar $(e_1,e_2,e_2)$ bergantung linear, tetapi
$e_1$ tidak berada dalam ruang pembangun dua vektor lainnya.

\begin{prop}[Menambah satu vektor dan keunikan koefisien]
    Jika $v_1,\ldots,v_r$ bebas linear, maka
    \[
        v_1,\ldots,v_r,w\text{ bebas linear}
        \iff w\notin\operatorname{span}\{v_1,\ldots,v_r\}.
    \]
    Selain itu, setiap anggota $\operatorname{span}\{v_1,\ldots,v_r\}$
    mempunyai tepat satu representasi sebagai kombinasi daftar itu.
\end{prop}
\begin{proof}
    Dalam hubungan $\sum a_jv_j+bw=0$, jika $b=0$ maka semua
    $a_j=0$. Jika $b\ne0$, hubungan itu menyatakan bahwa $w$
    berada dalam ruang pembangun. Untuk keunikan, kurangkan
    dua representasi dari vektor yang sama dan gunakan kebebasan linear.
\end{proof}

\subsection{Mengubah daftar menjadi sistem homogen}
Untuk $v_j\in\F^m$, bentuk matriks kolom
$A=[v_1\ \cdots\ v_r]$. Maka
\[
    v_1,\ldots,v_r\text{ bebas linear}
    \iff\ker A=\{0\}\iff\rank A=r.
\]
Jika $m=r$, kriteria ini ekuivalen dengan $\det A\ne0$.
Untuk vektor abstrak, matriks yang sama dapat dibentuk dari
koefisien terhadap suatu daftar yang sudah diketahui bebas linear.

\begin{exam}[Jumlah siklik: ganjil dan genap berbeda]
    Misalkan $v_1,\ldots,v_n$ bebas linear, $n\ge2$.
    Definisikan $w_j=v_j+v_{j+1}$ dengan $v_{n+1}=v_1$.
    Kapan $w_1,\ldots,w_n$ bebas linear?
\end{exam}
\begin{solution}
    Dalam hubungan $\sum_{j=1}^n a_jw_j=0$, koefisien $v_j$
    adalah $a_j+a_{j-1}$, dengan $a_0=a_n$.
    Jadi $a_j=-a_{j-1}$ dan $a_1=(-1)^na_1$.
    Jika $n$ ganjil, $2a_1=0$, sehingga seluruh koefisien nol.
    Jika $n$ genap, terdapat hubungan tak trivial
    \[
        w_1-w_2+w_3-\cdots-w_n=0.
    \]
    Dengan demikian, daftar itu bebas linear tepat ketika $n$ ganjil.
    Penggunaan $2a_1=0\Rightarrow a_1=0$ sah karena medannya
    $\R$ atau $\C$.
\end{solution}

\subsection{Evaluasi dan perbandingan koefisien}
Untuk polinom, persamaan vektor berarti kesamaan setiap koefisien.
Untuk fungsi, persamaan harus berlaku di \emph{seluruh} domain.
Evaluasi pada beberapa titik dapat membuktikan kebebasan jika
sistem koefisiennya hanya mempunyai solusi nol. Sebaliknya,
matriks evaluasi yang singular pada pilihan titik tertentu
belum membuktikan bahwa fungsi-fungsinya bergantung linear.

\begin{exam}[Derajat berbeda dan titik evaluasi]
    Buktikan dua pernyataan berikut.
    \begin{enumerate}[label=(\alph*),nosep]
        \item Polinom tak nol yang mempunyai derajat berbeda-beda
        selalu bebas linear.
        \item Untuk titik berbeda $a_0,\ldots,a_n\in\F$, polinom
        \[
            L_j(t)=\prod_{\substack{0\le k\le n\\k\ne j}}
                \frac{t-a_k}{a_j-a_k},\qquad 0\le j\le n,
        \]
        bebas linear.
    \end{enumerate}
\end{exam}
\begin{solution}
    Untuk (a), jika ada kombinasi nol tak trivial, pilih polinom
    berderajat terbesar yang koefisien kombinasinya tidak nol.
    Suku tertingginya tidak dapat dibatalkan polinom lain, kontradiksi.
    Untuk (b), sifat $L_j(a_i)=\delta_{ij}$ mengubah
    $\sum c_jL_j=0$ menjadi $c_i=0$ dengan mengevaluasi di $a_i$.
\end{solution}

\section{Basis dan Dimensi}

\begin{definisi}[][Basis dan ruang berdimensi hingga][bebas:basis]
    Basis bagi $V$ adalah himpunan bebas linear yang membangun $V$.
    Ruang $V$ disebut \emph{berdimensi hingga} jika mempunyai
    himpunan pembangun hingga. Jika tidak, $V$ disebut
    \emph{berdimensi tak hingga}.
\end{definisi}

\begin{teorema}[][Lema pertukaran][bebas:pertukaran]
    Misalkan $u_1,\ldots,u_r$ bebas linear dan
    $v_1,\ldots,v_s$ membangun $V$. Maka $r\le s$.
    Lebih lanjut, $r$ anggota daftar pembangun dapat diganti
    oleh $u_1,\ldots,u_r$ sehingga daftar baru tetap membangun $V$.
\end{teorema}
\begin{proof}
    Nyatakan $u_1$ dalam daftar $v_j$. Karena $u_1\ne0$, suatu
    koefisien tidak nol; selesaikan terhadap $v_j$ itu. Mengganti
    $v_j$ dengan $u_1$ tidak mengubah ruang pembangun.
    Setelah $u_1,\ldots,u_k$ dimasukkan, nyatakan $u_{k+1}$
    dalam daftar pembangun yang baru. Sedikitnya satu koefisien
    pada $v_j$ yang tersisa tidak nol, sebab jika semuanya nol,
    $u_{k+1}$ berada dalam ruang pembangun $u_1,\ldots,u_k$.
    Koefisien tersebut memungkinkan pertukaran berikutnya.
    Proses ini membutuhkan satu $v_j$ berbeda pada setiap tahap,
    sehingga $r\le s$ dan pernyataan pertukaran terbukti.
\end{proof}

\begin{teorema}[][Keberadaan basis dan kesamaan ukurannya][bebas:dimensi]
    Setiap ruang berdimensi hingga mempunyai basis. Semua
    basisnya mempunyai banyak anggota yang sama, disebut
    \emph{dimensi} $V$ dan ditulis $\dim_{\F}V$.
    Setiap daftar bebas linear dapat diperluas menjadi basis,
    dan dari setiap daftar pembangun hingga dapat dipilih suatu basis.
\end{teorema}
\begin{proof}
    Dari daftar pembangun hingga, buang satu vektor yang merupakan
    kombinasi lainnya selama daftar itu masih bergantung linear.
    Ruang pembangun tidak berubah, dan proses berhenti pada basis.
    Jika dua basis mempunyai $r$ dan $s$ anggota, lema pertukaran
    memberi $r\le s$ dan $s\le r$, sehingga $r=s$.
    Untuk memperluas daftar bebas linear, tambahkan anggota dari
    suatu daftar pembangun hingga yang belum berada dalam ruang
    pembangunnya. Daftar tetap bebas linear dan akhirnya membangun $V$.
\end{proof}

Untuk $V=\{0\}$, basisnya adalah himpunan kosong dan dimensinya
nol. Jika $\dim V=n$, diperoleh kriteria penting berikut:
\begin{itemize}[nosep]
    \item setiap daftar bebas linear mempunyai paling banyak $n$ anggota;
    \item setiap daftar pembangun mempunyai sedikitnya $n$ anggota;
    \item daftar yang terdiri atas tepat $n$ vektor merupakan basis
    jika bebas linear, atau jika membangun $V$.
\end{itemize}
Syarat \emph{tepat $n$} penting: satu vektor tak nol di $\R^3$
bebas linear, tetapi belum merupakan basis.

\subsection{Dimensi subruang dan contoh dasar}
Jika $U\subseteq V$ dan $\dim V=n$, proses memilih vektor
$u_{k+1}\in U\setminus\operatorname{span}\{u_1,\ldots,u_k\}$
berhenti setelah paling banyak $n$ langkah. Hasilnya adalah basis
bagi $U$. Maka $\dim U\le n$, dan
\[
    U\subseteq V,\quad \dim U=\dim V<\infty
    \quad\Longrightarrow\quad U=V.
\]
Kesamaan dimensi saja, tanpa inklusi, tidak menjamin kesamaan subruang.

Basis standar $\F^n$ adalah $(e_1,\ldots,e_n)$, sedangkan
$\Mat_{m\times n}(\F)$ mempunyai basis matriks satuan $E_{ij}$
dan dimensi $mn$. Notasi
\[
    \mathcal P_{\le n}(\F)=\{p\in\F[t]:\deg p\le n\}\cup\{0\}
\]
memiliki basis $(1,t,\ldots,t^n)$ dan dimensi $n+1$.
Ruang $\F[t]$ berdimensi tak hingga karena memuat daftar bebas
linear $1,t,\ldots,t^N$ untuk setiap $N$.

\subsection{Perubahan medan skalar}
Jika $(v_1,\ldots,v_n)$ basis kompleks bagi $V$, maka
\[
    (v_1,iv_1,\ldots,v_n,iv_n)
\]
merupakan basis real. Memisahkan bagian real dan imajiner dari
koefisien membuktikan sifat pembangun sekaligus kebebasannya.
Jadi $\dim_{\R}V=2\dim_{\C}V$.
Kebebasan atas $\C$ mengakibatkan kebebasan atas $\R$, tetapi
arah sebaliknya dapat gagal: $1,i$ bebas atas $\R$ dan bergantung atas $\C$.

\begin{exam}[Dimensi ruang matriks Hermitian]
    Tentukan basis real dan dimensi ruang $\mathcal H_n$
    dari semua matriks Hermitian berukuran $n$.
\end{exam}
\begin{solution}
    Entri diagonal matriks Hermitian real; pasangan entri
    di luar diagonal berbentuk $a+bi$ dan $a-bi$.
    Karena itu, suatu basis real diberikan oleh
    \[
        E_{jj}\ (1\le j\le n),\qquad
        E_{jk}+E_{kj},\ i(E_{jk}-E_{kj})\ (j<k).
    \]
    Setiap matriks Hermitian mempunyai ekspansi unik dalam daftar
    tersebut, dengan koefisien berupa entri diagonal serta bagian
    real dan imajiner entri di atas diagonal. Maka
    $\dim_{\R}\mathcal H_n=n+2\binom n2=n^2$.
    Ini dimensi real; $\mathcal H_n$ bukan subruang kompleks untuk $n\ge1$.
\end{solution}

\section{Koordinat dan Pergantian Basis}

\begin{definisi}[][Vektor koordinat][bebas:koordinat]
    Untuk basis terurut $\mathcal B=(b_1,\ldots,b_n)$, jika
    $v=x_1b_1+\cdots+x_nb_n$, definisikan
    \[
        [v]_{\mathcal B}=(x_1,\ldots,x_n)^T.
    \]
    Urutan basis merupakan bagian dari definisi koordinat.
\end{definisi}

Jika $\mathcal C=(c_1,\ldots,c_n)$ adalah basis lain, letakkan
\[
    P_{\mathcal B\leftarrow\mathcal C}
      =\bigl[[c_1]_{\mathcal B}\ \cdots\ [c_n]_{\mathcal B}\bigr].
\]
Kolom-kolomnya menyatakan vektor basis sumber dalam basis tujuan.
Dengan menyubstitusi ekspansi setiap $c_j$, diperoleh
\[
    [v]_{\mathcal B}
      =P_{\mathcal B\leftarrow\mathcal C}[v]_{\mathcal C},
    \qquad
    P_{\mathcal C\leftarrow\mathcal B}
      =P_{\mathcal B\leftarrow\mathcal C}^{-1}.
\]
Untuk tiga basis, urutan panah memberi aturan
$P_{\mathcal B\leftarrow\mathcal D}
=P_{\mathcal B\leftarrow\mathcal C}P_{\mathcal C\leftarrow\mathcal D}$.

\begin{exam}[Koordinat polinom bergeser]
    Pada $\mathcal P_{\le2}(\R)$, gunakan
    $\mathcal B=(1,t,t^2)$ dan $\mathcal C=(1,t-1,(t-1)^2)$.
    Tentukan $P_{\mathcal B\leftarrow\mathcal C}$ dan
    koordinat $p(t)=2+3t+t^2$ terhadap $\mathcal C$.
\end{exam}
\begin{solution}
    Ekspansi ketiga anggota $\mathcal C$ memberi
    \[
        P_{\mathcal B\leftarrow\mathcal C}
        =\begin{pmatrix}1&-1&1\\0&1&-2\\0&0&1\end{pmatrix}.
    \]
    Dengan $s=t-1$, polinom menjadi $6+5s+s^2$.
    Jadi $[p]_{\mathcal C}=(6,5,1)^T$, dan perkalian matriks
    di atas mengembalikannya menjadi $[p]_{\mathcal B}=(2,3,1)^T$.
\end{solution}

Secara umum, terhadap basis $(1,t-a,\ldots,(t-a)^n)$,
koordinat polinom $p\in\mathcal P_{\le n}(\F)$ adalah
\[
    p(a),\ p'(a),\ \frac{p''(a)}{2!},\ldots,
    \frac{p^{(n)}(a)}{n!}.
\]
Hal ini mengikuti ekspansi polinom dalam pangkat $t-a$ lalu
mendiferensialkan dan mengevaluasi di $a$.
Terhadap basis Lagrange $\mathcal L=(L_0,\ldots,L_n)$, diperoleh
\[
    [p]_{\mathcal L}=(p(a_0),\ldots,p(a_n))^T.
\]
Memang, selisih $p-\sum p(a_j)L_j$ berderajat paling besar $n$ dan mempunyai
$n+1$ akar berbeda, sehingga identik nol.

\section{Jumlah Subruang dan Rumus Dimensi}

\begin{definisi}[][Jumlah dan jumlah langsung][bebas:jumlah-langsung]
    Untuk subruang $U,W\subseteq V$, definisikan
    \[
        U+W=\{u+w:u\in U,\ w\in W\}.
    \]
    Jumlah ini disebut \emph{jumlah langsung}, ditulis $U\oplus W$,
    jika setiap anggotanya mempunyai representasi $u+w$ yang unik.
\end{definisi}

Ruang $U+W$ adalah subruang terkecil yang memuat $U\cup W$.
Berbeda dari gabungan, jumlah dua subruang selalu merupakan subruang.

\begin{teorema}[][Rumus dimensi dan kriteria jumlah langsung][bebas:grassmann]
    Untuk subruang berdimensi hingga $U,W$,
    \[
        \dim(U+W)=\dim U+\dim W-\dim(U\cap W).
    \]
    Selain itu, jumlahnya langsung jika dan hanya jika
    $U\cap W=\{0\}$, atau ekuivalen dengan
    $\dim(U+W)=\dim U+\dim W$.
\end{teorema}
\begin{proof}
    Ambil basis $(z_1,\ldots,z_k)$ bagi $U\cap W$.
    Perluas dengan $u_1,\ldots,u_a$ menjadi basis $U$, dan
    dengan $w_1,\ldots,w_b$ menjadi basis $W$.
    Daftar gabungan
    \[
        (z_1,\ldots,z_k,u_1,\ldots,u_a,w_1,\ldots,w_b)
    \]
    membangun $U+W$.
    Jika kombinasinya nol, kombinasi para $u_i$ berada dalam
    $U\cap W$; kebebasan basis $U$ memaksa semua koefisien
    $u_i$ nol. Kebebasan basis $W$ kemudian memaksa koefisien
    sisanya nol. Jadi daftar gabungan adalah basis, dan
    $\dim(U+W)=k+a+b=\dim U+\dim W-k$.

    Jika $U\cap W=\{0\}$ dan $u+w=u'+w'$, maka
    $u-u'=w'-w\in U\cap W$, sehingga representasi unik.
    Sebaliknya, untuk $z\in U\cap W$, persamaan
    $z+0=0+z$ memaksa $z=0$ jika representasi unik.
\end{proof}

Jika $U,W\subseteq V$ dan $\dim V=n$, rumus ini memberi
\[
    \dim(U\cap W)\ge\dim U+\dim W-n.
\]
Khususnya, jika $\dim U+\dim W>n$, keduanya beririsan tak trivial.
Untuk membuktikan kesamaan subruang, sering cukup membuktikan
satu inklusi lalu menyamakan dimensinya.

\begin{exam}[Irisan dan jumlah dua subruang di ruang berdimensi empat]
    Di $\R^4$, misalkan
    \[
        U=\{x:x_1+x_2+x_3+x_4=0\},\qquad
        W=\{x:x_1=x_2,\ x_3=x_4\}.
    \]
    Tentukan basis $U\cap W$ serta dimensi $U+W$.
\end{exam}
\begin{solution}
    Ruang $U$ berdimensi tiga, dengan basis
    $(e_1-e_4,e_2-e_4,e_3-e_4)$.
    Anggota $W$ berbentuk $(a,a,b,b)^T$, sehingga dimensinya dua.
    Syarat tambahan untuk berada di $U$ adalah $2a+2b=0$.
    Jadi $U\cap W=\operatorname{span}\{(1,1,-1,-1)^T\}$.
    Rumus dimensi memberi $\dim(U+W)=3+2-1=4$, maka $U+W=\R^4$.
    Jumlahnya bukan jumlah langsung.
\end{solution}

\subsection{Komplemen dan lebih dari dua subruang}
Setiap subruang $U$ dari ruang berdimensi hingga $V$ mempunyai
\emph{komplemen}: suatu $W$ dengan $V=U\oplus W$.
Perluas basis $U$ menjadi basis $V$, lalu ambil $W$ sebagai
ruang pembangun vektor-vektor tambahan. Komplemen tidak tunggal;
di $\R^2$, baik $\operatorname{span}\{e_2\}$ maupun
$\operatorname{span}\{e_1+e_2\}$ merupakan komplemen bagi
$\operatorname{span}\{e_1\}$.

Untuk $U_1,\ldots,U_k$, jumlahnya langsung jika hubungan
$u_1+\cdots+u_k=0$, dengan $u_j\in U_j$, memaksa setiap $u_j=0$.
Syarat $U_i\cap U_j=\{0\}$ untuk setiap $i\ne j$ saja belum cukup.
Tiga garis yang dibangun oleh $e_1,e_2,e_1+e_2$ di $\R^2$
beririsan sepasang-sepasang hanya di nol, tetapi jumlah ketiganya
tidak langsung. Uji yang tepat adalah
\[
    U_j\cap(U_1+\cdots+U_{j-1})=\{0\},\qquad 2\le j\le k.
\]

\section{Latihan Pilihan}

\begin{latihan}\label{bebas:lat-parameter}
    Misalkan $u,v,w$ bebas linear. Untuk $\lambda\in\F$, tentukan
    kapan $v-u,\lambda w-v,u-w$ bebas linear. Pada nilai parameter
    yang gagal, berikan basis ruang pembangunnya.
\end{latihan}

\begin{latihan}\label{bebas:lat-siklik}
    Misalkan $v_1,\ldots,v_n$ bebas linear, $n\ge2$.
    Definisikan $w_j=v_j+\lambda v_{j+1}$, dengan $v_{n+1}=v_1$.
    Buktikan bahwa $w_1,\ldots,w_n$ bebas linear tepat ketika
    $(-\lambda)^n\ne1$. Tentukan dimensi ruang pembangunnya
    jika $(-\lambda)^n=1$.
\end{latihan}

\begin{latihan}\label{bebas:lat-lokal}
    Berikan contoh tiga vektor yang setiap dua anggotanya bebas
    linear, tetapi ketiganya bergantung linear.
    Secara umum, jika setiap $r-1$ anggota dari daftar
    $v_1,\ldots,v_r$ bebas linear tetapi seluruh daftar bergantung
    linear, buktikan bahwa setiap hubungan tak trivialnya memakai
    koefisien yang semuanya tidak nol.
\end{latihan}

\begin{latihan}\label{bebas:lat-polinom}
    Untuk $n\ge2$, tentukan basis dan dimensi subruang
    \[
        U=\{p\in\mathcal P_{\le n}(\R):p(0)=p'(0)=0\},
        \quad
        W=\{p\in\mathcal P_{\le n}(\R):p(1)=0\}.
    \]
    Tentukan pula $\dim(U\cap W)$ dan $U+W$.
\end{latihan}

\begin{latihan}\label{bebas:lat-koordinat}
    Untuk basis $\mathcal B=(b_1,\ldots,b_n)$, letakkan
    $c_j=b_1+\cdots+b_j$. Buktikan bahwa
    $\mathcal C=(c_1,\ldots,c_n)$ basis.
    Jika $[v]_{\mathcal B}=(x_1,\ldots,x_n)^T$, tentukan
    $[v]_{\mathcal C}$.
\end{latihan}

\begin{latihan}\label{bebas:lat-eksponensial}
    Untuk bilangan real berbeda $a_1,\ldots,a_r$, buktikan bahwa
    fungsi $e^{a_1t},\ldots,e^{a_rt}$ bebas linear atas $\R$.
    Gunakan evaluasi turunan di $t=0$ dan determinan Vandermonde.
\end{latihan}

\begin{latihan}\label{bebas:lat-medan}
    Buktikan bahwa $(1,i)$ dan $(i,-1)$ bebas linear di $\C^2$
    sebagai ruang real, tetapi bergantung sebagai ruang kompleks.
    Tentukan basis real bagi ruang pembangun kompleks kedua vektor itu.
\end{latihan}

\begin{latihan}\label{bebas:lat-jumlah-nol}
    Untuk $n\ge2$, misalkan $S$ ruang matriks real $n\times n$
    yang jumlah entri pada setiap baris dan setiap kolomnya nol.
    Buktikan bahwa $\dim S=(n-1)^2$ dan berikan basis eksplisit.
\end{latihan}

\begin{latihan}\label{bebas:lat-tiga}
    Jika $U_1,U_2,U_3$ subruang dari ruang berdimensi $n$, buktikan
    \[
        \dim(U_1\cap U_2\cap U_3)
        \ge\dim U_1+\dim U_2+\dim U_3-2n.
    \]
    Berikan contoh yang mencapai kesamaan dengan ketiga subruang
    berbeda dan sejati.
\end{latihan}

\begin{latihan}\label{bebas:lat-selisih-dimensi}
    Misalkan $U,W$ subruang berdimensi hingga dan
    $\dim(U+W)=\dim(U\cap W)+1$.
    Buktikan bahwa $U\subseteq W$ atau $W\subseteq U$.
\end{latihan}

\begin{latihan}\label{bebas:lat-lagrange}
    Untuk titik berbeda $a_0,\ldots,a_n$, gunakan basis Lagrange
    untuk membuktikan
    \[
        \sum_{j=0}^{n}\frac{a_j^k}
        {\displaystyle\prod_{\ell\ne j}(a_j-a_\ell)}
        =\begin{cases}0,&0\le k<n,\\1,&k=n.\end{cases}
    \]
    Untuk $k=0$, pembilang ditafsirkan sebagai $1$ pada semua titik.
\end{latihan}

\begin{latihan}\label{bebas:lat-komplemen}
    Misalkan $U,W\subseteq V$, dengan $\dim V<\infty$ dan
    $\dim U=\dim W$. Buktikan bahwa ada subruang $K$ sehingga
    \[
        V=U\oplus K=W\oplus K.
    \]
    Apakah adanya komplemen yang sama memaksa $U=W$?
\end{latihan}

\ifpembahasan
\section{Pembahasan Latihan}

\subsection*{Latihan \ref{bebas:lat-parameter}: koefisien dalam daftar bebas}
Matriks koefisien terhadap $(u,v,w)$ adalah
\[
    A=\begin{pmatrix}-1&0&1\\1&-1&0\\0&\lambda&-1\end{pmatrix},
    \qquad \det A=\lambda-1.
\]
Jadi daftar bebas tepat ketika $\lambda\ne1$.
Untuk $\lambda=1$, jumlah ketiganya nol, sedangkan $v-u$ dan
$w-v$ bebas linear dengan membandingkan koefisien $u$ dan $w$.
Kedua vektor itu menjadi basis ruang pembangun yang berdimensi dua.

\subsection*{Latihan \ref{bebas:lat-siklik}: menutup hubungan rekursif}
Hubungan $\sum a_jw_j=0$ ekuivalen dengan
$a_j+\lambda a_{j-1}=0$, dengan $a_0=a_n$.
Maka $a_j=(-\lambda)^{j-1}a_1$ dan
$(1-(-\lambda)^n)a_1=0$.
Jika faktor itu tidak nol, semua $a_j=0$.
Jika faktornya nol, $\lambda\ne0$ dan pilihan $a_1=1$
menghasilkan hubungan tak trivial; setiap hubungan merupakan
kelipatannya. Selain itu, $w_1,\ldots,w_{n-1}$ bebas, sebab
hubungan yang koefisien $w_n$-nya nol memaksa semua koefisien nol.
Hubungan yang ditemukan menyatakan $w_n$ dari sisanya.
Jadi dimensi ruang pembangun tepat $n-1$.

\subsection*{Latihan \ref{bebas:lat-lokal}: pasangan bebas belum cukup}
Ambil $e_1,e_2,e_1+e_2$ di $\R^2$.
Setiap pasangan bebas linear, tetapi terdapat hubungan
$e_1+e_2-(e_1+e_2)=0$.
Secara umum, jika suatu koefisien dalam hubungan tak trivial
bernilai nol, hubungan itu menjadi hubungan tak trivial pada
paling banyak $r-1$ vektor, bertentangan dengan hipotesis.

\subsection*{Latihan \ref{bebas:lat-polinom}: faktor yang diwajibkan}
Syarat pada $U$ berarti $t^2$ membagi $p$; basisnya
$(t^2,t^3,\ldots,t^n)$ dan dimensinya $n-1$.
Syarat pada $W$ berarti $(t-1)$ membagi $p$; basisnya
$(t-1,t(t-1),\ldots,t^{n-1}(t-1))$ dan dimensinya $n$.
Anggota irisan habis dibagi $t^2(t-1)$.
Untuk $n\ge3$, basis irisannya
\[
    t^2(t-1),\ t^3(t-1),\ldots,t^{n-1}(t-1),
\]
dengan $n-2$ anggota. Untuk $n=2$, irisannya hanya nol dan
basisnya kosong, sehingga rumus dimensi yang sama tetap berlaku.
Maka $\dim(U+W)=(n-1)+n-(n-2)=n+1$ dan
$U+W=\mathcal P_{\le n}(\R)$.

\subsection*{Latihan \ref{bebas:lat-koordinat}: selisih koordinat berurutan}
Persamaan $b_1=c_1$ dan $b_j=c_j-c_{j-1}$ menunjukkan bahwa
$\mathcal C$ membangun $V$. Karena jumlahnya $n$, daftar itu basis.
Jika $v=\sum y_jc_j$, koefisien $b_k$ adalah
$x_k=y_k+\cdots+y_n$. Jadi
\[
    [v]_{\mathcal C}=(x_1-x_2,x_2-x_3,\ldots,x_{n-1}-x_n,x_n)^T.
\]

\subsection*{Latihan \ref{bebas:lat-eksponensial}: sistem Vandermonde}
Andaikan $\sum c_je^{a_jt}=0$ untuk setiap $t$.
Diferensialkan $k$ kali dan evaluasi di nol, untuk
$k=0,\ldots,r-1$, sehingga $\sum c_ja_j^k=0$.
Determinan matriks sistem ini adalah
$\prod_{i<j}(a_j-a_i)\ne0$. Maka semua $c_j=0$ dan
fungsi-fungsi tersebut bebas linear.

\subsection*{Latihan \ref{bebas:lat-medan}: koefisien real dan kompleks}
Tuliskan $u=(1,i)$; vektor kedua adalah $iu$.
Untuk $a,b\in\R$, hubungan $au+b(iu)=0$ berarti
$(a+bi)u=0$, sehingga $a=b=0$.
Di atas $\C$, hubungan $iu-(iu)=0$ memberi ketergantungan
antara kedua vektor dengan koefisien $i$ dan $-1$.
Ruang pembangun kompleksnya adalah $\{zu:z\in\C\}$;
sebagai ruang real, basisnya $(u,iu)$ dan dimensinya dua.

\subsection*{Latihan \ref{bebas:lat-jumlah-nol}: pilih blok bebas}
Pilih bebas entri $x_{ij}$ untuk $1\le i,j<n$.
Syarat jumlah nol memaksa
\[
    x_{in}=-\sum_{j<n}x_{ij},\qquad
    x_{nj}=-\sum_{i<n}x_{ij},\qquad
    x_{nn}=\sum_{i,j<n}x_{ij}.
\]
Rumus terakhir memenuhi baris terakhir sekaligus kolom terakhir;
jadi salah satu dari $2n$ persamaan semula berulang.
Basis eksplisitnya adalah
\[
    B_{ij}=E_{ij}-E_{in}-E_{nj}+E_{nn},\qquad 1\le i,j<n.
\]
Setiap $X\in S$ mempunyai ekspansi unik
$X=\sum_{i,j<n}x_{ij}B_{ij}$, sehingga $\dim S=(n-1)^2$.

\subsection*{Latihan \ref{bebas:lat-tiga}: dua kali batas irisan}
Terapkan batas dimensi irisan pertama pada $U_1,U_2$, kemudian
pada $U_1\cap U_2,U_3$:
\begin{align*}
    \dim(U_1\cap U_2\cap U_3)
    &\ge\dim(U_1\cap U_2)+\dim U_3-n\\
    &\ge\dim U_1+\dim U_2+\dim U_3-2n.
\end{align*}
Di $\R^3$, tiga bidang koordinat $x_1=0$, $x_2=0$, $x_3=0$
mencapai kesamaan: irisan ketiganya nol dan ruas kanan $2+2+2-6=0$.

\subsection*{Latihan \ref{bebas:lat-selisih-dimensi}: dua bilangan cacah}
Letakkan $k=\dim(U\cap W)$.
Rumus dimensi memberi $(\dim U-k)+(\dim W-k)=1$.
Kedua suku adalah bilangan cacah, sehingga salah satunya nol.
Jika $\dim U=k$, inklusi $U\cap W\subseteq U$ dan kesamaan
dimensi memberi $U=U\cap W\subseteq W$. Kasus lainnya serupa.

\subsection*{Latihan \ref{bebas:lat-lagrange}: bandingkan suku tertinggi}
Identitas interpolasi untuk $p(t)=t^k$ memberi
$t^k=\sum_{j=0}^n a_j^kL_j(t)$.
Koefisien $t^n$ pada $L_j$ adalah
$1/\prod_{\ell\ne j}(a_j-a_\ell)$.
Membandingkan koefisien $t^n$ pada kedua ruas memberi nilai
nol untuk $k<n$ dan satu untuk $k=n$.

\subsection*{Latihan \ref{bebas:lat-komplemen}: padankan bagian di luar irisan}
Pilih basis $(z_1,\ldots,z_k)$ dari $U\cap W$ dan perluas
menjadi $(z_1,\ldots,z_k,u_1,\ldots,u_r)$ pada $U$ serta
$(z_1,\ldots,z_k,w_1,\ldots,w_r)$ pada $W$.
Jumlah vektor tambahan sama karena $\dim U=\dim W$.
Daftar $(z_1,\ldots,z_k,u_1,\ldots,u_r,w_1,\ldots,w_r)$
adalah basis bagi $U+W$ sebagaimana pada bukti rumus dimensi.
Definisikan
\[
    K_0=\operatorname{span}\{u_1+w_1,\ldots,u_r+w_r\}.
\]
Jika $\sum a_j(u_j+w_j)\in U$, maka $\sum a_jw_j\in U\cap W$,
sehingga semua $a_j=0$. Jadi $K_0\cap U=\{0\}$; argumen
simetris memberi $K_0\cap W=\{0\}$.
Selain itu $w_j=(u_j+w_j)-u_j$ dan $u_j=(u_j+w_j)-w_j$,
sehingga $U+W=U\oplus K_0=W\oplus K_0$.
Ambil komplemen $L$ bagi $U+W$ di $V$ dan letakkan $K=K_0\oplus L$.
Maka $V=U\oplus K=W\oplus K$.
Kesamaan $U=W$ tidak wajib: di $\R^2$, pilih
$U=\operatorname{span}\{e_1\}$,
$W=\operatorname{span}\{e_1+e_2\}$, dan $K=\operatorname{span}\{e_2\}$.
\fi

% !TeX root = ../../ONMIPA Alin.tex
% Acuan: Zhang (2009), Bab 3, definisi dan fakta serta Soal 3.79--3.112;
% Zhang (2011), bagian 1.3. Contoh dan pembahasan disusun untuk handout ini.
\chapter{Transformasi Linear}
\label{ch:transformasi-linear}

Transformasi linear mempertahankan kombinasi linear. Karena itu,
perilakunya pada sebuah basis menentukan perilakunya di seluruh
ruang. Bab ini menghubungkan struktur ruang vektor dengan perhitungan
matriks, lalu memakai hubungan tersebut untuk menganalisis kernel,
citra, proyeksi, serta pangkat suatu operator.
Semua ruang dan transformasi menggunakan medan skalar yang sama,
yaitu $\F=\R$ atau $\C$, kecuali jika dinyatakan lain.

\section{Definisi dan Penentuan pada Basis}

\begin{definisi}[][Transformasi linear][trans:def]
    Pemetaan $T:V\to W$ disebut \emph{transformasi linear} jika
    untuk semua $u,v\in V$ dan $a,b\in\F$ berlaku
    \[
        T(au+bv)=aT(u)+bT(v).
    \]
    Transformasi linear dari $V$ ke dirinya sendiri disebut
    \emph{operator linear} pada $V$.
\end{definisi}

Definisi tersebut ekuivalen dengan aditivitas
$T(u+v)=T(u)+T(v)$ dan homogenitas $T(av)=aT(v)$.
Dengan mengambil koefisien nol, diperoleh $T(0_V)=0_W$.
Induksi kemudian memberi
\[
    T\left(\sum_{j=1}^r a_jv_j\right)
      =\sum_{j=1}^r a_jT(v_j).
\]
Jadi setiap hubungan linear antarvektor harus tetap berlaku
setelah transformasi diterapkan. Syarat $T(0)=0$ saja belum
cukup: pemetaan $T(x)=x^2$ pada $\R$ memenuhi syarat itu,
tetapi gagal pada penjumlahan.

\begin{exam}[Linear terhadap medan yang mana?]
    Pada $\C$, tinjau $T(z)=\overline z$ dan $S(z)=\operatorname{Re}z$,
    dengan nilai $S$ juga dipandang sebagai bilangan kompleks.
    Apakah keduanya linear atas $\R$? Atas $\C$?
\end{exam}
\begin{solution}
    Konjugasi dan bagian real mempertahankan penjumlahan dan
    perkalian oleh skalar real, sehingga keduanya linear atas $\R$.
    Namun, $T(i)=-i\ne iT(1)$ dan $S(i)=0\ne iS(1)$.
    Jadi keduanya tidak linear atas $\C$.
    Pilihan medan skalar merupakan bagian dari pernyataan linearitas.
\end{solution}

Pemetaan $x\mapsto Ax+b$ dari $\F^n$ ke $\F^m$ linear tepat
ketika $b=0$; syarat ini langsung diperoleh dari nilai di nol.
Pemetaan dengan $b\ne0$ disebut pemetaan afin.
Contoh linear lain mencakup evaluasi $p\mapsto p(a)$,
turunan $p\mapsto p'$, serta pemetaan matriks $X\mapsto AXB$
untuk $A,B$ tetap dengan ukuran yang sesuai.

\begin{teorema}[][Menentukan transformasi melalui basis][trans:basis]
    Misalkan $\mathcal B=(v_1,\ldots,v_n)$ basis bagi $V$.
    Untuk sembarang $w_1,\ldots,w_n\in W$, terdapat tepat satu
    transformasi linear $T:V\to W$ dengan $T(v_j)=w_j$.
    Transformasi tersebut diberikan oleh
    \[
        T\left(\sum_{j=1}^n a_jv_j\right)=\sum_{j=1}^n a_jw_j.
    \]
\end{teorema}
\begin{proof}
    Koordinat terhadap basis tunggal, sehingga rumus tersebut
    memberikan satu nilai yang tidak bergantung pada pilihan
    representasi. Rumusnya mempertahankan kombinasi linear.
    Setiap transformasi linear dengan nilai basis yang ditentukan
    harus memenuhi rumus itu, sehingga transformasinya unik.
\end{proof}

Jika nilai transformasi ditentukan pada daftar yang hanya
membangun ruang, semua hubungan pada daftar itu harus dihormati.
Secara tepat, aturan $v_j\mapsto w_j$ memperluas menjadi
transformasi pada $\operatorname{span}\{v_j\}$ jika dan hanya jika
\[
    \sum a_jv_j=0\quad\Longrightarrow\quad\sum a_jw_j=0.
\]
Syarat ini perlu karena linearitas, dan cukup karena dua
representasi dari vektor yang sama mempunyai selisih nol.

\begin{exam}[Data pada tiga vektor yang bukan basis standar]
    Carilah $T:\R^3\to\R^2$ yang memenuhi
    \[
        T(e_1+e_2)=(1,0),\quad
        T(e_2+e_3)=(0,1),\quad
        T(e_3+e_1)=(1,1).
    \]
\end{exam}
\begin{solution}
    Ketiga vektor sumber bebas linear, sebab determinan matriks
    kolomnya adalah $2$. Jadi transformasi ada dan unik.
    Tulis $a=T(e_1)$, $b=T(e_2)$, $c=T(e_3)$.
    Sistem $a+b=(1,0)$, $b+c=(0,1)$, $c+a=(1,1)$
    memberi $a=(1,0)$, $b=(0,0)$, $c=(0,1)$.
    Maka $T(x,y,z)=(x,z)$.
\end{solution}

\section{Kernel, Citra, dan Teorema Rank--Nulitas}

\begin{definisi}[][Kernel, citra, rank, dan nulitas][trans:kernel]
    Untuk transformasi linear $T:V\to W$, definisikan
    \[
        \ker T=\{v\in V:T(v)=0\},\qquad
        \operatorname{im}T=\{T(v):v\in V\}.
    \]
    Jika dimensinya hingga, tulis
    $\nul T=\dim\ker T$ dan $\rank T=\dim\operatorname{im}T$.
\end{definisi}

Kernel merupakan subruang \emph{domain} $V$, sedangkan citra
merupakan subruang \emph{kodomain} $W$.
Untuk kernel, $T(au+bv)=0$ jika $T(u)=T(v)=0$.
Untuk citra, kombinasi $aT(u)+bT(v)$ sama dengan $T(au+bv)$.
Kedua himpunan memuat nol, sehingga keduanya memenuhi uji subruang.

\begin{teorema}[][Teorema rank--nulitas][trans:rank-nulitas]
    Jika $T:V\to W$ linear dan $\dim V=n<\infty$, maka
    \[
        \rank T+\nul T=n.
    \]
    Lebih tepatnya, ambil basis $(u_1,\ldots,u_k)$ bagi $\ker T$,
    lalu perluas menjadi basis
    \[
        (u_1,\ldots,u_k,v_1,\ldots,v_{n-k})
    \]
    bagi $V$. Daftar $(T(v_1),\ldots,T(v_{n-k}))$ merupakan basis citra.
\end{teorema}
\begin{proof}
    Ekspansi vektor domain dalam basis yang diberikan menunjukkan
    bahwa citra dibangun oleh $T(v_1),\ldots,T(v_{n-k})$.
    Jika $\sum a_jT(v_j)=0$, maka $\sum a_jv_j\in\ker T$
    dan dapat dituliskan sebagai kombinasi $u_1,\ldots,u_k$.
    Kebebasan basis $V$ memaksa semua $a_j=0$.
    Maka dimensi citra adalah $n-k$, sedangkan nulitasnya $k$.
\end{proof}

\begin{prop}[Injektif, surjektif, dan kebebasan citra]
    Transformasi $T$ injektif jika dan hanya jika $\ker T=\{0\}$,
    dan surjektif jika dan hanya jika $\operatorname{im}T=W$.
    Jika $T(v_1),\ldots,T(v_r)$ bebas linear, maka $v_1,\ldots,v_r$
    bebas linear. Arah sebaliknya berlaku jika $T$ injektif.
\end{prop}
\begin{proof}
    Persamaan $T(u)=T(v)$ ekuivalen dengan $u-v\in\ker T$,
    sehingga kriteria injektivitas mengikuti langsung.
    Kriteria surjektivitas adalah definisi citra.
    Untuk kebebasan, terapkan $T$ pada suatu hubungan linear
    vektor sumber. Sebaliknya, jika citra suatu kombinasi nol
    dan $T$ injektif, kombinasi sumbernya juga nol.
\end{proof}

Jika $\dim V=\dim W=n<\infty$, teorema rank--nulitas memberi
ekuivalensi: $T$ injektif, $T$ surjektif, dan $T$ bijektif.
Untuk dimensi berbeda, injektivitas hanya mengharuskan
$\dim V\le\dim W$, sedangkan surjektivitas mengharuskan
$\dim V\ge\dim W$. Ketaksamaan dimensi saja tidak menjamin
sifat tersebut bagi transformasi yang diberikan.

\begin{exam}[Evaluasi polinom pada beberapa titik]
    Untuk titik berbeda $a_1,\ldots,a_m\in\F$, $m\ge1$, tinjau
    \[
        E:\mathcal P_{\le n}(\F)\to\F^m,\qquad
        E(p)=(p(a_1),\ldots,p(a_m))^T.
    \]
    Tentukan rank, nulitas, dan kapan $E$ injektif atau surjektif.
\end{exam}
\begin{solution}
    Kernel terdiri atas polinom yang habis dibagi
    $q(t)=\prod_{j=1}^m(t-a_j)$.
    Jika $m\le n$, kernel adalah
    $q\mathcal P_{\le n-m}(\F)$, berdimensi $n-m+1$.
    Jika $m>n$, polinom berderajat paling besar $n$ dengan
    $m$ akar berbeda harus nol. Jadi
    \[
        \nul E=\max\{n+1-m,0\},\qquad
        \rank E=\min\{n+1,m\}.
    \]
    Transformasi injektif tepat ketika $m\ge n+1$ dan surjektif
    tepat ketika $m\le n+1$. Untuk $m=n+1$, interpolasi
    Lagrange memberikan inversnya secara eksplisit.
\end{solution}

\subsection{Serat solusi dan prapeta subruang}
Jika $T(v_0)=w$, seluruh solusi $T(v)=w$ adalah
\[
    T^{-1}(\{w\})=v_0+\ker T.
\]
Memang, $T(v)=T(v_0)$ jika dan hanya jika $v-v_0\in\ker T$.
Jadi suatu persamaan linear konsisten tepat ketika
$w\in\operatorname{im}T$, dan jika konsisten, solusinya unik
tepat ketika kernel hanya nol.

Untuk subruang $Y\subseteq W$, prapeta
$T^{-1}(Y)=\{v:T(v)\in Y\}$ selalu subruang dari $V$.
Notasi prapeta ini tidak mensyaratkan bahwa $T$ mempunyai invers.
Pembatasan $T$ pada $T^{-1}(Y)$ mempunyai kernel $\ker T$
dan citra $Y\cap\operatorname{im}T$. Maka, jika $V$ berdimensi hingga,
\[
    \dim T^{-1}(Y)=\nul T+\dim(Y\cap\operatorname{im}T).
\]
Demikian pula, untuk subruang $U\subseteq V$,
$\dim T(U)=\dim U-\dim(U\cap\ker T)$.

\begin{exam}[Jumlah dimensi tidak menjamin jumlah langsung]
    Pada $\R^2$, definisikan $T(x,y)=(y,0)$.
    Tentukan kernel, citra, dan jumlah keduanya.
\end{exam}
\begin{solution}
    Kernel dan citra keduanya $\operatorname{span}\{e_1\}$.
    Walaupun $\nul T+\rank T=1+1=2$, jumlah kedua subruang
    hanya garis $\operatorname{span}\{e_1\}$, bukan seluruh $\R^2$.
    Teorema rank--nulitas menyatakan kesamaan bilangan dimensi;
    ia tidak menyatakan bahwa domain adalah jumlah langsung
    kernel dan citra. Untuk domain dan kodomain berbeda,
    penjumlahan kernel dan citra bahkan belum tentu bermakna.
\end{solution}

\section{Representasi Matriks dan Pergantian Basis}

\begin{definisi}[][Matriks suatu transformasi][trans:matriks]
    Misalkan $\mathcal B=(b_1,\ldots,b_n)$ basis $V$ dan
    $\mathcal C=(c_1,\ldots,c_m)$ basis $W$.
    Matriks transformasi $T:V\to W$ relatif terhadap kedua basis adalah
    \[
        [T]_{\mathcal C\leftarrow\mathcal B}
        =\bigl[[T(b_1)]_{\mathcal C}\ \cdots\
        [T(b_n)]_{\mathcal C}\bigr]\in\Mat_{m\times n}(\F).
    \]
\end{definisi}

Dengan menuliskan $v=\sum x_jb_j$, linearitas memberi
\[
    [T(v)]_{\mathcal C}
      =[T]_{\mathcal C\leftarrow\mathcal B}[v]_{\mathcal B}.
\]
Jadi kolom ke-$j$ berisi \emph{koordinat citra vektor basis
ke-$j$}, bukan barisnya. Untuk operator dengan basis domain dan
kodomain sama, notasi disingkat menjadi $[T]_{\mathcal B}$.
Rank dan nulitas matriks ini sama dengan rank dan nulitas $T$,
karena persamaan koordinat mempertahankan hubungan linear.

\begin{exam}[Turunan dalam dua basis polinom]
    Pada $V=\mathcal P_{\le3}(\R)$, letakkan $D(p)=p'$.
    Tentukan matriks $D$ dalam basis
    $\mathcal B=(1,t,t^2,t^3)$ dan
    $\mathcal C=(1,t,t^2/2,t^3/6)$.
\end{exam}
\begin{solution}
    Turunan anggota basis langsung memberi
    \[
        [D]_{\mathcal B}
        =\begin{pmatrix}0&1&0&0\\0&0&2&0\\0&0&0&3\\0&0&0&0\end{pmatrix},
        \qquad
        [D]_{\mathcal C}
        =\begin{pmatrix}0&1&0&0\\0&0&1&0\\0&0&0&1\\0&0&0&0\end{pmatrix}.
    \]
    Kedua matriks menyatakan operator yang sama.
    Basis kedua menyerap faktor dari penurunan pangkat.
    Kernel dibangun oleh $1$, citra adalah $\mathcal P_{\le2}(\R)$,
    dan $D^4=0$ pada ruang ini.
\end{solution}

\begin{teorema}[][Rumus pergantian basis][trans:ganti-basis]
    Tulis $A=[T]_{\mathcal C\leftarrow\mathcal B}$,
    $P=P_{\mathcal B\leftarrow\mathcal B'}$, dan
    $Q=P_{\mathcal C\leftarrow\mathcal C'}$.
    Maka
    \[
        [T]_{\mathcal C'\leftarrow\mathcal B'}=Q^{-1}AP.
    \]
    Khusus untuk operator dengan perubahan basis yang sama
    pada domain dan kodomain, matriks baru adalah $P^{-1}AP$.
\end{teorema}
\begin{proof}
    Untuk setiap $v$, lakukan perubahan koordinat berurutan:
    \[
        [T(v)]_{\mathcal C'}
        =Q^{-1}[T(v)]_{\mathcal C}
        =Q^{-1}A[v]_{\mathcal B}
        =Q^{-1}AP[v]_{\mathcal B'}.
    \]
    Keunikan matriks representasi memberikan rumus yang dinyatakan.
\end{proof}

Rumus ini menjelaskan mengapa keserupaan matriks muncul pada
operator linear: dua matriks serupa merepresentasikan operator
yang sama dalam basis berbeda. Akibatnya, determinan, trace,
polinom karakteristik, dan nilai eigen operator tidak bergantung
pada pilihan basis. Untuk transformasi antar dua ruang dengan
perubahan basis terpisah, bentuk umumnya tetap $Q^{-1}AP$;
bentuk itu tidak harus merupakan keserupaan.

\section{Komposisi, Invers, dan Ruang Transformasi}

Jika $T:V\to W$ dan $S:W\to Z$ linear, komposisi
$ST=S\circ T$ linear dan berarti \emph{menerapkan $T$ lebih dahulu}.
Dengan basis yang sesuai,
\[
    [ST]_{\mathcal D\leftarrow\mathcal B}
      =[S]_{\mathcal D\leftarrow\mathcal C}
       [T]_{\mathcal C\leftarrow\mathcal B}.
\]
Urutan perkalian mengikuti urutan komposisi; dua operator pada
ruang yang sama tidak harus komutatif.

\begin{teorema}[][Rank komposisi][trans:komposisi]
    Untuk $T:V\to W$ dan $S:W\to Z$ dengan semua ruang
    berdimensi hingga,
    \[
        \rank(ST)=\rank T-\dim(\operatorname{im}T\cap\ker S).
    \]
    Akibatnya,
    \[
        \rank S+\rank T-\dim W
        \le\rank(ST)\le\min\{\rank S,\rank T\}.
    \]
\end{teorema}
\begin{proof}
    Batasi $S$ pada $\operatorname{im}T$.
    Kernel pembatasannya adalah $\operatorname{im}T\cap\ker S$,
    sedangkan citranya $\operatorname{im}(ST)$.
    Teorema rank--nulitas memberikan rumus utama.
    Karena dimensi irisan paling besar $\dim\ker S
    =\dim W-\rank S$, diperoleh batas bawah.
    Batas atas oleh $\rank T$ mengikuti rumus, dan batas oleh
    $\rank S$ mengikuti $\operatorname{im}(ST)\subseteq\operatorname{im}S$.
\end{proof}

Jika $T$ bijektif, inversnya linear: untuk $w_j=T(v_j)$,
persamaan $T(av_1+bv_2)=aw_1+bw_2$ memberi
$T^{-1}(aw_1+bw_2)=av_1+bv_2$.
Transformasi linear bijektif disebut \emph{isomorfisme}.
Dua ruang berdimensi hingga isomorfik jika dan hanya jika
dimensinya sama: kirim basis ruang pertama ke basis ruang kedua.

\begin{exam}[Invers sebagai jumlah hingga turunan]
    Buktikan bahwa $I-D$ invertibel pada $\mathcal P_{\le n}(\F)$,
    dengan $D(p)=p'$, dan selesaikan $p-p'=q$ di ruang tersebut.
\end{exam}
\begin{solution}
    Karena $D^{n+1}=0$, identitas deret geometri operator memberi
    \[
        (I-D)(I+D+\cdots+D^n)=I-D^{n+1}=I.
    \]
    Perkalian dalam urutan terbalik juga bernilai $I$.
    Jadi $(I-D)^{-1}=I+D+\cdots+D^n$ dan satu-satunya solusi
    adalah $p=q+q'+\cdots+q^{(n)}$.
    Secara umum, jika $N^k=0$, maka
    $(I-N)^{-1}=I+N+\cdots+N^{k-1}$.
\end{solution}

Pada ruang berdimensi hingga yang sama, $ST=I$ mengakibatkan
$TS=I$: persamaan pertama membuat $T$ injektif, sehingga
$T$ bijektif dan $S=T^{-1}$.
Untuk ruang tak hingga, injektif dan surjektif tidak ekuivalen.
Pada $\F[t]$, operator turunan surjektif tetapi tidak injektif,
sedangkan perkalian oleh $t$ injektif tetapi tidak surjektif.

\subsection{Ruang semua transformasi linear}
Himpunan $\mathcal L(V,W)$ menjadi ruang linear dengan operasi
\[
    (S+T)(v)=S(v)+T(v),\qquad (aT)(v)=aT(v).
\]
Jika $\dim V=n$ dan $\dim W=m$, maka
$\dim\mathcal L(V,W)=mn$.
Untuk membuktikannya, pilih basis $(b_j)$ dan $(c_i)$, lalu
definisikan $E_{ij}^{\mathcal L}$ melalui
\[
    E_{ij}^{\mathcal L}(b_k)=\begin{cases}c_i,&k=j,\\0,&k\ne j.\end{cases}
\]
Setiap $T$ mempunyai ekspansi unik
$T=\sum_{i,j}a_{ij}E_{ij}^{\mathcal L}$, dengan $a_{ij}$ entri
matriks representasinya. Ini memberi basis yang beranggotakan $mn$.
Tidak diperlukan hubungan khusus antara vektor basis domain
dan vektor basis kodomain.

\section{Proyeksi dan Subruang Invarian}

\begin{definisi}[][Proyeksi dan subruang invarian][trans:proyeksi]
    Operator $P:V\to V$ disebut \emph{proyeksi} jika $P^2=P$.
    Subruang $U\subseteq V$ disebut \emph{invarian} terhadap
    operator $T$ jika $T(U)\subseteq U$.
\end{definisi}

Proyeksi di sini tidak memerlukan hasil kali dalam. Istilah
proyeksi ortogonal memuat syarat tambahan yang akan bergantung
pada struktur hasil kali dalam.

\begin{teorema}[][Proyeksi dan dekomposisi ruang][trans:dekomposisi-proyeksi]
    Jika $P^2=P$, maka
    \[
        V=\operatorname{im}P\oplus\ker P,\qquad
        \operatorname{im}(I-P)=\ker P.
    \]
    Sebaliknya, jika $V=U\oplus W$, terdapat tepat satu
    proyeksi dengan citra $U$ dan kernel $W$, yaitu $P(u+w)=u$.
\end{teorema}
\begin{proof}
    Setiap $v$ dapat ditulis sebagai $Pv+(v-Pv)$.
    Suku pertama berada di citra, sedangkan suku kedua berada
    di kernel karena $P(v-Pv)=Pv-P^2v=0$.
    Jika $x=Py\in\ker P$, maka $0=Px=P^2y=Py=x$;
    jadi jumlahnya langsung.
    Selain itu, setiap $(I-P)v$ berada di kernel, dan untuk
    $x\in\ker P$ berlaku $x=(I-P)x$.
    Arah sebaliknya mengikuti keunikan dekomposisi $u+w$.
    Proyeksi dengan citra $U$ bertindak sebagai identitas pada $U$
    dan bernilai nol pada $W$, sehingga rumusnya juga unik.
\end{proof}

Jika $\dim V=n$ dan $\rank P=r$, gabungan basis citra dan
kernel memberi matriks $\diag(I_r,0)$.
Maka $\tr P=r$ dan $\rank(I-P)=n-r$.
Untuk operator umum, kernel dan citra memang invarian terhadap
operator itu sendiri, tetapi tidak harus saling komplemen.

\begin{exam}[Memisahkan bagian simetris suatu matriks]
    Pada $\Mat_n(\F)$, tinjau $P(X)=(X+X^T)/2$.
    Tentukan citra, kernel, dan rank dari $P$.
\end{exam}
\begin{solution}
    Transpose linear atas $\F$, sehingga $P$ linear.
    Matriks $P(X)$ simetris, dan untuk matriks simetris $S$
    berlaku $P(S)=S$. Jadi $P^2=P$ dan citranya adalah ruang
    matriks simetris. Kernelnya ruang matriks antisimetris,
    yaitu $X^T=-X$. Dengan menghitung entri bebas,
    \[
        \rank P=\frac{n(n+1)}2,\qquad
        \nul P=\frac{n(n-1)}2.
    \]
    Dekomposisi yang dihasilkan adalah
    $X=(X+X^T)/2+(X-X^T)/2$.
    Jika transpose diganti transpose konjugat, pemetaannya
    linear atas $\R$, tetapi umumnya tidak linear atas $\C$.
\end{solution}

\subsection{Basis yang disesuaikan dengan subruang}
Jika $U$ invarian terhadap $T$, perluas basis $U$ menjadi basis
$V$. Dalam basis ini, matriks $T$ berbentuk
\[
    \begin{pmatrix}A&B\\0&C\end{pmatrix}.
\]
Blok nol muncul karena citra anggota basis $U$ tidak mempunyai
komponen di luar $U$.
Jika $V=U\oplus W$ dan kedua subruang invarian, matriksnya
berbentuk blok diagonal $\diag(A,C)$.
Komplemen suatu subruang selalu ada dalam dimensi hingga,
tetapi komplemen yang invarian belum tentu ada.

Jika $ST=TS$, maka $S$ mempertahankan $\ker T$ dan
$\operatorname{im}T$. Untuk $v\in\ker T$, berlaku
$T(Sv)=S(Tv)=0$; untuk $v=Tu$, berlaku $Sv=T(Su)$.
Argumen yang sama menunjukkan bahwa setiap ruang eigen
$\ker(T-\lambda I)$ invarian terhadap $S$.

\section{Operator Nilpoten dan Stabilisasi Kernel}

Operator $N$ disebut nilpoten jika $N^k=0$ untuk suatu $k\ge1$.
Rantai citra suatu vektor sering menghasilkan daftar bebas linear
tanpa menghitung determinan.

\begin{teorema}[][Rantai nilpoten][trans:rantai]
    Jika $N^kv=0$ dan $N^{k-1}v\ne0$, maka
    \[
        v,Nv,\ldots,N^{k-1}v
    \]
    bebas linear. Khususnya, indeks nilpotensi operator pada
    ruang berdimensi $n\ge1$ paling besar $n$.
\end{teorema}
\begin{proof}
    Andaikan $\sum_{j=0}^{k-1}a_jN^jv=0$ mempunyai koefisien
    tidak nol, dan pilih indeks terkecil $j_0$ dengan $a_{j_0}\ne0$.
    Terapkan $N^{k-1-j_0}$ pada hubungan itu.
    Suku dengan indeks lebih besar lenyap, sehingga
    $a_{j_0}N^{k-1}v=0$, kontradiksi.
    Jika $k$ adalah indeks nilpotensi, ada $v$ dengan
    $N^{k-1}v\ne0$, sehingga $k\le n$.
\end{proof}

\begin{exam}[Basis dari satu rantai turunan]
    Pada $\mathcal P_{\le n}(\F)$, ambil $v=t^n$ dan $D(p)=p'$.
    Tunjukkan bahwa $v,Dv,\ldots,D^nv$ basis dan tentukan
    $\nul D^j$ untuk $0\le j\le n+1$.
\end{exam}
\begin{solution}
    Karena $D^nv=n!\ne0$ dan $D^{n+1}v=0$, teorema rantai
    memberi $n+1$ vektor bebas, sehingga semuanya membentuk basis.
    Untuk $1\le j\le n+1$, kernel $D^j$ terdiri atas polinom
    berderajat paling besar $j-1$, maka $\nul D^j=j$.
    Untuk $j=0$, $D^0=I$ dan nulitasnya nol.
\end{solution}

Untuk operator umum $T:V\to V$ dengan $\dim V=n$, selalu ada
rantai inklusi
\[
    \ker T^0\subseteq\ker T\subseteq\ker T^2\subseteq\cdots,
    \qquad
    V\supseteq\operatorname{im}T\supseteq\operatorname{im}T^2\supseteq\cdots.
\]
Rantai kernel tidak dapat membesar secara ketat lebih dari $n$ kali.
Jika $\ker T^k=\ker T^{k+1}$, maka kesamaan berlaku untuk semua
pangkat selanjutnya. Memang, jika $T^{k+2}v=0$, maka
$Tv\in\ker T^{k+1}=\ker T^k$, sehingga $T^{k+1}v=0$.
Ulangi argumen ini secara induktif.

\begin{teorema}[][Dekomposisi setelah kernel stabil][trans:stabil]
    Misalkan $T$ operator pada $V$ dengan $\dim V=n<\infty$.
    Jika $\ker T^k=\ker T^{k+1}$ untuk suatu $k\ge1$, maka
    \[
        V=\ker T^k\oplus\operatorname{im}T^k.
    \]
    Kedua subruang invarian; pembatasan $T$ pada yang pertama
    nilpoten, dan pada yang kedua invertibel.
    Untuk $\dim V=n\ge1$, pilihan $k=n$ selalu dapat digunakan.
\end{teorema}
\begin{proof}
    Stabilisasi memberi $\ker T^{2k}=\ker T^k$.
    Jika $x=T^ky\in\ker T^k$, maka $T^{2k}y=0$, sehingga
    $y\in\ker T^k$ dan $x=0$.
    Jadi irisan kernel dan citra $T^k$ hanya nol; rank--nulitas
    menyatakan jumlah dimensinya $n$, sehingga jumlahnya seluruh $V$.
    Keduanya invarian karena $T$ komutatif dengan pangkatnya.
    Pada kernel, pangkat ke-$k$ bernilai nol.
    Pada citra, jika $Tx=0$, maka $x\in\ker T^k$ juga,
    sehingga $x=0$. Pembatasan ini injektif pada ruang berdimensi
    hingga, maka invertibel.
    Rantai kernel telah stabil paling lambat pada pangkat $n$.
\end{proof}

Kasus $k=1$ memberi kriteria yang sering berguna:
\[
    V=\ker T\oplus\operatorname{im}T
    \iff\rank T^2=\rank T.
\]
Satu arah mengikuti teorema dan rank--nulitas. Untuk arah lainnya,
jika jumlahnya langsung dan $T^2v=0$, maka
$Tv\in\ker T\cap\operatorname{im}T=\{0\}$, sehingga
$\ker T^2=\ker T$.
Sifat ini tidak memaksa $T$ idempoten: $T=2I$ memenuhi
dekomposisi tersebut, tetapi $T^2\ne T$ pada ruang tak nol.

\section{Latihan Pilihan}

\begin{latihan}\label{trans:lat-data}
    Tentukan semua transformasi linear
    $T:\mathcal P_{\le2}(\R)\to\R^2$ yang memenuhi
    \[
        T(1)=(1,0),\quad T(t)=(0,1),\quad T(1+t)=(a,b).
    \]
    Untuk parameter yang memungkinkan, apakah $T$ unik?
    Tentukan rank dan nulitas setiap transformasi yang mungkin.
\end{latihan}

\begin{latihan}\label{trans:lat-jet}
    Definisikan $T:\mathcal P_{\le4}(\R)\to\R^3$ melalui
    $T(p)=(p(0),p'(0),p(1))^T$.
    Tentukan matriks terhadap basis standar, basis kernel,
    serta rank. Carilah semua $p$ dengan $T(p)=(\alpha,\beta,\gamma)^T$.
\end{latihan}

\begin{latihan}\label{trans:lat-basis}
    Pada $\R^2$, definisikan $T(x,y)=(2x+y,x+2y)$.
    Gunakan basis $\mathcal C=((1,1)^T,(1,-1)^T)$ untuk
    menentukan $[T]_{\mathcal C}$ dan rumus $T^k(x,y)$
    untuk setiap bilangan bulat $k$, termasuk $k<0$.
\end{latihan}

\begin{latihan}\label{trans:lat-euler}
    Untuk $a\in\F$, tinjau operator
    $T_a(p)=tp'-ap$ pada $\mathcal P_{\le n}(\F)$.
    Tentukan kernel, citra, dan syarat invertibilitasnya.
    Jika invertibel, tuliskan $T_a^{-1}$ secara eksplisit.
\end{latihan}

\begin{latihan}\label{trans:lat-komutator}
    Ambil $A=\diag(1,1,2)$ dan definisikan
    $T:\Mat_3(\F)\to\Mat_3(\F)$ melalui $T(X)=AX-XA$.
    Tentukan basis kernel dan citra serta rank dan nulitasnya.
    Generalisasikan hasilnya untuk matriks diagonal dengan
    nilai diagonal berbeda yang bermultiplisitas $m_1,\ldots,m_s$.
\end{latihan}

\begin{latihan}\label{trans:lat-kuadrat-nol}
    Misalkan $T^2=0$ pada ruang berdimensi $n$.
    Buktikan bahwa $\rank T\le\lfloor n/2\rfloor$.
    Tunjukkan bahwa setiap rank bulat antara nol dan batas itu
    dapat dicapai. Kapan $\operatorname{im}T=\ker T$?
\end{latihan}

\begin{latihan}\label{trans:lat-dua-proyeksi}
    Untuk proyeksi $P,Q$ pada ruang yang sama, buktikan bahwa
    $P+Q$ proyeksi jika dan hanya jika $PQ=QP=0$.
    Dalam keadaan itu, buktikan
    \[
        \operatorname{im}(P+Q)=\operatorname{im}P\oplus\operatorname{im}Q.
    \]
\end{latihan}

\begin{latihan}\label{trans:lat-invers-sepihak}
    Misalkan $T:V\to W$ dan $S:W\to V$ linear dengan
    $ST=I_V$, serta kedua ruang berdimensi hingga.
    Buktikan bahwa $TS$ proyeksi pada $W$ dan
    \[
        W=\operatorname{im}T\oplus\ker S.
    \]
    Tentukan kapan $TS=I_W$.
\end{latihan}

\begin{latihan}\label{trans:lat-irisan}
    Untuk transformasi $T:V\to W$ dan subruang $U_1,U_2\subseteq V$,
    buktikan
    \[
        T(U_1+U_2)=T(U_1)+T(U_2),\qquad
        T(U_1\cap U_2)\subseteq T(U_1)\cap T(U_2).
    \]
    Berikan contoh inklusi kedua yang sejati, dan buktikan bahwa
    kesamaan berlaku jika $T$ injektif.
\end{latihan}

\begin{latihan}\label{trans:lat-faktorisasi}
    Misalkan $A:V\to W$ dan $B:V\to Z$ linear, dengan semua
    ruang berdimensi hingga. Buktikan bahwa
    \[
        \ker A\subseteq\ker B
        \iff B=CA\text{ untuk suatu transformasi linear }C:W\to Z.
    \]
    Tunjukkan bahwa $C$ unik jika $A$ surjektif.
\end{latihan}

\begin{latihan}\label{trans:lat-komplemen-invarian}
    Pada $\R^2$, definisikan $N(x,y)=(y,0)$ dan
    $U=\operatorname{span}\{e_1\}$.
    Buktikan bahwa $U$ invarian, tetapi tidak mempunyai
    komplemen invarian terhadap $N$.
\end{latihan}

\begin{latihan}\label{trans:lat-komutator-identitas}
    Pada $\F[t]$, definisikan $D(p)=p'$ dan $M(p)=tp$.
    Buktikan $DM-MD=I$ dan, untuk $k\ge1$,
    \[
        D^kM-MD^k=kD^{k-1}.
    \]
    Buktikan pula bahwa pada ruang tak nol berdimensi hingga
    atas $\R$ atau $\C$, tidak ada operator $S,T$ dengan $ST-TS=I$.
\end{latihan}

\ifpembahasan
\section{Pembahasan Latihan}

\subsection*{Latihan \ref{trans:lat-data}: nilai pada daftar bergantung}
Linearitas memaksa $T(1+t)=T(1)+T(t)=(1,1)$,
sehingga transformasi ada tepat ketika $(a,b)=(1,1)$.
Nilai $T(t^2)=(c,d)$ dapat dipilih sembarang.
Seluruh transformasi yang memenuhi adalah
\[
    T(x+yt+zt^2)=(x+cz,y+dz),\qquad c,d\in\R.
\]
Jadi transformasi tidak unik. Citra memuat $(1,0)$ dan $(0,1)$,
sehingga rank selalu dua dan nulitas satu.
Kernelnya dibangun oleh $t^2-c-dt$.

\subsection*{Latihan \ref{trans:lat-jet}: nilai dan turunan sebagai kendala}
Matriksnya terhadap $(1,t,t^2,t^3,t^4)$ dan basis standar $\R^3$ adalah
\[
    [T]=\begin{pmatrix}1&0&0&0&0\\0&1&0&0&0\\1&1&1&1&1\end{pmatrix}.
\]
Kernel terdiri atas kelipatan $t^2(t-1)$ yang berderajat paling
besar empat. Basisnya $(t^2(t-1),t^3(t-1))$, sehingga nulitas
dua dan rank tiga. Solusi khusus untuk target yang diberikan adalah
\[
    p_0(t)=\alpha(1-t^2)+\beta(t-t^2)+\gamma t^2.
\]
Memeriksa nilai dan turunannya menunjukkan $T(p_0)=(\alpha,\beta,\gamma)^T$.
Semua solusi diperoleh dengan menambahkan kernel:
$p=p_0+t^2(t-1)(c+dt)$, dengan $c,d\in\R$.

\subsection*{Latihan \ref{trans:lat-basis}: pangkat dalam basis yang tepat}
Vektor $c_1=(1,1)^T$ dan $c_2=(1,-1)^T$ memenuhi
$T(c_1)=3c_1$ dan $T(c_2)=c_2$. Jadi $[T]_{\mathcal C}=\diag(3,1)$.
Karena kedua entri diagonal tak nol, $T$ invertibel dan
pangkat bilangan bulat terdefinisi. Ekspansi
\[
    (x,y)^T=\frac{x+y}{2}c_1+\frac{x-y}{2}c_2
\]
memberi
\[
    T^k(x,y)=\left(
      \frac{3^k(x+y)+x-y}{2},\
      \frac{3^k(x+y)-x+y}{2}\right).
\]

\subsection*{Latihan \ref{trans:lat-euler}: bekerja pada setiap monomial}
Karena $T_a(t^j)=(j-a)t^j$, matriksnya terhadap basis monomial
adalah $\diag(-a,1-a,\ldots,n-a)$.
Jika $a\notin\{0,1,\ldots,n\}$, semua entri tak nol,
kernel hanya nol, citra seluruh ruang, dan
\[
    T_a^{-1}\left(\sum_{j=0}^n b_jt^j\right)
      =\sum_{j=0}^n\frac{b_j}{j-a}t^j.
\]
Jika $a=r\in\{0,\ldots,n\}$, kernel dibangun oleh $t^r$ dan
citra terdiri atas polinom yang koefisien $t^r$-nya nol.
Ranknya $n$, nulitasnya satu, dan operator tidak invertibel.

\subsection*{Latihan \ref{trans:lat-komutator}: matriks satuan sebagai basis}
Untuk $A=\diag(\lambda_1,\ldots,\lambda_n)$, berlaku
\[
    T(E_{ij})=(\lambda_i-\lambda_j)E_{ij}.
\]
Pada kasus $\diag(1,1,2)$, basis kernel adalah
$(E_{11},E_{12},E_{21},E_{22},E_{33})$, sedangkan basis citra
adalah $(E_{13},E_{23},E_{31},E_{32})$.
Jadi nulitas lima dan rank empat.
Secara umum, kernel dibangun oleh $E_{ij}$ dengan
$\lambda_i=\lambda_j$ dan citra oleh sisanya. Kelompok nilai
diagonal bermultiplisitas $m_\ell$ menyumbang $m_\ell^2$ anggota
basis kernel. Maka
\[
    \nul T=\sum_{\ell=1}^s m_\ell^2,\qquad
    \rank T=n^2-\sum_{\ell=1}^s m_\ell^2.
\]

\subsection*{Latihan \ref{trans:lat-kuadrat-nol}: citra termuat dalam kernel}
Dari $T^2=0$, diperoleh $\operatorname{im}T\subseteq\ker T$.
Jika $r=\rank T$, maka $r\le n-r$, sehingga $r\le\lfloor n/2\rfloor$.
Untuk setiap $r$ dalam rentang itu, pilih basis $(e_1,\ldots,e_n)$
dan tetapkan $T(e_{2j})=e_{2j-1}$ untuk $1\le j\le r$,
serta $T(e_k)=0$ untuk semua indeks lainnya.
Operator ini berpangkat dua nol dan citranya dibangun oleh
$e_1,e_3,\ldots,e_{2r-1}$, sehingga ranknya $r$.
Inklusi citra dalam kernel menjadi kesamaan tepat ketika
dimensinya sama, yaitu $n=2r$.

\subsection*{Latihan \ref{trans:lat-dua-proyeksi}: dua perkalian harus lenyap}
Jika $P+Q$ idempoten, ekspansi kuadrat memberi $PQ+QP=0$.
Kalikan di kiri dan di kanan oleh $P$, secara terpisah:
\[
    PQ+PQP=0,\qquad PQP+QP=0.
\]
Mengurangkan keduanya memberi $PQ=QP$.
Karena $PQ+QP=0$ dan medan berkarakteristik nol, diperoleh
$PQ=QP=0$. Arah sebaliknya langsung dari ekspansi kuadrat.
Jika $x\in\operatorname{im}P\cap\operatorname{im}Q$, maka
$Px=x$, tetapi $x=Qy$ memberi $Px=PQy=0$; jadi $x=0$.
Untuk $u=Pa$ dan $v=Qb$ berlaku $(P+Q)(u+v)=u+v$.
Maka setiap anggota jumlah citra berada di citra $P+Q$;
inklusi sebaliknya langsung dari rumus $P+Q$.

\subsection*{Latihan \ref{trans:lat-invers-sepihak}: proyeksi pada ruang yang lebih besar}
Persamaan $ST=I_V$ membuat $T$ injektif dan $S$ surjektif.
Selain itu, $(TS)^2=T(ST)S=TS$, sehingga $TS$ proyeksi.
Citranya sama dengan citra $T$: satu inklusi jelas, dan
untuk $Tv$ berlaku $TS(Tv)=Tv$.
Kernelnya sama dengan kernel $S$, sebab $TSw=0$ dan injektivitas
$T$ mengakibatkan $Sw=0$.
Teorema proyeksi memberi $W=\operatorname{im}T\oplus\ker S$.
Selanjutnya, $TS=I_W$ tepat ketika $T$ surjektif, atau ekuivalen
dengan $S$ injektif. Dalam dimensi hingga dan dengan $ST=I_V$,
ini terjadi tepat ketika $\dim V=\dim W$.

\subsection*{Latihan \ref{trans:lat-irisan}: prapeta yang berbeda}
Linearitas memberi $T(u_1+u_2)=T(u_1)+T(u_2)$, sehingga
kesamaan jumlah berlaku. Jika $v\in U_1\cap U_2$, citranya
berada dalam kedua citra, membuktikan inklusi irisan.
Untuk inklusi sejati, ambil $T:\R^2\to\R$,
$T(x,y)=x+y$, serta $U_1=\operatorname{span}\{e_1\}$ dan
$U_2=\operatorname{span}\{e_2\}$.
Irisan sumber hanya nol, tetapi kedua citra sama dengan $\R$.
Jika $T$ injektif dan $w=T(u_1)=T(u_2)$ dengan $u_j\in U_j$,
maka $u_1=u_2\in U_1\cap U_2$, sehingga inklusi balik berlaku.

\subsection*{Latihan \ref{trans:lat-faktorisasi}: definisikan dahulu pada citra}
Jika $B=CA$, setiap $v\in\ker A$ memenuhi $Bv=C(Av)=0$.
Sebaliknya, andaikan $\ker A\subseteq\ker B$.
Definisikan $C_0:\operatorname{im}A\to Z$ melalui
$C_0(Av)=Bv$.
Jika $Av=Av'$, maka $v-v'\in\ker A\subseteq\ker B$,
sehingga $Bv=Bv'$; jadi definisinya sah.
Linearitas mengikuti dari $A$ dan $B$.
Perluas basis citra $A$ menjadi basis $W$, lalu tentukan
nilai $C$ pada vektor tambahan secara sembarang, misalnya nol.
Teorema penentuan pada basis memberi perluasan linear dengan $B=CA$.
Jika $A$ surjektif, semua anggota $W$ berbentuk $Av$ dan
nilainya sudah dipaksa oleh $C(Av)=Bv$, sehingga $C$ unik.

\subsection*{Latihan \ref{trans:lat-komplemen-invarian}: komplemen tidak selalu terpelihara}
Karena $N(e_1)=0$, subruang $U$ invarian.
Setiap komplemennya berupa garis yang dibangun oleh
$w=(a,b)$ dengan $b\ne0$.
Jika garis ini invarian, harus ada $\lambda\in\R$ dengan
$Nw=(b,0)=\lambda(a,b)$.
Koordinat kedua memaksa $\lambda=0$, sedangkan koordinat
pertama kemudian memaksa $b=0$, kontradiksi.

\subsection*{Latihan \ref{trans:lat-komutator-identitas}: perbedaan dimensi hingga dan tak hingga}
Aturan turunan hasil kali memberi
$(DM-MD)p=(tp)'-tp'=p$, sehingga $DM-MD=I$.
Untuk $k\ge1$, gunakan identitas teleskopik
\[
    D^kM-MD^k
      =\sum_{j=0}^{k-1}D^j(DM-MD)D^{k-1-j}
      =kD^{k-1}.
\]
Pada ruang berdimensi $n\ge1$, representasi matriks akan
memberikan $AB-BA=I_n$ jika operator demikian ada.
Mengambil trace memberi $0=\tr(AB)-\tr(BA)=n$, kontradiksi
di atas $\R$ atau $\C$.
Tidak ada pertentangan dengan contoh polinom: $\F[t]$
berdimensi tak hingga, dan perkalian oleh $t$ tidak mempertahankan
subruang $\mathcal P_{\le n}(\F)$.
\fi


\part{Hasil Kali Dalam}{% !TeX root = ../../ONMIPA Alin.tex
Struktur ruang vektor memungkinkan kita membentuk kombinasi linear,
memilih basis, dan menghitung dimensi. Hasil kali dalam menambahkan
gagasan geometris: panjang, jarak, dan ortogonalitas. Dengan struktur
ini, kita dapat menentukan vektor terdekat dalam suatu subruang
serta memilih basis yang menyederhanakan perhitungan.

Pembahasan dimulai dari aksioma dan ketaksamaan dasar, dilanjutkan
dengan basis ortonormal, proyeksi, serta pendekatan kuadrat terkecil.
Pada bagian akhir, hasil kali dalam digunakan untuk mempelajari
adjoin, isometri, dan operator adjoin diri. Pembaca diharapkan telah
memahami basis, dimensi, kernel, citra, dan representasi matriks
dari transformasi linear.
}
% !TeX root = ../../ONMIPA Alin.tex
% Acuan: Zhang (2009), Bab 5, definisi serta Soal 5.1--5.21;
% Zhang (2011), bagian 1.4. Konvensi di sini linear pada argumen kedua,
% sesuai x^*y pada bab matriks; semua rumus konjugasi disesuaikan.
\chapter{Ruang Hasil Kali Dalam}
\label{ch:hasil-kali-dalam}

Dimensi dan kebebasan linear tidak menentukan panjang suatu vektor
atau sudut antarvektor. Kedua gagasan itu memerlukan struktur
tambahan berupa hasil kali dalam. Struktur yang sama juga berlaku
pada ruang polinom dan ruang matriks, sehingga argumen geometris
dapat digunakan untuk menyelesaikan persoalan yang tampaknya aljabar.
Sepanjang bagian ini, medan skalar adalah $\F=\R$ atau $\C$.

\section{Definisi dan Konvensi Konjugasi}

\begin{definisi}[][Hasil kali dalam][hkd:def]
    Hasil kali dalam pada ruang linear $V$ adalah pemetaan
    $\langle\cdot,\cdot\rangle:V\times V\to\F$ yang memenuhi
    \begin{enumerate}[label=(\roman*),nosep]
        \item $\langle x,x\rangle$ real dan nonnegatif, serta
        $\langle x,x\rangle=0$ tepat ketika $x=0$;
        \item $\langle x,y\rangle=\overline{\langle y,x\rangle}$;
        \item $\langle x,ay+bz\rangle
        =a\langle x,y\rangle+b\langle x,z\rangle$.
    \end{enumerate}
    Syarat-syarat tersebut berlaku untuk semua vektor dan skalar
    yang sesuai. Ruang $V$ bersama hasil kali dalamnya disebut
    \emph{ruang hasil kali dalam}.
\end{definisi}

Konvensi handout ini adalah \textbf{linear pada argumen kedua}.
Simetri konjugat kemudian memberi
\[
    \langle ax+by,z\rangle
      =\overline a\langle x,z\rangle+\overline b\langle y,z\rangle.
\]
Jadi pada ruang kompleks, hasil kali dalam bersifat konjugat-linear
pada argumen pertama. Secara umum,
\[
    \left\langle\sum_i a_ix_i,\sum_j b_jy_j\right\rangle
      =\sum_{i,j}\overline{a_i}b_j\langle x_i,y_j\rangle.
\]
Pada ruang real, konjugasi tidak mengubah skalar sehingga bentuknya
bilinear dan simetris. Selain itu, $\langle0,x\rangle
=\langle x,0\rangle=0$.

Hasil kali dalam standar pada $\F^n$ adalah
$\langle x,y\rangle=x^*y=\sum_j\overline{x_j}y_j$.
Untuk vektor real, rumus ini menjadi $x^Ty$.

\begin{exam}[Konjugasi tidak boleh dihilangkan]
    Apakah $B(x,y)=x_1y_1+x_2y_2$ merupakan hasil kali dalam
    pada $\C^2$?
\end{exam}
\begin{solution}
    Tidak. Vektor tak nol $x=(1,i)^T$ memenuhi
    $B(x,x)=1+i^2=0$, sehingga syarat positif definit gagal.
    Bahkan $B((i,0)^T,(i,0)^T)=-1$.
    Pada hasil kali dalam standar, kedua norma kuadrat itu
    masing-masing bernilai dua dan satu.
\end{solution}

Jika $U\subseteq V$ subruang, pembatasan hasil kali dalam pada
$U\times U$ kembali menjadi hasil kali dalam. Namun, suatu
bentuk yang hanya nonnegatif dapat mempunyai vektor tak nol
dengan nilai kuadrat nol; bentuk seperti itu belum menjadi hasil kali dalam.

\section{Hasil Kali Dalam pada Polinom dan Matriks}

\begin{exam}[Hasil kali dalam integral]
    Pada $\mathcal P_{\le n}(\F)$, buktikan bahwa
    \[
        \langle p,q\rangle=\int_{-1}^{1}\overline{p(t)}q(t)\,dt
    \]
    merupakan hasil kali dalam.
\end{exam}
\begin{solution}
    Linearitas dan simetri konjugat mengikuti sifat integral.
    Nilai $\langle p,p\rangle=\int_{-1}^1|p(t)|^2\,dt$ nonnegatif.
    Jika integralnya nol, fungsi kontinu nonnegatif $|p|^2$
    harus nol di seluruh interval: nilai positif di satu titik
    akan memberikan integral positif pada interval kecil di sekitarnya.
    Maka polinom $p$ mempunyai tak hingga banyak akar dan identik nol.
    Arah sebaliknya langsung.
\end{solution}

Pada ruang matriks berukuran $m\times n$, hasil kali dalam Frobenius adalah
\[
    \langle A,B\rangle_F=\tr(A^*B)
      =\sum_{i=1}^m\sum_{j=1}^n\overline{a_{ij}}b_{ij}.
\]
Kesamaan terakhir menunjukkan bahwa ini tepat hasil kali dalam
standar pada daftar seluruh entri. Norma yang dihasilkannya adalah
$\|A\|_F=(\sum_{i,j}|a_{ij}|^2)^{1/2}$.
Matriks satuan $E_{ij}$ saling ortogonal dan masing-masing bernorma satu.

\begin{teorema}[][Representasi oleh matriks positif definit][hkd:matriks-metrik]
    Pada $\F^n$, rumus $\langle x,y\rangle_G=x^*Gy$ merupakan
    hasil kali dalam jika dan hanya jika $G$ Hermitian positif definit.
    Setiap hasil kali dalam pada $\F^n$ mempunyai representasi
    tersebut secara tunggal.
\end{teorema}
\begin{proof}
    Jika $G$ Hermitian positif definit, ketiga aksioma segera
    mengikuti perkalian matriks dan $x^*Gx>0$ untuk $x\ne0$.
    Sebaliknya, letakkan $g_{ij}=\langle e_i,e_j\rangle$.
    Ekspansi kedua argumen dalam basis standar memberi
    $\langle x,y\rangle=x^*Gy$.
    Simetri konjugat memberi $G^*=G$, dan aksioma kepositifan
    memberi sifat positif definit. Setiap entri dipaksa oleh
    hasil kali dua vektor basis, sehingga $G$ unik.
\end{proof}

Untuk basis terurut sembarang $\mathcal B=(b_1,\ldots,b_n)$,
matriks $G_{\mathcal B}=(\langle b_i,b_j\rangle)$ memenuhi
\[
    \langle x,y\rangle=[x]_{\mathcal B}^*G_{\mathcal B}[y]_{\mathcal B}.
\]
Jika $P=P_{\mathcal B\leftarrow\mathcal C}$, maka
$G_{\mathcal C}=P^*G_{\mathcal B}P$.
Jadi perubahan basis pada hasil kali dalam menghasilkan
\emph{kongruensi}, bukan rumus keserupaan $P^{-1}GP$.

\begin{exam}[Parameter dalam hasil kali dalam berbobot]
    Untuk $z\in\C$ dan $a\in\R$, kapan matriks
    \[
        G=\begin{pmatrix}1&z\\\overline z&a\end{pmatrix}
    \]
    mendefinisikan hasil kali dalam pada $\C^2$?
\end{exam}
\begin{solution}
    Untuk $x=(x_1,x_2)^T$, lengkapi kuadrat:
    \[
        x^*Gx=|x_1+zx_2|^2+(a-|z|^2)|x_2|^2.
    \]
    Nilainya positif untuk setiap $x\ne0$ tepat ketika $a>|z|^2$.
    Jika $a=|z|^2$, vektor $(-z,1)^T$ menghasilkan nilai nol.
    Jika $a<|z|^2$, vektor yang sama menghasilkan nilai negatif.
\end{solution}

\section{Norma dan Ketaksamaan Cauchy--Schwarz}

\begin{definisi}[][Norma dan jarak][hkd:norma]
    Norma yang diinduksi oleh hasil kali dalam didefinisikan oleh
    $\|x\|=\sqrt{\langle x,x\rangle}$.
    Jarak antara $x$ dan $y$ adalah $d(x,y)=\|x-y\|$.
    Vektor dengan norma satu disebut vektor satuan.
\end{definisi}

\begin{teorema}[][Ketaksamaan Cauchy--Schwarz][hkd:cauchy-schwarz]
    Untuk setiap $x,y$,
    \[
        |\langle x,y\rangle|^2\le\|x\|^2\|y\|^2.
    \]
    Kesamaan berlaku jika dan hanya jika $x,y$ bergantung linear.
\end{teorema}
\begin{proof}
    Kasus $y=0$ langsung. Untuk $y\ne0$, letakkan
    $c=\langle y,x\rangle/\|y\|^2$ dan $r=x-cy$.
    Maka $\langle y,r\rangle=0$ dan
    \[
        0\le\|r\|^2
          =\|x\|^2-\frac{|\langle y,x\rangle|^2}{\|y\|^2}.
    \]
    Ini memberi ketaksamaan yang diminta.
    Kesamaan ekuivalen dengan $r=0$, yakni $x=cy$.
    Dengan memasukkan kembali kasus vektor nol, kondisi itu
    tepat ketergantungan linear kedua vektor.
\end{proof}

\begin{prop}[Ketaksamaan segitiga dan kesamaannya]
    Norma memenuhi $\|ax\|=|a|\|x\|$ serta
    \[
        \|x+y\|\le\|x\|+\|y\|,\qquad
        \bigl|\|x\|-\|y\|\bigr|\le\|x-y\|.
    \]
    Jika $x,y$ tak nol, kesamaan pada ketaksamaan segitiga
    pertama berlaku tepat ketika $x=cy$ untuk suatu real $c>0$.
    Jika salah satu vektor nol, kesamaan juga berlaku.
\end{prop}
\begin{proof}
    Homogenitas norma mengikuti $\langle ax,ax\rangle
    =|a|^2\langle x,x\rangle$.
    Selanjutnya,
    \begin{align*}
        \|x+y\|^2
        &=\|x\|^2+\|y\|^2+2\operatorname{Re}\langle x,y\rangle\\
        &\le\|x\|^2+\|y\|^2+2|\langle x,y\rangle|\\
        &\le(\|x\|+\|y\|)^2.
    \end{align*}
    Untuk vektor tak nol, kedua ketaksamaan menjadi kesamaan
    jika keduanya bergantung linear dan hasil kali dalamnya
    real positif. Ini memaksa kelipatan real positif.
    Ketaksamaan kedua diperoleh dari
    $\|x\|\le\|x-y\|+\|y\|$, lalu menukar $x,y$.
\end{proof}

Ketergantungan linear saja tidak cukup untuk kesamaan segitiga.
Di $\C$, vektor $1$ dan $i$ bergantung linear, tetapi
$|1+i|=\sqrt2<2=|1|+|i|$.

\begin{exam}[Ketaksamaan berbobot dan kondisi kesamaan]
    Untuk $w_j>0$ dan $a_j\in\C$, buktikan
    \[
        \left|\sum_{j=1}^n a_j\right|^2
        \le\left(\sum_{j=1}^n w_j|a_j|^2\right)
            \left(\sum_{j=1}^n\frac1{w_j}\right).
    \]
\end{exam}
\begin{solution}
    Gunakan Cauchy--Schwarz pada
    $x=(1/\sqrt{w_j})_j$ dan $y=(\sqrt{w_j}a_j)_j$.
    Kesamaan berlaku tepat ketika $y=cx$, yaitu
    $a_j=c/w_j$ untuk suatu konstanta kompleks $c$ yang sama.
    Kondisi kesamaan ini penting ketika ketaksamaan digunakan
    untuk menentukan nilai minimum atau maksimum.
\end{solution}

\section{Identitas Polarisasi dan Ortogonalitas}

\begin{teorema}[][Jajargenjang dan polarisasi][hkd:polarisasi]
    Norma dari hasil kali dalam memenuhi
    \[
        \|x+y\|^2+\|x-y\|^2=2\|x\|^2+2\|y\|^2.
    \]
    Pada ruang real,
    $4\langle x,y\rangle=\|x+y\|^2-\|x-y\|^2$.
    Pada ruang kompleks dengan konvensi handout ini,
    \begin{align*}
        4\langle x,y\rangle
        ={}&\|x+y\|^2-\|x-y\|^2\\
            &-i\|x+iy\|^2+i\|x-iy\|^2.
    \end{align*}
\end{teorema}
\begin{proof}
    Ekspansi $\|x\pm y\|^2$ memberikan suku silang
    $\pm2\operatorname{Re}\langle x,y\rangle$.
    Menjumlahkan memberi identitas jajargenjang, sedangkan
    mengurangkan memberi empat kali bagian real.
    Karena $\langle x,iy\rangle=i\langle x,y\rangle$,
    selisih $\|x+iy\|^2-\|x-iy\|^2$ adalah
    $-4\operatorname{Im}\langle x,y\rangle$.
    Gabungkan bagian real dan imajiner untuk memperoleh rumus kompleks.
\end{proof}

Norma menentukan hasil kali dalamnya secara tunggal melalui
polarisasi. Namun, tidak setiap norma berasal dari hasil kali dalam.
Sebagai contoh, norma $\|(x_1,x_2)\|_1=|x_1|+|x_2|$ memberi
\[
    \|e_1+e_2\|_1^2+\|e_1-e_2\|_1^2=8,
    \qquad 2\|e_1\|_1^2+2\|e_2\|_1^2=4,
\]
sehingga gagal memenuhi identitas jajargenjang.

Dua vektor disebut \emph{ortogonal}, ditulis $x\perp y$, jika
$\langle x,y\rangle=0$. Ortogonalitas mengakibatkan identitas
Pythagoras $\|x+y\|^2=\|x\|^2+\|y\|^2$.
Pada ruang real, arah sebaliknya juga benar.

\begin{exam}[Pythagoras belum cukup di ruang kompleks]
    Di $\C$ dengan hasil kali dalam standar, ambil $x=1$ dan $y=i$.
    Maka $\|x+y\|^2=2=\|x\|^2+\|y\|^2$, tetapi
    $\langle x,y\rangle=i\ne0$.
    Jadi identitas Pythagoras hanya memaksa bagian real hasil
    kali dalam bernilai nol. Pada ruang kompleks, syarat yang tepat adalah
    \[
        x\perp y\iff
        \begin{cases}
            \|x+y\|=\|x-y\|,\\
            \|x+iy\|=\|x-iy\|.
        \end{cases}
    \]
    Kedua persamaan menghilangkan bagian real dan imajiner secara terpisah.
\end{exam}

\section{Matriks Gram dan Ketergantungan Linear}

\begin{definisi}[][Matriks Gram][hkd:gram]
    Untuk vektor $v_1,\ldots,v_m$, matriks Gram adalah
    \[
        G=(g_{ij})_{i,j=1}^m,\qquad g_{ij}=\langle v_i,v_j\rangle.
    \]
\end{definisi}

\begin{teorema}[][Kepositifan dan rank matriks Gram][hkd:rank-gram]
    Matriks Gram Hermitian semidefinit positif dan memenuhi
    \[
        \rank G=\dim\operatorname{span}\{v_1,\ldots,v_m\}.
    \]
    Matriks $G$ positif definit tepat ketika daftar vektor itu
    bebas linear; ekuivalen dengan $\det G>0$.
\end{teorema}
\begin{proof}
    Untuk $c=(c_1,\ldots,c_m)^T$, berlaku
    \[
        c^*Gc=\left\|\sum_{j=1}^m c_jv_j\right\|^2\ge0.
    \]
    Simetri konjugat memberi $G^*=G$.
    Definisikan $S:\F^m\to V$ melalui $Sc=\sum c_jv_j$.
    Jika $Gc=0$, persamaan di atas memberi $Sc=0$.
    Sebaliknya, $Sc=0$ memberi $(Gc)_i=\langle v_i,Sc\rangle=0$.
    Maka $\ker G=\ker S$; teorema rank--nulitas memberi
    $\rank G=\rank S=\dim\operatorname{im}S$.
    Kepositifan definit ekuivalen dengan $Sc\ne0$ untuk setiap
    $c\ne0$, yakni kebebasan linear. Karena semua nilai eigen
    $G$ nonnegatif, keadaan ini juga ekuivalen dengan $\det G>0$.
\end{proof}

Dalam koordinat ortonormal, jika kolom matriks $A$ adalah para
$v_j$, maka $G=A^*A$. Sebaliknya, setiap matriks Hermitian
semidefinit positif merupakan matriks Gram: ambil akar kuadrat
Hermitian $B=G^{1/2}$ sehingga $G=B^*B$, lalu gunakan kolom-kolom $B$.
Hasil mengenai akar kuadrat telah dibahas pada bab matriks khusus.

\begin{exam}[Tiga hasil kali dalam tidak dapat dipilih sembarang]
    Misalkan $u,v,w$ vektor satuan dalam ruang real dan
    $a=\langle u,v\rangle$, $b=\langle u,w\rangle$,
    $c=\langle v,w\rangle$. Buktikan
    \[
        (c-ab)^2\le(1-a^2)(1-b^2).
    \]
\end{exam}
\begin{solution}
    Matriks Gram ketiganya semidefinit positif, sehingga
    \[
        0\le\det\begin{pmatrix}1&a&b\\a&1&c\\b&c&1\end{pmatrix}
          =1+2abc-a^2-b^2-c^2.
    \]
    Menata ulang memberi ketaksamaan yang diminta.
    Kesamaan berarti determinan Gram nol, sehingga ketiga
    vektor bergantung linear. Jika $|a|,|b|<1$, batas untuk $c$ adalah
    $ab\pm\sqrt{(1-a^2)(1-b^2)}$.
\end{solution}

\section{Latihan Pilihan}

\begin{latihan}\label{hkd:lat-tarik-balik}
    Misalkan $A:V\to W$ linear dan $W$ ruang hasil kali dalam.
    Buktikan bahwa $[x,y]=\langle Ax,Ay\rangle$ merupakan hasil
    kali dalam pada $V$ tepat ketika $A$ injektif.
    Nyatakan syaratnya jika $A$ adalah matriks $m\times n$.
\end{latihan}

\begin{latihan}\label{hkd:lat-evaluasi}
    Pada $\mathcal P_{\le n}(\F)$, ambil titik real berbeda
    $a_1,\ldots,a_m$ dan bobot $w_j>0$. Tentukan kapan
    \[
        [p,q]=\sum_{j=1}^m w_j\overline{p(a_j)}q(a_j)
    \]
    merupakan hasil kali dalam.
\end{latihan}

\begin{latihan}\label{hkd:lat-turunan}
    Pada $\mathcal P_{\le n}(\F)$, buktikan bahwa
    \[
        [p,q]=\overline{p(0)}q(0)+\int_0^1\overline{p'(t)}q'(t)\,dt
    \]
    merupakan hasil kali dalam. Jelaskan mengapa menghapus
    suku pertama membuat sifat positif definit gagal.
\end{latihan}

\begin{latihan}\label{hkd:lat-kesamaan}
    Di $\C^2$, ambil $x=(1,i)^T$ dan $y=(a,b)^T$.
    Tentukan semua $a,b$ yang memberikan kesamaan pada
    Cauchy--Schwarz, lalu tentukan yang juga memenuhi
    $\|x+y\|=\|x\|+\|y\|$.
\end{latihan}

\begin{latihan}\label{hkd:lat-trace}
    Untuk $A,B\in\Mat_n(\C)$, buktikan
    \[
        |\tr A|^2\le n\tr(A^*A),\qquad
        |\tr(AB)|\le\|A\|_F\|B\|_F.
    \]
    Tentukan kondisi kesamaan pada ketaksamaan pertama.
\end{latihan}

\begin{latihan}\label{hkd:lat-invers-positif}
    Jika $G$ Hermitian positif definit dan $x\ne0$, buktikan
    \[
        (x^*Gx)(x^*G^{-1}x)\ge(x^*x)^2.
    \]
    Buktikan bahwa kesamaan berlaku tepat ketika $x$ vektor eigen $G$.
\end{latihan}

\begin{latihan}\label{hkd:lat-korelasi}
    Untuk $m\ge2$ dan $a\in\R$, tentukan kapan terdapat
    vektor satuan real $v_1,\ldots,v_m$ dengan
    $\langle v_i,v_j\rangle=a$ untuk setiap $i\ne j$.
    Tentukan dimensi terkecil ruang yang dapat memuat vektor-vektor itu.
\end{latihan}

\begin{latihan}\label{hkd:lat-negatif}
    Dalam ruang hasil kali dalam real berdimensi $n\ge1$,
    misalkan $v_1,\ldots,v_m$ memenuhi
    $\langle v_i,v_j\rangle<0$ untuk setiap $i\ne j$.
    Buktikan $m\le n+1$ dan tunjukkan bahwa batas tersebut dapat dicapai.
\end{latihan}

\ifpembahasan
\section{Pembahasan Latihan}

\subsection*{Latihan \ref{hkd:lat-tarik-balik}: kepositifan ditentukan oleh kernel}
Linearitas pada argumen kedua dan simetri konjugat selalu berlaku.
Selanjutnya $[x,x]=\|Ax\|^2=0$ tepat ketika $Ax=0$.
Jadi sifat positif definit ekuivalen dengan $\ker A=\{0\}$.
Untuk matriks $m\times n$, ini berarti $\rank A=n$;
khususnya, diperlukan $m\ge n$.

\subsection*{Latihan \ref{hkd:lat-evaluasi}: banyak titik dibandingkan derajat}
Karena semua bobot positif, $[p,p]=0$ tepat ketika
$p(a_1)=\cdots=p(a_m)=0$.
Jika $m\ge n+1$, polinom berderajat paling besar $n$ dengan
akar sebanyak itu harus nol, sehingga bentuknya hasil kali dalam.
Jika $m\le n$, polinom tak nol $\prod_{j=1}^m(t-a_j)$
berada di ruang yang diberikan dan bernilai nol pada semua titik.
Jadi syaratnya tepat $m\ge n+1$.

\subsection*{Latihan \ref{hkd:lat-turunan}: memulihkan konstanta yang hilang}
Jika $[p,p]=0$, kedua suku nonnegatifnya nol.
Integral $|p'|^2$ yang nol memaksa $p'=0$, sehingga $p$ konstan.
Syarat $p(0)=0$ kemudian memaksa konstanta itu nol.
Aksioma lainnya mengikuti sifat konjugasi dan integral.
Tanpa suku pertama, setiap polinom konstan tak nol mempunyai
nilai kuadrat nol, sehingga bentuknya bukan hasil kali dalam.

\subsection*{Latihan \ref{hkd:lat-kesamaan}: bergantung linear dan searah}
Karena $x\ne0$, kesamaan Cauchy--Schwarz berarti $y=cx$,
yakni $a=c$ dan $b=ic$ untuk suatu $c\in\C$.
Kesamaan segitiga mengharuskan $c$ real nonnegatif, termasuk
$c=0$ yang menghasilkan $y=0$.
Kelipatan kompleks dengan fase berbeda umumnya tidak memenuhi
kesamaan segitiga.

\subsection*{Latihan \ref{hkd:lat-trace}: pandang matriks sebagai vektor}
Gunakan $\langle I,A\rangle_F=\tr A$ dan $\|I\|_F^2=n$.
Cauchy--Schwarz memberi ketaksamaan pertama, dengan kesamaan
tepat ketika $A=cI$ untuk suatu $c\in\C$.
Untuk yang kedua, gunakan
$\langle A^*,B\rangle_F=\tr(AB)$ serta $\|A^*\|_F=\|A\|_F$.

\subsection*{Latihan \ref{hkd:lat-invers-positif}: gunakan metrik yang sesuai}
Gunakan Cauchy--Schwarz pada $x,G^{-1}x$ untuk hasil kali dalam
$\langle u,v\rangle_G=u^*Gv$. Ketiga besaran yang muncul adalah
\[
    \langle x,G^{-1}x\rangle_G=x^*x,\quad
    \|x\|_G^2=x^*Gx,\quad
    \|G^{-1}x\|_G^2=x^*G^{-1}x.
\]
Kesamaan berarti $G^{-1}x=cx$, dengan $c>0$ karena $G^{-1}$
positif definit. Ini ekuivalen dengan $x$ vektor eigen $G$.

\subsection*{Latihan \ref{hkd:lat-korelasi}: spektrum matriks Gram}
Matriks Gram yang diminta adalah
$G=(1-a)I_m+a\boldsymbol{1}\boldsymbol{1}^T$.
Pada arah $\boldsymbol{1}$, nilai eigennya $1+(m-1)a$;
pada komplemen ortogonalnya, nilai eigennya $1-a$.
Jadi $G$ semidefinit positif tepat ketika
\[
    -\frac1{m-1}\le a\le1.
\]
Kondisi ini cukup karena setiap matriks real simetris semidefinit
positif dapat diwujudkan sebagai Gram vektor real.
Dimensi terkecil sama dengan rank Gram: bernilai $m$ di bagian
dalam interval, $m-1$ pada ujung kiri, dan $1$ pada ujung kanan.
Kecukupan dimensi rank diperoleh dengan mengambil hanya baris
tak nol pada faktorisasi spektral $G=B^TB$.

\subsection*{Latihan \ref{hkd:lat-negatif}: pisahkan koefisien positif dan negatif}
Andaikan $m>n+1$. Vektor $(v_j,1)$ pada $V\times\R$
bergantung linear, sehingga ada real $c_j$ tidak semuanya nol dengan
\[
    \sum_jc_jv_j=0,\qquad \sum_jc_j=0.
\]
Ada koefisien positif dan negatif. Letakkan
$u=\sum_{c_j>0}c_jv_j=\sum_{c_k<0}(-c_k)v_k$.
Menggunakan dua ekspresi tersebut pada kedua argumen memberi
\[
    \|u\|^2=\sum_{c_j>0,\ c_k<0}c_j(-c_k)\langle v_j,v_k\rangle<0,
\]
kontradiksi. Jadi $m\le n+1$.
Untuk mencapai batas, di hiperbidang $\sum x_j=0$ pada $\R^{n+1}$
yang berdimensi $n$, ambil
$v_j=e_j-\boldsymbol{1}/(n+1)$.
Untuk $j\ne k$, hasil kali dalamnya $-1/(n+1)<0$.
\fi

% !TeX root = ../../ONMIPA Alin.tex
% Acuan: Zhang (2009), Bab 5, Soal 5.6--5.8, 5.19--5.23,
% 5.26--5.35, dan 5.50; Zhang (2011), bagian 1.4, khususnya
% Gram--Schmidt, komplemen ortogonal, dan ketaksamaan Bessel.
\chapter{Ortogonalitas dan Proyeksi}
\label{ch:ortogonalitas}

Basis ortonormal mengubah pencarian koordinat menjadi perhitungan
hasil kali dalam. Dari basis ini, kita memperoleh dekomposisi
ortogonal dan proyeksi, yang menjawab pertanyaan: vektor mana
dalam suatu subruang yang paling dekat dengan vektor yang diberikan?
Kecuali dinyatakan lain, ruang dalam bab ini berdimensi hingga.
Hasil kali dalam tetap linear pada argumen kedua.

\section{Sistem Ortogonal dan Ortonormal}

\begin{definisi}[][Daftar ortogonal dan ortonormal][ortho:def]
    Daftar $v_1,\ldots,v_m$ disebut ortogonal jika
    $\langle v_i,v_j\rangle=0$ untuk $i\ne j$.
    Daftar itu disebut ortonormal jika, sebagai tambahan,
    setiap vektornya bernorma satu; ekuivalen dengan
    $\langle v_i,v_j\rangle=\delta_{ij}$.
\end{definisi}

Setiap daftar ortogonal yang semua anggotanya tak nol bebas linear.
Memang, dari $\sum c_jv_j=0$, ambil hasil kali dalam dengan
$v_i$ pada argumen pertama untuk memperoleh $c_i\|v_i\|^2=0$.
Selain itu, ekspansi langsung memberikan
\[
    \left\|\sum_jc_jv_j\right\|^2=\sum_j|c_j|^2\|v_j\|^2.
\]
Syarat tak nol tidak boleh diabaikan dalam kesimpulan kebebasan linear.

\begin{teorema}[][Koordinat ortonormal dan identitas Parseval][ortho:parseval]
    Jika $e_1,\ldots,e_n$ basis ortonormal bagi $V$, maka
    \[
        x=\sum_{j=1}^n\langle e_j,x\rangle e_j,
        \qquad \|x\|^2=\sum_{j=1}^n|\langle e_j,x\rangle|^2.
    \]
    Lebih umum,
    \[
        \langle x,y\rangle
        =\sum_{j=1}^n\overline{\langle e_j,x\rangle}\langle e_j,y\rangle.
    \]
\end{teorema}
\begin{proof}
    Tulis $x=\sum c_je_j$. Mengambil $\langle e_i,x\rangle$
    memberi $c_i$. Ekspansi hasil kali dalam dua kombinasi
    tersebut hanya menyisakan suku dengan indeks yang sama,
    sehingga kedua identitas lainnya mengikuti.
\end{proof}

Jika $(v_1,\ldots,v_n)$ basis ortogonal, koefisien $v_j$ dalam
ekspansi $x$ adalah $\langle v_j,x\rangle/\|v_j\|^2$.
Penyebut boleh dihilangkan hanya ketika basisnya ortonormal.

\begin{teorema}[][Ketaksamaan Bessel][ortho:bessel]
    Jika $e_1,\ldots,e_m$ ortonormal, maka untuk setiap $x$,
    \[
        \sum_{j=1}^m|\langle e_j,x\rangle|^2\le\|x\|^2.
    \]
    Kesamaan berlaku tepat ketika
    $x\in\operatorname{span}\{e_1,\ldots,e_m\}$.
\end{teorema}
\begin{proof}
    Letakkan $p=\sum_j\langle e_j,x\rangle e_j$ dan $r=x-p$.
    Untuk setiap $i$, $\langle e_i,r\rangle=0$, sehingga $p\perp r$.
    Maka
    \[
        \|x\|^2=\|p\|^2+\|r\|^2
          =\sum_j|\langle e_j,x\rangle|^2+\|r\|^2.
    \]
    Kesamaan berlaku tepat ketika $r=0$.
\end{proof}

\section{Proses Gram--Schmidt}

\begin{teorema}[][Ortonormalisasi Gram--Schmidt][ortho:gram-schmidt]
    Untuk daftar bebas linear $v_1,\ldots,v_m$, definisikan
    \[
        u_1=v_1,\quad e_1=\frac{u_1}{\|u_1\|},\qquad
        u_k=v_k-\sum_{j=1}^{k-1}\langle e_j,v_k\rangle e_j,
        \quad e_k=\frac{u_k}{\|u_k\|}.
    \]
    Setiap $u_k$ tak nol, daftar $(e_1,\ldots,e_m)$ ortonormal,
    dan untuk setiap $k$ berlaku
    \[
        \operatorname{span}\{e_1,\ldots,e_k\}
        =\operatorname{span}\{v_1,\ldots,v_k\}.
    \]
\end{teorema}
\begin{proof}
    Gunakan induksi. Jika $e_1,\ldots,e_{k-1}$ sudah ortonormal,
    maka $\langle e_i,u_k\rangle=0$ untuk $i<k$.
    Jika $u_k=0$, vektor $v_k$ berada dalam ruang pembangun
    vektor-vektor sebelumnya, bertentangan dengan kebebasan linear.
    Normalisasi memberi vektor satuan baru yang ortogonal
    terhadap semua vektor sebelumnya.
    Rumus $u_k$ dan bentuk baliknya menyatakan bahwa kedua
    daftar membangun ruang yang sama pada setiap tahap.
\end{proof}

Akibatnya, setiap ruang hasil kali dalam berdimensi hingga mempunyai
basis ortonormal. Daftar ortonormal yang belum lengkap dapat
diperluas menjadi basis ortonormal: perluas dahulu menjadi basis
sembarang, lalu terapkan Gram--Schmidt. Vektor ortonormal yang
sudah ada pada awal daftar tidak berubah.

Jika daftar awal bergantung linear, suatu sisa $u_k$ dapat nol.
Sisa tersebut harus dilewati, bukan dibagi dengan normanya.
Proses pada vektor berikutnya tetap dapat dilanjutkan untuk
memperoleh basis ortonormal ruang pembangun daftar awal.

\begin{exam}[Ortonormalisasi tiga vektor real]
    Terapkan Gram--Schmidt pada
    $v_1=(1,1,0)^T$, $v_2=(1,0,1)^T$, $v_3=(0,1,1)^T$.
\end{exam}
\begin{solution}
    Langkah pertama memberi $e_1=(1,1,0)^T/\sqrt2$.
    Selanjutnya,
    \begin{align*}
        u_2&=v_2-\tfrac12(1,1,0)^T=(1/2,-1/2,1)^T,\\
        e_2&=(1,-1,2)^T/\sqrt6,\\
        u_3&=v_3-\tfrac12(1,1,0)^T-\tfrac16(1,-1,2)^T
             =\tfrac23(-1,1,1)^T.
    \end{align*}
    Jadi $e_3=(-1,1,1)^T/\sqrt3$.
    Ketiganya bernorma satu dan hasil kali setiap dua vektor
    berbeda nol, sehingga menjadi basis ortonormal $\R^3$.
\end{solution}

\begin{exam}[Basis ortonormal polinom]
    Pada $\mathcal P_{\le2}(\R)$ dengan
    $\langle p,q\rangle=\int_{-1}^1p(t)q(t)\,dt$,
    ortonormalkan daftar $(1,t,t^2)$.
\end{exam}
\begin{solution}
    Simetri interval membuat fungsi ganjil berintegral nol.
    Maka $1\perp t$, dengan norma kuadrat $2$ dan $2/3$.
    Sisa untuk $t^2$ adalah $t^2-1/3$, yang bernorma kuadrat $8/45$.
    Hasil akhirnya adalah
    \[
        e_0=\frac1{\sqrt2},\qquad
        e_1=\sqrt{\frac32}\,t,\qquad
        e_2=\sqrt{\frac58}\,(3t^2-1).
    \]
    Ortogonalitas di sini ditentukan oleh integral, bukan oleh
    hasil kali daftar koefisien dalam basis monomial.
\end{solution}

\section{Komplemen Ortogonal dan Dekomposisi}

\begin{definisi}[][Komplemen ortogonal][ortho:komplemen]
    Untuk $S\subseteq V$, definisikan
    \[
        S^\perp=\{x\in V:\langle s,x\rangle=0
            \text{ untuk setiap }s\in S\}.
    \]
\end{definisi}

Himpunan $S^\perp$ selalu subruang, bahkan jika $S$ bukan subruang.
Selain itu, $S^\perp=(\operatorname{span}S)^\perp$.
Jika $S\subseteq T$, maka $T^\perp\subseteq S^\perp$;
semakin banyak syarat ortogonal, semakin kecil ruang yang tersisa.

\begin{teorema}[][Dekomposisi ortogonal][ortho:dekomposisi]
    Untuk subruang $W$ dari ruang hasil kali dalam $V$ berdimensi $n$,
    \[
        V=W\oplus W^\perp,\qquad
        \dim W^\perp=n-\dim W,\qquad (W^\perp)^\perp=W.
    \]
    Untuk himpunan sembarang $S\subseteq V$,
    $(S^\perp)^\perp=\operatorname{span}S$.
\end{teorema}
\begin{proof}
    Pilih basis ortonormal $e_1,\ldots,e_r$ dari $W$ dan perluas
    menjadi basis ortonormal $e_1,\ldots,e_n$ dari $V$.
    Ekspansi koordinat menunjukkan bahwa $W^\perp$ dibangun
    oleh $e_{r+1},\ldots,e_n$. Ini membuktikan dekomposisi dan
    rumus dimensi. Vektor yang ortogonal terhadap semua anggota
    $W^\perp$ tepat kombinasi $e_1,\ldots,e_r$, sehingga
    komplemen dua kali mengembalikan $W$.
    Pernyataan terakhir mengikuti $S^\perp=(\operatorname{span}S)^\perp$.
\end{proof}

Untuk subruang $U,W$, berlaku
\[
    (U+W)^\perp=U^\perp\cap W^\perp,\qquad
    (U\cap W)^\perp=U^\perp+W^\perp.
\]
Identitas pertama mengikuti langsung dari definisi.
Untuk yang kedua, terapkan identitas pertama pada $U^\perp,W^\perp$
dan ambil komplemen ortogonal sekali lagi.
Penggunaan komplemen ganda di sini memakai asumsi dimensi hingga.

\begin{exam}[Ruang matriks bertrace nol]
    Dalam $\Mat_n(\C)$ dengan hasil kali dalam Frobenius,
    tentukan komplemen ortogonal dari
    $W=\{X:\tr X=0\}$ dan dekomposisi ortogonal setiap matriks $A$.
\end{exam}
\begin{solution}
    Karena $\langle I,X\rangle_F=\tr X$, diperoleh
    $W=\operatorname{span}\{I\}^\perp$ dan $W^\perp=\operatorname{span}\{I\}$.
    Dekomposisinya adalah
    \[
        A=\left(A-\frac{\tr A}{n}I\right)+\frac{\tr A}{n}I.
    \]
    Suku pertama bertrace nol dan suku kedua matriks skalar.
    Jarak $A$ ke $W$ adalah $|\tr A|/\sqrt n$.
\end{solution}

\section{Proyeksi dan Vektor Terdekat}

\begin{definisi}[][Proyeksi ortogonal][ortho:proyeksi]
    Jika $x=p+r$ dengan $p\in W$ dan $r\in W^\perp$,
    definisikan $P_Wx=p$.
    Operator $P_W$ disebut proyeksi ortogonal ke $W$.
\end{definisi}

Keunikan dekomposisi membuat definisi ini sah. Jika $(e_1,\ldots,e_r)$
basis ortonormal $W$, maka
\[
    P_Wx=\sum_{j=1}^r\langle e_j,x\rangle e_j.
\]
Rumus ini juga membuktikan linearitas, $P_W^2=P_W$,
$\operatorname{im}P_W=W$, dan $\ker P_W=W^\perp$.
Selain itu, $I-P_W=P_{W^\perp}$, serta
\[
    \|x\|^2=\|P_Wx\|^2+\|x-P_Wx\|^2.
\]
Dengan demikian $\|P_Wx\|\le\|x\|$, dengan kesamaan tepat ketika $x\in W$.

\begin{teorema}[][Karakterisasi pendekatan terbaik][ortho:terdekat]
    Untuk $p\in W$, pernyataan berikut ekuivalen:
    \begin{enumerate}[label=(\roman*),nosep]
        \item $p=P_Wx$;
        \item $x-p\in W^\perp$;
        \item $\|x-p\|\le\|x-w\|$ untuk semua $w\in W$.
    \end{enumerate}
    Vektor terdekat ini unik. Bahkan, untuk setiap $w\in W$,
    \[
        \|x-w\|^2=\|x-P_Wx\|^2+\|P_Wx-w\|^2.
    \]
\end{teorema}
\begin{proof}
    Dua pernyataan pertama adalah keunikan dekomposisi ortogonal.
    Untuk setiap $w\in W$, tulis $x-w=(x-P_Wx)+(P_Wx-w)$.
    Kedua suku ortogonal, sehingga Pythagoras memberi identitas
    yang dinyatakan. Nilainya minimum tepat ketika suku kedua
    nol, yaitu $w=P_Wx$. Ini membuktikan pernyataan ketiga dan keunikan.
\end{proof}

\begin{exam}[Proyeksi ke hiperbidang jumlah koordinat nol]
    Pada $\F^n$, letakkan $W=\{x:\sum_jx_j=0\}$.
    Tentukan $P_Wx$ dan jarak $x$ ke $W$.
\end{exam}
\begin{solution}
    Komplemen ortogonal $W$ dibangun oleh $\boldsymbol{1}=(1,\ldots,1)^T$.
    Dengan $\mu=\frac1n\sum_jx_j$, diperoleh
    \[
        P_Wx=x-\mu\boldsymbol{1},\qquad
        d(x,W)=\sqrt n\,|\mu|
          =\frac{|\sum_jx_j|}{\sqrt n}.
    \]
    Jadi proyeksi mengurangi setiap koordinat dengan rata-ratanya.
\end{solution}

\begin{exam}[Polinom linear terdekat dengan polinom kubik]
    Dalam $\mathcal P_{\le3}(\R)$ dengan hasil kali dalam integral
    pada $[-1,1]$, carilah polinom di $W=\operatorname{span}\{1,t\}$
    yang paling dekat dengan $t^3$.
\end{exam}
\begin{solution}
    Basis $(1,t)$ ortogonal. Koefisien proyeksi adalah
    \[
        \frac{\langle1,t^3\rangle}{\|1\|^2}=0,\qquad
        \frac{\langle t,t^3\rangle}{\|t\|^2}
          =\frac{2/5}{2/3}=\frac35.
    \]
    Jadi pendekatan terbaik adalah $3t/5$ dan galat kuadratnya
    $\|t^3\|^2-\|3t/5\|^2=2/7-6/25=8/175$.
\end{solution}

\section{Rumus Matriks dan Kuadrat Terkecil}

Jika $(w_1,\ldots,w_r)$ basis $W$ yang belum ortonormal,
tulis $P_Wx=\sum_jc_jw_j$. Syarat sisa ortogonal memberi
\[
    \sum_{j=1}^r\langle w_i,w_j\rangle c_j=\langle w_i,x\rangle,
    \qquad 1\le i\le r.
\]
Jadi $Gc=b$, dengan $G$ matriks Gram basis dan
$b_i=\langle w_i,x\rangle$. Karena basis bebas linear, $G$
positif definit dan $c=G^{-1}b$.

Untuk $W=\operatorname{im}A\subseteq\F^m$ dengan kolom-kolom
$A\in\Mat_{m\times r}(\F)$ bebas linear, rumus tersebut menjadi
\[
    P_W=A(A^*A)^{-1}A^*.
\]
Jika hasil kali dalamnya $x^*Hy$ dengan $H$ positif definit,
rumus yang sesuai adalah $P_W=A(A^*HA)^{-1}A^*H$.
Invers pada kedua rumus mensyaratkan kolom-kolom $A$ bebas linear.

\begin{exam}[Proyeksi berubah ketika hasil kali dalam berubah]
    Pada $\R^2$, gunakan $H=\begin{pmatrix}2&1\\1&2\end{pmatrix}$
    dan $\langle x,y\rangle_H=x^THy$.
    Tentukan proyeksi $(0,1)^T$ ke garis yang dibangun oleh $e_1$.
\end{exam}
\begin{solution}
    Koefisiennya adalah
    $\langle e_1,(0,1)^T\rangle_H/\langle e_1,e_1\rangle_H=1/2$.
    Maka proyeksinya $(1/2,0)^T$.
    Sisa $(-1/2,1)^T$ memang ortogonal terhadap $e_1$ dalam
    metrik $H$, meskipun tidak ortogonal dalam hasil kali dalam standar.
\end{solution}

\begin{teorema}[][Persamaan normal kuadrat terkecil][ortho:normal]
    Untuk $A\in\Mat_{m\times n}(\F)$ dan $b\in\F^m$, vektor
    $c$ meminimumkan $\|Ac-b\|$ jika dan hanya jika
    \[
        A^*Ac=A^*b.
    \]
    Solusi minimum selalu ada. Vektor $Ac$ yang dihasilkan unik,
    tetapi koefisien $c$ unik tepat ketika $\ker A=\{0\}$.
\end{teorema}
\begin{proof}
    Vektor $Ac$ terdekat dengan $b$ tepat ketika
    $b-Ac\perp\operatorname{im}A$.
    Ini ekuivalen dengan $A^*(b-Ac)=0$, yakni persamaan normal.
    Proyeksi ke $\operatorname{im}A$ selalu ada, dan menurut
    definisi citra proyeksi itu dapat ditulis sebagai $Ac$.
    Keunikan vektor proyeksi sudah diketahui. Jika $c_0$ salah
    satu solusi, seluruh solusi adalah $c_0+\ker A$.
\end{proof}

\begin{exam}[Garis yang paling sesuai dengan tiga data]
    Carilah garis $y=a+bt$ yang meminimumkan jumlah kuadrat
    galat pada data $(0,1)$, $(1,2)$, dan $(2,2)$.
\end{exam}
\begin{solution}
    Matriks desain dan vektor data adalah
    \[
        A=\begin{pmatrix}1&0\\1&1\\1&2\end{pmatrix},\qquad
        y=\begin{pmatrix}1\\2\\2\end{pmatrix}.
    \]
    Persamaan normal memberi
    \[
        \begin{pmatrix}3&3\\3&5\end{pmatrix}
        \begin{pmatrix}a\\b\end{pmatrix}
        =\begin{pmatrix}5\\6\end{pmatrix}.
    \]
    Maka $a=7/6$, $b=1/2$.
    Sisa $y-A(a,b)^T=(-1/6,1/3,-1/6)^T$ ortogonal terhadap
    kedua kolom $A$, dan jumlah kuadrat galat minimum adalah $1/6$.
\end{solution}

\section{Latihan Pilihan}

\begin{latihan}\label{ortho:lat-lengkapi}
    Lengkapi $u_1=(1,1,1,1)^T/2$ dan $u_2=(1,1,-1,-1)^T/2$
    menjadi basis ortonormal $\R^4$.
\end{latihan}

\begin{latihan}\label{ortho:lat-kompleks}
    Pada $\C^2$, letakkan $W=\operatorname{span}\{(1,i)^T\}$.
    Tentukan $W^\perp$, proyeksi $x=(1,1)^T$ ke $W$, dan jaraknya.
\end{latihan}

\begin{latihan}\label{ortho:lat-polinom}
    Lanjutkan ortonormalisasi $(1,t,t^2,t^3)$ terhadap hasil
    kali dalam integral pada $[-1,1]$. Berikan basis ortonormal
    lengkap bagi $\mathcal P_{\le3}(\R)$.
\end{latihan}

\begin{latihan}\label{ortho:lat-kuadratik}
    Tentukan real $a,b$ yang meminimumkan
    $\int_{-1}^1(t^2-a-bt)^2\,dt$, beserta nilai minimumnya.
\end{latihan}

\begin{latihan}\label{ortho:lat-kendala}
    Di antara $x\in\R^3$ yang memenuhi
    $x_1+x_2+x_3=1$ dan $x_1-x_3=2$, tentukan vektor bernorma
    terkecil dan norma kuadrat minimumnya.
\end{latihan}

\begin{latihan}\label{ortho:lat-dimensi}
    Dalam ruang berdimensi $n$, misalkan $u_1,\ldots,u_p$
    ortogonal dan semuanya tak nol. Jika $v_1,\ldots,v_q$
    masing-masing ortogonal terhadap semua $u_i$ dan $p+q>n$,
    buktikan bahwa daftar $v_1,\ldots,v_q$ bergantung linear.
\end{latihan}

\begin{latihan}\label{ortho:lat-dua-proyeksi}
    Misalkan $P,Q$ matriks proyeksi ortogonal pada $\C^n$.
    Buktikan bahwa $PQ$ merupakan matriks proyeksi ortogonal
    jika dan hanya jika $PQ=QP$. Dalam keadaan itu, tentukan citranya.
\end{latihan}

\begin{latihan}\label{ortho:lat-kontraksi}
    Misalkan $P^2=P$ pada ruang hasil kali dalam berdimensi hingga.
    Buktikan bahwa $P$ proyeksi ortogonal jika dan hanya jika
    $\|Px\|\le\|x\|$ untuk setiap $x$.
\end{latihan}

\ifpembahasan
\section{Pembahasan Latihan}

\subsection*{Latihan \ref{ortho:lat-lengkapi}: pola tanda yang saling meniadakan}
Ambil $u_3=(1,-1,1,-1)^T/2$ dan $u_4=(1,-1,-1,1)^T/2$.
Norma kuadrat masing-masing adalah empat kali $1/4$, yaitu satu.
Pada hasil kali dua vektor berbeda, terdapat dua suku $1/4$
dan dua suku $-1/4$, sehingga hasilnya nol.
Empat vektor ortonormal di $\R^4$ membentuk basis.

\subsection*{Latihan \ref{ortho:lat-kompleks}: konjugasi pada koefisien proyeksi}
Untuk $u=(1,i)^T$, persamaan $\langle u,(a,b)^T\rangle=a-ib=0$
memberi $W^\perp=\operatorname{span}\{(1,-i)^T\}$.
Karena $\langle u,x\rangle=1-i$ dan $\|u\|^2=2$,
\[
    P_Wx=\frac{1-i}{2}\begin{pmatrix}1\\i\end{pmatrix}
        =\begin{pmatrix}(1-i)/2\\(1+i)/2\end{pmatrix}.
\]
Norma kuadrat proyeksinya satu, sedangkan $\|x\|^2=2$.
Jadi jarak $x$ ke $W$ adalah satu.

\subsection*{Latihan \ref{ortho:lat-polinom}: pisahkan paritas polinom}
Tiga vektor pertama adalah $e_0,e_1,e_2$ pada contoh sebelumnya.
Karena $t^3$ ganjil, proyeksinya ke $e_0,e_2$ nol.
Sisa setelah mengurangi komponen pada $t$ adalah $t^3-3t/5$,
dengan norma kuadrat $8/175$. Jadi vektor terakhir dapat dipilih
\[
    e_3=\sqrt{\frac78}(5t^3-3t).
\]
Basis lengkapnya ialah
$1/\sqrt2$, $\sqrt{3/2}\,t$, $\sqrt{5/8}(3t^2-1)$,
dan $\sqrt{7/8}(5t^3-3t)$.

\subsection*{Latihan \ref{ortho:lat-kuadratik}: syarat sisa ortogonal}
Proyeksi $t^2$ ke $\operatorname{span}\{1,t\}$ mempunyai
koefisien $a=(2/3)/2=1/3$ dan $b=0$.
Nilai minimum adalah
\[
    \int_{-1}^1(t^2-1/3)^2\,dt=\frac25-\frac29=\frac8{45}.
\]
Keunikan proyeksi menjamin tidak ada pasangan lain yang mencapai minimum.

\subsection*{Latihan \ref{ortho:lat-kendala}: pilih solusi yang ortogonal terhadap kernel}
Dua vektor baris kendala, $(1,1,1)^T$ dan $(1,0,-1)^T$,
saling ortogonal. Solusi yang berada dalam ruang pembangun
keduanya adalah
\[
    x_0=\tfrac13(1,1,1)^T+(1,0,-1)^T
       =(4/3,1/3,-2/3)^T.
\]
Seluruh solusi lain berbentuk $x_0+t(1,-2,1)^T$.
Karena $x_0\perp(1,-2,1)^T$, norma kuadratnya adalah
$\|x_0\|^2+6t^2=7/3+6t^2$.
Minimum $7/3$ dicapai hanya pada $t=0$.

\subsection*{Latihan \ref{ortho:lat-dimensi}: hitung ruang yang masih tersedia}
Daftar $u_i$ bebas linear, sehingga ruang pembangunnya berdimensi $p$.
Semua $v_j$ berada dalam komplemen ortogonal yang berdimensi $n-p$.
Karena $q>n-p$, daftar tersebut pasti bergantung linear.

\subsection*{Latihan \ref{ortho:lat-dua-proyeksi}: simetri dan idempotensi}
Matriks proyeksi ortogonal memenuhi $P^*=P$ dan $P^2=P$,
dan hal yang sama berlaku untuk $Q$.
Sifat Hermitian dapat dilihat dari rumus
$P=EE^*$ dengan kolom-kolom $E$ basis ortonormal citranya.
Sebaliknya, matriks Hermitian idempoten mempunyai kernel
ortogonal terhadap citra, sehingga merupakan proyeksi ortogonal.
Jika $PQ$ proyeksi ortogonal, maka
$PQ=(PQ)^*=QP$.
Jika $PQ=QP$, maka $(PQ)^*=PQ$ dan $(PQ)^2=P^2Q^2=PQ$,
sehingga $PQ$ proyeksi ortogonal.
Citranya adalah $\operatorname{im}P\cap\operatorname{im}Q$:
setiap $PQx$ tetap oleh $P$ dan $Q$, dan setiap vektor pada
irisan tersebut tetap oleh $PQ$.

\subsection*{Latihan \ref{ortho:lat-kontraksi}: uji komponen kernel}
Arah perlu telah dibuktikan dengan Pythagoras.
Untuk arah cukup, ambil $u\in\operatorname{im}P$ dan $v\in\ker P$.
Untuk setiap skalar $t$, berlaku $P(u+tv)=u$, sehingga
$\|u\|^2\le\|u+tv\|^2$.
Jika $v\ne0$, pilih $t=-\langle v,u\rangle/\|v\|^2$.
Ekspansi norma kemudian memberi
\[
    \|u\|^2\le\|u\|^2-\frac{|\langle v,u\rangle|^2}{\|v\|^2},
\]
sehingga $\langle v,u\rangle=0$.
Kasus $v=0$ otomatis. Jadi citra dan kernel ortogonal.
Karena $P^2=P$, keduanya merupakan komplemen langsung;
maka $P$ tepat proyeksi ortogonal ke citranya.
\fi

% !TeX root = ../../ONMIPA Alin.tex
% Acuan: Zhang (2009), Bab 5, Soal 5.33, 5.36--5.51;
% Zhang (2011), bagian 1.4, hubungan kernel--citra dan adjoin.
% Seluruh operator dibahas dalam dimensi hingga; tidak memakai SVD, QR, atau LU.
\chapter[Operator pada Ruang Hasil Kali Dalam]{Operator pada Ruang\\Hasil Kali Dalam}
\label{ch:operator-hkd}

Hasil kali dalam memungkinkan sebuah transformasi dipindahkan
dari satu argumen ke argumen lainnya melalui operator adjoin.
Konsep ini menghubungkan kernel dengan komplemen ortogonal citra,
serta menjelaskan sifat geometris operator uniter dan adjoin diri.
Semua ruang pada bab ini berdimensi hingga atas $\R$ atau $\C$.
Konvensi hasil kali dalam tetap linear pada argumen kedua.

\section{Representasi Fungsional Linear}

Fungsional linear adalah transformasi linear $\ell:V\to\F$.
Untuk setiap $z\in V$, aturan $x\mapsto\langle z,x\rangle$
merupakan fungsional linear. Dalam dimensi hingga, semua
fungsional linear dapat diperoleh dengan cara tersebut.

\begin{teorema}[][Representasi Riesz dalam dimensi hingga][op:riesz]
    Untuk setiap fungsional linear $\ell:V\to\F$, terdapat
    tepat satu $z\in V$ dengan
    \[
        \ell(x)=\langle z,x\rangle\qquad\text{untuk semua }x\in V.
    \]
    Jika $V\ne\{0\}$, maka
    $\max_{\|x\|=1}|\ell(x)|=\|z\|$.
\end{teorema}
\begin{proof}
    Ambil basis ortonormal $e_1,\ldots,e_n$ dan letakkan
    $a_j=\ell(e_j)$. Untuk $x=\sum x_je_j$,
    $\ell(x)=\sum a_jx_j$.
    Vektor $z=\sum\overline{a_j}e_j$ memenuhi
    $\langle z,x\rangle=\sum a_jx_j$.
    Jika $z'$ juga merepresentasikan $\ell$, maka
    $\langle z-z',x\rangle=0$ untuk semua $x$.
    Pilih $x=z-z'$ untuk memperoleh $z=z'$.
    Cauchy--Schwarz memberi $|\ell(x)|\le\|z\|$ pada vektor satuan.
    Jika $z\ne0$, batas dicapai pada $x=z/\|z\|$; jika $z=0$,
    semua nilai fungsional nol.
\end{proof}

Pada ruang kompleks, pemetaan $z\mapsto\langle z,\cdot\rangle$
bersifat konjugat-linear terhadap $z$. Konjugasi pada koefisien
representer dalam bukti di atas tidak boleh dihapus.

\begin{exam}[Kendala linear sebagai hiperbidang ortogonal]
    Pada $\C^3$, ambil $\ell(x)=2x_1-ix_2+3x_3$.
    Tentukan representer, nilai maksimum $|\ell(x)|$ pada
    vektor satuan, dan solusi bernorma minimum dari $\ell(x)=b$.
\end{exam}
\begin{solution}
    Representernya $z=(2,i,3)^T$, dengan $\|z\|^2=14$.
    Maka maksimum yang diminta adalah $\sqrt{14}$ dan
    $\ker\ell=\operatorname{span}\{z\}^\perp$.
    Solusi minimum adalah $x_0=(b/14)z$ karena
    $\langle z,x_0\rangle=b$.
    Semua solusi lain berbentuk $x_0+w$ dengan $w\perp z$,
    sehingga norma kuadratnya $|b|^2/14+\|w\|^2$.
\end{solution}

\section{Adjoin dan Matriks Representasinya}

\begin{definisi}[][Adjoin transformasi linear][op:adjoin]
    Adjoin dari $T:V\to W$ adalah transformasi $T^*:W\to V$
    yang memenuhi
    \[
        \langle Tx,y\rangle_W=\langle x,T^*y\rangle_V
        \qquad(x\in V,\ y\in W).
    \]
\end{definisi}

\begin{teorema}[][Keberadaan dan keunikan adjoin][op:ada-adjoin]
    Setiap transformasi linear antara dua ruang hasil kali dalam
    berdimensi hingga mempunyai tepat satu adjoin.
    Jika matriks $T$ dalam basis ortonormal adalah $A$, matriks
    adjoin dalam arah sebaliknya adalah $A^*$.
\end{teorema}
\begin{proof}
    Pilih basis ortonormal pada kedua ruang, dan tulis koordinat
    $x,y$ sebagai $u,v$. Identitas
    \[
        (Au)^*v=u^*(A^*v)
    \]
    menunjukkan bahwa transformasi bermatriks $A^*$ memenuhi
    definisi adjoin. Untuk keunikan, jika $S$ juga memenuhi definisi,
    maka $\langle x,(S-T^*)y\rangle=0$ untuk semua $x,y$.
    Dengan memilih $x=(S-T^*)y$, diperoleh $(S-T^*)y=0$
    untuk semua $y$, sehingga $S=T^*$.
\end{proof}

Jika basis yang digunakan tidak ortonormal, misalkan matriks
hasil kali dalam domain dan kodomain adalah $G_V$ dan $G_W$.
Untuk $A=[T]_{\mathcal C\leftarrow\mathcal B}$, identitas adjoin memberi
\[
    [T^*]_{\mathcal B\leftarrow\mathcal C}=G_V^{-1}A^*G_W.
\]
Memang, kedua ruas identitas hasil kali dalam dalam koordinat
adalah $u^*A^*G_Wv$ dan $u^*G_V[T^*]v$.
Matriks adjoin baru sama dengan transpose konjugat matriks asal
ketika basis-basisnya ortonormal, atau dalam keadaan khusus yang sesuai.

\begin{exam}[Adjoin diri dalam hasil kali dalam berbobot]
    Pada $\R^2$ dengan $G=\diag(1,2)$, tinjau operator
    bermatriks $A=\begin{pmatrix}0&2\\1&0\end{pmatrix}$.
    Tentukan matriks adjoinnya.
\end{exam}
\begin{solution}
    Perhitungan memberi
    \[
        G^{-1}A^TG
        =\begin{pmatrix}1&0\\0&1/2\end{pmatrix}
         \begin{pmatrix}0&1\\2&0\end{pmatrix}
         \begin{pmatrix}1&0\\0&2\end{pmatrix}=A.
    \]
    Jadi operator sama dengan adjoinnya walaupun $A$ tidak
    simetris dalam basis standar tersebut.
    Vektor eigen $(\sqrt2,1)^T$ dan $(-\sqrt2,1)^T$
    ortogonal dalam metrik $G$, karena hasil kali dalamnya $-2+2=0$.
\end{solution}

\begin{prop}[Aljabar adjoin]
    Untuk transformasi dengan domain dan kodomain yang sesuai,
    \[
        (T^*)^*=T,\qquad
        (aS+bT)^*=\overline aS^*+\overline bT^*,\qquad
        (ST)^*=T^*S^*.
    \]
    Jika $T$ invertibel, maka $(T^{-1})^*=(T^*)^{-1}$.
\end{prop}
\begin{proof}
    Masing-masing rumus dapat diperiksa melalui identitas
    pendefinisi adjoin dan keunikannya.
    Misalnya,
    $\langle STx,y\rangle=\langle Tx,S^*y\rangle
    =\langle x,T^*S^*y\rangle$.
    Rumus invers mengikuti pengambilan adjoin pada
    $TT^{-1}=I$ dan $T^{-1}T=I$.
\end{proof}

\begin{exam}[Adjoin operator pada ruang matriks]
    Pada ruang matriks dengan hasil kali dalam Frobenius,
    tentukan adjoin $L(X)=AXB$, dengan ukuran matriks yang sesuai.
\end{exam}
\begin{solution}
    Sifat siklik trace memberi
    \[
        \langle AXB,Y\rangle_F
        =\tr(B^*X^*A^*Y)
        =\tr(X^*A^*YB^*)=\langle X,A^*YB^*\rangle_F.
    \]
    Jadi $L^*(Y)=A^*YB^*$.
    Secara khusus, adjoin dari $X\mapsto AX-XA$ adalah
    $Y\mapsto A^*Y-YA^*$.
\end{solution}

\section{Hubungan Kernel, Citra, dan Ortogonalitas}

\begin{teorema}[][Empat ruang fundamental melalui adjoin][op:empat-ruang]
    Untuk $T:V\to W$,
    \[
        \ker T^*=(\operatorname{im}T)^\perp,\qquad
        \operatorname{im}T^*=(\ker T)^\perp.
    \]
    Akibatnya,
    \[
        V=\ker T\oplus\operatorname{im}T^*,\qquad
        W=\operatorname{im}T\oplus\ker T^*,
    \]
    dengan kedua jumlah ortogonal, dan $\rank T=\rank T^*$.
\end{teorema}
\begin{proof}
    Vektor $y\in W$ berada di $(\operatorname{im}T)^\perp$
    tepat ketika $\langle Tx,y\rangle=0$ untuk semua $x$.
    Dengan identitas adjoin, hal ini ekuivalen dengan
    $\langle x,T^*y\rangle=0$ untuk semua $x$, yakni $T^*y=0$.
    Terapkan hasil ini pada $T^*$ dan ambil komplemen ortogonal
    untuk memperoleh identitas kedua.
    Dekomposisi mengikuti teorema komplemen ortogonal.
    Terakhir, $\rank T^*=\dim V-\dim\ker T=\rank T$.
\end{proof}

Selain itu,
\[
    \langle x,T^*Tx\rangle=\|Tx\|^2,
    \qquad \ker(T^*T)=\ker T.
\]
Satu inklusi kernel jelas; inklusi lainnya mengikuti bahwa
$T^*Tx=0$ mengakibatkan $\|Tx\|^2=0$.
Jadi $\rank(T^*T)=\rank T$.
Dalam masalah kuadrat terkecil, hubungan ini menjelaskan mengapa
matriks persamaan normal invertibel tepat ketika kolom matriks
asal bebas linear.

Jika $T$ surjektif, setiap persamaan $Tx=y$ mempunyai satu
solusi yang berada dalam $(\ker T)^\perp$; solusi tersebut
adalah solusi dengan norma minimum. Keunikannya mengikuti
dekomposisi ortogonal, sebab solusi lain berbeda sebesar
vektor dari $\ker T$.

\section{Isometri dan Operator Uniter}

\begin{definisi}[][Isometri linear dan operator uniter][op:isometri]
    Transformasi linear $T:V\to W$ disebut isometri jika
    $\langle Tx,Ty\rangle_W=\langle x,y\rangle_V$ untuk semua $x,y$.
    Isometri bijektif pada ruang yang sama disebut operator
    \emph{uniter} pada ruang kompleks dan \emph{ortogonal} pada ruang real.
\end{definisi}

\begin{teorema}[][Karakterisasi isometri linear][op:uji-isometri]
    Untuk transformasi linear $T:V\to W$, pernyataan berikut ekuivalen:
    \begin{enumerate}[label=(\roman*),nosep]
        \item $T$ mempertahankan hasil kali dalam;
        \item $\|Tx\|=\|x\|$ untuk setiap $x$;
        \item $T^*T=I_V$;
        \item citra suatu basis ortonormal $V$ adalah daftar ortonormal di $W$.
    \end{enumerate}
    Isometri selalu injektif. Jika $\dim V=\dim W$, ia bijektif
    dan $T^{-1}=T^*$.
\end{teorema}
\begin{proof}
    Pernyataan pertama memberi kedua dengan mengambil $y=x$.
    Arah balik memakai identitas polarisasi dan linearitas $T$.
    Identitas $\langle Tx,Ty\rangle=\langle x,T^*Ty\rangle$
    membuktikan ekuivalensi dengan pernyataan ketiga.
    Memeriksa pada pasangan anggota basis memberi pernyataan
    keempat; ekspansi koordinat membuktikan arah baliknya.
    Pelestarian norma memaksa kernel hanya nol.
    Jika dimensi sama, injektivitas memberi bijektivitas dan
    $T^*T=I$ kemudian memberi $T^*=T^{-1}$.
\end{proof}

\begin{exam}[Pelestarian norma tanpa linearitas]
    Pemetaan $C(z)=\overline z$ pada $\C$ mempertahankan norma
    dan jarak, tetapi bukan linear kompleks: $C(i)=-i\ne iC(1)$.
    Pemetaan translasi $x\mapsto x+a$ pada $\R^n$, dengan $a\ne0$,
    mempertahankan jarak tetapi tidak linear karena tidak
    memetakan nol ke nol. Jadi asumsi linearitas dalam teorema
    pelestarian norma tidak boleh dihilangkan.
\end{exam}

Dalam basis ortonormal, isometri mempunyai matriks $A$ dengan
$A^*A=I$. Jika matriks persegi, ini ekuivalen dengan uniter.
Pergantian antara dua basis ortonormal juga diberikan oleh
matriks uniter, sehingga representasi suatu operator berubah
melalui $A\mapsto U^*AU$.
Untuk basis dengan matriks metrik $G$, operator bermatriks $A$
mempertahankan hasil kali dalam tepat ketika $A^*GA=G$.

\section{Operator Adjoin Diri, Normal, dan Positif}

\begin{definisi}[][Kelas operator berdasarkan adjoin][op:kelas]
    Operator $T:V\to V$ disebut \emph{adjoin diri} jika $T^*=T$,
    \emph{antiadjoin diri} jika $T^*=-T$, dan \emph{normal}
    jika $T^*T=TT^*$.
    Operator adjoin diri disebut semidefinit positif jika
    $\langle x,Tx\rangle\ge0$ untuk semua $x$, dan positif definit
    jika ketaksamaan itu ketat untuk $x\ne0$.
\end{definisi}

Dalam basis ortonormal, sifat-sifat ini tepat sifat Hermitian,
anti-Hermitian, normal, serta positif pada matriks representasinya.
Operator adjoin diri dan operator uniter selalu normal.
Untuk operator adjoin diri, $\langle x,Tx\rangle$ real karena
sama dengan konjugatnya.

\begin{teorema}[][Teorema spektral untuk operator][op:spektral]
    Setiap operator adjoin diri mempunyai basis ortonormal
    yang terdiri atas vektor eigen, dengan nilai eigen real.
    Pada ruang kompleks, setiap operator normal mempunyai
    basis ortonormal vektor eigen.
\end{teorema}
\begin{proof}
    Pilih basis ortonormal awal. Matriks operator menjadi
    Hermitian pada kasus adjoin diri, atau normal pada kasus kedua.
    Terapkan teorema spektral matriks pada
    Teorema~\ref{eigen:teorema-spektral}.
    Kolom-kolom matriks perubahan basis uniter memberikan
    koordinat basis ortonormal vektor eigen operator.
    Pada ruang real dan operator adjoin diri, matriksnya simetris
    real sehingga basis eigen dapat dipilih real.
\end{proof}

Kerealan nilai eigen adjoin diri juga dapat dilihat langsung.
Jika $Tx=\lambda x$ dan $x\ne0$, maka
$\lambda\|x\|^2=\langle x,Tx\rangle$ real.
Jika $Tu=\lambda u$ dan $Tv=\mu v$ dengan $\lambda\ne\mu$,
identitas adjoin memberi $(\lambda-\mu)\langle u,v\rangle=0$;
jadi ruang eigen berbeda saling ortogonal.

\begin{exam}[Normal real belum tentu mempunyai basis eigen real]
    Matriks $R=\begin{pmatrix}0&-1\\1&0\end{pmatrix}$ ortogonal,
    sehingga normal. Namun, polinom karakteristiknya $t^2+1$
    tidak mempunyai akar real. Jadi tidak ada basis eigen real.
    Setelah ruang diperluas ke $\C^2$, nilai eigennya $i,-i$
    dan teorema spektral kompleks berlaku.
\end{exam}

Jika $T$ normal, ekspansi dalam basis eigen ortonormal memberi
$\|Tx\|=\|T^*x\|$ dan $\ker T=\ker T^*$.
Pada ruang real, kesamaan norma langsung mengikuti
$\langle x,(T^*T-TT^*)x\rangle=0$.
Oleh karena itu, operator normal memenuhi
$V=\ker T\oplus\operatorname{im}T$ secara ortogonal.
Operator normal yang nilpoten harus nol: pada ruang kompleks,
semua nilai eigen nol dan diagonalisasi memberi $T=0$;
untuk matriks real, argumen yang sama dapat diterapkan setelah
dipandang sebagai matriks kompleks.

\section{Hasil Bagi Rayleigh dan Dimensi Subruang Positif}

\begin{teorema}[][Batas hasil bagi Rayleigh][op:rayleigh]
    Misalkan $T$ adjoin diri pada ruang tak nol, dengan nilai
    eigen terkecil $\lambda_{\min}$ dan terbesar $\lambda_{\max}$.
    Untuk $x\ne0$,
    \[
        \lambda_{\min}\le
        \frac{\langle x,Tx\rangle}{\|x\|^2}
        \le\lambda_{\max}.
    \]
    Kesamaan pada masing-masing batas berlaku tepat ketika
    $x$ berada dalam ruang eigen nilai batas tersebut.
\end{teorema}
\begin{proof}
    Tulis $x=\sum c_je_j$ dalam basis eigen ortonormal.
    Maka hasil bagi yang dimaksud adalah
    \[
        \frac{\sum_j\lambda_j|c_j|^2}{\sum_j|c_j|^2},
    \]
    yakni rata-rata berbobot nilai-nilai eigen.
    Batas atas tercapai tepat ketika setiap koefisien untuk
    nilai eigen yang lebih kecil nol. Argumen batas bawah serupa.
\end{proof}

\begin{exam}[Maksimum bentuk kuadratik pada lingkaran satuan]
    Tentukan rentang $2x^2+2xy+2y^2$ untuk real $x^2+y^2=1$.
\end{exam}
\begin{solution}
    Matriks bentuk kuadratiknya $A=\begin{pmatrix}2&1\\1&2\end{pmatrix}$
    mempunyai nilai eigen $1$ dan $3$, masing-masing pada arah
    $(1,-1)^T$ dan $(1,1)^T$.
    Jadi minimum satu dicapai ketika $x=-y=\pm1/\sqrt2$,
    dan maksimum tiga ketika $x=y=\pm1/\sqrt2$.
    Semua nilai antara keduanya dicapai dengan mengubah
    koefisien dalam basis eigen ortonormal.
\end{solution}

\begin{teorema}[][Banyak nilai eigen positif][op:indeks-positif]
    Jika operator adjoin diri $T$ mempunyai tepat $r$ nilai
    eigen positif, dihitung dengan multiplisitas, maka $r$
    adalah dimensi terbesar subruang $U$ yang memenuhi
    \[
        \langle x,Tx\rangle>0\qquad\text{untuk setiap }x\in U\setminus\{0\}.
    \]
\end{teorema}
\begin{proof}
    Jumlah ruang eigen untuk nilai positif mempunyai dimensi $r$
    dan memenuhi ketaksamaan tersebut.
    Ruang $N$ yang dibangun oleh vektor eigen untuk nilai
    nonpositif berdimensi $n-r$.
    Jika $\dim U>r$, rumus dimensi memberi
    $\dim(U\cap N)\ge\dim U+(n-r)-n>0$.
    Vektor tak nol pada irisan menghasilkan nilai kuadrat
    positif menurut definisi $U$ dan nonpositif menurut definisi
    $N$, kontradiksi.
\end{proof}

Dengan mengganti $T$ oleh $-T$, diperoleh pernyataan serupa
untuk banyaknya nilai eigen negatif. Hasil ini sering mengubah
informasi tentang suatu subruang menjadi batas jumlah nilai
eigen tanpa menghitung polinom karakteristik.

\section{Latihan Pilihan}

\begin{latihan}\label{op:lat-riesz}
    Pada $\mathcal P_{\le2}(\R)$ dengan hasil kali dalam integral
    pada $[-1,1]$, tentukan representer dari $\ell(p)=p(0)$.
    Carilah konstanta terkecil $C$ sehingga
    \[
        |p(0)|^2\le C\int_{-1}^1|p(t)|^2\,dt
    \]
    untuk semua polinom pada ruang tersebut.
\end{latihan}

\begin{latihan}\label{op:lat-adjoin-berbobot}
    Pada $\R^2$, gunakan metrik $G=\begin{pmatrix}2&1\\1&2\end{pmatrix}$.
    Jika matriks $T$ dalam basis standar adalah
    $A=\begin{pmatrix}1&1\\0&2\end{pmatrix}$, tentukan matriks
    $T^*$. Apakah $T$ adjoin diri?
\end{latihan}

\begin{latihan}\label{op:lat-komutator}
    Untuk matriks Hermitian $A\in\Mat_n(\C)$, definisikan
    $D_A(X)=AX-XA$ pada ruang matriks dengan hasil kali dalam
    Frobenius. Buktikan bahwa $D_A$ adjoin diri.
    Buktikan pula bahwa $D_A$ semidefinit positif tepat ketika
    $A$ merupakan matriks skalar.
\end{latihan}

\begin{latihan}\label{op:lat-titik-tetap}
    Jika $U$ uniter, buktikan dekomposisi ortogonal
    \[
        V=\ker(U-I)\oplus\operatorname{im}(U-I).
    \]
    Tafsirkan kedua subruang sebagai vektor tetap dan selisih
    vektor dengan citranya.
\end{latihan}

\begin{latihan}\label{op:lat-gram-sama}
    Misalkan $v_1,\ldots,v_m$ dan $w_1,\ldots,w_m$ dua daftar
    dalam ruang hasil kali dalam $V$ dengan
    $\langle v_i,v_j\rangle=\langle w_i,w_j\rangle$ untuk semua $i,j$.
    Buktikan bahwa ada operator uniter atau ortogonal $U$ dengan
    $Uv_j=w_j$ untuk setiap $j$. Daftar tidak harus bebas linear.
\end{latihan}

\begin{latihan}\label{op:lat-norma-sama}
    Untuk transformasi linear $A,B:V\to W$, andaikan
    $\|Ax\|=\|Bx\|$ untuk semua $x$.
    Buktikan bahwa ada operator uniter atau ortogonal $U$ pada $W$
    sehingga $A=UB$.
\end{latihan}

\begin{latihan}\label{op:lat-kuadratik-real}
    Pada ruang hasil kali dalam kompleks, buktikan bahwa
    $\langle x,Tx\rangle$ real untuk semua $x$ jika dan hanya jika
    $T$ adjoin diri. Jelaskan mengapa pernyataan ini tidak menjadi
    kriteria adjoin diri pada ruang real.
\end{latihan}

\begin{latihan}\label{op:lat-isotropik}
    Misalkan $T$ adjoin diri mempunyai $p$ nilai eigen positif,
    $q$ negatif, dan $z$ nol, dihitung dengan multiplisitas.
    Tentukan dimensi terbesar subruang $W$ yang memenuhi
    $\langle x,Tx\rangle=0$ untuk setiap $x\in W$.
\end{latihan}

\begin{latihan}\label{op:lat-pilih-metrik}
    Buktikan bahwa suatu operator pada ruang linear berdimensi
    hingga dapat dibuat adjoin diri dengan memilih suatu hasil
    kali dalam tepat ketika operator itu dapat didiagonalkan
    dengan semua nilai eigen real.
    Simpulkan bahwa setiap operator idempoten dapat dibuat
    menjadi proyeksi ortogonal untuk suatu hasil kali dalam.
\end{latihan}

\ifpembahasan
\section{Pembahasan Latihan}

\subsection*{Latihan \ref{op:lat-riesz}: evaluasi sebagai hasil kali dalam}
Tuliskan representer $r(t)=a+bt+ct^2$.
Syarat $\langle r,1\rangle=1$, $\langle r,t\rangle=0$,
dan $\langle r,t^2\rangle=0$ memberi
\[
    2a+\frac23c=1,\qquad \frac23b=0,\qquad
    \frac23a+\frac25c=0.
\]
Maka $r(t)=9/8-15t^2/8$.
Norma kuadrat representer adalah $\langle r,r\rangle=\ell(r)=9/8$.
Jadi konstanta terbaik $C=9/8$, dengan kesamaan untuk
polinom yang merupakan kelipatan $r$, termasuk polinom nol.

\subsection*{Latihan \ref{op:lat-adjoin-berbobot}: transpose saja belum cukup}
Rumus adjoin dalam basis ini memberi
\[
    [T^*]=G^{-1}A^TG
       =\frac13\begin{pmatrix}2&-1\\-1&2\end{pmatrix}
         \begin{pmatrix}2&1\\4&5\end{pmatrix}
       =\begin{pmatrix}0&-1\\2&3\end{pmatrix}.
\]
Matriks tersebut berbeda dari $A$, sehingga $T$ bukan adjoin diri.
Ia juga berbeda dari $A^T$, sebab basis standar tidak ortonormal
untuk metrik $G$.

\subsection*{Latihan \ref{op:lat-komutator}: selisih nilai eigen muncul berpasangan}
Rumus adjoin pada ruang matriks memberi $D_A^*=D_{A^*}=D_A$.
Pilih basis eigen ortonormal $u_1,\ldots,u_n$ bagi $A$,
dengan nilai eigen real $\lambda_j$.
Matriks $u_iu_j^*$ adalah vektor eigen $D_A$ untuk nilai
$\lambda_i-\lambda_j$.
Jika $D_A$ semidefinit positif, semua selisih itu nonnegatif.
Menukar $i,j$ memaksa semuanya nol, sehingga semua nilai
eigen $A$ sama. Teorema spektral memberi $A=cI$.
Sebaliknya, $A=cI$ menghasilkan $D_A=0$, yang semidefinit positif.

\subsection*{Latihan \ref{op:lat-titik-tetap}: kernel adjoin sama dengan kernel asal}
Karena $U^*=U^{-1}$, persamaan $U^*x=x$ ekuivalen dengan $Ux=x$.
Jadi $\ker((U-I)^*)=\ker(U-I)$.
Identitas kernel--citra memberi
\[
    \bigl(\operatorname{im}(U-I)\bigr)^\perp=\ker(U-I).
\]
Dekomposisi ortogonal mengikuti dengan mengambil komplemen.
Kernel berisi vektor yang tetap oleh $U$, sedangkan citra
berisi semua $Ux-x$. Mengganti tandanya memberi himpunan
yang sama, sehingga dapat pula ditulis sebagai semua $x-Ux$.

\subsection*{Latihan \ref{op:lat-gram-sama}: perluas isometri pada ruang pembangun}
Definisikan dahulu $U_0(\sum c_jv_j)=\sum c_jw_j$.
Kesamaan matriks Gram memberi
\[
    \left\|\sum c_jv_j\right\|^2
      =\left\|\sum c_jw_j\right\|^2.
\]
Jika dua representasi sumber sama, selisihnya bernorma nol;
persamaan tersebut memaksa selisih citranya nol juga.
Jadi $U_0$ terdefinisi dengan baik, linear, dan merupakan
isometri dari ruang pembangun $v_j$ ke ruang pembangun $w_j$.
Pilih basis ortonormal pada ruang pertama; citranya menjadi
basis ortonormal pada ruang kedua. Perluas kedua basis ke
basis ortonormal $V$, lalu kirim vektor tambahan secara berpasangan.
Operator yang dihasilkan uniter atau ortogonal dan memenuhi semua syarat.

\subsection*{Latihan \ref{op:lat-norma-sama}: cocokkan kedua citra}
Definisikan $U_0(Bx)=Ax$ pada $\operatorname{im}B$.
Jika $Bx=By$, hipotesis pada $x-y$ memberi
$\|A(x-y)\|=\|B(x-y)\|=0$, sehingga definisinya sah.
Pemetaan ini linear, mempertahankan norma, dan surjektif ke
$\operatorname{im}A$. Dengan polarisasi, ia mempertahankan
hasil kali dalam. Perluas isometri ini ke operator uniter
atau ortogonal pada $W$ dengan melengkapi basis ortonormal
kedua citra seperti pada latihan sebelumnya. Maka $U(Bx)=Ax$
untuk setiap $x$, yakni $A=UB$.

\subsection*{Latihan \ref{op:lat-kuadratik-real}: pisahkan dua bagian adjoin diri}
Arah perlu mengikuti simetri konjugat dan $T^*=T$.
Untuk arah cukup, tulis
\[
    T=H+iK,\qquad H=\frac{T+T^*}{2},\qquad
    K=\frac{T-T^*}{2i}.
\]
Keduanya adjoin diri. Karena
$\langle x,Tx\rangle=\langle x,Hx\rangle+i\langle x,Kx\rangle$
real, diperoleh $\langle x,Kx\rangle=0$ untuk semua $x$.
Uji pada basis eigen ortonormal $K$ menunjukkan bahwa setiap
nilai eigennya nol, sehingga $K=0$ dan $T=H$ adjoin diri.
Pada ruang real, setiap hasil kali dalam bernilai real secara
otomatis. Rotasi $\begin{pmatrix}0&-1\\1&0\end{pmatrix}$,
misalnya, tidak adjoin diri walaupun semua nilai kuadratnya real.

\subsection*{Latihan \ref{op:lat-isotropik}: batas dimensi dan konstruksi}
Tuliskan dekomposisi spektral ortogonal
$V=V_+\oplus V_-\oplus V_0$, dengan dimensi $p,q,z$.
Jika bentuk kuadratik nol pada $W$, maka
$W\cap V_+=W\cap V_-=\{0\}$.
Rumus dimensi memberi $\dim W\le q+z$ dan $\dim W\le p+z$.
Jadi batas atasnya $z+\min\{p,q\}$.
Untuk mencapai batas, pilih vektor eigen ortonormal $e_i$ untuk
nilai positif $\lambda_i$ dan $f_i$ untuk nilai negatif $-\mu_i$,
dengan $\lambda_i,\mu_i>0$.
Gabungkan basis $V_0$ dengan
\[
    w_i=\frac{e_i}{\sqrt{\lambda_i}}+
        \frac{f_i}{\sqrt{\mu_i}},\qquad 1\le i\le\min\{p,q\}.
\]
Daftar itu bebas linear. Untuk setiap kombinasinya, kontribusi
positif dan negatif pada bentuk kuadratik saling meniadakan,
sedangkan komponen $V_0$ tidak menyumbang apa pun.
Jadi dimensi terbesar tepat $z+\min\{p,q\}$.

\subsection*{Latihan \ref{op:lat-pilih-metrik}: jadikan basis eigen ortonormal}
Jika operator adjoin diri untuk suatu hasil kali dalam,
teorema spektral memberi basis eigen dan nilai eigen real.
Sebaliknya, pilih basis eigen dengan nilai eigen real dan
definisikan hasil kali dalam dengan menjadikan basis itu ortonormal.
Matriks operator dalam basis tersebut diagonal real, sehingga
operator adjoin diri.
Jika matriks kolom basis eigen dalam koordinat awal adalah $S$,
pilihan metriknya secara eksplisit
$G=(S^{-1})^*S^{-1}$.
Operator idempoten mempunyai dekomposisi
$V=\operatorname{im}P\oplus\ker P$, dan matriksnya dalam basis
yang sesuai adalah $\diag(I,0)$.
Jadikan basis tersebut ortonormal. Citra dan kernel sekarang
ortogonal, sehingga $P$ menjadi proyeksi ortogonal.
\fi


\nocite{zhang2009,zhang2011}
\bibliographystyle{plainnat}
\bibliography{references}
\end{document}
